%% parts/p10/algebraic-and-categorical-logic.tex --- Part 10
\chapter{Algebraic and Categorical Logic}
\label{sec:p10-algebraic-and-categorical-logic}

Algebraic logic replaces formulas by equivalence classes and logical
operations by algebraic operations; categorical logic studies syntax,
semantics, and substitution through universal mappings. These are not merely
changes of notation. Each translation below states the language, algebra or
category, and preservation claim that makes the comparison meaningful.

\section{Universal algebra and equational consequence}
\label{sec:p10-universal-algebra}

\begin{hbdefinition}{def:p10-algebra}{Algebras, homomorphisms, and congruences}
For a finitary signature \(\tau\), a \(\tau\)-algebra \(A\) is a nonempty set
equipped with an \(n\)-ary operation \(f^A\) for each \(n\)-ary symbol \(f\).
A homomorphism \(h:A\to B\) satisfies
\(h(f^A(\bar a))=f^B(h\bar a)\). A congruence is an equivalence relation
compatible with every basic operation, and \(A/\theta\) carries the induced
quotient operations. These names are necessary to distinguish operation-
preserving maps and quotient-compatible equivalences from arbitrary functions
and relations. They stand as algebraic objects and relations and are used in
the representation and completeness theorems below.
\index{algebra}\index{homomorphism!algebraic}\index{congruence}
\end{hbdefinition}

Subalgebras, products, and homomorphic images are denoted by the closure
operators \(\mathsf S\), \(\mathsf P\), and \(\mathsf H\). A class of similar
algebras is a \emph{variety} when it is the class of all models of a set of
equations.

\begin{hbtheorem}{thm:p10-hsp}{cited=Birkhoff1935}{Birkhoff HSP theorem}
A class of finitary \(\tau\)-algebras is a variety if and only if it is closed
under homomorphic images, subalgebras, and arbitrary direct products, where
arbitrary products include the empty product (the one-element
\(\tau\)-algebra). Thus an HSP-closed class is nonempty under this convention.
\end{hbtheorem}

%% Terms, identities, equational deduction, and free algebras.

\subsection{Terms, identities, and equational deduction}
\label{sec:p10-equational-logic}

\begin{hbdefinition}{def:p10-term-algebra}{Term algebra and evaluation}
Fix a finitary signature \(\tau\). For a set \(X\) disjoint from the
operation symbols, \(T_\tau(X)\) is the set of finite terms generated
by \(X\) and the symbols of \(\tau\) with their declared arities.
An \(n\)-ary basic operation on this set sends
\((t_1,\ldots,t_n)\) to the term \(f(t_1,\ldots,t_n)\).
In particular, a nullary symbol is itself a term.
Under the nonempty-algebra convention, this term algebra exists when
\(X\ne\varnothing\) or \(\tau\) has a nullary symbol.
If neither condition holds, there are no terms and the empty set is
not an algebra in the present convention.

For a \(\tau\)-algebra \(A\) and a map \(v:X\to A\), evaluation is
defined recursively by
\[
 x^A(v)=v(x),\qquad
 f(t_1,\ldots,t_n)^A(v)
       =f^A(t_1^A(v),\ldots,t_n^A(v)).
\]
A substitution \(\sigma:X\to T_\tau(Y)\) replaces each variable
simultaneously and extends recursively to a map
\(\widehat\sigma:T_\tau(X)\to T_\tau(Y)\).
Write \(t\sigma=\widehat\sigma(t)\).
No variable is bound in a term.
These constructions use the term-algebra presentation of
\cite[II\S10, Definitions~10.1--10.4]{BurrisSankappanavar2012}.
\index{term algebra}\index{substitution!equational}
\end{hbdefinition}

\begin{hbtheorem}{thm:p10-term-universal}{proved=annex:p10.eq.1}{Universal property and substitution law for terms}
If \(T_\tau(X)\) is nonempty, every \(v:X\to A\) extends to a unique
homomorphism \(T_\tau(X)\to A\), namely term evaluation.
For \(\sigma:X\to T_\tau(Y)\) and \(v:Y\to A\),
\[
 (t\sigma)^A(v)=t^A(w),
 \qquad w(x)=(\sigma(x))^A(v).
\]
Every algebra homomorphism \(h:A\to B\) also satisfies
\[
 h(t^A(v))=t^B(h\circ v).
\]
\end{hbtheorem}

\begin{hbdefinition}{def:p10-equational-logic}{Equational consequence}
Fix a variable set \(V\) for which \(T_\tau(V)\) is nonempty.
An equation \(s\approx t\) is an \emph{identity of \(A\)} when
\[
 A\models s\approx t
 \quad\Longleftrightarrow\quad
 \forall v:V\to A\quad s^A(v)=t^A(v).
\]
Write \(A\models\Gamma\) when every equation in \(\Gamma\) is an
identity of \(A\). Equational consequence is
\[
 \Gamma\models_{\mathrm{alg}}s\approx t
 \quad\Longleftrightarrow\quad
 \forall A\,
 (A\models\Gamma\ \Longrightarrow\ A\models s\approx t),
\]
where \(A\) ranges over nonempty \(\tau\)-algebras.
Thus each premise is required under every valuation.

The judgment \(\Gamma\vdash_{\mathrm{eq}}s\approx t\) means that a
finite derivation uses equations from \(\Gamma\), reflexivity,
symmetry, transitivity, compatibility with the basic operations,
and simultaneous term substitution. The complete rule forms are on
Sheet~\ref{sheet:equational-logic}.
Compatibility is equivalent to replacing one occurrence of a term
by an equal term inside a larger term; the proof below establishes
this correspondence with the source's replacement rule.
The semantic and deductive presentations are those of
\cite[II\S11, Definition~11.1; II\S14, Definitions~14.11,
14.16, 14.18]{BurrisSankappanavar2012}.
\index{equational logic}\index{identity!algebraic}
\end{hbdefinition}

\showcalculus{equational-logic}

\begin{hbexample}{def:p10-identities-valuations}{An identity and a local equation}
In the empty signature, the identity \(x\approx y\), for distinct
variables \(x,y\), requires an algebra to have exactly one element.
Consequently it implies every identity.
In contrast, the single valuation \(v(x)=v(y)=v(u)=0\), \(v(z)=1\)
in the set \(\{0,1\}\) satisfies \(x\approx y\) and fails
\(u\approx z\), where the displayed variables are distinct.
This local implication is a different consequence relation:
unrestricted substitution of a fixed local premise would not preserve it.
All uses of \(\models_{\mathrm{alg}}\) below have the universal
identity semantics just defined.
\end{hbexample}

\begin{hbdefinition}{def:p10-deductive-congruence}{The congruence of provable identities}
On \(T_\tau(V)\), put
\[
 s\mathrel{\theta_\Gamma}t
 \quad\Longleftrightarrow\quad
 \Gamma\vdash_{\mathrm{eq}}s\approx t.
\]
A congruence \(\theta\) on a term algebra is \emph{fully invariant}
when every endomorphism preserves it:
\[
 s\mathrel\theta t\ \Longrightarrow\
 \widehat\sigma(s)\mathrel\theta\widehat\sigma(t)
 \quad\text{for every }\sigma:V\to T_\tau(V).
\]
The deduction rules make \(\theta_\Gamma\) the least fully invariant
congruence containing the pairs of \(\Gamma\).
The exact leastness argument and the quotient-model construction are
given in the following proof
\cite[II\S14, Theorems~14.12 and~14.17]{BurrisSankappanavar2012}.
\index{congruence!fully invariant}
\end{hbdefinition}

\begin{hbtheorem}{thm:p10-equational-completeness}{proved=annex:p10.eq.2}{Completeness of equational logic}
For every set \(\Gamma\) of equations over \(V\) and every equation
\(s\approx t\) over \(V\),
\[
 \Gamma\models_{\mathrm{alg}}s\approx t
 \quad\Longleftrightarrow\quad
 \Gamma\vdash_{\mathrm{eq}}s\approx t.
\]
The nonempty quotient algebra \(T_\tau(V)/\theta_\Gamma\) satisfies
every identity in \(\Gamma\).
If the displayed consequence fails, the valuation \(x\mapsto[x]\)
in this quotient falsifies \(s\approx t\).
Moreover, consequence is finitary: for each consequence \(s\approx t\),
some finite \(\Gamma_0\subseteq\Gamma\) already implies \(s\approx t\).
\end{hbtheorem}

This is Birkhoff's completeness theorem, with a direct quotient proof
below; the source's deductive-closure proof is
\cite[II\S14, Theorem~14.19, p.~95]{BurrisSankappanavar2012}.
The locators in this section refer to the 2012 update of the
Millennium Edition, whose main-text pagination was aligned with the
1981 edition.

\subsection{Free algebras in an equational class}
\label{sec:p10-free-algebras}

\begin{hbdefinition}{def:p10-free-algebra}{Free algebra on a set of generators}
For a variety \(\mathcal V\), a free algebra on \(X\) is a pair
\((F_{\mathcal V}(X),\eta_X)\), where
\(F_{\mathcal V}(X)\in\mathcal V\) and
\(\eta_X:X\to F_{\mathcal V}(X)\), such that every map \(v:X\to A\)
with \(A\in\mathcal V\) factors uniquely as
\[
 v=\bar v\circ\eta_X,\qquad
 \bar v:F_{\mathcal V}(X)\longrightarrow A
 \quad\text{an algebra homomorphism}.
\]
The universal property is stated for the map \(\eta_X\);
injectivity is not part of this definition.
For the variety consisting of one-element algebras, all generators
have the same image.
\index{free algebra}\index{universal property}
\end{hbdefinition}

\begin{hbtheorem}{thm:p10-free-algebra}{proved=annex:p10.eq.3}{Construction and uniqueness of free algebras}
Let \(\mathcal V\) be the class of all nonempty models of a set
\(\Gamma\) of identities. If \(T_\tau(X)\ne\varnothing\), define
\[
 s\mathrel{\theta_{\mathcal V,X}}t
 \quad\Longleftrightarrow\quad
 \forall A\in\mathcal V\ \forall v:X\to A\quad
                    s^A(v)=t^A(v).
\]
Then
\[
 F_{\mathcal V}(X)=T_\tau(X)/\theta_{\mathcal V,X},
 \qquad \eta_X(x)=[x],
\]
is free in \(\mathcal V\) on \(X\).
Any two free algebras on \(X\) are uniquely isomorphic by an
isomorphism commuting with the generator maps
\cite[II\S10, Theorems~10.7, 10.10, 10.12;
II\S11, Theorem~11.4]{BurrisSankappanavar2012}.
\end{hbtheorem}

The hypothesis on \(T_\tau(X)\) matters when \(X=\varnothing\).
With a nullary operation, ground terms provide the construction.
Without nullary operations, the empty term set does not provide a
nonempty algebra. No uniform existence theorem for an empty set of
generators is then available under the nonempty-domain convention.
For example, algebras of the empty signature are nonempty sets.
An initial nonempty set \(F\) would have a unique map to \(\{0,1\}\),
but its two constant maps are distinct.
This example does not exclude special varieties that have an initial
algebra for other reasons.

For the signature consisting of a binary multiplication and a unit
constant, the monoid identities give the free monoid \(X^*\).
The empty word is the unit, concatenation is multiplication, and
\(\eta_X(x)\) is the one-letter word.
For \(v:X\to M\), the required extension sends
\(x_1\cdots x_n\) to \(v(x_1)\cdots v(x_n)\) and sends the empty
word to \(1_M\).
Associativity and the unit laws make it a homomorphism; every
homomorphism extending \(v\) must take these values.
In particular, the free monoid on the empty set is the one-element
monoid.

%% Subdirect representation for nonempty finitary algebras.

\subsection{Subdirect products and irreducible quotients}
\label{sec:p10-subdirect-products}

Fix a finitary signature \(\tau\); its set of operation symbols need not
be finite.  All algebras have nonempty carriers, and all product indices
form sets.  We work in set theory with the axiom of choice; its use
through Zorn's lemma in the representation theorem will be explicit.
The source for this toolkit is Chapter~II, \S8 of the 2012 Millennium
Edition of \textcite{BurrisSankappanavar2012}, with the quotient
correspondence from Chapter~II, Theorem~6.20.

Write \(\operatorname{Con}A\) for the congruences on \(A\), ordered by
inclusion, and put
\[
 \Delta_A=\{(a,a):a\in A\},\qquad \nabla_A=A\times A.
\]
Intersections of congruences are congruences; the empty intersection is
\(\nabla_A\).  Thus a family has meet given by intersection and join
given by the intersection of all its common congruence upper bounds.
In particular, \(\theta\vee\phi\) is the least congruence containing
both \(\theta\) and \(\phi\).  The \emph{principal congruence}
\(\operatorname{Cg}_A(a,b)\) is the intersection of all congruences
containing \((a,b)\).  The quotient map is
\(q_\theta:A\to A/\theta\), \(a\mapsto[a]_\theta\), and
\(\ker h=\{(a,b):h(a)=h(b)\}\) for a homomorphism \(h\).

\begin{hbdefinition}{def:p10-subdirect-product}{Subdirect products and embeddings}
For an indexed family \((A_i)_{i\in I}\) of \(\tau\)-algebras, its
direct product has the Cartesian product as carrier and coordinatewise
operations:
\[
 f^{\prod_{i\in I}A_i}(x_1,\ldots,x_n)(i)
       =f^{A_i}(x_1(i),\ldots,x_n(i)).
\]
This includes nullary operations.  For \(I=\varnothing\), the product
is the one-element \(\tau\)-algebra whose element is the empty tuple.
For \(i\in I\), let \(\pi_i\) be the coordinate projection.

A subalgebra \(B\leq\prod_{i\in I}A_i\) is a nonempty subset closed
under the basic operations, with those operations restricted to \(B\).
It is a \emph{subdirect product} when every restricted projection
\(\pi_i|_B:B\to A_i\) is surjective.
A \emph{subdirect embedding} is an injective homomorphism
\(e:A\to\prod_{i\in I}A_i\) whose image is a subdirect product.
Equivalently, \(e\) is injective and every \(\pi_i\circ e\) is onto.
The image may be a proper subalgebra of the full product.
For an empty index set a subdirect product is exactly the one-element
product, and an algebra embeds subdirectly into it exactly when that
algebra is trivial, meaning that its carrier has one element.
\index{subdirect product}\index{subdirect embedding}
\end{hbdefinition}

\begin{hbtheorem}{thm:p10-subdirect-kernels}{proved=annex:p10.subdirect.1}{The congruence-intersection criterion}
Let \(h_i:A\to A_i\) be homomorphisms, for \(i\in I\), and define
\[
 e:A\longrightarrow\prod_{i\in I}A_i,\qquad
 e(a)=(h_i(a))_{i\in I}.
\]
Then \(e\) is a homomorphism and
\[
 \ker e=\bigcap_{i\in I}\ker h_i.
\]
It is an embedding exactly when that intersection is \(\Delta_A\);
if each \(h_i\) is onto, this is exactly the condition for \(e\) to
be subdirect.

In particular, for congruences \(\theta_i\in\operatorname{Con}A\),
put \(\theta=\bigcap_{i\in I}\theta_i\).  The map
\[
 A/\theta\longrightarrow\prod_{i\in I}A/\theta_i,\qquad
 [a]_\theta\longmapsto([a]_{\theta_i})_{i\in I}
\]
is a subdirect embedding.  The diagonal quotient map with domain
\(A\) is a subdirect embedding exactly when \(\theta=\Delta_A\).
These statements include \(I=\varnothing\), with
\(\theta=\nabla_A\) and \(A/\theta\) trivial.
\end{hbtheorem}

\begin{hbtheorem}{thm:p10-congruence-correspondence}{proved=annex:p10.subdirect.2}{Congruences above a quotient kernel}
For \(\theta\in\operatorname{Con}A\), the assignments
\[
\begin{aligned}
 \phi&\longmapsto
 \phi/\theta
   =\{([a]_\theta,[b]_\theta):(a,b)\in\phi\},
       &&\theta\subseteq\phi\in\operatorname{Con}A,\\
 \psi&\longmapsto
 \{(a,b):([a]_\theta,[b]_\theta)\in\psi\},
       &&\psi\in\operatorname{Con}(A/\theta)
\end{aligned}
\]
are inverse order isomorphisms between
\([\theta,\nabla_A]\) and \(\operatorname{Con}(A/\theta)\).
They preserve all meets and joins in these lattices and send
\(\theta\) to \(\Delta_{A/\theta}\).
\end{hbtheorem}

\begin{hbdefinition}{def:p10-subdirect-irreducible}{Subdirect irreducibility and the monolith}
An algebra \(A\) is \emph{subdirectly irreducible} when, for every
subdirect embedding
\[
 e:A\longrightarrow\prod_{i\in I}A_i
 \quad\text{with }I\ne\varnothing,
\]
at least one coordinate map \(\pi_i\circ e:A\to A_i\) is an
isomorphism.  The nonempty-index condition includes trivial algebras
in this convention; empty products remain allowed in
Definition~\ref{def:p10-subdirect-product}.

For a nontrivial algebra, a least congruence strictly above
\(\Delta_A\), when it exists, is called its \emph{monolith}.
\index{subdirectly irreducible algebra}\index{monolith}
\end{hbdefinition}

\begin{hbtheorem}{thm:p10-subdirect-monolith}{proved=annex:p10.subdirect.3}{Characterizing subdirect irreducibility}
An algebra \(A\) is subdirectly irreducible exactly when every
nonempty family \((\theta_i)_{i\in I}\) of congruences satisfying
\(\bigcap_i\theta_i=\Delta_A\) contains a member equal to
\(\Delta_A\).

Every trivial algebra has this property.  For nontrivial \(A\),
the following conditions are equivalent:
\[
\begin{gathered}
 A\text{ is subdirectly irreducible};\\
 \operatorname{Con}A\setminus\{\Delta_A\}
                 \text{ has a least member }\mu;\\
 \bigcap\bigl(\operatorname{Con}A\setminus\{\Delta_A\}\bigr)
                       \ne\Delta_A.
\end{gathered}
\]
In this case the intersection is the monolith \(\mu\).
For every \((a,b)\in\mu\) with \(a\ne b\),
\(\mu=\operatorname{Cg}_A(a,b)\).
\end{hbtheorem}

\begin{hbtheorem}{thm:p10-subdirect-representation}{proved=annex:p10.subdirect.4}{Birkhoff's subdirect representation theorem}
Assuming the axiom of choice, every nonempty finitary algebra \(A\)
is isomorphic to a subdirect product of subdirectly irreducible
homomorphic images of \(A\).

For nontrivial \(A\), one can take one quotient \(A/\theta_{a,b}\)
for every ordered pair \(a\ne b\), where \(\theta_{a,b}\) is maximal
among the congruences omitting \((a,b)\).  Then
\[
 a\longmapsto([a]_{\theta_{u,v}})_{u\ne v}
\]
is the required embedding.  A trivial algebra has a one-factor
representation by itself; alternatively it is isomorphic to the
empty product.

If \(A\) is finite and nontrivial, at most
\(|A|(|A|-1)\) finite quotient factors suffice.  This finite case
does not require the axiom of choice.
\end{hbtheorem}

The congruence-intersection criterion and the monolith characterization
correspond to Lemma~8.2 and Theorem~8.4 of the cited edition;
the representation theorem and its finite case are Theorem~8.6 and
Corollary~8.7 \cite[Chapter~II, \S8, pp.~57--59]{BurrisSankappanavar2012}.
The annex proves the chain-union and quotient steps explicitly.
Maximality here means maximal while separating a specified pair:
it need not mean maximal among all proper congruences.

\paragraph{A proper subdirect product.}
Let \(C_3=\{0<c<1\}\) be the three-element bounded lattice, and
\(C_2=\{0<1\}\) the two-element bounded lattice.
The maps with value lists
\[
 h_1=(0,0,1),\qquad h_2=(0,1,1)
\]
on \(0,c,1\) are monotone and preserve the endpoints.
Monotone maps between chains preserve binary minima and maxima,
so these maps preserve meets, joins, and the two constants.
Both are onto, and their kernels intersect in \(\Delta_{C_3}\).
They embed \(C_3\) as
\(\{(0,0),(0,1),(1,1)\}\leq C_2\times C_2\).
This is a three-element subdirect product inside a four-element
direct product.  Each factor is subdirectly irreducible because
its only congruences are equality and the universal relation.

%% toc: P10 139 139.5
\section{Varieties and Equational Representation by HSP}
\label{sec:p10-hsp-equational-representation}

The term and quotient constructions in
\cref{sec:p10-equational-logic,sec:p10-free-algebras} establish
equational completeness and free algebras for an equational class.
The converse representation required by
\cref{thm:p10-hsp} has a further task: starting with a class \(K\),
construct a product from members of \(K\), a subalgebra of that
product and an onto map to any algebra satisfying all identities
of \(K\). The following constructions supply those objects and maps.
The kernel route follows the corrected 2012 author edition of
Burris--Sankappanavar \cite[II\S\S9--11]{BurrisSankappanavar2012}.
Birkhoff's original 1935 all-valuations route is recorded separately
\cite[\S\S7--10, pp.\ 437--441]{Birkhoff1935}.

\subsection{Source conventions and identity preservation}
\label{sec:p10-hsp-conventions}

\begin{hbconvention}{def:p10-hsp-conventions}{Fixed type, nonempty carriers and two product conventions}
Work in ZFC with a set-sized finitary signature.
The book calls the signature \(\mathcal F\), with \(\mathcal F_n\) its \(n\)-ary
symbols; it is the current reader's \(\tau\).
All algebras have nonempty carriers and total basic operations.
Nullary operations are included. Write \(\mathbf A\) for
an algebra with carrier \(A\), and preserve the book's
\(T(X)\) for terms, \(\mathbf T(X)\) for their term algebra.
It exists exactly when \(X\cup \mathcal F_0\ne\varnothing\).
Classes considered below are sets or definable classes of similar
algebras, so separation and collection apply to their set-indexed
witness conditions.

The book's operators \(I,S,H,P,V\) have their native roles:
isomorphic copies, subalgebras, onto homomorphic images,
products of \emph{nonempty} set-indexed families, and least
nonempty H/S/P-closed class.
In particular its Definition~9.3 requires a variety to be nonempty.
The host's \(\mathsf H,\mathsf S,\mathsf P\) have the first
two corresponding roles, but \(\mathsf P\) also allows an
empty family, whose product is the one-element algebra.
These conventions are not silently identified at an empty class.

Fix a countably infinite variable set \(W\).
Preserve the book's identity glyph \(\approx\), which is the
same equality-of-all-valuations judgment as the current reader.
For any variable set \(X\), write
\[
 \operatorname{Id}_K(X)=
 \{p\approx q:p,q\in T(X),\
       \mathbf A\models p\approx q\text{ for every }\mathbf A\in K\}.
 \tag{HSP.1}
\]
The source writes \(\operatorname{Id}_K(X)\) and
\(M(\Sigma)\) for the class satisfying an identity set
\(\Sigma\), in II Definitions 11.1 and 11.7.
For a definable proper class \(K\), (HSP.1) remains a subset of
the set \(T(X)^2\), not a proper-class set of equations.
\end{hbconvention}

\begin{hblemma}{thm:p10-hsp-preservation}{proved=annex:p10-hsp-preservation}{Identities survive H, S and P}
If all members of \(K\) satisfy an identity, then so do their
nonempty subalgebras, homomorphic images and set-indexed
products, including the host's empty product.
Consequently
\[
 \operatorname{Id}_{\mathsf H\mathsf S\mathsf P(K)}(W)
       =\operatorname{Id}_K(W).
 \tag{HSP.2}
\]
For \(K\ne\varnothing\), the same equality holds with the book's
native \(HSP\). Each individual identity uses only finitely
many variables. Preservation under an onto homomorphism requires
only their finite family of lifts, not a section of the entire
quotient map.
\end{hblemma}

\subsection{The relative-free algebra and its separating product}
\label{sec:p10-hsp-kernel-construction}

\begin{hbdefinition}{def:p10-hsp-relative-free}{The source's set of quotient kernels}
Suppose \(\mathbf T(X)\) exists. Keep the native definitions
of II Definition~10.9:
\[
 \begin{split}
 \Phi_K(X)&=\{\phi\in\operatorname{Con}\mathbf T(X):
               \mathbf T(X)/\phi\in IS(K)\},\\
 \theta_K(X)&=\bigcap\Phi_K(X),\\
 \bar X&=X/\theta_K(X)=\{[x]_{\theta_K(X)}:x\in X\},\\
 \mathbf F_K(\bar X)&=\mathbf T(X)/\theta_K(X),\qquad
 \eta_X(x)=[x]_{\theta_K(X)}.
 \end{split}
 \tag{HSP.3}
\]
An empty congruence intersection means the universal congruence.
The book also uses \(\bar X\), the set of generator classes,
and \(\bar p=[p]_{\theta_K(X)}\) for a term's class.
The generator map \(X\to\bar X\) need not be injective.
In a trivial class all its values coincide.
\end{hbdefinition}

\begin{hbtheorem}{thm:p10-hsp-witness-product}{proved=annex:p10-hsp-witness-product}{Explicit relative-free embedding and extensions}
Let \(K\ne\varnothing\) and \(T(X)\ne\varnothing\).
Then \(\Phi_K(X)\) is a nonempty set, and
\[
 (p,q)\in\theta_K(X)
 \quad\Longleftrightarrow\quad
 h(p)=h(q)\text{ for every }h:\mathbf T(X)\to\mathbf A,\
                                             \mathbf A\in K.
 \tag{HSP.4}
\]
For each \(\phi\in\Phi_K(X)\), choose a member
\(\mathbf A_\phi\in K\) and an embedding
\[
 j_\phi:\mathbf T(X)/\phi\longrightarrow\mathbf A_\phi.
 \]
These choices are over a set of congruences; for a definable
proper class, collection first bounds a set of possible witnesses.
Put
\[
 \begin{split}
 h_\phi&=j_\phi\circ q_\phi,\\
 d([p]_{\theta_K(X)})&=(h_\phi(p))_{\phi\in\Phi_K(X)},\\
 \mathbf B&=
   d\bigl(\mathbf F_K(\bar X)\bigr)
      \leq\prod_{\phi\in\Phi_K(X)}\mathbf A_\phi.
 \end{split}
 \tag{HSP.5}
\]
Then \(d\) is a well-defined embedding, and \(\mathbf B\)
is generated by the tuples \(d(\eta_X(x))\), with constants
interpreted coordinatewise. Thus \(\mathbf F_K(\bar X)\in ISP(K)\).

Every \(v:X\to A\), for \(\mathbf A\in K\), extends
uniquely to a homomorphism
\(\bar v:\mathbf F_K(\bar X)\to\mathbf A\), by
\(\bar v([p])=p^{\mathbf A}(v)\).
Moreover \(\mathbf F_K(\bar X)\) satisfies all identities of \(K\)
over every variable set, not just equalities at the named generators.
This is the construction behind II Theorems 10.10,10.12 and 11.4;
no appeal to the HSP theorem is used to obtain its embedding.
\end{hbtheorem}

\subsection{From all identities to an actual HSP representation}
\label{sec:p10-hsp-representation}

\begin{hbmethod}{def:p10-hsp-target-map}{Generators indexed by the target's elements}
Let \(K\ne\varnothing\) and
\(\mathbf A\models\operatorname{Id}_K(W)\).
Take \(X=\{x_a:a\in A\}\), a disjointly tagged copy of the
nonempty target carrier. The evaluation
\[
 e_A:\mathbf T(X)\longrightarrow\mathbf A,\qquad
 e_A(x_a)=a
 \tag{HSP.6}
\]
is onto. The full identity hypothesis, using a countably infinite
alphabet, will imply \(\theta_K(X)\subseteq\ker e_A\).
Only finitely many variables occur in any particular pair of
terms, so their identity can be renamed into \(W\) and back.
Define
\[
 \begin{aligned}
 r_A:\mathbf F_K(\bar X)&\longrightarrow\mathbf A,
          &r_A([p])&=e_A(p),\\
 b_A:\mathbf B&\longrightarrow\mathbf A,
          &b_A(d([p]))&=e_A(p).
 \end{aligned}
 \tag{HSP.7}
\]
The proof verifies independence of representatives and surjectivity,
so the required object is a subalgebra of the actual product
in (HSP.5), followed by an actual onto map, not an unspecified
application of a representation theorem.
\end{hbmethod}

\begin{hbtheorem}{thm:p10-hsp-representation}{proved=annex:p10-hsp-representation}{HSP generation and the variety characterization}
For every class \(K\) under \cref{def:p10-hsp-conventions},
\[
 M(\operatorname{Id}_K(W))=
             \mathsf H\mathsf S\mathsf P(K).
 \tag{HSP.8}
\]
For \(K\ne\varnothing\), the right side also equals the book's
\(HSP(K)\). If a nonempty algebra lies on the left, the maps
(HSP.5)--(HSP.7) give its representation when \(K\ne\varnothing\).
For \(K=\varnothing\), both sides under the host convention are
exactly all one-element algebras.

A class of nonempty finitary algebras is equational exactly when
it is closed under homomorphic images, nonempty subalgebras and
all set-indexed products including the empty product.
Equivalently, with the book's product convention it is equational
exactly when it is nonempty and H/S/P-closed.
Isomorphic copies are automatically included by closure under
homomorphic images.
\end{hbtheorem}

\begin{hbwarning}{def:p10-hsp-empty-boundary}{Empty generators and empty generating classes are different}
When \(X=\varnothing\), the term algebra is available if \(\mathcal F_0\)
has a constant symbol; otherwise this construction has no terms
and no nonempty free object is asserted. The representation
(HSP.6) always uses a nonempty \(X\), regardless of constants.

The book's \(HSP(\varnothing)=\varnothing\), because its product
families are nonempty. Its least \emph{nonempty} variety containing
\(\varnothing\) is instead the class of one-element algebras.
Thus the formula \(V(K)=HSP(K)\), with these native definitions,
has the necessary nonempty-\(K\) scope.
The host permits an empty product and obtains
\(\mathsf H\mathsf S\mathsf P(\varnothing)\) equal to that
trivial variety. Every identity is true in a one-element algebra,
so no equational class of nonempty algebras is empty.
\end{hbwarning}

\subsection{Birkhoff's original all-valuations product}
\label{sec:p10-hsp-original-product}

\begin{hbdefinition}{def:p10-hsp-original-route}{The 1935 finite-species construction}
In the original, a \emph{uniform} operator has the entire
power of its carrier as its proper domain, and a \emph{species}
fixes the indices of the operators. A symbolic \emph{function}
is an expression formed from primitive symbols by operators;
a \emph{law} is an equation true under every element substitution.
The original permits ordinal indices and transfinite ranks.
The construction here is explicitly its finitary specialization:
finite-arity species, finite terms and arbitrary set-sized
sets of primitive symbols. Partial-domain algebras from its
opening definition are not included.
The current nonempty-carrier convention is also a hypothesis
of this specialization, not inferred from the original's
opening definition.
Keep its \(\Sigma\) for a species, distinct in type from the
book's identity set \(\Sigma\).
The original writes \(\mathfrak F(\mathfrak A)\) for the
algebra family generated from an algebra basis,
\(\Phi(\mathfrak A)\) for all its common laws, and
\(\Phi(B)\) for the deductive family generated from a law basis.
It also writes \(\mathfrak F(B)\) for the algebra family obtained
from the corresponding free quotients.
The algebra-basis and law-basis operands distinguish these
genuine source overloads; \(\Phi(B)\) is not the book's
congruence-kernel set \(\Phi_K(X)\).

For a nonempty set \(\mathfrak A\) of such algebras, let \(X\)
be a nonempty set of primitive symbols. Index every map
\(\xi:X\to A_i\), for \(\mathbf A_i\in\mathfrak A\).
Let \(\mathbf S_{i,\xi}\leq\mathbf A_i\) be generated
by its values, and form
\[
 \mathbf S^*=\prod_{(i,\xi)}\mathbf S_{i,\xi},\qquad
 z_x=(\xi(x))_{(i,\xi)}\quad(x\in X).
 \tag{HSP.9}
\]
The original writes \(F(\mathfrak A,n)\) for the subalgebra
generated by these diagonally assigned primitive symbols,
with \(n\) their number. Here put
\(\mathbf F(\mathfrak A,X)=
       \operatorname{Sg}_{\mathbf S^*}\{z_x:x\in X\}\)
to specify a possibly infinite set without treating cardinalities
as literal term alphabets.
For finite \(X=\{x_1,\ldots,x_n\}\), this is the original
\(F(\mathfrak A,n)\).
\end{hbdefinition}

\begin{hbtheorem}{thm:p10-hsp-original-product}{proved=annex:p10-hsp-original-product}{The original product realizes exactly the common laws}
Term evaluation into (HSP.9) gives an onto homomorphism to
\(\mathbf F(\mathfrak A,X)\) whose kernel is precisely
equality under every valuation into every basis algebra.
Each coordinate is onto its generated \(\mathbf S_{i,\xi}\),
so the image is a subdirect product of those generated
subalgebras in the current terminology. It need not be a
subdirect product of the full basis algebras, and no
subdirect-irreducibility assertion follows.
Consequently it induces an isomorphism
\[
 \mathbf T(X)/\theta_{\mathfrak A}(X)
             \ \cong\ \mathbf F(\mathfrak A,X).
 \tag{HSP.10}
\]
For a law basis \(B\), the original's \(F(B,n)\) is the quotient
by its equated functions. For the modern equality/deduction bridge,
choose a variable alphabet \(U\) containing \(X\) and all variables
of \(B\), and let \(\theta_{B,U}\) be the deductive congruence
on \(\mathbf T(U)\).
Put \(R_{B,X}=\theta_{B,U}\cap T(X)^2\).
Then the corresponding quotient is
\(\mathbf T(X)/R_{B,X}\), using
\cref{def:p10-deductive-congruence,thm:p10-equational-completeness}.
It is the relative-free algebra for \(M(B)\); the number of free
generators does not limit the number of variables permitted in \(B\).
Thus (HSP.10) gives the finite-term content of the original
equation 10-1: the all-valuations construction and the quotient
by all common laws agree. Every algebra obeying the common laws
is an onto image of one of these relative-free algebras on a
sufficiently large generator set.
In the original notation the isomorphism reads
\[
 F(\mathfrak A,n)\sim F(\Phi(\mathfrak A),n),
 \tag{1935.10-1}
\]
where its \(\sim\) means algebra isomorphism.
Under the explicitly modern equality/deduction bridge,
algebra families and deductively closed law families are
mutually determined and reverse inclusion:
\[
 \mathfrak F(\Phi(\mathfrak F))=\mathfrak F,\qquad
 \Phi(\mathfrak F(\Phi))=\Phi.
 \tag{HSP.11}
\]
Consequently their lattices are dually isomorphic, the finitary
family/law form of the original Theorem~10 and Corollary~1.
For a finite nonempty basis
\(\mathfrak A=\{\mathbf A_1,\ldots,\mathbf A_s\}\)
with finite carrier sizes \(a_i\), every algebra generated
by \(m\geq1\) elements and obeying its laws has at most
\[
 \prod_{i=1}^{s} a_i^{\,a_i^{\,m}}
 \tag{HSP.12}
\]
elements. This is the finite-basis case of the original
Corollary~2, obtained from the actual valuation-indexed product.
\end{hbtheorem}

The original's Theorem~7 displays operator compatibility and
uniform function substitution; it does not separately print the
three elementary equality rules as a modern calculus sheet.
The comparison with \(\theta_B\) explicitly uses the reader's
reflexive, symmetric and transitive equality deductions in addition
to those two operations. It is not a claim that this modern rule
listing occurs verbatim in the 1935 article.
The original's family/law correspondence in Theorems 9--10 is the
source of the two directions; the 1944 subdirectly-irreducible
representation theorem is a different result.

\begin{hbexample}{def:p10-hsp-two-element-set}{Two generators in the empty signature}
For a basis containing the two-element set \(\{0,1\}\) and no
operation symbols, there are four assignments to \(x,y\).
In their order \(00,01,10,11\), the diagonal generator tuples are
\[
 z_x=(0,0,1,1),\qquad z_y=(0,1,0,1).
\]
The generated subalgebra consists of these two tuples.
The coordinate-generated carriers have sizes \(1,2,2,1\),
so the original product has four elements, not sixteen;
the ambient product of four full two-element algebras has sixteen.
Neither product is confused with its two-element generated subalgebra.
More generally, the class of all finite nonempty sets is H/S-closed
and closed under finite products, but not under all products.
It has the same identities as all nonempty sets. Finite products
therefore do not suffice in the HSP characterization.
\end{hbexample}

\begin{proofboundary}{def:p10-hsp-choice-and-scope}{Choice and presentation limits}
The relative-free witness selection and arbitrary nonempty products
are taken in ZFC. Preservation of each finitary identity requires
only finite lifts, and the all-valuations diagonal itself supplies
a point of its product whenever a term is available.
No construction here is a decision procedure for class membership,
identity validity or choice of separating witnesses.
An infinite signature is allowed; the crucial restriction is
finite arity and hence finite variable support of each term.
The original ordinal-arity/rank framework requires its own
infinitary syntax and lifting conditions and is not silently
identified with this finitary theorem.
\end{proofboundary}


\section{Boolean algebras and Stone duality}
\label{sec:p10-stone-duality}

\begin{hbdefinition}{def:p10-boolean-algebra}{Boolean algebra, filter, ideal, and ultrafilter}
A Boolean algebra is a bounded distributive lattice
\((B,\wedge,\vee,\neg,0,1)\) satisfying
\(a\wedge\neg a=0\) and \(a\vee\neg a=1\).
The one-element algebra with \(0=1\) is allowed.
The order is \(a\leq b\) exactly when \(a\wedge b=a\).
A filter \(F\subseteq B\) contains \(1\), is upward closed, and is
closed under binary meets; it is proper when \(0\notin F\).
An ideal contains \(0\), is downward closed, and is closed under
binary joins; it is proper when it excludes \(1\).
An ultrafilter is a maximal proper filter. Equivalently, it is a
filter containing exactly one of \(a,\neg a\) for every \(a\), as
proved in Lemma~\ref{thm:p10-ultrafilter-separation}. These terms are
needed because algebraic truth assignments, deductively closed theories, and
prime points of a dual space occupy different levels. They stand as algebraic
structures and subsets and are used in Stone representation and the
Lindenbaum--Tarski construction.
\index{Boolean algebra}\index{filter!Boolean algebra}\index{ultrafilter}\index{ideal}
\end{hbdefinition}

An atom is a minimal nonzero element. Atomicity means every nonzero element
lies above an atom; it does not mean that the algebra is finite. The Boolean
prime ideal theorem asserts that every proper filter extends to an ultrafilter.
It is weaker than the axiom of choice over ZF but is not provable in ZF alone;
proofs using ultrafilter extension must state that choice principle
\cite{Halmos1963,Koppelberg1989}.

%% Stone representation, reconstruction, and functorial duality.

\subsection{Ultrafilters and Stone representation}
\label{sec:p10-stone-representation}

The proofs in this subsection work in \(\mathsf{ZF}\) together with
the Boolean Prime Ideal Theorem (BPIT), in its proper-filter-extension
form stated above. This assumption suffices for the arbitrary Boolean
algebras considered here. The complete source presentation is
\cite[IV\S\S3--4]{BurrisSankappanavar2012}.

\begin{hblemma}{thm:p10-ultrafilter-separation}{proved=annex:p10.stone.1}{Ultrafilter dichotomy and separation}
A proper filter \(U\) is an ultrafilter if and only if, for every
\(a\in B\), exactly one of \(a,\neg a\) belongs to \(U\).
Consequently,
\[
 a\vee b\in U\quad\Longleftrightarrow\quad a\in U\text{ or }b\in U.
\]
Assuming BPIT, if \(F\) is a filter and \(b\notin F\), there is an
ultrafilter \(U\supseteq F\) with \(b\notin U\).
In particular, each nonzero element belongs to an ultrafilter,
and distinct Boolean elements are separated by an ultrafilter.
The trivial Boolean algebra has no ultrafilters.
\end{hblemma}

\begin{hbdefinition}{def:p10-stone-space}{Stone space of a Boolean algebra}
A \emph{Stone space} is a compact Hausdorff topological space with
a basis of sets that are both open and closed, called \emph{clopen}
sets. The empty space is allowed.
Compact means that every open cover has a finite subcover.
For a Boolean algebra \(B\), set
\[
 X_B=\{U:U\text{ is an ultrafilter of }B\},\qquad
 \widehat a=\{U\in X_B:a\in U\}.
\]
Give \(X_B\) the topology whose opens are unions of the sets
\(\widehat a\).
The filter laws give
\(\widehat1=X_B\) and
\(\widehat a\cap\widehat b=\widehat{a\wedge b}\), so these sets
form a basis.
The theorem below proves that this is a Stone space.

For any topological space \(X\), the set
\(\operatorname{Clop}(X)\) is a Boolean algebra under intersection,
union, and complement relative to \(X\), with bounds
\(\varnothing,X\).
Thus \(\operatorname{Clop}(\varnothing)\) is the trivial Boolean
algebra. The topology on \(X_B\) is a structure on the set of
ultrafilters, distinct from the underlying set of \(B\).
\index{Stone space}\index{clopen set}\index{Stone duality}
\end{hbdefinition}

\begin{hbtheorem}{thm:p10-stone-representation}{proved=annex:p10.stone.2}{Stone representation theorem}
Assuming BPIT, \(X_B\) is a Stone space and
\[
 \beta_B:B\longrightarrow\operatorname{Clop}(X_B),
 \qquad \beta_B(a)=\widehat a,
\]
is a Boolean-algebra isomorphism.
In particular, it is an injective Boolean homomorphism onto all
clopen subsets, not merely onto a Boolean subalgebra of them.
\end{hbtheorem}

The proof establishes compactness before using it to show that a
clopen set is a finite union of basic clopens. It applies to arbitrary
open covers, without choosing a basic neighborhood independently at
every point.
The representation and the inverse construction below are
\cite[IV\S4, Theorem~4.6, pp.~136--138]{BurrisSankappanavar2012};
see also \textcite{Stone1936,Johnstone1982,GivantHalmos2009}.

\subsection{Reconstruction of spaces and action on maps}
\label{sec:p10-stone-reconstruction}

\begin{hbtheorem}{thm:p10-stone-space-reconstruction}{proved=annex:p10.stone.3}{A Stone space is recovered from its clopen algebra}
For every Stone space \(X\), the map
\[
 \varepsilon_X:X\longrightarrow X_{\operatorname{Clop}(X)},
 \qquad
 \varepsilon_X(x)=\{C\in\operatorname{Clop}(X):x\in C\},
\]
is a homeomorphism.
Every ultrafilter of \(\operatorname{Clop}(X)\) consists of exactly
the clopen neighborhoods of one point of \(X\).
For \(X=\varnothing\), both sides of the displayed map are empty.
\end{hbtheorem}

\begin{hbdefinition}{def:p10-stone-dual-maps}{The inverse-image maps of Stone duality}
Let \(\mathsf{BA}\) have Boolean algebras as objects and maps
preserving \(\wedge,\vee,\neg,0,1\) as arrows.
Let \(\mathsf{Stone}\) have Stone spaces as objects and continuous
maps as arrows. In both cases identities and composition are those
of functions.

For a Boolean homomorphism \(h:B\to C\), define
\[
 X(h):X_C\longrightarrow X_B,\qquad X(h)(U)=h^{-1}[U].
\]
For a continuous map \(g:X\to Y\), define
\[
 \operatorname{Clop}(g):
 \operatorname{Clop}(Y)\longrightarrow\operatorname{Clop}(X),
 \qquad \operatorname{Clop}(g)(D)=g^{-1}[D].
\]
Both assignments reverse the direction of arrows.
A contravariant functor preserves identities and reverses the order
of composition; the theorem verifies these laws for the displayed
assignments. Naturality of the reconstruction maps means the two
displayed commuting equations in that theorem.
These are the concrete categorical data behind the contravariant
equivalence, with the general category terminology formalized in
Definition~\ref{def:p10-category}.
\end{hbdefinition}

\begin{hbtheorem}{thm:p10-stone-duality}{proved=annex:p10.stone.4}{Functorial Stone duality}
Assuming BPIT, the inverse-image assignments are well-defined
contravariant functors
\[
 \mathsf{BA}\ \xrightarrow{\ X\ }\ \mathsf{Stone},
 \qquad
 \mathsf{Stone}\ \xrightarrow{\ \operatorname{Clop}\ }\ \mathsf{BA}.
\]
The Boolean isomorphisms \(\beta_B\) and homeomorphisms
\(\varepsilon_X\) are natural: for \(h:B\to C\) and \(g:X\to Y\),
\[
\begin{aligned}
 \operatorname{Clop}(X(h))\circ\beta_B&=\beta_C\circ h,\\
 X(\operatorname{Clop}(g))\circ\varepsilon_X&=\varepsilon_Y\circ g.
\end{aligned}
\]
They satisfy the identities
\[
\begin{aligned}
 X(\beta_B)\circ\varepsilon_{X_B}&=1_{X_B},\\
 \operatorname{Clop}(\varepsilon_X)\circ
                 \beta_{\operatorname{Clop}(X)}
      &=1_{\operatorname{Clop}(X)}.
\end{aligned}
\]
Thus the categories are contravariantly equivalent.
Every continuous map \(f:X_C\to X_B\) arises from a unique Boolean
homomorphism \(h:B\to C\) by \(f=X(h)\).
\end{hbtheorem}

The source places the action on homomorphisms and continuous maps
in exercises following the object-level representation theorem
\cite[IV\S4, Exercises~9--10, p.~140]{BurrisSankappanavar2012}.
The annex supplies the well-definedness, composition laws,
naturality calculations, and inverse map on arrows in full.

\paragraph{Finite algebras and finite spaces.}
For \(B=\mathcal P(S)\) with \(S\) finite, every ultrafilter is
\(U_s=\{A\subseteq S:s\in A\}\) for a unique \(s\in S\).
Indeed, the finite join of the singleton atoms is \(1\), so finite-join
primeness forces an ultrafilter to contain a singleton; properness
prevents it from containing two different ones.
Then upward closure and the complement criterion determine the entire
ultrafilter. Each \(\widehat{\{s\}}\) is the singleton \(\{U_s\}\),
so the dual topology is discrete.
For \(S=\varnothing\), the algebra is trivial and its dual is empty.

%% The countable atomless Boolean algebra and its Stone space.

\subsection{The countable atomless Boolean algebra and Cantor space}
\label{sec:p10-countable-atomless}

The finite-splitting and countable constructions below are proved in
\(\mathsf{ZF}\). The final correspondence for arbitrary Boolean
algebras states its additional BPIT assumption explicitly.

\begin{hbdefinition}{def:p10-atomless-countability}{Atomlessness and countability}
In a Boolean algebra \(B\), an \emph{atom} is an element \(p\ne0\)
such that \(0<b\leq p\) implies \(b=p\).
The algebra is \emph{atomless} if
\[
 \forall b\ne0\ \exists c\quad 0<c<b.
\]
It is \emph{nontrivial} when \(0\ne1\).
Here \emph{countable} means finite or countably infinite.
Every finite nontrivial Boolean algebra has an atom, so a nontrivial
countable atomless algebra is necessarily countably infinite.
The nontriviality hypothesis excludes the one-element algebra, which
is atomless vacuously and has empty Stone space.

All Boolean subalgebras contain the same \(0,1\) as their ambient
algebra.  Atoms of a finite subalgebra are minimal nonzero elements
of that subalgebra; they can split further in an atomless ambient
algebra.
\index{atomless Boolean algebra}\index{countable atomless Boolean algebra}
\end{hbdefinition}

\begin{hbtheorem}{thm:p10-atomless-extension}{proved=annex:p10.atomless.1}{Extending an isomorphism of finite Boolean subalgebras}
Let \(B,C\) be nontrivial Boolean algebras, with \(C\) atomless.
Suppose \(F\leq F'\leq B\) and \(G\leq C\) are finite subalgebras,
and \(h:F\to G\) is an isomorphism.
Then there are a finite subalgebra \(G'\) of \(C\) containing \(G\)
and an isomorphism \(h':F'\to G'\) extending \(h\).

More explicitly, if an atom \(p\) of \(F\) is the join of \(k_p\)
atoms of \(F'\), split \(h(p)\) into \(k_p\) pairwise disjoint
nonzero elements of \(C\), and map those atoms bijectively to the
chosen pieces.
\end{hbtheorem}

\begin{hbtheorem}{thm:p10-countable-atomless-uniqueness}{proved=annex:p10.atomless.2}{Uniqueness of the countable atomless Boolean algebra}
Any two nontrivial countable atomless Boolean algebras are isomorphic.
More generally, every isomorphism between finite subalgebras of two
such algebras extends to an isomorphism of the whole algebras.
\end{hbtheorem}

The extension theorem is used alternately in the two algebras:
one step includes the next element in the domain, and the next
includes the next element in the range.
This is the back-and-forth argument suggested in
\cite[Chapter~IV, \S1, Exercise~5, p.~121]{BurrisSankappanavar2012};
the annex supplies the splitting and union proofs.

\begin{hbdefinition}{def:p10-cantor-cylinder-algebra}{Cantor space and its cylinder algebra}
Put \(K=2^{\mathbb N}=\{0,1\}^{\mathbb N}\), with the product
topology of the discrete two-element space.
Coordinates are numbered \(0,1,2,\ldots\).
For \(s\in\{0,1\}^n\), its \emph{cylinder} is
\[
 [s]=\{x\in K:x|_{\{0,\ldots,n-1\}}=s\}.
\]
The cylinder of the empty word is \(K\).
Write \(x|_n\) for the initial segment of \(x\) of length \(n\).
These cylinders form a clopen basis; this product space is called
\emph{Cantor space}.
Define
\[
 \mathcal B_n
   =\left\{\ \bigcup_{s\in E}[s]:E\subseteq\{0,1\}^n\right\},
 \qquad
 \mathcal B_{\rm cyl}=\bigcup_{n\in\mathbb N}\mathcal B_n.
\]
The Boolean operations are union, intersection, and complement
relative to \(K\), with constants \(\varnothing,K\).
\index{Cantor space}\index{cylinder set}\index{Boolean algebra!cylinder algebra}
\end{hbdefinition}

\begin{hbtheorem}{thm:p10-cantor-cylinder-algebra}{proved=annex:p10.atomless.3}{Existence through clopen cylinders}
Cantor space is nonempty, compact, and Hausdorff.
Every clopen subset of \(K\) is a finite union of cylinders of
one common length, so
\[
 \operatorname{Clop}(K)=\mathcal B_{\rm cyl}.
\]
This is a nontrivial countably infinite atomless Boolean algebra.
The finite subalgebras \(\mathcal B_n\) have
\(2^{2^n}\) elements; their inclusions split every cylinder atom
\([s]\) into \([s0]\) and \([s1]\).
Their union is also the direct limit of this inclusion sequence.
Consequently there exists exactly one nontrivial countable atomless
Boolean algebra up to isomorphism.
\end{hbtheorem}

\begin{hbtheorem}{thm:p10-countable-atomless-stone}{proved=annex:p10.atomless.4}{The Stone space is Cantor space}
Let \(B\) be a nontrivial countable atomless Boolean algebra.
Choose an isomorphism
\(\alpha:B\to\mathcal B_{\rm cyl}=\operatorname{Clop}(K)\).
Then
\[
 H:K\longrightarrow X_B,\qquad
 H(x)=\{b\in B:x\in\alpha(b)\}
\]
is a homeomorphism, with
\[
 H^{-1}(\widehat b)=\alpha(b).
\]
Thus every ultrafilter of \(B\) corresponds to exactly one infinite
binary sequence.
\end{hbtheorem}

The Cantor-space connection is stated for a countably infinite free Boolean
algebra in \cite[Chapter~IV, \S4, Exercise~3, p.~140]{BurrisSankappanavar2012}.
Here existence, algebraic uniqueness, and the displayed point map
give the conclusion directly for every nontrivial countable
atomless Boolean algebra.

\begin{hbtheorem}{thm:p10-atoms-isolated}{proved=annex:p10.atomless.5}{Atoms and isolated Stone points}
Assume the Boolean prime ideal theorem, as in
Theorem~\ref{thm:p10-stone-representation}.
For every Boolean algebra \(B\), its atoms correspond bijectively
to the isolated points of \(X_B\).
The atom \(p\) corresponds to the principal ultrafilter
\[
 U_p=\{b\in B:p\leq b\},\qquad \widehat p=\{U_p\}.
\]
In particular, \(B\) is atomless exactly when \(X_B\) has no
isolated points.  If \(B\) is finite and nontrivial with \(k\)
atoms, then \(X_B\) is a discrete \(k\)-point space and
\(|B|=2^k\).
\end{hbtheorem}

This is the correspondence in
\cite[Chapter~IV, \S4, Exercise~11, p.~141]{BurrisSankappanavar2012}.
It also explains why the finite nontrivial cases cannot have
Cantor space as their Stone space: every cylinder in \(K\)
contains two smaller nonempty cylinders, so \(K\) has no isolated
points.


\begin{hbdefinition}{def:p10-lindenbaum-algebra}{Lindenbaum--Tarski algebra}
For a propositional calculus \(S\) whose provable equivalence is a
congruence for its connectives, define
\(\varphi\equiv_S\psi\) when
\(S\vdash\varphi\leftrightarrow\psi\). The quotient of formulas by this
congruence, with connectives inducing operations, is the
\emph{Lindenbaum--Tarski algebra} of \(S\). The name is needed to distinguish a
syntactic quotient from an arbitrary Boolean algebra. It stands as an algebra
presented by formulas, differs because equality is provable equivalence, and
is used to turn maximal consistent theories into ultrafilters.
\index{Lindenbaum--Tarski algebra}
\end{hbdefinition}

For classical propositional logic this quotient is the free Boolean algebra on
the propositional variables. If \(\varphi\) is unprovable, its class
\(a=[\varphi]\) is different from \(1\), so \(\neg a\ne0\).
The principal filter
\[
 F=\{b:\neg a\leq b\}
\]
is proper: it contains \(1\), is upward closed and closed under meets,
and cannot contain \(0\).
By the Boolean Prime Ideal Theorem, extend \(F\) to an ultrafilter \(U\).
Then \(\neg a\in U\) and \(a\notin U\).
Its characteristic map \(\chi_U:B\to\{0,1\}\) is a Boolean
homomorphism, since ultrafilters preserve finite meets and select
exactly one of each element and its complement.
Assigning a propositional variable \(p\) the value \(\chi_U([p])\)
therefore evaluates every formula \(\psi\) to \(\chi_U([\psi])\),
by induction on formulas. In particular it falsifies \(\varphi\).
This proves the required countermodel direction of CPL completeness;
the converse follows by induction on a derivation.
For arbitrary sets of propositional variables, the ultrafilter step
uses the stated Boolean Prime Ideal Theorem. The countable-language
construction is given separately in Part~II.
Compactness becomes topological compactness of the Stone space of complete
theories, and model-theoretic type spaces are Stone spaces whose basic clopens
are sets of types containing a formula.

\section{Distributive, Heyting, and modal algebras}
\label{sec:p10-nonclassical-algebras}

%% Prime-filter separation and bounded distributive-lattice duality.

\subsection{Prime filters of bounded distributive lattices}
\label{sec:p10-prime-filters}

\begin{hbdefinition}{def:p10-distributive-prime-filter}{Distributive-lattice filters and prime filters}
A bounded distributive lattice \(D\) has finite meet and join
operations, least element \(0\), and greatest element \(1\).
The one-element lattice is permitted.
A filter contains \(1\), is upward closed, and is closed under
binary meets. An ideal contains \(0\), is downward closed, and
is closed under binary joins.
A filter is proper when it excludes \(0\).
A \emph{prime filter} is a proper filter \(P\) such that
\[
 a\vee b\in P\quad\Longrightarrow\quad a\in P\text{ or }b\in P.
\]
A bounded lattice homomorphism preserves \(0,1,\wedge,\vee\).
The principal filter and ideal are
\(\uparrow a=\{d:a\leq d\}\) and
\(\downarrow b=\{d:d\leq b\}\).
The characteristic map of a prime filter is a bounded lattice
homomorphism \(D\to\{0,1\}\), and every such homomorphism arises
this way. This assertion concerns lattice operations.
Preservation of any additional implication operation is a separate
condition.
\index{prime filter!distributive lattice}
\end{hbdefinition}

\begin{hbtheorem}{thm:p10-prime-filter-separation}{proved=annex:p10.priestley.1}{Prime-filter separation}
Assume BPIT. If \(F\) is a filter and \(I\) an ideal of \(D\)
with \(F\cap I=\varnothing\), there is a prime filter
\(P\supseteq F\) disjoint from \(I\).
In particular, if \(a\not\leq b\), a prime filter contains \(a\)
and omits \(b\).
If \(D\) is countable, these conclusions hold in \(\mathsf{ZF}\)
without BPIT.
\end{hbtheorem}

The proof first handles finite distributive lattices, then reduces
the general case to arbitrary-variable propositional compactness.
For the countable case it supplies a separate enumeration construction.
The prime-filter theorem and its Heyting application are treated in
\cite[Prime Filter Theorem and Corollary~2.14, p.~41;
Theorem~2.26, pp.~47--48]{Moraschini2023AlgebraLogicNotes}.

\subsection{Priestley spaces and representation}
\label{sec:p10-priestley-duality}

\begin{hbdefinition}{def:p10-priestley-space}{Priestley space and clopen upsets}
For a partially ordered set \(X\), an \emph{upset} \(U\) satisfies
\(x\in U,\ x\leq y\Rightarrow y\in U\).
A Priestley space is a compact Hausdorff space \(X\) with a partial
order such that
\[
 x\not\leq y\quad\Longrightarrow\quad
 \text{some clopen upset }U\text{ contains }x\text{ and omits }y.
\]
The empty ordered space is allowed.
Write \(\operatorname{ClopUp}(X)\) for its clopen upsets, a bounded
distributive lattice under intersection and union, with bounds
\(\varnothing,X\).
For a bounded distributive lattice \(D\), let
\[
 X_D=\{P:P\text{ is a prime filter of }D\},\qquad
 \widehat a=\{P\in X_D:a\in P\}.
\]
Order \(X_D\) by inclusion and give it the topology generated by
the sets \(\widehat a\) and their complements.
Its basic sets can be written
\(\widehat a\setminus\widehat b\):
finite positive memberships combine by meet and finite negative
memberships by join, with empty meet \(1\) and empty join \(0\).
\index{Priestley space}\index{Priestley duality}\index{clopen upset}
\end{hbdefinition}

\begin{hbtheorem}{thm:p10-priestley-representation}{proved=annex:p10.priestley.2}{Priestley representation}
Assuming BPIT, \(X_D\) is a Priestley space, and
\[
 \eta_D:D\longrightarrow\operatorname{ClopUp}(X_D),
 \qquad \eta_D(a)=\widehat a,
\]
is a bounded lattice isomorphism.
It preserves and reflects the lattice order.
The trivial lattice corresponds to the empty Priestley space.
\end{hbtheorem}

\begin{hbtheorem}{thm:p10-priestley-reconstruction}{proved=annex:p10.priestley.3}{Recovering a Priestley space}
For a Priestley space \(X\), the evaluation map
\[
 \varepsilon_X:X\longrightarrow
      X_{\operatorname{ClopUp}(X)},\qquad
 \varepsilon_X(x)=\{U\in\operatorname{ClopUp}(X):x\in U\},
\]
is an order isomorphism and a homeomorphism.
In particular, every prime filter of its clopen-upset lattice
is the membership filter of exactly one point.
The assertion includes \(X=\varnothing\).
\end{hbtheorem}

\begin{hbtheorem}{thm:p10-priestley-functors}{proved=annex:p10.priestley.4}{Functorial Priestley duality}
Assume BPIT.
Bounded distributive lattices with bounded lattice homomorphisms
are contravariantly equivalent to Priestley spaces with continuous
order-preserving maps.
A homomorphism \(h:D\to E\) gives the map
\[
 X(h):X_E\to X_D,\qquad P\mapsto h^{-1}[P].
\]
A continuous order-preserving \(g:X\to Y\) gives the lattice
homomorphism
\[
 \operatorname{ClopUp}(g):
 \operatorname{ClopUp}(Y)\to\operatorname{ClopUp}(X),
 \qquad U\mapsto g^{-1}[U].
\]
The maps \(\eta_D\) and \(\varepsilon_X\) are natural, and the
two reconstruction composites are identities.
\end{hbtheorem}

These statements give the ordered-space representation of
\textcite{Priestley1970}, with full proofs of the maps and
reconstruction supplied here.
Priestley's later author presentation identifies the hom-functors
and evaluation maps explicitly
\cite[slides~17--20]{Priestley2019Slides}.
The Heyting specialization and its order-topological conventions
are also described in
\cite[\S2.3.3, pp.~33--37]{Bezhanishvili2006Thesis}.

\paragraph{Why the order is part of the representation.}
The three-element chain \(D=\{0<a<1\}\) has exactly two prime
filters, \(P=\{1\}\) and \(Q=\{a,1\}\), with \(P\subsetneq Q\).
The topology is discrete: \(\widehat a=\{Q\}\) and its complement
is \(\{P\}\).
There are exactly three clopen upsets,
\(\varnothing,\{Q\},\{P,Q\}\), reproducing \(D\).
The underlying discrete space has four clopens, including
\(\{P\}\), which is not an upset.
The clopen-upset restriction records the original non-Boolean
lattice structure.

%% Countably many dense requirements in a partially ordered set.

\subsection{The Rasiowa--Sikorski dense-set lemma}
\label{sec:p10-rasiowa-sikorski}

\begin{hbdefinition}{def:p10-dense-generic-filter}{Dense sets and generic filters in a poset}
Let \((P,\leq)\) be a nonempty partially ordered set.
The convention is that \(q\leq p\) means that \(q\) is a
\emph{stronger} condition than \(p\).
A subset \(D\subseteq P\) is \emph{downward dense} when
\[
 \forall p\in P\ \exists q\in D\quad q\leq p.
\]
A \emph{filter in the poset} is a nonempty \(G\subseteq P\) that
is upward closed and downward directed:
\[
\begin{aligned}
 q\in G,\ q\leq p&\quad\Longrightarrow\quad p\in G,\\
 p,q\in G&\quad\Longrightarrow\quad
       \exists r\in G\ (r\leq p\text{ and }r\leq q).
\end{aligned}
\]
For a family \(\mathcal D\) of downward dense subsets of \(P\),
the filter is \emph{\(\mathcal D\)-generic} when
\[
 G\cap D\ne\varnothing\qquad(D\in\mathcal D).
\]
Density requires a stronger member of \(D\) below each condition;
it imposes no closure condition on \(D\).
These are the order and genericity conventions of
\cite[\S8.1, pp.~130--132]{Ciesielski1997}.
\index{dense set!in a poset}\index{generic filter}\index{filter!poset}
\end{hbdefinition}

\begin{hbdefinition}{def:p10-dependent-choice}{Dependent choice for a prescribed start}
We use the following form of the axiom of dependent choice,
\(\mathsf{DC}\): if \(S\) is a nonempty set and
\(R\subseteq S\times S\) satisfies
\[
 \forall s\in S\ \exists t\in S\quad sRt,
\]
then for every prescribed \(s_0\in S\) there is a sequence
\((s_n)_{n\in\mathbb N}\) starting at that element and satisfying
\(s_nRs_{n+1}\) for every \(n\).
The next proof uses this principle to make successive choices of
stronger conditions, where each choice depends on the preceding one.
\index{dependent choice}
\end{hbdefinition}

\begin{hblemma}{thm:p10-rasiowa-sikorski}{proved=annex:p10.rs.1}{Rasiowa--Sikorski lemma}
Let \((P,\leq)\) be a nonempty poset, let \(p\in P\), and let
\(\mathcal D\) be a countable family of downward dense subsets of
\(P\), where countable includes finite.
Under \(\mathsf{ZF}+\mathsf{DC}\), there is a
\(\mathcal D\)-generic filter \(G\) containing \(p\).

If \(\mathcal D=\{D_n:n\in\mathbb N\}\) is given by a sequence,
one can take
\[
\begin{gathered}
 p_0=p,\qquad p_{n+1}\leq p_n,\qquad p_{n+1}\in D_n,\\
 G=\{q\in P:\exists n\in\mathbb N\quad p_n\leq q\}.
\end{gathered}
\]
If \(P\) has a fixed enumeration \(e:\mathbb N\to P\) onto it,
the same conclusion is provable in \(\mathsf{ZF}\) alone, by always
using the first enumerated suitable extension.
For a finite family \(\mathcal D\), \(\mathsf{ZF}\) suffices even
when \(P\) is arbitrary.  If \(P\) is finite, a filter containing
\(p\) can meet every downward dense subset of \(P\) at once in
\(\mathsf{ZF}\).
\end{hblemma}

The countable dense-set theorem and its descending-sequence construction
are \cite[Theorem~8.1.2, pp.~132--133]{Ciesielski1997}; the proof here
makes the sufficient choice assumptions explicit.
This theorem constructs a directed poset filter meeting specified dense
sets.  Theorem~\ref{thm:p10-prime-filter-separation} instead extends a
lattice filter to a prime filter while avoiding an ideal.
The arbitrary-poset conclusion above is proved using \(\mathsf{DC}\);
no implication from BPIT alone is asserted.

%% Heyting algebras, IPC, and prime-filter Kripke completeness.

\subsection{Heyting algebras and intuitionistic completeness}
\label{sec:p10-heyting-completeness}

\begin{hbdefinition}{def:p10-heyting-algebra}{Heyting algebra}
A Heyting algebra is a nonempty bounded distributive lattice
\(H=(H,\wedge,\vee,0,1)\), equipped with an operation
\(\Rightarrow\) satisfying
\[
 c\leq(a\Rightarrow b)
       \quad\Longleftrightarrow\quad c\wedge a\leq b
       \qquad(a,b,c\in H).
\]
This is the residuation law: \(a\Rightarrow b\) is the greatest
element whose meet with \(a\) lies below \(b\).
It determines the operation uniquely.
Write \(\neg a=a\Rightarrow0\).
The one-element algebra, with \(0=1\), is allowed.
A Heyting homomorphism preserves the bounds and all three
binary operations \(\wedge,\vee,\Rightarrow\).
The arrow \(\Rightarrow\) here acts on algebra elements; the
object-language connective is \(\to\), as in
Definition~\ref{def:p08-heyting-algebra}.
The source notation \(a\to b\) for algebra elements denotes this
same uniquely determined residual; thus \(a\to b=a\Rightarrow b\)
when the operands are elements of a Heyting algebra.
\index{Heyting algebra}\index{residuation!Heyting algebra}
\end{hbdefinition}

\begin{hbtheorem}{thm:p10-heyting-examples}{proved=annex:p10.heyting.1}{Open-set and upset Heyting algebras}
For every topological space \(X\), its open subsets form a
Heyting algebra \(\mathcal O(X)\), with finite meets given by
intersection, finite joins by union, bounds \(\varnothing,X\),
and
\[
 U\Rightarrow V=\operatorname{int}_X((X\setminus U)\cup V).
\]
For every preordered set \((W,\leq)\), the upsets form a
Heyting algebra \(\operatorname{Up}(W)\), with the same
set operations and
\[
 U\Rightarrow V
   =\{w\in W:\forall v\geq w\ (v\in U\Longrightarrow v\in V)\}.
\]
These statements include empty \(X\) or \(W\), when the respective
algebra is trivial.  A nonempty \(W\) gives a nontrivial upset
algebra and supports the intuitionistic Kripke semantics of
Part~VIII.
\end{hbtheorem}

The residuation definition appears in
\cite[Definition~2.2.4, p.~21]{Bezhanishvili2006Thesis};
the upset construction and prime-filter order are given there
in \S2.2.3, pp.~26--27.  The proofs below also cover preorders.

\begin{hbdefinition}{def:p10-ipc-lindenbaum}{IPC derivability and its Lindenbaum algebra}
Let \(\mathsf{Prop}\) be any set of propositional variables.
Use the language \(\bot,\land,\lor,\to\) of
Section~\ref{sec:intuitionistic-logic-01}, with
\(\neg\varphi=(\varphi\to\bot)\) and
\(\top=(\bot\to\bot)\).
Write \(\Gamma\vdash_{\mathsf{IPC}}\varphi\) for that
finite-support \(\mathsf{IPL}\) derivability relation.
Thus IPC and IPL denote the same logic, and every proof here
uses the propositional fragment of
\hyperref[sheet:lj]{\texttt{sheet:lj}}.
Uniform substitution transforms a derivation of
\(\Gamma\vdash\varphi\) into one of
\(\sigma[\Gamma]\vdash\sigma(\varphi)\); it does not permit
free substitution in a conclusion while leaving its premises fixed.

Define
\[
 \varphi\sim\psi
 \quad\Longleftrightarrow\quad
 \{\varphi\}\vdash_{\mathsf{IPC}}\psi
       \text{ and }\{\psi\}\vdash_{\mathsf{IPC}}\varphi.
\]
Let \(H_{\mathsf{IPC}}\) have formula classes \([\varphi]\),
with bounds \([\bot],[\top]\) and operations induced by
\(\land,\lor,\to\).
Its order is
\[
 [\varphi]\leq[\psi]
          \quad\Longleftrightarrow\quad
          \{\varphi\}\vdash_{\mathsf{IPC}}\psi.
\]
The next theorem proves that these definitions are well-defined
and give a Heyting algebra.
A valuation \(v:\mathsf{Prop}\to H\) into a Heyting algebra extends
recursively, interpreting the formula connectives by
\(0,\wedge,\vee,\Rightarrow\).
\end{hbdefinition}

\begin{hbtheorem}{thm:p10-heyting-completeness}{proved=annex:p10.heyting.2}{Algebraic completeness of IPC}
The algebra \(H_{\mathsf{IPC}}\) is a nontrivial Heyting algebra.
For every formula \(\varphi\),
\[
 \vdash_{\mathsf{IPC}}\varphi
 \quad\Longleftrightarrow\quad
 [\varphi]=1
 \quad\Longleftrightarrow\quad
 \llbracket\varphi\rrbracket_v^H=1_H
       \text{ for every Heyting algebra }H\text{ and valuation }v.
\]
The canonical valuation \(v_0(p)=[p]\) gives
\(\llbracket\varphi\rrbracket_{v_0}=[\varphi]\), and therefore
refutes every non-theorem.
More generally, for a finite premise set \(\Gamma\),
\[
 \Gamma\vdash_{\mathsf{IPC}}\varphi
 \quad\Longleftrightarrow\quad
 \bigwedge_{\gamma\in\Gamma}\llbracket\gamma\rrbracket_v^H
       \leq\llbracket\varphi\rrbracket_v^H
       \text{ for every }H,v.
\]
The empty meet is \(1\).
\end{hbtheorem}

The theoremhood statement is
\cite[Theorem~2.2.16, p.~25]{Bezhanishvili2006Thesis}.
The annex proves soundness directly for the displayed LJ rules,
then constructs the quotient order and its residual.

\begin{hbdefinition}{def:p10-heyting-prime-frame}{The prime-filter order of a Heyting algebra}
A subset \(P\subseteq H\) is a prime filter when it is a lattice filter containing \(1\),
closed upward and under finite meets, excluding \(0\), and
satisfying
\[
 a\vee b\in P\Longrightarrow a\in P\text{ or }b\in P.
\]
Let \(\operatorname{Pf}(H)\) be all such filters, ordered by inclusion.
For \(a\in H\), put
\[
 \beta_H(a)=\{P\in\operatorname{Pf}(H):a\in P\}.
\]
This is an upset of the prime-filter order.
The filters need not be maximal.
\end{hbdefinition}

\begin{hbtheorem}{thm:p10-heyting-prime-representation}{proved=annex:p10.heyting.3}{Prime-filter representation and the implication witness}
Assume BPIT.  For every Heyting algebra \(H\), the map
\[
 \beta_H:H\longrightarrow
       \operatorname{Up}(\operatorname{Pf}(H),\subseteq)
\]
is an injective Heyting homomorphism.
In particular, for \(a,b\in H\) and \(P\in\operatorname{Pf}(H)\),
\[
 a\Rightarrow b\in P
 \quad\Longleftrightarrow\quad
 \forall Q\supseteq P\
       (Q\in\operatorname{Pf}(H)\text{ and }a\in Q
                     \Longrightarrow b\in Q).
\]
If \(a\Rightarrow b\notin P\), a witnessing prime filter
\(Q\supseteq P\) contains \(a\) and omits \(b\).
If \(H\) is nontrivial, its prime-filter order is nonempty.
For trivial \(H\), it is empty and the displayed map is the
isomorphism between two one-element algebras.
For countable \(H\), these conclusions hold in ZF using the
countable case of Theorem~\ref{thm:p10-prime-filter-separation}.
\end{hbtheorem}

The prime-filter construction is the underlying ordered
representation in
\cite[Theorem~2.3.3 and its proof sketch, p.~28]{Bezhanishvili2006Thesis}.
The image is a Heyting algebra of upsets; it is not asserted
to contain every upset.

\begin{hbtheorem}{thm:p10-ipc-prime-kripke-completeness}{proved=annex:p10.heyting.4}{Strong local Kripke completeness of IPC}
Assume BPIT for an arbitrary variable set.
For every formula set \(\Gamma\) and formula \(\varphi\),
\[
 \Gamma\vdash_{\mathsf{IPC}}\varphi
       \quad\Longleftrightarrow\quad
 \Gamma\entails_{\mathsf{IPL}}\varphi
\]
with the persistent valuations, nonempty preorders, and local
forcing convention of
Definitions~\ref{def:p08-ipl-kripke-structure}--\ref{def:p08-ipl-consequence}.
The same equivalence holds if frames are required to be partial orders.
Every failure of derivability is witnessed on the prime filters
of \(H_{\mathsf{IPC}}\), ordered by inclusion, with
\[
 V(p)=\beta_{H_{\mathsf{IPC}}}([p]).
\]
If the variable set is countable, the entire argument works in
ZF by the countable prime-filter-separation theorem.
\end{hbtheorem}

%% Esakia spaces, Heyting homomorphisms, and the order back condition.

\subsection{Esakia duality}
\label{sec:p10-esakia-duality}

\begin{hbdefinition}{def:p10-esakia-space}{Esakia spaces}
For a subset \(S\) of a partially ordered set, put
\[
 \downarrow S=\{x:\exists s\in S\ (x\leq s)\},\qquad
 \uparrow x=\{y:x\leq y\}.
\]
An Esakia space is a Priestley space \(X\) such that
\(\downarrow C\) is clopen for every clopen subset \(C\) of \(X\).
The condition applies to all clopens, including those that are
not upsets.
The empty Esakia space is allowed, corresponding to the trivial
Heyting algebra.
For a Heyting algebra \(H\), the underlying order of \(X_H\) is
\((\operatorname{Pf}(H),\subseteq)\) from
Definition~\ref{def:p10-heyting-prime-frame}.
As point sets, \(\eta_H(a)=\beta_H(a)=\widehat a\); the Priestley
topology identifies these sets as clopen upsets.
This is the modern Priestley-space formulation in
\cite[Definition~2.3.20, p.~34]{Bezhanishvili2006Thesis}.
\index{Esakia space}\index{Esakia duality}
\end{hbdefinition}

\begin{hbtheorem}{thm:p10-esakia-objects}{proved=annex:p10.esakia.1}{Heyting algebras and Esakia spaces}
Assume BPIT.
The Priestley space \(X_H\) of a Heyting algebra \(H\) is Esakia,
and
\[
 \downarrow(\widehat a\setminus\widehat b)
       =X_H\setminus\widehat{a\Rightarrow b}.
\]
Conversely, the clopen upsets of an Esakia space form a Heyting
algebra with
\[
 U\Rightarrow V=X\setminus\downarrow(U\setminus V).
\]
The map \(\eta_H:a\mapsto\widehat a\) is a Heyting isomorphism
onto \(\operatorname{ClopUp}(X_H)\).
The evaluation map
\(x\mapsto\{U\in\operatorname{ClopUp}(X):x\in U\}\)
recovers \(X\) as an ordered topological space.
\end{hbtheorem}

\begin{hbdefinition}{def:p10-esakia-morphism}{Esakia morphisms}
An Esakia morphism \(g:X\to Y\) is a continuous order-preserving
map satisfying the order back condition
\[
 g(x)\leq y\quad\Longrightarrow\quad
       \exists z\geq x\quad g(z)=y.
\]
Equivalently, \(g[\uparrow x]=\uparrow g(x)\) for every \(x\).
Order preservation gives one inclusion and back gives the other.
These are continuous p-morphisms of the underlying partial orders
\cite[\S2.3.3, pp.~35--37]{Bezhanishvili2006Thesis}.
\end{hbdefinition}

\begin{hblemma}{thm:p10-heyting-prime-lift}{proved=annex:p10.esakia.2}{Lifting prime filters along a Heyting homomorphism}
Assume BPIT.
Let \(h:H\to K\) be a Heyting homomorphism, \(F\) a prime filter
of \(K\), and \(G\) a prime filter of \(H\), with
\(h^{-1}[F]\subseteq G\).
There is a prime filter \(F'\supseteq F\) of \(K\) satisfying
\[
 h^{-1}[F']=G.
\]
Thus the inverse-image map \(X_K\to X_H\) has the order back
condition.
\end{hblemma}

\begin{hbtheorem}{thm:p10-esakia-functors}{proved=annex:p10.esakia.3}{Functorial Esakia duality}
Assume BPIT.
Heyting algebras with Heyting homomorphisms are contravariantly
equivalent to Esakia spaces with Esakia morphisms.
The functors are the prime-filter and clopen-upset constructions.
On arrows they act by inverse image:
\[
\begin{aligned}
 h:H\to K&\quad\longmapsto\quad
     X(h):X_K\to X_H,\quad F\mapsto h^{-1}[F],\\
 g:X\to Y&\quad\longmapsto\quad
     \operatorname{ClopUp}(g):
     \operatorname{ClopUp}(Y)\to\operatorname{ClopUp}(X).
\end{aligned}
\]
The Priestley reconstruction isomorphisms are natural for these
restricted objects and arrows.
The trivial algebra and empty space are included.
\end{hbtheorem}

Esakia's original paper formulates the duality using perfect
topological Kripke models and strongly isotone maps
\cite[Definition~2, p.~298; Theorem~3, case \(\alpha=1\), p.~300]{Esakia1974}.
Its representation on p.~301 uses clopen upper cones, with
implication determined by their future-point condition.
The explicit modern prime-filter and Priestley-space
presentation above follows
\cite[Theorems~2.3.22 and~2.3.25, pp.~35, 37]{Bezhanishvili2006Thesis};
the annex supplies the full filter-lifting and morphism proofs.
The original paper and that thesis use nonempty spaces.
The empty-space case is included explicitly here to retain
the trivial Heyting algebra.

\paragraph{A monotone map need not preserve implication dually.}
Let \(X=\{x\}\) and \(Y=\{0<1\}\), both with the discrete topology.
These finite ordered spaces are Esakia.
The map \(g(x)=0\) is continuous and order-preserving, but does
not satisfy back, since \(1\geq g(x)\) has no preimage.
For the clopen upsets \(U=\{1\}\), \(V=\varnothing\) of \(Y\),
\[
 U\Rightarrow V=\varnothing,\qquad
 g^{-1}[U]\Rightarrow g^{-1}[V]=X.
\]
Thus inverse image under this map does not preserve Heyting
implication. The order back condition is part of the arrow
correspondence.

%% Intermediate theorem logics and varieties of Heyting algebras.

\subsection{Intermediate logics and Heyting varieties}
\label{sec:p10-intermediate-algebras}

Work in the fixed algebraic signature
\(\tau_H=(\wedge,\vee,\Rightarrow,0,1)\).
The residual \(\Rightarrow\) is the operation of
Definition~\ref{def:p10-heyting-algebra}.
Write \(\mathsf{HA}\) for all Heyting algebras, including the trivial
one, and \(\mathsf{BA}\) for those whose bounded lattice reduct is
Boolean, with \(a\Rightarrow b=\neg a\vee b\).
Thus both classes are considered in the same signature.

\begin{hbtheorem}{thm:p10-heyting-equational}{proved=annex:p10.intermediate.1}{The equational class of Heyting algebras}
A bounded distributive lattice equipped with a binary operation
\(\Rightarrow\) is a Heyting algebra exactly when it satisfies
\[
\begin{aligned}
 a\Rightarrow a&=1,\\
 a\wedge(a\Rightarrow b)&=a\wedge b,\\
 b\wedge(a\Rightarrow b)&=b,\\
 a\Rightarrow(b\wedge c)&=(a\Rightarrow b)\wedge(a\Rightarrow c).
\end{aligned}
\]
Together with the bounded distributive lattice identities, these
are a finite equational axiomatization of \(\mathsf{HA}\).
Inside \(\mathsf{HA}\), the additional identity
\[
 a\vee(a\Rightarrow0)=1
\]
defines precisely \(\mathsf{BA}\).
\end{hbtheorem}

The finite residual axioms are
\cite[Theorem~2.2.6, pp.~21--22]{Bezhanishvili2006Thesis};
the Boolean characterization is Proposition~2.2.9 on p.~23.
The proof here includes both directions, so the use of varieties does
not depend on treating residuation as an equational axiom without
justification.

\begin{hbdefinition}{def:p10-intermediate-logic}{Intermediate theorem logics and their quotients}
Fix a countably infinite variable set \(X=\{p_0,p_1,\ldots\}\).
Use formulas built from \(X\) by \(\bot,\wedge,\vee,\to\), with
\(\neg\varphi=(\varphi\to\bot)\), \(\top=(\bot\to\bot)\), and
\(\varphi\leftrightarrow\psi=(\varphi\to\psi)\wedge(\psi\to\varphi)\).
IPC is the theorem set of Definition~\ref{def:p10-ipc-lindenbaum},
restricted to \(X\), and CPL is the classical theorem set in this
same grammar.
An \emph{intermediate logic} is a set of formulas \(L\) such that
\[
 \mathsf{IPC}\subseteq L\subseteq\mathsf{CPL},
\]
closed under modus ponens and uniform substitution.
Write \(\vdash_L\varphi\) for \(\varphi\in L\).
These are closure rules for theorem sets.
If local premises are used, a derivation from \(\Gamma\) is a finite
sequence using \(\Gamma\), theorems of \(L\), and modus ponens;
substitution does not act freely on premise-dependent conclusions.

Put
\[
 \varphi\equiv_L\psi\quad\Longleftrightarrow\quad
                    \vdash_L\varphi\leftrightarrow\psi,
 \qquad H_L=\operatorname{Fm}(X)/{\equiv_L}.
\]
The proposed operations on \(H_L\) are induced by the formula
connectives, with algebraic \(\Rightarrow\) induced by \(\to\), and
the proposed order is
\[
 [\varphi]_L\leq[\psi]_L
       \quad\Longleftrightarrow\quad\vdash_L\varphi\to\psi.
\]
The following theorem verifies these constructions.
\index{intermediate logic}\index{Lindenbaum--Tarski algebra!intermediate logic}
\end{hbdefinition}

\begin{hbdefinition}{def:p10-intermediate-variety-maps}{The logic--variety assignments}
For an intermediate logic \(L\), define
\[
 \mathsf V(L)=\{H\in\mathsf{HA}:H\models\lambda
                                      \text{ for every }\lambda\in L\},
\]
where validity means value \(1\) under every valuation \(X\to H\).
For a class \(\mathcal V\subseteq\mathsf{HA}\), put
\[
 \operatorname{Log}(\mathcal V)
    =\{\varphi:H\models\varphi\text{ for every }H\in\mathcal V\}.
\]
A subvariety of \(\mathsf{HA}\) is the class of Heyting algebras
satisfying some set of identities in \(\tau_H\).
Identities may be taken over the fixed infinite variable set \(X\):
every identity uses only finitely many variables, which can be renamed
to \(p_0,p_1,\ldots\) in order of first occurrence.
\index{variety!Heyting algebras}\index{algebraic semantics!intermediate logic}
\end{hbdefinition}

\begin{hbtheorem}{thm:p10-intermediate-algebra-completeness}{proved=annex:p10.intermediate.2}{Algebraic completeness of an intermediate logic}
For every intermediate \(L\), \(\equiv_L\) is a congruence,
\(H_L\) is a nontrivial Heyting algebra, and \(H_L\in\mathsf V(L)\).
For every formula \(\varphi\),
\[
 \varphi\in L
 \quad\Longleftrightarrow\quad [\varphi]_L=1
 \quad\Longleftrightarrow\quad H_L\models\varphi
 \quad\Longleftrightarrow\quad\mathsf V(L)\models\varphi.
\]
The valuation \(p\mapsto[p]_L\) refutes every non-theorem.
\end{hbtheorem}

\begin{hbtheorem}{thm:p10-intermediate-variety-correspondence}{proved=annex:p10.intermediate.3}{The lattice correspondence for intermediate logics}
The assignments \(L\mapsto\mathsf V(L)\) and
\(\mathcal V\mapsto\operatorname{Log}(\mathcal V)\) are mutually
inverse order-reversing bijections between intermediate logics and
the interval of subvarieties
\[
 \mathsf{BA}\subseteq\mathcal V\subseteq\mathsf{HA}.
\]
In particular,
\[
 \operatorname{Log}(\mathsf V(L))=L,\qquad
 \mathsf V(\operatorname{Log}(\mathcal V))=\mathcal V.
\]
The endpoints correspond as
\[
 \mathsf V(\mathsf{IPC})=\mathsf{HA},\qquad
 \mathsf V(\mathsf{CPL})=\mathsf{BA},
 \qquad
 \operatorname{Log}(\mathsf{HA})=\mathsf{IPC},\quad
 \operatorname{Log}(\mathsf{BA})=\mathsf{CPL}.
\]
These are complete lattices.  On the logic side, nonempty meets are
intersections and joins are the least MP- and substitution-closed
extensions of IPC containing the union.  The empty meet is CPL and
the empty join is IPC.
On the variety side, nonempty meets are intersections and joins are
the least subvarieties containing \(\mathsf{BA}\) and the union.
The empty meet is \(\mathsf{HA}\) and the empty join is \(\mathsf{BA}\).
The correspondence interchanges all these meets and joins.
\end{hbtheorem}

The source defines intermediate logics in Definition~2.1.4, p.~12,
states their algebraic completeness in Theorem~2.2.18, p.~25, and
gives the lattice anti-isomorphism in Theorem~2.2.19, p.~26
\cite{Bezhanishvili2006Thesis}.
The interval above makes the consistent-logic convention explicit:
the trivial algebra belongs to every variety in it, but the variety
consisting only of trivial algebras is outside this interval.
These are algebraic completeness statements; no assertion of ordinary
Kripke completeness for every intermediate logic follows from them.


Locales retain the frame of opens without assuming enough points to recover a
space, so pointfree topology is a genuine change of semantic carrier rather
than another spelling of ordinary topology \cite{Johnstone1982}.

%% Finitary Boolean operators, canonical extensions, and descriptive duality.

\subsection{Boolean operators and their relational representation}
\label{sec:p10-boolean-operators}

This development uses the held 19 October 2005 preprint of Venema's
\emph{Algebras and Coalgebras}; its local page and result numbers
are cited as preprint locators \cite{Venema2005Preprint}.
The published 2007 chapter has its own bibliography entry
\cite{Venema2007}.
For the published literature on varieties of complex algebras,
see also \textcite{Goldblatt1989}.
The general representation theorem is attributed to
\textcite{JonssonTarski1951}.
The proofs below assume \(\mathsf{ZF}\) plus BPIT, as in the
preceding Stone construction.

\begin{hbdefinition}{def:p10-bao}{Boolean algebra with operators}
Fix a set \(\Lambda\) of additional operation symbols, with a finite
arity \(n_f\in\mathbb N\) for every \(f\in\Lambda\).
A Boolean algebra with these operators is
\(\mathbf B=(B,(f^B)_{f\in\Lambda})\), whose Boolean reduct is \(B\),
such that for every positive-arity \(f\) and each coordinate,
\[
\begin{aligned}
 f^B(b_1,\ldots,0,\ldots,b_n)&=0,\\
 f^B(b_1,\ldots,a\vee c,\ldots,b_n)
 &=f^B(b_1,\ldots,a,\ldots,b_n)
       \vee f^B(b_1,\ldots,c,\ldots,b_n).
\end{aligned}
\]
These are coordinatewise normality and additivity.
They imply monotonicity in each coordinate.
For \(n_f=0\), \(f^B()\) is a designated Boolean element; the
coordinate conditions are vacuous and do not require it to be zero.
A homomorphism preserves the Boolean operations and every \(f^B\),
including nullary ones.

A normal modal algebra has one unary operator \(\Diamond\).
Its dual is \(\Box a=\neg\Diamond\neg a\), which preserves finite
meets and \(1\).
The general definition allows an arbitrary set of symbols of
different finite arities
\cite[preprint Definition~3.11, p.~9]{Venema2005Preprint}.
\index{Boolean algebra with operators}\index{modal algebra}
\end{hbdefinition}

\begin{hbdefinition}{def:p10-bao-complex-algebra}{Relational frames and complex algebras}
A relational \(\Lambda\)-frame
\(\mathfrak F=(W,(R_f)_{f\in\Lambda})\) has
\(R_f\subseteq W^{n_f+1}\).
Empty carriers are permitted in this algebraic construction, so
the trivial Boolean algebra can be represented. A pointed modal
model still requires a chosen point.
The complex algebra \(\mathfrak F^+\) has Boolean reduct
\(\mathcal P(W)\) and operations
\[
 f^{\mathfrak F^+}(A_1,\ldots,A_n)
 =\{w:\exists w_1\in A_1\cdots\exists w_n\in A_n\
                     R_f(w,w_1,\ldots,w_n)\}.
\]
For \(n=0\), this is \(\{w:R_f(w)\}\).
For positive arity these operations preserve arbitrary unions
in every coordinate, including empty unions.
This is the complex-operator construction of
\cite[preprint Definition~3.5, p.~8]{Venema2005Preprint},
with the empty-frame convention explicitly included here.
\end{hbdefinition}

\begin{hbdefinition}{def:p10-bao-ultrafilter-frame}{The ultrafilter frame of a Boolean algebra with operators}
Let \(X_B\) be the Stone space of the Boolean reduct.
Define \(\mathbf B^\bullet=(X_B,(R_f)_{f\in\Lambda})\) by
\[
 R_f(U,V_1,\ldots,V_n)
 \quad\Longleftrightarrow\quad
 \forall b_1\in V_1\cdots\forall b_n\in V_n\
                      f^B(b_1,\ldots,b_n)\in U.
\]
For a nullary symbol, this reads \(R_f(U)\) iff \(f^B()\in U\).
No witness choices enter the definition.
The source calls this the ultrafilter frame or canonical structure
\cite[preprint Definition~5.12, p.~19]{Venema2005Preprint}.
\end{hbdefinition}

\begin{hbexample}{def:p10-bao-binary-nullary}{Binary meet and a nullary operator}
For the binary operator \(f(a,b)=a\wedge b\), the canonical relation is
\[
 R_f(U,V_1,V_2)\quad\Longleftrightarrow\quad U=V_1=V_2.
\]
Indeed, setting the other argument to \(1\) in the defining relation
gives \(V_i\subseteq U\), hence equality of the ultrafilters.
Conversely, a single ultrafilter is closed under meets.
The complex operation on subsets of \(X_B\) is therefore intersection.

For a nullary operator with value \(c\in B\), the relation is
\(R_f(U)\) iff \(c\in U\), and the complex nullary value is
\(\widehat c\).
The choices \(c=1\) and \(c=0\) give the whole Stone space and the
empty set, respectively.
\end{hbexample}

\begin{hblemma}{thm:p10-bao-ultrafilter-witness}{proved=annex:p10.bao.1}{Simultaneous ultrafilter witnesses}
Let \(f:B^n\to B\) be an operator, \(U\in X_B\), and
\(F_1,\ldots,F_n\) proper filters such that
\[
 f(a_1,\ldots,a_n)\in U
              \quad\text{whenever }a_i\in F_i\ (1\leq i\leq n).
\]
There are ultrafilters \(V_i\supseteq F_i\) such that
\(R_f(U,V_1,\ldots,V_n)\).
In particular, if \(f(a_1,\ldots,a_n)\in U\), the \(V_i\) can be
chosen to contain the respective \(a_i\).
For \(n=0\), the conclusion has the empty witness tuple and is
exactly the assumption \(f()\in U\).
\end{hblemma}

Normality is needed for the element-witness conclusion.
The unary constant-one operation on the two-element Boolean
algebra preserves binary joins and sends \(0\) into its ultrafilter.
It cannot have a witness ultrafilter containing that input \(0\).

\begin{hbtheorem}{thm:p10-jonsson-tarski}{proved=annex:p10.bao.2}{J\'onsson--Tarski representation}
The Stone map
\[
 \beta_B:B\longrightarrow(\mathbf B^\bullet)^+,
 \qquad b\longmapsto\widehat b,
\]
is an embedding of Boolean algebras with operators:
\[
 \widehat{f^B(b_1,\ldots,b_n)}
      =f^{(\mathbf B^\bullet)^+}(\widehat b_1,\ldots,\widehat b_n).
\]
This includes every finite arity and every nullary operation.
The image of the Boolean reduct is
\(\operatorname{Clop}(X_B)\).
For a trivial \(\mathbf B\), its ultrafilter frame is empty and
its complex algebra is trivial.
\end{hbtheorem}

The source states this theorem for the full finitary signature but
gives its operator-preservation proof for one unary diamond
\cite[preprint Theorem~5.15, p.~20]{Venema2005Preprint}.
The annex proves the simultaneous filter-extension lemma and all
finite-arity cases.

\subsection{Canonical extensions}
\label{sec:p10-bao-canonical-extension}

\begin{hbdefinition}{def:p10-bao-canonical-extension}{The canonical extension and its Boolean completion}
The canonical extension of \(\mathbf B\) is
\[
 \mathbf B^\sigma=(\mathbf B^\bullet)^+,
\]
with the embedding \(\beta_B\).
This is the source's canonical-extension construction
\cite[preprint Definition~5.16, p.~21]{Venema2005Preprint}.
Its underlying Boolean algebra is complete and atomic.

For a Boolean embedding \(e:B\to D\) into a complete Boolean
algebra, call a meet of elements of \(e[B]\) \emph{closed} and a
join of them \emph{open}, including empty meets and joins.
The embedding is \emph{dense} when every element of \(D\) is
both a join of closed elements and a meet of open elements.
It is \emph{compact} when
\[
 \bigwedge e[S]\leq\bigvee e[T]\quad(S,T\subseteq B)
\]
implies the corresponding inequality for some finite
\(S_0\subseteq S\), \(T_0\subseteq T\).
Equivalently, an inequality between the meet of a family of
closed elements and the join of a family of open elements has
a witness using finitely many members of those two families.
The annex proves the equivalence of these compactness formulations.
All infinite meets and joins here are taken in the completion \(D\).
These are order-theoretic completion conditions
\cite[preprint \S7.2, Definition~7.1]{Venema2005Preprint}.
\index{canonical extension}\index{Boolean completion}
\end{hbdefinition}

\begin{hbtheorem}{thm:p10-bao-canonical-completion}{proved=annex:p10.bao.3}{The dense and compact canonical completion}
The Boolean embedding
\(\beta_B:B\to\mathcal P(X_B)\) is dense and compact.
Any other dense and compact Boolean completion of \(B\) is
isomorphic to it by a unique Boolean isomorphism commuting with
the embeddings.
This uniqueness concerns the Boolean completion.
The operators of \(\mathbf B^\sigma\) are the explicitly defined
complex operators; each preserves arbitrary joins in every
positive-arity coordinate.
\end{hbtheorem}

The existence and uniqueness statement is the source's
Theorem~7.2, with the Boolean atomicity argument supplied in
the annex \cite[preprint pp.~38--39]{Venema2005Preprint}.
The embedding need not preserve an infinite join that already
exists in \(B\).
For example, let \(B\) be the Boolean algebra of finite and
cofinite subsets of \(\mathbb N\).
The join of all singleton elements is \(\mathbb N=1\) in \(B\).
The cofinite sets form an ultrafilter \(U_\infty\), since exactly
one of any finite/cofinite set and its complement is cofinite.
Then \(U_\infty\in\widehat1\) but
\(U_\infty\notin\bigcup_n\widehat{\{n\}}\).
Thus that union, the join of the singleton images in the
completion, is strictly below the image of their join in \(B\).

The embedding transfers a valuation in \(\mathbf B\) to one whose
values lie in \(\beta_B[B]\).
Consequently every algebraic equation valid in \(\mathbf B\)
holds for those transported valuations.
Validity of an additional equation for all subsets of \(X_B\)
requires a further argument; it is the issue addressed by
canonicity in the following modal-completeness development.

\subsection{Admissible sets and descriptive general frames}
\label{sec:p10-bao-general-frames}

\begin{hbdefinition}{def:p10-bao-general-frame}{A general frame for a finitary operator signature}
A general frame is
\(\mathcal G=(W,(R_f)_{f\in\Lambda},\mathcal A)\), where
\(\mathcal A\subseteq\mathcal P(W)\) is a Boolean subalgebra
closed under every complex operator \(f^{\mathfrak F^+}\).
Its elements are the admissible proposition sets.
Valuations in a general frame assign proposition variables to
members of \(\mathcal A\); recursive formula evaluation remains
in \(\mathcal A\).
Write \(\mathcal G^+\) for the Boolean algebra with operators
whose carrier is precisely \(\mathcal A\).
The underlying relational frame's full complex algebra has
carrier \(\mathcal P(W)\).
For the unary nonempty case this agrees with
Definition~\ref{def:p08-general-frame}
\cite[preprint Definition~5.24, p.~22]{Venema2005Preprint}.
\end{hbdefinition}

\begin{hbdefinition}{def:p10-bao-descriptive}{Descriptive general frames}
A general frame is \emph{differentiated} when its admissible sets
separate distinct points.
It is \emph{compact} when every family of admissible sets with
nonempty finite intersections has nonempty total intersection.
Here the empty finite intersection is \(W\); this convention
also permits the empty frame.
It is \emph{tight} when for every \(f\) and tuple,
\[
 R_f(w,w_1,\ldots,w_n)
 \quad\Longleftrightarrow\quad
 \forall A_i\in\mathcal A\,
 \left(\bigwedge_{i=1}^n w_i\in A_i
       \ \Longrightarrow\
       w\in f^{\mathcal G^+}(A_1,\ldots,A_n)\right).
\]
The nullary case has no \(A_i\) to quantify over.
A \emph{descriptive} general frame has all three properties.
Its topology consists of unions of admissible sets
\cite[preprint Definition~5.25, p.~22]{Venema2005Preprint}.
\end{hbdefinition}

\begin{hbtheorem}{thm:p10-bao-descriptive-representation}{proved=annex:p10.bao.4}{Descriptive representation and reconstruction}
For every \(\mathbf B\),
\[
 \mathbf B_*=(X_B,(R_f),\{\widehat b:b\in B\})
\]
is descriptive, and \(\beta_B:\mathbf B\to(\mathbf B_*)^+\)
is an isomorphism of Boolean algebras with operators.
Conversely, for every descriptive \(\mathcal G\), its induced
topology is a Stone topology, its admissible sets are exactly
the clopens, and
\[
 w\longmapsto\{A\in\mathcal A:w\in A\}
\]
is an isomorphism of general frames
\(\mathcal G\to(\mathcal G^+)_*\).
It is a homeomorphism, preserves and reflects every relation,
and identifies the admissible families.
\end{hbtheorem}

\begin{hbdefinition}{def:p10-bao-general-morphism}{Bounded morphisms respecting admissible sets}
A morphism \(g:\mathcal G\to\mathcal H\) of general frames is a
function on their carriers such that preimages of admissible
sets are admissible, and for each \(f\) it has both properties:
\[
\begin{aligned}
 R_f^{\mathcal G}(w,u_1,\ldots,u_n)
   &\Longrightarrow
     R_f^{\mathcal H}(g(w),g(u_1),\ldots,g(u_n)),\\
 R_f^{\mathcal H}(g(w),v_1,\ldots,v_n)
   &\Longrightarrow
     \exists u_1,\ldots,u_n\
       \left(R_f^{\mathcal G}(w,u_1,\ldots,u_n)
                 \land\bigwedge_i g(u_i)=v_i\right).
\end{aligned}
\]
For a nullary symbol these say
\(R_f^{\mathcal G}(w)\) iff \(R_f^{\mathcal H}(g(w))\).
The admissible-preimage condition makes the map continuous for
the induced topologies
\cite[preprint Definition~5.27, p.~23]{Venema2005Preprint}.
\end{hbdefinition}

\begin{hbtheorem}{thm:p10-bao-descriptive-duality}{proved=annex:p10.bao.5}{Duality with descriptive general frames}
Boolean algebras with the fixed \(\Lambda\)-operators and their
homomorphisms are contravariantly equivalent to descriptive
\(\Lambda\)-general frames with the preceding morphisms.
A homomorphism \(h:\mathbf A\to\mathbf B\) gives
\[
 h_*:\mathbf B_*\to\mathbf A_*,
 \qquad h_*(U)=h^{-1}[U],
\]
and a general-frame morphism gives the inverse-image
homomorphism on its admissible algebras.
The reconstruction maps of
Theorem~\ref{thm:p10-bao-descriptive-representation} are natural.
These statements include trivial algebras and empty descriptive
frames \cite[preprint Proposition~5.18 and Theorem~5.28,
pp.~21, 23]{Venema2005Preprint}.
\end{hbtheorem}

%% Unary normal modal Lindenbaum algebras and general-frame completeness.

\subsection{Normal modal logics, Lindenbaum algebras, and general frames}
\label{sec:p10-modal-algebra-completeness}

\begin{hbdefinition}{def:p10-normal-modal-algebra-logic}{The fixed modal language and its algebras}
Fix the propositional-variable set \(\mathsf{Prop}\) of Part~VIII.
Formulas are built using \(\bot,\neg,\wedge,\Box\), with the other
Boolean connectives abbreviated as usual and
\(\Diamond\varphi=\neg\Box\neg\varphi\).
A \emph{normal modal logic} \(L\) is a set of formulas containing
all classical propositional tautology instances and
\[
 \Box(\varphi\to\psi)\to(\Box\varphi\to\Box\psi),
\]
and closed under modus ponens, uniform substitution, and
theorem necessitation: \(\varphi\in L\) implies \(\Box\varphi\in L\).
Write \(\vdash_L\varphi\) for \(\varphi\in L\).
The logic is consistent when \(\bot\notin L\).
For premises \(\Gamma\), a derivation uses members of \(\Gamma\),
members of \(L\), and modus ponens.  Necessitation applies only
to theorems, whose boxes already belong to \(L\); uniform
substitution is likewise a closure rule for \(L\), not a rule
for freely substituting in premises.
This extends the local convention of
Section~\ref{sec:modal-temporal-epistemic-01} to an arbitrary normal \(L\).

A \emph{normal modal algebra} is a Boolean algebra \(A\), possibly
trivial, with a unary operation \(\Box_A\) satisfying
\[
 \Box_A1=1,\qquad
 \Box_A(a\wedge b)=\Box_Aa\wedge\Box_Ab.
\]
Its dual \(\Diamond_Aa=\neg\Box_A\neg a\) preserves \(0\) and
binary joins, so this is the unary case of a Boolean algebra
with operators.
A valuation \(v:\mathsf{Prop}\to A\) extends recursively to
\(\llbracket\varphi\rrbracket_v^A\).
Write \(A\models\varphi\) when that value is \(1_A\) for every
valuation, and let
\[
 \operatorname{Alg}(L)
   =\{A:A\text{ is a normal modal algebra and }
                   A\models\lambda\text{ for every }\lambda\in L\}.
\]
Thus membership in \(\operatorname{Alg}(L)\) requires validity
under all algebra valuations.
\end{hbdefinition}

\begin{hbdefinition}{def:p10-modal-lindenbaum}{The normal modal Lindenbaum algebra}
Put
\[
 \varphi\equiv_L\psi
       \quad\Longleftrightarrow\quad
       \vdash_L\varphi\leftrightarrow\psi.
\]
The following theorem verifies that this is a congruence.
Let \(F_L\) be the quotient of the formula algebra by this relation,
with classes \([\varphi]_L\) and induced operations
\[
\begin{gathered}
 0=[\bot]_L,\quad 1=[\top]_L,\quad
 \neg[\varphi]_L=[\neg\varphi]_L,\\
 [\varphi]_L\wedge[\psi]_L=[\varphi\wedge\psi]_L,\qquad
 \Box[\varphi]_L=[\Box\varphi]_L.
\end{gathered}
\]
This is the \emph{normal modal Lindenbaum algebra} of \(L\).
\index{Lindenbaum--Tarski algebra!normal modal logic}
\end{hbdefinition}

\begin{hbtheorem}{thm:p10-modal-algebra-completeness}{proved=annex:p10.modal-algebra.1}{Algebraic completeness for every normal modal logic}
The relation \(\equiv_L\) is a congruence, \(F_L\) is a normal
modal algebra, and \(F_L\in\operatorname{Alg}(L)\).
For every formula \(\varphi\),
\[
 \vdash_L\varphi
 \quad\Longleftrightarrow\quad
 [\varphi]_L=1
 \quad\Longleftrightarrow\quad
 F_L\models\varphi
 \quad\Longleftrightarrow\quad
 \operatorname{Alg}(L)\models\varphi.
\]
If \(\varphi\notin L\), the valuation \(p\mapsto[p]_L\)
refutes \(\varphi\) in \(F_L\).
If \(L\) is inconsistent, it contains every formula and
\(\operatorname{Alg}(L)\) consists exactly of the trivial modal
algebras.
\end{hbtheorem}

The quotient and algebraic completeness are the unary normal case of
\cite[Definitions~3.14 and~3.16, Proposition~3.15, and
Theorems~3.13 and~3.17, pp.~10--11]{Venema2005Preprint}.
For an arbitrary valuation in the quotient, the annex chooses
representatives only for the finitely many variables in the formula
being evaluated.

\begin{hbdefinition}{def:p10-modal-canonical-general-frame}{The canonical general frame and its underlying frame}
Suppose \(L\) is consistent.
Apply the unary case of Theorem~\ref{thm:p10-jonsson-tarski} to
\(F_L\), using the Boolean prime ideal theorem.
Let \(X_L\) be the ultrafilters of its Boolean reduct, and put
\[
 \widehat a=\{U\in X_L:a\in U\},\qquad
 U R_L V
   \quad\Longleftrightarrow\quad
       \forall a\in V\ \Diamond a\in U.
\]
Equivalently,
\[
 U R_L V
   \quad\Longleftrightarrow\quad
       \{a:\Box a\in U\}\subseteq V.
\]
Define
\[
 \mathfrak G_L=(X_L,R_L,\mathcal A_L),\qquad
 \mathcal A_L=\{\widehat a:a\in F_L\},\qquad
 \mathfrak F_L=(X_L,R_L).
\]
The admissible sets \(\mathcal A_L\) are exactly the Stone clopens
and are closed under Boolean and modal operations, as in
Definition~\ref{def:p10-bao-general-frame}.
The \emph{canonical general frame} \(\mathfrak G_L\) restricts
variable valuations to \(\mathcal A_L\).
Its \emph{underlying canonical frame} \(\mathfrak F_L\) instead
permits every subset of \(X_L\) as a variable value.
These are the unary constructions of
\cite[Example~5.13 and Definition~5.24, pp.~20, 22]{Venema2005Preprint}.
\end{hbdefinition}

\begin{hbtheorem}{thm:p10-modal-general-frame-completeness}{proved=annex:p10.modal-algebra.2}{General-frame completeness and the canonical truth lemma}
Assume the Boolean prime ideal theorem.
For consistent normal \(L\), the frame \(\mathfrak G_L\) validates
every member of \(L\) under every admissible valuation.
For the canonical valuation \(v_0(p)=\widehat{[p]_L}\),
\[
 \mathfrak G_L,v_0,U\models\varphi
           \quad\Longleftrightarrow\quad[\varphi]_L\in U.
\]
Consequently \(\mathfrak G_L\models\varphi\) exactly when
\(\varphi\in L\).
The map
\[
 U\longmapsto\{\varphi:[\varphi]_L\in U\}
\]
identifies \(\mathfrak F_L\) with the maximal-consistent-set
canonical frame of Method~\ref{def:p08-modal-canonical-model}.
\end{hbtheorem}

\begin{hbtheorem}{thm:p10-modal-general-frame-strong}{proved=annex:p10.modal-algebra.3}{Strong completeness for local general-frame consequence}
Assume the Boolean prime ideal theorem.
Let \(\operatorname{Gen}(L)\) be all nonempty general frames
validating every member of \(L\) under every admissible valuation.
For every set \(\Gamma\) of formulas,
\[
 \Gamma\vdash_L\varphi
       \quad\Longleftrightarrow\quad
 \Gamma\models^\ell_{\operatorname{Gen}(L)}\varphi.
\]
Here derivations have the theorem-only necessitation convention of
Definition~\ref{def:p10-normal-modal-algebra-logic}, and the right side
requires all premises at the same point.
For consistent \(L\), every failure of derivability already has
a counterexample in \(\mathfrak G_L\) under \(v_0\).
For inconsistent \(L\), \(\operatorname{Gen}(L)\) is empty and
both sides hold for all \(\Gamma,\varphi\).
\end{hbtheorem}

\begin{hbdefinition}{def:p10-modal-canonicity}{Canonicity and unrestricted valuations}
For a consistent normal logic in the fixed language, the
\emph{canonical-frame validity condition} is
\[
 \mathfrak F_L\models\lambda\quad\text{for every }\lambda\in L,
\]
where valuations range over all subsets of \(X_L\).
This is an additional condition beyond validity on
\(\mathfrak G_L\).

When varying the set of variables to a set \(Y\), interpret \(L\)
schematically: let \(L(Y)\) be the least normal logic over \(Y\)
containing all substitution instances of members of \(L\).
In the all-variable-set convention, \(L\) is \emph{canonical}
when the canonical frame of \(L(Y)\) validates \(L(Y)\) for
every \(Y\).
This is the convention in
\cite[Definition~6.10, p.~31]{Venema2005Preprint}.
The next theorem needs only the fixed-language condition.
\index{canonicity!normal modal logic}
\end{hbdefinition}

\begin{hbtheorem}{thm:p10-modal-canonicity-completeness}{proved=annex:p10.modal-algebra.4}{Completeness from canonical-frame validity}
Assume the Boolean prime ideal theorem.
If consistent normal \(L\) satisfies the canonical-frame validity
condition, then \(L\) is strongly complete for local consequence
over the class \(\operatorname{Fr}(L)\) of nonempty ordinary
Kripke frames validating \(L\).

A sufficient algebraic condition is that
\(\operatorname{Alg}(L)\) be closed under
\[
 A\longmapsto A^\sigma
       =\text{the full complex algebra of the ultrafilter frame of }A.
\]
This closure property is called canonicity of the algebra class.
\end{hbtheorem}

The representation embedding \(A\hookrightarrow A^\sigma\)
preserves evaluations whose variable values lie in its image.
Equations true in a subalgebra need not hold under all valuations
in the larger algebra.  Thus the embedding alone supplies
truth preservation on admissible values.  Together with the
Lindenbaum construction this gives general-frame completeness;
the additional validity or
closure condition above supplies the ordinary-frame conclusion.
The algebraic closure formulation is
\cite[Definition~6.9 and Theorem~6.12, p.~31]{Venema2005Preprint}.
The equivalence with logical canonicity in that source's
Theorem~6.11 uses its all-variable-set convention.
The Kripke-incomplete normal logics of
Theorem~\ref{thm:p08-kripke-incomplete-logics} illustrate why the
ordinary-frame conclusion needs its own hypotheses.

%% Finitary normal modal logics and descriptive general-frame completeness.

\subsection{Normal logics for arbitrary finitary operator signatures}
\label{sec:p10-polyadic-modal-algebras}

The algebraic and local completeness arguments extend to the full
signature of Definition~\ref{def:p10-bao}.  The number of operator
symbols need not be finite; each individual operator has finite arity.
The following language and normality conventions follow
\cite[preprint \S2, pp.~4, 6]{Venema2005Preprint}.

\begin{hbdefinition}{def:p10-polyadic-normal-logic}{Finitary normal modal language and logic}
Fix a set \(\Lambda\) with an arity \(n_f\in\mathbb N\) for each
\(f\in\Lambda\), and a countably infinite set \(X\) of proposition
variables disjoint from the connective symbols.
The set \(\operatorname{Fm}_\Lambda(X)\) consists of the finite formulas
\[
 \varphi ::= p\ (p\in X)\mid\bot\mid\neg\varphi
       \mid(\varphi\wedge\psi)\mid f(\varphi_1,\ldots,\varphi_{n_f}).
\]
The other Boolean connectives are abbreviations.
For \(n_f=0\), the expression \(f()\) is a formula constant.
Substitution replaces variables simultaneously by formulas and leaves
every connective, including a nullary one, fixed.

For \(n_f>0\), fix a coordinate \(1\leq i\leq n_f\) and formulas in the other
coordinates, and write \(f_i(\chi)\) for the resulting formula with
\(\chi\) in coordinate \(i\).
A \emph{normal \(\Lambda\)-logic} is a set
\(L\subseteq\operatorname{Fm}_\Lambda(X)\) containing all classical
propositional tautology instances and, for every such context, the axioms
\[
 \neg f_i(\bot),\qquad
 f_i(\varphi\vee\psi)\to
                 \bigl(f_i(\varphi)\vee f_i(\psi)\bigr).
\]
It is closed under modus ponens and uniform substitution.
It also has \emph{theorem-level coordinate monotonicity}: whenever
\(\varphi\to\psi\) belongs to \(L\), so does
\(f_i(\varphi)\to f_i(\psi)\), for every choice of the other arguments.
Write \(\vdash_L\varphi\) for \(\varphi\in L\), and call \(L\)
consistent when \(\bot\notin L\).

The reverse additivity implication follows from monotonicity applied
to \(\varphi\to\varphi\vee\psi\) and
\(\psi\to\varphi\vee\psi\).
The source states monotonicity for simultaneous changes of all arguments;
the coordinate form is equivalent by finite iteration, using reflexive
implications for unchanged arguments.  Its normality axiom puts
\(\top\) in the other coordinates.  Monotonicity in those coordinates
gives the arbitrary-context version displayed here, and the converse is
an instance of that version.
No coordinate rule or normality axiom is imposed when \(n_f=0\).
In particular, normality does not assert \(f()\leftrightarrow\bot\).

A derivation from a premise set \(\Gamma\) is a finite sequence using
members of \(\Gamma\), members of \(L\), and modus ponens.
Write \(\Gamma\vdash_L\varphi\) when such a sequence ends in
\(\varphi\).  Substitution and coordinate monotonicity apply to
theorems of \(L\), not freely to premise-dependent lines.
For one unary symbol \(\Diamond\), putting
\(\Box\varphi=\neg\Diamond\neg\varphi\) gives the normal rules
and local-premise convention of
Definition~\ref{def:p10-normal-modal-algebra-logic}; the proof below
verifies this comparison.
\index{normal modal logic!finitary signature}\index{polyadic modal logic}
\end{hbdefinition}

\begin{hbdefinition}{def:p10-polyadic-modal-algebra}{Algebraic interpretation and the finitary Lindenbaum algebra}
In a Boolean algebra \(\mathbf A\) with \(\Lambda\)-operators,
a valuation \(v:X\to A\) extends recursively to
\(\llbracket\varphi\rrbracket_v^{\mathbf A}\), interpreting each
\(f\) by \(f^A\).  For a nullary formula its value is \(f^A()\).
Write \(\mathbf A\models\varphi\) when every valuation gives value
\(1_A\), and put
\[
 \operatorname{Alg}_\Lambda(L)
   =\{\mathbf A:\mathbf A\text{ is a }\Lambda\text{-BAO and }
                     \mathbf A\models\lambda\text{ for every }\lambda\in L\}.
\]
Trivial algebras are included.
Define
\[
 \varphi\equiv_L\psi
       \quad\Longleftrightarrow\quad\vdash_L\varphi\leftrightarrow\psi.
\]
The next theorem proves this is a congruence.  Its quotient
\(F_{\Lambda,L}=\operatorname{Fm}_\Lambda(X)/{\equiv_L}\) has the
usual Boolean operations on classes and
\[
 f^{F_{\Lambda,L}}([\varphi_1]_L,\ldots,[\varphi_n]_L)
      =[f(\varphi_1,\ldots,\varphi_n)]_L,
 \qquad f^{F_{\Lambda,L}}()=[f()]_L\quad(n_f=0).
\]
\end{hbdefinition}

\begin{hbtheorem}{thm:p10-polyadic-algebra-completeness}{proved=annex:p10.polyadic.1}{Algebraic completeness for a finitary normal logic}
For every normal \(\Lambda\)-logic \(L\), the relation
\(\equiv_L\) is a congruence and
\(F_{\Lambda,L}\in\operatorname{Alg}_\Lambda(L)\).
For every formula \(\varphi\),
\[
 \vdash_L\varphi
 \quad\Longleftrightarrow\quad [\varphi]_L=1
 \quad\Longleftrightarrow\quad F_{\Lambda,L}\models\varphi
 \quad\Longleftrightarrow\quad
                  \operatorname{Alg}_\Lambda(L)\models\varphi.
\]
A non-theorem is refuted by \(p\mapsto[p]_L\) in this quotient.
If \(L\) is inconsistent, it contains every formula, its quotient
is trivial, and \(\operatorname{Alg}_\Lambda(L)\) consists exactly
of the trivial \(\Lambda\)-BAOs.
\end{hbtheorem}

This is the full finitary normal case of
\cite[preprint Definition~3.11, p.~9, and Theorem~3.13,
pp.~10--11]{Venema2005Preprint}.
The proof uses representatives only for the finitely many variables
in the formula under consideration, even when \(\Lambda\) is infinite.

\begin{hbdefinition}{def:p10-polyadic-canonical-frame}{The descriptive canonical frame}
Assume BPIT and let \(L\) be consistent.
Write \(\operatorname{Ult}(A)\) for the ultrafilters of the Boolean
reduct of \(A\).
On the ultrafilters of
the Boolean reduct of \(F_{\Lambda,L}\), put
\[
\begin{gathered}
 W_L=\operatorname{Ult}(F_{\Lambda,L}),\qquad
 \widehat a=\{U\in W_L:a\in U\},\qquad
 \mathcal A_L=\{\widehat a:a\in F_{\Lambda,L}\},\\
 R_f(U,U_1,\ldots,U_n)
 \quad\Longleftrightarrow\quad
 \forall a_1\in U_1\cdots\forall a_n\in U_n\quad
                f^{F_{\Lambda,L}}(a_1,\ldots,a_n)\in U.
\end{gathered}
\]
For \(n_f=0\), this means \(R_f(U)\) iff \([f()]_L\in U\).
Write \(\mathcal D_L=(W_L,(R_f),\mathcal A_L)\).
Its canonical valuation is \(v_0(p)=\widehat{[p]_L}\).
These are the constructions of
\cite[preprint Definition~5.12, p.~19, and Definition~5.24,
p.~22]{Venema2005Preprint}.

In any general \(\Lambda\)-frame, valuations take their variable
values in the admissible algebra.  Boolean satisfaction has its usual
clauses, and
\[
 \mathcal G,v,w\models f(\varphi_1,\ldots,\varphi_n)
 \quad\Longleftrightarrow\quad
 \exists w_1,\ldots,w_n\quad
 \left(R_f(w,w_1,\ldots,w_n)\ \land\
       \bigwedge_{i=1}^n\mathcal G,v,w_i\models\varphi_i\right).
\]
For \(n=0\), the right side is simply \(R_f(w)\).
Validity quantifies over every admissible valuation and every point.
\end{hbdefinition}

\begin{hbtheorem}{thm:p10-polyadic-canonical-truth}{proved=annex:p10.polyadic.2}{Descriptive representation and canonical truth}
Assume BPIT.  For consistent normal \(L\), the general frame
\(\mathcal D_L\) is nonempty and descriptive.
It validates every member of \(L\) under every admissible valuation,
and
\[
 \mathcal D_L,v_0,U\models\varphi
       \quad\Longleftrightarrow\quad [\varphi]_L\in U.
\]
Consequently \(\mathcal D_L\models\varphi\) iff \(\varphi\in L\).
\end{hbtheorem}

\begin{hbtheorem}{thm:p10-polyadic-local-completeness}{proved=annex:p10.polyadic.3}{Strong local completeness over descriptive general frames}
Assume BPIT, and let \(\operatorname{DGen}_\Lambda(L)\) consist of
all nonempty descriptive general \(\Lambda\)-frames validating
every member of \(L\) under every admissible valuation.
For every set \(\Gamma\subseteq\operatorname{Fm}_\Lambda(X)\) and
formula \(\varphi\),
\[
 \Gamma\vdash_L\varphi
       \quad\Longleftrightarrow\quad
 \Gamma\models^\ell_{\operatorname{DGen}_\Lambda(L)}\varphi.
\]
The right side means that, at every point of every such frame and
under every admissible valuation, truth of all premises at that
point implies truth of the conclusion there.
For consistent \(L\), every failure of derivability is witnessed
in \(\mathcal D_L\) under \(v_0\).
For inconsistent \(L\), the frame class is empty and both sides
hold for every \(\Gamma,\varphi\).
The same equivalence holds if all nonempty general frames validating
\(L\) are used instead of only the descriptive ones.
\end{hbtheorem}

The algebraic statement of Venema's Theorem~3.13 is a theoremhood
result.  The strong local result here uses the finite deduction
property and a Boolean filter containing the premise values and
the complement of the conclusion.  The general-frame representation
is descriptive by Theorem~\ref{thm:p10-bao-descriptive-representation},
corresponding to \cite[preprint Theorem~5.28, p.~23]{Venema2005Preprint}.
No claim that the underlying ordinary frame validates \(L\) under
all subset valuations is needed.


Under topological semantics, \(\Box\) is interior and \(\Diamond\) is
closure. S4 is sound and complete for validity in all topological spaces.
The McKinsey--Tarski space-specific construction belongs to
\emph{The Algebra of Topology} (1944): every nonempty dense-in-itself
separable metrizable space, including the real line, has S4 as its logic.
The original separability assumption and its later removal are distinct
results \cite{McKinseyTarski1944Topology,BezhanishviliEtAl2017NewMT}.

\section{Many-valued and substructural algebras}
\label{sec:p10-many-substructural}

\begin{hbdefinition}{def:p10-mv-algebra}{MV-algebra}
An MV-algebra is a nonempty algebra \((A,\oplus,\neg,0)\) satisfying, for all
\(x,y,z\in A\),
\[
\begin{gathered}
x\oplus(y\oplus z)=(x\oplus y)\oplus z,\qquad
x\oplus y=y\oplus x,\qquad x\oplus0=x,\\
\neg\neg x=x,\qquad x\oplus\neg0=\neg0,\\
\neg(\neg x\oplus y)\oplus y
=\neg(\neg y\oplus x)\oplus x.
\end{gathered}
\]
Put \(1=\neg0\), define \(x\to y=\neg x\oplus y\), put
\(x\le y\) iff \(x\to y=1\), and
\(x\odot y=\neg(\neg x\oplus\neg y)\). The standard algebra on \([0,1]\) has
\(x\oplus y=\min(1,x+y)\) and \(\neg x=1-x\). The name is needed to identify
the exact algebraic semantics of infinite-valued \L{}ukasiewicz logic rather
than an arbitrary many-valued algebra. It stands as an equational class,
differs through the Chang identities and truncated addition, and is used in
standard completeness and the ordered-group correspondence.
\index{MV-algebra}
\end{hbdefinition}

Chang completeness says that validity in the standard MV-algebra agrees with
derivability in infinite-valued \L{}ukasiewicz logic \cite{Chang1958,Chang1959}.
Mundici's theorem gives an equivalence between MV-algebras and abelian lattice-
ordered groups with a strong unit \cite{Mundici1986}. BL- and MTL-algebras
replace the particular MV operations by residuated monoidal operations and
additional prelinearity/divisibility equations; their standard-completeness
claims depend on the chosen algebra class \cite{Hajek1998,
CignoliDOttavianoMundici2000}.

\begin{hbdefinition}{def:p10-residuated-lattice}{Residuated lattice}
A residuated lattice is a lattice-ordered monoid
\((A,\wedge,\vee,\otimes,1)\) with residuals satisfying
\[
 a\otimes b\le c
 \quad\Longleftrightarrow\quad
 b\le a\backslash c
 \quad\Longleftrightarrow\quad
 a\le c/b.
\]
The two residuals are necessary when exchange is absent; in a commutative
algebra they coincide. This stands as the algebraic semantics of Full
Lambek-style calculi, differs from a Heyting algebra by separating monoidal
product from lattice meet, and is used to track weakening, contraction, and
exchange as algebraic equations \cite{OnoKomori1985,
GalatosJipsenKowalskiOno2007}.
\index{residuated lattice}\index{FL-algebra}
\end{hbdefinition}

Cylindric algebras encode first-order Boolean operations together with
cylindrifications for existential quantification and diagonal elements for
equality. Representable cylindric algebras form a proper subclass of abstract
cylindric algebras in dimensions at least three, and their representation
problems carry substantial undecidability phenomena
\cite{HenkinMonkTarski1971}. Relation algebras similarly algebraize binary
relations with Boolean operations, converse, identity, and relational
composition; not every abstract relation algebra is representable
\cite{Tarski1935,Maddux2006}. These are pointers to distinct algebraic
programs, not claims that first-order logic reduces to a finite Boolean
signature without representability conditions.

Blok--Pigozzi algebraizability requires translations between formulas and
equations that make logical consequence correspond to equational consequence
and interderivability correspond to equality in the equivalent algebraic
semantics \cite{BlokPigozzi1989}. This criterion is stronger than merely having
some sound algebra-valued semantics.

\section{Categories, adjunctions, and internal logic}
\label{sec:p10-categorical-logic}

%% Elementary categorical definitions and complete proof-annex statements.

\subsection{Categories, functors, and natural transformations}
\label{sec:p10-categorical-toolkit}

\begin{hbconvention}{def:p10-categorical-size}{Size conventions for categorical constructions}
The fixed universe \(\mathcal U\) of small sets is transitive,
has \(\mathbb N\in\mathcal U\), and is closed under pairing, powersets,
and unions of small-indexed families:
\[
 \begin{gathered}
 a\in\mathcal U\ \Longrightarrow\ a\subseteq\mathcal U,\qquad
 a,b\in\mathcal U\ \Longrightarrow\ \{a,b\}\in\mathcal U,\\
 a\in\mathcal U\ \Longrightarrow\ \mathcal P(a)\in\mathcal U,\\
 I\in\mathcal U,\quad A_i\in\mathcal U\ (i\in I)
       \ \Longrightarrow\ \bigcup_{i\in I}A_i\in\mathcal U .
 \end{gathered}
\]
These closures also keep subsets, function sets, quotients, and
small-indexed products within \(\mathcal U\).
The ambient collection of all sets is a permitted default:
in that reading small means set-sized, with the ordinary
set/class distinction, and the closures are ordinary set-theoretic
properties.  If a set-sized \(\mathcal U\) is used for relativization,
it is required to be a Grothendieck universe with the stated
closures.  The existence of such a set universe, or a larger
containing one, is an additional size hypothesis when invoked;
no universe-existence assertion is part of the categorical proofs.

A category is small when its object and arrow sets are small.
A diagram is small when its indexing category is small; its
values may be objects of a large category.
Families indexed by a large category are handled in the containing
ambient framework.  Statements about class-sized natural
transformations use their displayed component families and
parametrizations; they do not assert a set of all such families.
A small set of parameters need not make every entire component
family a small object.
\end{hbconvention}

\begin{hbdefinition}{def:p10-category}{Category, functor, and natural transformation}
A category \(\mathcal C\) consists of objects, arrows with specified
domains and codomains, an identity \(1_A:A\to A\) for each object,
and a composite \(g\circ f:A\to C\) whenever \(f:A\to B\) and
\(g:B\to C\).  For composable arrows,
\[
 h\circ(g\circ f)=(h\circ g)\circ f,\qquad
 1_B\circ f=f=f\circ1_A.
\]
Write \(\mathcal C(A,B)\) for the arrows \(A\to B\).
The category is locally small if each such collection is a small
set, and small if its objects and all its arrows form small sets.
An isomorphism \(f:A\to B\) has an arrow \(g:B\to A\) with
\(g f=1_A\) and \(f g=1_B\).

A functor \(F:\mathcal C\to\mathcal D\) assigns objects and arrows
with the indicated domains and codomains, and satisfies
\[
 F(1_A)=1_{F A},\qquad F(g\circ f)=F(g)\circ F(f).
\]
The identity functor is \(1_{\mathcal C}\).
The opposite category \(\mathcal C^{\mathrm{op}}\) reverses every
arrow and the order of composition.  A contravariant set-valued
functor is an ordinary functor
\(F:\mathcal C^{\mathrm{op}}\to\mathsf{Set}\): thus \(f:A\to B\)
induces \(F(f):F(B)\to F(A)\).

For functors \(F,G:\mathcal C\to\mathcal D\), a natural
transformation \(\alpha:F\Rightarrow G\) consists of arrows
\(\alpha_A:F A\to G A\) satisfying
\[
                  G(f)\circ\alpha_A=\alpha_B\circ F(f)
                  \quad(f:A\to B).
\]
It is a natural isomorphism when every component is an
isomorphism.  The naturality equation then also holds for the
componentwise inverses.  Transformations compose componentwise;
identity components give the identity transformation.
\index{category}\index{functor}\index{natural transformation}
\index{opposite category}\index{locally small category}
\end{hbdefinition}

The category \(\mathsf{Set}\) has small sets as objects and functions
as arrows.  A small partially ordered set gives a small category with one
arrow \(a\to b\) when \(a\leq b\), and none otherwise.  Its identity
and composition laws are reflexivity and transitivity; antisymmetry
makes isomorphic objects equal.

\begin{hbdefinition}{def:p10-category-equivalence}{Fullness, faithfulness, and equivalence}
A functor \(F:\mathcal C\to\mathcal D\) is full when each map
\(\mathcal C(A,B)\to\mathcal D(F A,F B)\) is surjective, and faithful
when each is injective.  It is essentially surjective when every
object of \(\mathcal D\) is isomorphic to some \(F A\).
An equivalence consists of functors \(F:\mathcal C\to\mathcal D\)
and \(G:\mathcal D\to\mathcal C\), together with natural
isomorphisms \(1_{\mathcal C}\Rightarrow GF\) and
\(FG\Rightarrow1_{\mathcal D}\).
The functor \(G\) is a quasi-inverse.  The particular equivalence
constructed below supplies these functors and isomorphisms
explicitly.
\index{equivalence!of categories}\index{full functor}
\index{faithful functor}\index{essentially surjective functor}
\end{hbdefinition}

These definitions use the elementary conventions of
\cite[Chapter~1, especially Definitions~1.1.1, 1.3.1, 1.4.1,
and 1.4.3]{Riehl2026ContextOnline}.  The classical references
\cite{MacLane1998,Awodey2010} give separate book presentations.
The Riehl locators in this section refer to the author-hosted
second-edition text retrieved on 6 September 2026.

\subsection{Limits and colimits}

\begin{hbdefinition}{def:p10-limit}{Diagrams, cones, limits, and colimits}
Let \(J\) be a small category and \(D:J\to\mathcal C\) a functor,
called a diagram of shape \(J\).
A cone with vertex \(X\) is a family \(p_j:X\to D j\) satisfying
\[
                         D(u)\circ p_j=p_k
                         \quad(u:j\to k\text{ in }J).
\]
A limit is a cone \((L,(\ell_j)_j)\) such that every cone
\((X,(p_j)_j)\) has exactly one mediating arrow \(a:X\to L\)
with \(\ell_j a=p_j\) for all \(j\).
Equivalently, it is a terminal object in the category of cones,
whose arrows preserve every cone leg.

A cocone with vertex \(X\) has arrows \(q_j:D j\to X\) satisfying
\(q_k D(u)=q_j\).  A colimit is a cocone \((K,(k_j)_j)\) such that
every cocone \(q_j:D j\to X\) has exactly one arrow \(b:K\to X\)
with \(b k_j=q_j\).  It is an initial object in the category of
cocones.  The empty diagram has no legs or commutativity conditions.
Its limits are terminal objects and its colimits are initial
objects.
\index{limit!categorical}\index{colimit}\index{cone!categorical}
\index{terminal object}\index{initial object}
\end{hbdefinition}

\begin{hbdefinition}{def:p10-products-equalizers}{Products, equalizers, and their duals}
For a small family \((A_i)_{i\in I}\), a product is an object
\(\prod_i A_i\) with projections \(\pi_i\) for which every family
\(f_i:X\to A_i\) has a unique arrow
\(\langle f_i\rangle_i:X\to\prod_i A_i\) satisfying
\(\pi_i\langle f_i\rangle_i=f_i\).
It is the limit of the discrete diagram on \(I\).
A coproduct \(\coprod_i A_i\) has injections \(\iota_i\) and the
dual unique-arrow property for families \(A_i\to X\).
The empty product is terminal; the empty coproduct is initial.

For parallel arrows \(f,g:A\to B\), an equalizer is an arrow
\(e:E\to A\) with \(f e=g e\), universal among arrows with this
property: every \(a:X\to A\) with \(f a=g a\) factors uniquely
as \(a=e h\).
A coequalizer is \(q:B\to Q\) with \(q f=q g\), and every
\(b:B\to X\) with \(b f=b g\) factors uniquely as \(b=h q\).

A pullback of \(f:A\to C\) and \(g:B\to C\) consists of
\(p:P\to A\), \(q:P\to B\), with \(f p=g q\), such that every
commuting pair \(a:X\to A,b:X\to B\) factors uniquely through
\((p,q)\).  A pushout is the dual construction for
\(A\leftarrow C\to B\).  These are limits and colimits of the
indicated three-object diagrams.

An arrow \(m:A\to B\) is a monomorphism when
\(m r=m s\) implies \(r=s\) for every pair \(r,s:X\to A\).
An epimorphism satisfies the dual right-cancellation property.
\index{product!categorical}\index{coproduct}\index{equalizer}
\index{coequalizer}\index{pullback}\index{pushout}
\index{monomorphism}\index{epimorphism}
\end{hbdefinition}

For chosen binary products, write
\[
 f\times g=\langle f\pi_1,g\pi_2\rangle:
           A\times B\longrightarrow A'\times B'
 \quad(f:A\to A',\ g:B\to B').
\]
The projections distinguish arrows into a product.  Thus equations
between such arrows can be established by composing with both
projections.

\begin{hbtheorem}{thm:p10-limits-unique}{proved=annex:p10.cat.1}{Uniqueness of universal constructions}
Two limiting cones for the same diagram have a unique isomorphism
that preserves every cone leg.  Two colimiting cocones have the
dual unique isomorphism.
Every equalizer arrow is monic, and every coequalizer arrow is epic.
For products, pairing commutes with precomposition and product
arrows preserve identities and composition:
\[
 \langle f,g\rangle a=\langle f a,g a\rangle,\qquad
 1_A\times1_B=1_{A\times B},\qquad
 (f'\times g')(f\times g)=(f'f)\times(g'g).
\]
The uniqueness assertion concerns isomorphisms compatible with
the specified legs; it imposes no uniqueness on arbitrary
isomorphisms of the underlying objects.
\end{hbtheorem}

\begin{hbtheorem}{thm:p10-limits-products-equalizers}{proved=annex:p10.cat.2}{Constructing limits and colimits}
A locally small category with all small products and equalizers
has all small limits.  A locally small category with all small
coproducts and coequalizers has all small colimits.
Terminal objects, binary products, and equalizers suffice for
finite limits; initial objects, binary coproducts, and
coequalizers suffice for finite colimits.

For a small diagram \(D:J\to\mathsf{Set}\), a limit is the set
\[
 \left\{(x_j)_j\in\prod_{j\in\operatorname{Ob}J}D j:
       D(u)(x_j)=x_k\text{ for every }u:j\to k\right\}.
\]
A colimit is the disjoint union of the \(D j\)'s modulo the least
equivalence relation identifying
\((j,x)\) with \((k,D(u)(x))\) for every \(u:j\to k\).
The empty-diagram cases give a singleton and the empty set.
\end{hbtheorem}

Cones and their universal properties are presented in
\cite[Section~3.1, Definitions~3.1.1--3.1.6 and
Proposition~3.1.7]{Riehl2026ContextOnline}.
The general product/equalizer and coproduct/coequalizer
constructions are Theorem~3.5.11, pp.~112--113, of that text.

\subsection{Representable functors and Yoneda}

For an object \(A\) of a locally small category, the representable
functor \(\mathcal C(-,A):\mathcal C^{\mathrm{op}}\to\mathsf{Set}\)
sends \(X\) to \(\mathcal C(X,A)\), and \(u:Y\to X\) to
precomposition \(f\mapsto f u\).
Write \(\operatorname{Nat}(H,F)\) for natural transformations
between functors with the same domain and codomain.

\begin{hbtheorem}{thm:p10-yoneda}{proved=annex:p10.cat.3}{Yoneda lemma}
For a locally small category \(\mathcal C\), a functor
\(F:\mathcal C^{\mathrm{op}}\to\mathsf{Set}\), and an object \(A\),
evaluation at \(1_A\) gives a bijection
\[
 \operatorname{Nat}(\mathcal C(-,A),F)\longrightarrow F(A),
 \qquad \alpha\longmapsto\alpha_A(1_A).
\]
The inverse sends \(x\in F(A)\) to the transformation
\(\alpha^x_X(f)=F(f)(x)\).
This bijection is natural in \(A\) and \(F\).
It gives a canonical parametrization of the displayed family of
transformations by the small set \(F(A)\).
When \(\mathcal C\) is large, this does not assert that each
transformation's entire component family belongs to the fixed
universe.

Every natural transformation
\(\mathcal C(-,A)\Rightarrow\mathcal C(-,B)\) has the form
\(f\mapsto h f\) for a unique \(h:A\to B\).
Thus the Yoneda assignment \(A\mapsto\mathcal C(-,A)\) is full
and faithful on these hom-sets, and isomorphic representable
functors have isomorphic representing objects.
\index{Yoneda lemma}\index{representable functor}
\end{hbtheorem}

This is the contravariant form of
\cite[Theorem~2.2.4, pp.~61--63; Exercise~2.2.i,
p.~66]{Riehl2026ContextOnline}, with its proof supplied here.

\subsection{Adjunctions}

\begin{hbdefinition}{def:p10-adjunction}{Adjunction by natural hom-set bijections}
Let \(F:\mathcal C\to\mathcal D\) and \(G:\mathcal D\to\mathcal C\)
be functors between locally small categories.
An adjunction \(F\dashv G\) is a family of bijections
\[
 \Phi_{A,B}:\mathcal D(F A,B)\longrightarrow\mathcal C(A,G B)
\]
natural in both variables.  Explicitly, for
\(u:A'\to A\), \(f:F A\to B\), and \(v:B\to B'\),
\[
 \Phi_{A',B'}(v f F(u))=G(v)\,\Phi_{A,B}(f)\,u.
\]
The functor \(F\) is the left adjoint and \(G\) the right adjoint.
The image \(\Phi(f)\) and its inverse image are called
transposes with respect to this adjunction.
\index{adjunction}\index{left adjoint}\index{right adjoint}
\end{hbdefinition}

\begin{hbtheorem}{thm:p10-adjunction-data}{proved=annex:p10.cat.4}{Units, counits, and triangle identities}
Adjunction bijections \(\Phi\) determine natural transformations
\[
 \eta:1_{\mathcal C}\Rightarrow GF,\qquad
 \varepsilon:FG\Rightarrow1_{\mathcal D}
\]
by
\[
 \eta_A=\Phi_{A,F A}(1_{F A}),\qquad
 \varepsilon_B=\Phi^{-1}_{G B,B}(1_{G B}).
\]
They satisfy the triangle identities
\[
 \varepsilon_{F A}\circ F(\eta_A)=1_{F A},\qquad
 G(\varepsilon_B)\circ\eta_{G B}=1_{G B}.
\]
Conversely, natural transformations \(\eta,\varepsilon\) of these
types satisfying these identities determine an adjunction by the
mutually inverse maps
\[
 f\longmapsto G(f)\circ\eta_A,\qquad
 g\longmapsto\varepsilon_B\circ F(g).
\]
The constructions recover the original bijections and the original
unit and counit.
\index{unit!of an adjunction}\index{counit}
\index{triangle identities}
\end{hbtheorem}

\begin{hbtheorem}{thm:p10-adjoint-preservation}{proved=annex:p10.cat.5}{Adjoints preserve universal constructions}
If \(F\dashv G\), then \(F\) sends every existing small colimiting
cocone in \(\mathcal C\) to a colimiting cocone in \(\mathcal D\),
and \(G\) sends every existing small limiting cone in
\(\mathcal D\) to a limiting cone in \(\mathcal C\).
Consequently left adjoints preserve initial objects, coproducts,
coequalizers, and pushouts whenever these exist; right adjoints
preserve terminal objects, products, equalizers, and pullbacks.
The conclusion identifies the image cone or cocone by its
universal property, independently of any different chosen
representative of the same limit or colimit.
\end{hbtheorem}

The hom-set and unit/counit presentations are
\cite[Definition~4.1.1, p.~132; Definition~4.2.5 and
Theorem~4.2.7, pp.~140--142]{Riehl2026ContextOnline}.
Preservation of existing limits and colimits is Theorem~4.5.2,
pp.~158--159, of that text.

\subsection{Cartesian closed categories}

\begin{hbdefinition}{def:p10-ccc}{Cartesian closed category}
A cartesian closed category is a locally small category with a
terminal object, binary products, and exponentials.
When interpreting syntax, specify choices of these constructions.
For objects \(A,B\), an exponential consists of an object \(B^A\)
and an evaluation arrow
\[
                    \operatorname{ev}_{A,B}:B^A\times A\to B
\]
such that every \(f:X\times A\to B\) has a unique transpose
\(\Lambda(f):X\to B^A\) satisfying
\[
 \operatorname{ev}_{A,B}\circ(\Lambda(f)\times1_A)=f.
\]
Equivalently, the two maps
\[
 \mathcal C(X\times A,B)\ \longleftrightarrow\
 \mathcal C(X,B^A),\qquad
 f\longmapsto\Lambda(f),\quad
 k\longmapsto\operatorname{ev}_{A,B}\circ(k\times1_A)
\]
are inverse.  The inverse laws are the categorical \(\beta\)
and \(\eta\) equations.
In particular,
\(\Lambda(\operatorname{ev}_{A,B}\circ(k\times1_A))=k\).
\index{cartesian closed category}\index{exponential!categorical}
\index{evaluation!of an exponential}
\index{Curry--Howard--Lambek correspondence}
\end{hbdefinition}

\begin{hbtheorem}{thm:p10-exponential-adjunction}{proved=annex:p10.cat.6}{Exponential adjunction and examples}
In a cartesian closed category, transposition is natural in its
parameter:
\[
 \Lambda(f\circ(s\times1_A))=\Lambda(f)\circ s
       \quad(f:X\times A\to B,\ s:Y\to X).
\]
For fixed \(A\), the assignments \(X\mapsto X\times A\) and
\(B\mapsto B^A\) extend to functors with
\[
                         (-)\times A\dashv(-)^A.
\]
The category \(\mathsf{Set}\) is cartesian closed, with a singleton
as terminal object, Cartesian products as products, and the set
of functions \(A\to B\) as \(B^A\).
A Heyting algebra, regarded as a category under its order, is
cartesian closed with terminal object \(1\), product \(a\wedge b\),
and exponential \(b^a=(a\Rightarrow b)\).
\end{hbtheorem}

Function-type semantics uses precisely these evaluation and
transposition operations
\cite[Section~3.3, pp.~22--24]{Pitts1995Draft}.
The next construction applies them to the registered arrow-only
typed lambda calculus; the following internal-language construction
gives an explicit equivalence for a named small category.

%% The context category of the registered arrow-only STLC.

\subsection{Simply typed terms as a cartesian closed category}
\label{sec:p10-cartesian-lambda}

Use exactly the types and terms of the registered calculus
\hyperref[sheet:stlc]{\(\mathsf{STLC}\)}: types are generated from
\(\mathsf{Base}\) by \(\alpha\mathbin{\to_{\mathrm{ty}}}\beta\),
and terms are variables, typed abstractions, and applications.
The products and terminal object below are constructions on
contexts.  A tuple of terms is a metalinguistic list.
No product or unit constructor is added to this term language.

\begin{hbdefinition}{def:p10-stlc-beta-eta}{Typed \(\beta\eta\) equality and simultaneous substitution}
Write \(\Theta\vdash_{\beta\eta}M=N:\alpha\) for the least
equivalence relation on terms of type \(\alpha\) in the finite
typing context \(\Theta\) that contains consistent bound
\(\alpha\)-renaming, preserves application and typed abstraction,
and contains the following typed equation schemas:
\[
 (\lambda x{:}\alpha.M)N = M[N/x]:\beta,
 \qquad
 \lambda x{:}\alpha.Mx=M:
                  \alpha\mathbin{\to_{\mathrm{ty}}}\beta.
\]
In the first schema,
\(\Theta,x{:}\alpha\derives_{\mathsf{STLC}}M:\beta\) and
\(\Theta\derives_{\mathsf{STLC}}N:\alpha\).
In the second,
\(\Theta\derives_{\mathsf{STLC}}M:
 \alpha\mathbin{\to_{\mathrm{ty}}}\beta\), and \(x\) is fresh
for \(\Theta\) and \(M\).
Abstraction congruence carries an equation in
\(\Theta,x{:}\alpha\) to an equation in \(\Theta\).
These schemas define an equational metatheory of the registered
syntax; the reduction relation on the calculus sheet is
\(\beta\)-reduction.

A typed simultaneous substitution \(\sigma:\Theta\to\Psi\),
where \(\Psi=(y_1{:}\beta_1,\ldots,y_n{:}\beta_n)\), is a tuple
of terms \(\Theta\derives_{\mathsf{STLC}}\sigma_j:\beta_j\).
For a term \(M\) in \(\Psi\), the notation \(M[\sigma]\)
replaces its free \(y_j\)'s simultaneously by \(\sigma_j\).
Inserted terms are not substituted into again during this
operation.  The operation is capture-avoiding, following
Definition~\ref{def:p07-alpha-substitution}, with binder types
retained.  Under an abstraction, first rename its bound variable
to be fresh for the finitely many terms and contexts involved.
\end{hbdefinition}

\begin{hbtheorem}{thm:p10-stlc-equational-substitution}{proved=annex:p10.lambda.1}{Typing and equality under substitution}
If \(\Psi\derives_{\mathsf{STLC}}M:\alpha\) and
\(\sigma:\Theta\to\Psi\), then
\(\Theta\derives_{\mathsf{STLC}}M[\sigma]:\alpha\).
Identity substitutions act as identities, and for
\(\tau:\Omega\to\Theta\),
\[
 (M[\sigma])[\tau]=_{\alpha}M[\sigma[\tau]],
 \qquad
 \sigma[\tau]=(\sigma_1[\tau],\ldots,\sigma_n[\tau]).
\]
Here \(=_{\alpha}\) denotes bound \(\alpha\)-equivalence.
Both substitution into a typed \(\beta\eta\) equation and
componentwise replacement of a substitution by a
\(\beta\eta\)-equal substitution preserve typed \(\beta\eta\)
equality.  Weakening by unused declarations and consistent
renaming of declared variables preserve typing and equality.
\end{hbtheorem}

\begin{hbdefinition}{def:p10-stlc-context-category}{The category of typed contexts}
Let \(\mathcal S_{\beta\eta}\) have as objects finite type lists
\(\Gamma=(\alpha_1,\ldots,\alpha_m)\), including the empty list
\(()\).  Present \(\Gamma\) by a typing context
\(\Theta_\Gamma=(x_1{:}\alpha_1,\ldots,x_m{:}\alpha_m)\)
with distinct names.  A fixed enumeration of the variable supply
can choose these names once and for all.
An arrow from \(\Gamma\) to
\(\Delta=(\beta_1,\ldots,\beta_n)\) is a componentwise
\(\beta\eta\)-equivalence class of tuples \((M_1,\ldots,M_n)\)
with \(\Theta_\Gamma\derives_{\mathsf{STLC}}M_j:\beta_j\).
Changing the declared names consistently by position changes
only the presentation; this renaming of free context variables
is distinct from bound \(\alpha\)-renaming.

The identity of \(\Gamma\) is the tuple \((x_1,\ldots,x_m)\).
If \(M:\Gamma\to\Delta\) and \(N:\Delta\to\Xi\), define
\[
                 N\circ M=(N_j[M_1/y_1,\ldots,M_n/y_n])_j,
\]
where the \(y_j\)'s name the slots of \(\Delta\).
Fresh names may be used for intermediate presentations.
\end{hbdefinition}

\begin{hbtheorem}{thm:p10-stlc-context-ccc}{proved=annex:p10.lambda.2}{Cartesian closed structure on contexts}
The preceding data form a category.
Its terminal object is \(()\); its product
\(\Gamma\times\Delta\) is the concatenation \(\Gamma\cdot\Delta\),
with block projections and concatenation of tuples as pairing.

For \(\Gamma=(\alpha_1,\ldots,\alpha_m)\), define the
right-associated type
\[
 c_\Gamma(\beta)=
 \alpha_1\mathbin{\to_{\mathrm{ty}}}\cdots
 \mathbin{\to_{\mathrm{ty}}}\alpha_m
 \mathbin{\to_{\mathrm{ty}}}\beta,
 \qquad c_{()}(\beta)=\beta.
\]
If \(\Delta=(\beta_1,\ldots,\beta_n)\), an exponential is
\[
       \Delta^\Gamma=(c_\Gamma(\beta_1),\ldots,c_\Gamma(\beta_n)).
\]
Its evaluation
\(\Delta^\Gamma\cdot\Gamma\to\Delta\) is represented by
\((f_jx_1\cdots x_m)_{j=1}^n\).
The transpose of \(T:\Xi\cdot\Gamma\to\Delta\) is represented by
\[
       \bigl(\lambda x_1{:}\alpha_1\cdots
                       \lambda x_m{:}\alpha_m.T_j\bigr)_{j=1}^n.
\]
For \(m=0\), application and abstraction over the empty list
do nothing.  For \(n=0\), these tuples are empty.
Evaluation and transposition give inverse bijections
\[
 \mathcal S_{\beta\eta}(\Xi\cdot\Gamma,\Delta)
       \cong
 \mathcal S_{\beta\eta}(\Xi,\Delta^\Gamma),
\]
natural in \(\Xi\).  Hence \(\mathcal S_{\beta\eta}\) is
cartesian closed in the sense of Definition~\ref{def:p10-ccc}.
\end{hbtheorem}

\begin{hbdefinition}{def:p10-stlc-ccc-interpretation}{Interpretation of arrow-only STLC in a CCC}
Fix a cartesian closed category \(\mathcal C\) with chosen
terminal object \(1\), binary products, and exponentials, and
assign an object to every base type.  Interpret types by
\[
 \llbracket\alpha\mathbin{\to_{\mathrm{ty}}}\beta\rrbracket
       =\llbracket\beta\rrbracket^{\llbracket\alpha\rrbracket}.
\]
Interpret context lists recursively by
\(\llbracket()\rrbracket=1\) and
\(\llbracket\Gamma\cdot(\alpha)\rrbracket
 =\llbracket\Gamma\rrbracket\times\llbracket\alpha\rrbracket\).
Thus a singleton context is interpreted as
\(1\times\llbracket\alpha\rrbracket\); the construction does
not require this object to be literally
\(\llbracket\alpha\rrbracket\).
A judgment \(\Theta_\Gamma\derives_{\mathsf{STLC}}M:\alpha\)
denotes an arrow
\(\llbracket M\rrbracket:
 \llbracket\Gamma\rrbracket\to\llbracket\alpha\rrbracket\).
Variables denote the corresponding product projections,
application denotes
\[
 \llbracket MN\rrbracket
      =\operatorname{ev}\circ
           \langle\llbracket M\rrbracket,\llbracket N\rrbracket\rangle,
\]
and abstraction denotes
\(\llbracket\lambda x{:}\alpha.M\rrbracket
  =\Lambda(\llbracket M\rrbracket)\), where \(\Lambda\) is
exponential transposition.
The interpretation of a typed substitution is the unique
arrow into the interpreted target context with the specified
term components.
\end{hbdefinition}

\begin{hbtheorem}{thm:p10-stlc-ccc-completeness}{proved=annex:p10.lambda.3}{Substitution, soundness, and completeness for CCC equality}
Interpretation respects bound \(\alpha\)-renaming and satisfies
\[
             \llbracket M[\sigma]\rrbracket
                 =\llbracket M\rrbracket\circ\llbracket\sigma\rrbracket.
\]
For terms \(M,N\) of the same type \(\alpha\) in the same
context \(\Theta\), the following are equivalent:
\[
 \Theta\vdash_{\beta\eta}M=N:\alpha
 \quad\Longleftrightarrow\quad
 \llbracket M\rrbracket=\llbracket N\rrbracket
 \text{ in every CCC interpretation of the base types.}
\]
In particular, the context category itself detects every
failure of typed \(\beta\eta\) equality.
\end{hbtheorem}

Pitts derives the interpretation of function types and the
categorical \(\beta\) and \(\eta\) equations in
\cite[\S3.3, pp.~22--24, equations (3.19)--(3.22)]{Pitts1995Draft}.
His classifying-context construction uses tuples and
simultaneous substitution, with finite products given by
concatenation
\cite[\S4.2, pp.~30--33, Proposition~4.2.1]{Pitts1995Draft}.
The curried-list exponential is also given in
\cite[\S4.4, p.~38]{Pitts1995Draft}.
Its use for this handbook's arrow-only syntax is proved in the annex.
All page locators here refer to the draft dated 18 May 1995.
\index{lambda calculus!cartesian closed interpretation}
\index{cartesian closed category!typed context category}

%% A small CCC recovered from its named full equational diagram.

\subsection{Recovering a category from its internal language}
\label{sec:p10-internal-cartesian-language}

Fix a small cartesian closed category \(\mathcal C\) with given terminal
object \(1\), binary products, and exponentials \(B^A\) with evaluation
\(\operatorname{ev}_{A,B}:B^A\times A\to B\).
Small means that its objects and arrows form sets in the chosen universe.
The chosen universal constructions are part of the data.
We construct a particular equational language for this category.

\begin{hbdefinition}{def:p10-internal-ccc-language}{The named internal language of a small CCC}
For each object \(A\) introduce a base type \(\underline A\), using
distinct formal symbols for distinct objects.  Form all other types by
the arrow constructor of Definition~\ref{def:p10-stlc-beta-eta}.
Use the same variables, application, and typed abstraction, and add
the following closed constants:
\[
\begin{array}{ll}
 c_f:\underline A\mathbin{\to_{\mathrm{ty}}}\underline B
       &(f:A\to B),\\
 u:\underline1,&\\
 \mathsf{pair}_{A,B}:\underline A\mathbin{\to_{\mathrm{ty}}}
        \underline B\mathbin{\to_{\mathrm{ty}}}\underline{A\times B},&\\
 e_{A,B}:\underline{B^A}\mathbin{\to_{\mathrm{ty}}}
        (\underline A\mathbin{\to_{\mathrm{ty}}}\underline B),&\\
 l_{A,B}:(\underline A\mathbin{\to_{\mathrm{ty}}}\underline B)
        \mathbin{\to_{\mathrm{ty}}}\underline{B^A}.&
\end{array}
\]
Each constant has its declared type in any context; substitution leaves
it fixed.  No product or unit type constructor is added.
The structural constants compare named objects with contexts and arrow
types; \(\underline{B^A}\) and
\(\underline A\mathbin{\to_{\mathrm{ty}}}\underline B\) are distinct
type expressions.

Interpret \(\underline A\) by \(A\), and write \(A_T=\llbracket T\rrbracket\)
for the recursively interpreted object of a type \(T\).
As in Definition~\ref{def:p10-stlc-ccc-interpretation}, put
\[
 P_{()}=1,\qquad P_{\Gamma\cdot(T)}=P_\Gamma\times A_T.
\]
In particular, \(P_{(\underline A)}=1\times A\), not \(A\).
Terms and substitutions use the categorical interpretation of that
definition, extended by the declared constants.

To specify a closed constant of function type \(S\mathbin{\to_{\mathrm{ty}}}T\)
by an arrow \(a:A_S\to A_T\), use the closed denotation
\(\Lambda(a\circ\pi_2):1\to A_T^{A_S}\), where
\(\pi_2:1\times A_S\to A_S\).
In this sense \(c_f\) denotes \(f\), while \(e_{A,B}\) and
\(l_{A,B}\) denote the identity of \(B^A\).
The constant \(u\) denotes \(1_1:1\to1\).
The denotation of \(\mathsf{pair}_{A,B}\) is obtained by first
transposing in \(B\), then in \(A\), the arrow
\[
 ((1\times A)\times B)\longrightarrow A\times B,
 \qquad \langle\pi_2\circ\pi_1,\pi_2\rangle.
\]
Thus its applications denote categorical pairing.
The signature is a set because \(\mathcal C\) is small.
\index{internal language!cartesian closed category}
\end{hbdefinition}

\begin{hbdefinition}{def:p10-internal-ccc-diagram}{The full true equational diagram}
Let \(\operatorname{Diag}(\mathcal C)\) contain every well-typed equation
\[
 \Theta_\Gamma\vdash M=N:T
\]
whose two denotations are equal arrows \(P_\Gamma\to A_T\).
This is a set of equations-in-context.
Let \(=_{\mathcal C}\) be the least typed equational theory containing
this diagram and the \(\beta\eta\) equations, closed under equivalence,
application and abstraction congruence, weakening, renaming, and typed
capture-avoiding simultaneous substitution.
Define \(\mathcal S_{\mathcal C}\) by the context and tuple construction
of Definition~\ref{def:p10-stlc-context-category}, now using these
constants and componentwise \(=_{\mathcal C}\) equality.

Some equations in the diagram are listed below. Each line is read in
its own typing context: \(f,g\) are composable, \(!_A:A\to1\), the
product projections have domain \(A\times B\), and their variables
have the corresponding named types.
\[
\begin{gathered}
 c_{1_A}x=x,\qquad c_g(c_fx)=c_{g\circ f}x,\qquad c_{!_A}x=u,\\
 c_{\pi_1}(\mathsf{pair}_{A,B}xy)=x,\qquad
 c_{\pi_2}(\mathsf{pair}_{A,B}xy)=y,\\
 \mathsf{pair}_{A,B}(c_{\pi_1}z)(c_{\pi_2}z)=z
                     \quad(z:\underline{A\times B}),\\
 e_{A,B}(l_{A,B}h)=h,\qquad l_{A,B}(e_{A,B}z)=z,\\
 (e_{A,B}z)a
   =c_{\operatorname{ev}_{A,B}}(\mathsf{pair}_{B^A,A}za).
\end{gathered}
\]
In the exponential lines,
\(h:\underline A\mathbin{\to_{\mathrm{ty}}}\underline B\),
\(z:\underline{B^A}\), and \(a:\underline A\).
Terminal uniqueness is \(z=u\) in context \(z:\underline1\).
For \(f:D\times A\to B\), the diagram also records transposition:
\[
 (e_{A,B}(c_{\Lambda f}d))a
       =c_f(\mathsf{pair}_{D,A}da).
\]
These displayed equations describe part of the diagram; the theory
contains all true equations-in-context, not just this subfamily.
No effective enumeration or finite axiomatization of that theory is
asserted.
\index{equational diagram!internal language}
\end{hbdefinition}

The naming and structural equations follow the internal-language method
of \cite[\S2.3, pp.~10--11, Proposition~2.3.1]{Pitts1995Draft}.
The full semantic diagram is the construction in \S4.3, pp.~35--36;
the context category and its type structure are in \S\S4.2--4.4,
pp.~30--38.  These locators refer to the draft dated 18 May 1995.

\begin{hbtheorem}{thm:p10-internal-ccc-comparisons}{proved=annex:p10.internal.1}{Named type and context comparisons}
For every type \(T\), there are closed terms
\[
 \kappa_T:T\mathbin{\to_{\mathrm{ty}}}\underline{A_T},
 \qquad
 \iota_T:\underline{A_T}\mathbin{\to_{\mathrm{ty}}}T
\]
such that applying either comparison to a well-typed argument preserves
that argument's denotation in every context.
Their two composites equal identity terms in \(=_{\mathcal C}\).
For every context \(\Gamma\), there are inverse context arrows
\[
 \operatorname{pack}_\Gamma:\Gamma\longrightarrow(\underline{P_\Gamma}),
 \qquad
 \operatorname{unpack}_\Gamma:(\underline{P_\Gamma})\longrightarrow\Gamma.
\]
The component packing term denotes \(1_{P_\Gamma}\), while the packing
\emph{context arrow} denotes
\(\langle!,1\rangle:P_\Gamma\to1\times P_\Gamma\).
The unpacking context arrow denotes
\(\pi_2:1\times P_\Gamma\to P_\Gamma\).
These comparisons include the empty context.
\end{hbtheorem}

\begin{hbtheorem}{thm:p10-internal-ccc-equivalence}{proved=annex:p10.internal.2}{Evaluation recovers the given category}
Evaluation defines a full and faithful functor
\[
 E:\mathcal S_{\mathcal C}\longrightarrow\mathcal C,
 \qquad E(\Gamma)=P_\Gamma.
\]
It is essentially surjective.
An explicit quasi-inverse \(J\) sends \(A\) to the singleton context
\((\underline A)\) and \(f:A\to B\) to the term tuple \((c_fx)\).
The unit is \(\eta_\Gamma=\operatorname{pack}_\Gamma\), and the counit is
\[
 \varepsilon_A=\pi_2:EJ(A)=1\times A\longrightarrow A.
\]
They are natural isomorphisms and satisfy both triangle identities.
\end{hbtheorem}

\begin{hbtheorem}{thm:p10-internal-ccc-structure}{proved=annex:p10.internal.3}{Preservation of cartesian closed structure}
The category \(\mathcal S_{\mathcal C}\) is cartesian closed, using
the empty context, concatenation, and the curried-list exponential
\(\Delta^\Gamma\) of Theorem~\ref{thm:p10-stlc-context-ccc}.
Evaluation preserves the terminal object and has the canonical product
comparison isomorphism
\[
 j_{\Gamma,\Delta}:P_{\Gamma\cdot\Delta}\longrightarrow
 P_\Gamma\times P_\Delta,
 \qquad
 j_{\Gamma,\Delta}=\langle E(\pi_1),E(\pi_2)\rangle.
\]
If \(\Omega=\Delta^\Gamma\) and
\(\operatorname{ev}_{\rm syn}:\Omega\cdot\Gamma\to\Delta\) is its
syntactic evaluation, the exponential comparison is
\[
 \chi_{\Gamma,\Delta}
  =\Lambda\bigl(E(\operatorname{ev}_{\rm syn})
                   \circ j_{\Omega,\Gamma}^{-1}\bigr):
 P_\Omega\longrightarrow P_\Delta^{P_\Gamma}.
\]
It is an isomorphism.  Thus the equivalence preserves cartesian closed
structure up to these specified comparisons, including empty contexts.
\end{hbtheorem}

This reconstructs the given category from its named true equational
theory.  The bare \(\beta\eta\) context category instead has the
universal semantic property of
Theorem~\ref{thm:p10-stlc-ccc-completeness}; it need not be equivalent
to an arbitrary category in which its terms can be interpreted.


The categorical part of the Curry--Howard--Lambek correspondence
identifies typed substitution with composition and function types with
exponentials.  Theorems~\ref{thm:p10-stlc-ccc-completeness}
and~\ref{thm:p10-internal-ccc-equivalence} state the two precise
conclusions for the bare calculus and the named internal theory.
The earlier proof-theoretic correspondence supplies the interpretation
of typed terms as proofs; the categorical constructions specify the
objects, arrows, equations, and comparison isomorphisms
\cite{Lambek1980,LambekScott1986}.

%% One-sorted algebraic theories and the equivalence with equational algebras.

\subsection{Lawvere theories and equational algebras}
\label{sec:p10-lawvere-theories}

Throughout the Lawvere-theory development, signatures, equation sets,
algebra carriers, and generator sets belong to the fixed universe
used by \(\mathsf{Set}\) in Section~\ref{sec:p10-categorical-toolkit}.
The same constructions may be relativized to any larger universe
satisfying the size hypotheses of
Convention~\ref{def:p10-categorical-size}, when such a universe is fixed.

A finitary algebraic presentation can be recorded by its operations
in every finite list of variables and by the ways those operations
compose under substitution.  A Lawvere theory packages these data
as a category.  The number of variables is part of an arrow's type.

Lawvere's 1963 thesis uses the coproduct convention: the structural
functor goes from a category of finite sets to a theory category,
an \(n\)-ary operation goes from the generator to its \(n\)-fold
coproduct, and a model is a product-preserving functor from the
opposite theory to \(\mathsf{Set}\).
This is the presentation in Chapter~II, Section~1, pp.~61--63,
and Chapter~III, Section~1, pp.~74--75, of the authorized reprint
\cite{Lawvere2004FunctorialReprint}.
Passing to the opposite theory gives the product convention used
in Lawvere's separate 1968 article, Section~1, pp.~109--110 of that
same reprint, and in the Adamek--Rosicky--Vitale treatment below.
The two arrow conventions are related by this explicit reversal.

\begin{hbdefinition}{def:p10-lawvere-theory}{A one-sorted Lawvere theory}
Let \(\mathbb F\) be the category of finite ordinals
\(n=\{0,\ldots,n-1\}\) and all functions between them.
The empty ordinal is \(0\).
A one-sorted Lawvere theory consists of a small category \(T\)
whose objects are these natural numbers, and an identity-on-objects
functor
\[
                         j:\mathbb F^{\mathrm{op}}\longrightarrow T
\]
preserving finite products.
For the selector \(s_i:1\to n\), \(s_i(0)=i\), put
\(p_i^n=j(s_i^{\mathrm{op}}):n\to1\).
These projections exhibit \(n\) as the product of \(n\) copies
of the distinguished object \(1\).
In particular \(p_0^1=1_1\).
The object \(0\) is terminal; the distinguished generator \(1\)
is not by its name the terminal object.

For each family \(f_i:k\to1\), write
\(\langle f_i\rangle_{i<n}:k\to n\) for its unique product tuple.
For a function \(u:m\to n\) in \(\mathbb F\),
\[
                  j(u^{\mathrm{op}})
                     =\langle p_{u(i)}^n\rangle_{i<m}:n\to m.
\]
Thus structural maps select, repeat, or discard variables.
The functor \(j\) is not required to be faithful.
An \(n\)-ary operation is an arrow \(n\to1\); an arrow \(n\to m\)
is an \(m\)-tuple of \(n\)-ary operations.
A nullary operation is an arrow \(0\to1\).
\index{Lawvere theory}\index{algebraic theory!one-sorted}
\index{operation!as an arrow}\index{nullary operation}
\end{hbdefinition}

This is Definition~11.3 and Remarks~11.2--11.6, pp.~97--99,
of the author draft dated 1 February 2010
\cite{AdamekRosickyVitale2010Draft}.  The draft is a separately
identified version of the later published book.

\begin{hbdefinition}{def:p10-lawvere-presentation}{Equational presentations with specified finite contexts}
Fix a small finitary signature \(\Sigma\), allowing operation symbols
of arity zero.  Put \(X_n=(x_0,\ldots,x_{n-1})\), and let
\(T_\Sigma(n)=T_\Sigma(X_n)\) be its set of terms.
Here the term set is permitted to be empty: if \(n=0\) and there
are no nullary symbols, \(T_\Sigma(0)=\varnothing\).
This extends the carrier convention for the present construction;
Definition~\ref{def:p10-term-algebra} states the earlier nonempty
convention.

An equational presentation is a set \(E\) of triples
\((k,s,t)\), where \(s,t\in T_\Sigma(k)\).
The displayed context \(X_k\) records the variables over which
the identity is universally asserted.
A \(\Sigma\)-algebra with arbitrary carrier is a set \(A\) with
operations \(\sigma^A:A^k\to A\) of the declared arities.
The power \(A^k\) consists of functions from \(k\) to \(A\),
also written as \(k\)-tuples, and \(A^0\) is always a singleton.
The algebra satisfies \((k,s,t)\) when
\[
                     s^A(a)=t^A(a)\qquad\text{for every }a\in A^k.
\]
Evaluation uses the recursion of Definition~\ref{def:p10-term-algebra}.
An empty carrier admits all positive-arity operations uniquely,
but no nullary operation.
Write \(\operatorname{Alg}_{\mathrm{all}}(\Sigma,E)\) for these
models and their operation-preserving functions, and
\(\operatorname{Alg}_{\mathrm{ne}}(\Sigma,E)\) for the full
subcategory on nonempty carriers.

For every \(n\geq0\), let \(\equiv_E^n\) be the least family of
equivalence relations on \(T_\Sigma(n)\) that is compatible with
all operations, contains \(s\equiv_E^k t\) for \((k,s,t)\in E\),
and satisfies
\[
 s\equiv_E^k t
 \ \Longrightarrow\
 s[r_0/x_0,\ldots,r_{k-1}/x_{k-1}]
       \equiv_E^n
 t[r_0/x_0,\ldots,r_{k-1}/x_{k-1}]
       \quad(r_i\in T_\Sigma(n)).
\]
All substitutions are simultaneous.
The least family is the intersection of all families with these
properties; the family of full relations supplies one candidate.
It records the contextual form of equational deduction, including
instances of axioms written in other finite contexts.
\index{equational presentation}\index{algebra!empty carrier}
\index{equational consequence!finite contexts}
\end{hbdefinition}

\begin{hbtheorem}{thm:p10-lawvere-term-theory}{proved=annex:p10.lawvere.1}{The Lawvere theory of an equational presentation}
Define
\[
 T_{\Sigma,E}(n,m)=
             \bigl(T_\Sigma(n)/{\equiv_E^n}\bigr)^m.
\]
The identity of \(n\) is \(([x_i])_{i<n}\), and composition
substitutes the first arrow's tuple into each component of the
second arrow.
These data form a Lawvere theory.
Its structural functor sends \(u:m\to n\) to
\(([x_{u(i)}])_{i<m}:n\to m\).
Its terminal object is \(0\), and its binary product of \(n,m\)
is \(n+m\), with the two variable blocks as projections.
These assertions include empty source or target contexts.

For an algebra \(A\models E\), the formula
\[
 \begin{aligned}
 M_A(n)&=A^n,\\
 M_A(([t_i])_{i<m})(a)&=(t_i^A(a))_{i<m}
 \end{aligned}
\]
defines a finite-product-preserving functor
\(M_A:T_{\Sigma,E}\to\mathsf{Set}\).
An algebra homomorphism \(h:A\to B\) induces a natural
transformation \(M_h:M_A\Rightarrow M_B\) with components
\((M_h)_n=h^n\).
\end{hbtheorem}

The term theory and its evaluation are Notation~13.2 and
Lemma~13.3, p.~122, in
\cite{AdamekRosickyVitale2010Draft}; equations give the quotient
theory in Definition~13.10 and Theorem~13.11, p.~125.
The explicit context-category construction and its product
property are also given by
\cite[Section~4.2, pp.~30--35, especially
Proposition~4.2.1]{Pitts1995Draft}.

\begin{hbdefinition}{def:p10-lawvere-model}{Models and their underlying sets}
For a Lawvere theory \(T\), let
\[
                \operatorname{Mod}(T)
                    =\operatorname{FP}(T,\mathsf{Set})
\]
be the category of finite-product-preserving functors, with all
natural transformations as arrows.
The underlying-set functor is evaluation at the distinguished
generator:
\[
                 U(M)=M(1),\qquad U(\alpha)=\alpha_1.
\]
For every model the canonical comparison
\[
 c_n^M:M(n)\longrightarrow M(1)^n,\qquad
 c_n^M(z)=(M(p_i^n)(z))_{i<n},
\]
is a bijection.  For \(n=0\), this means that \(M(0)\) is
a singleton.  For \(n=1\), the comparison sends \(a\) to the
singleton tuple \((a)\).
Products are preserved through these comparisons; no strict
identity \(M(n)=M(1)^n\) is part of the definition.
\index{functorial semantics}\index{model!of an algebraic theory}
\index{underlying-set functor}
\end{hbdefinition}

\begin{hbtheorem}{thm:p10-lawvere-models}{proved=annex:p10.lawvere.2}{Equational algebras and functorial models}
Term evaluation gives an equivalence
\[
 \operatorname{Alg}_{\mathrm{all}}(\Sigma,E)
           \simeq\operatorname{Mod}(T_{\Sigma,E}).
\]
For the reverse construction, put \(A_M=M(1)\).
For a symbol \(\sigma\) of arity \(k\), let
\(\bar\sigma:k\to1\) be the one-component term arrow
\(([\sigma(x_0,\ldots,x_{k-1})])\), and define
\[
                     \sigma^{A_M}
                         =M(\bar\sigma)\circ(c_k^M)^{-1}.
\]
Then \(A_M\models E\), and a natural transformation
\(\alpha:M\Rightarrow N\) gives the algebra homomorphism
\(\alpha_1:A_M\to A_N\).
Every component of \(\alpha\) is determined by \(\alpha_1\):
\[
                     c_n^N\circ\alpha_n
                        =(\alpha_1)^n\circ c_n^M.
\]
The equivalence has explicit natural comparison isomorphisms.
On an algebra \(A\), the comparison with \(A_{M_A}=A^1\)
is \(a\mapsto(a)\); on a model \(M\), it is the family \(c_n^M\).

The equivalence restricts to
\[
 \operatorname{Alg}_{\mathrm{ne}}(\Sigma,E)
 \simeq
 \{M\in\operatorname{Mod}(T_{\Sigma,E}):M(1)\ne\varnothing\},
\]
where the right side is a full subcategory.
It therefore recovers the manuscript's earlier nonempty-algebra
convention without excluding empty set-valued models from the
unrestricted theorem.
\end{hbtheorem}

This functorial form, including all homomorphisms and natural
transformations, is
\cite[Theorem~13.11, p.~125; Proposition~11.8, p.~99]
{AdamekRosickyVitale2010Draft}
and \cite[Corollary~4.2.6(ii), p.~35]{Pitts1995Draft}.
Lawvere gives the original transformation-to-homomorphism
argument in Chapter~III, Section~1, p.~74 of the 1963 thesis
\cite{Lawvere2004FunctorialReprint}.

\begin{hbexample}{def:p10-lawvere-structural-collapse}{A structural functor can identify projections}
Take the empty signature and the equation
\((2,x_0,x_1)\).
The corresponding algebras have at most one element when empty
carriers are allowed.
In their term theory the two different selectors \(1\to2\) in
\(\mathbb F\) become the same projection \(2\to1\), because
\([x_0]=[x_1]\).
Hence the structural functor need not be faithful.
The object \(0\) still has no arrow to \(1\), since there are
no terms in the empty context.  This will distinguish its empty
initial model from the nonempty singleton models.
\end{hbexample}

%% An explicit equational basis presenting every one-sorted Lawvere theory.

\subsection{An equational presentation of every Lawvere theory}
\label{sec:p10-lawvere-presentations}

Let \(\mathcal T\) be a Lawvere theory in the sense of
Definition~\ref{def:p10-lawvere-theory}.
Write \(p_i^n:n\to1\), \(0\leq i<n\), for its selected projections.
The structural functor fixes \(p_0^1=1_1\).
For arrows \(g_i:n\to1\), write
\(\langle g_0,\ldots,g_{k-1}\rangle:n\to k\) for their unique
tuple.  An empty tuple is the unique arrow \(!_n:n\to0\).
Let \(X_n=(x_0,\ldots,x_{n-1})\), and abbreviate this variable
list by \(\mathbf x_n\).

\begin{hbdefinition}{def:p10-lawvere-named-presentation}{The named operation presentation}
The signature \(\Sigma_{\mathcal T}\) has one \(n\)-ary symbol
\(\sigma_f\) for every arrow \(f:n\to1\) in \(\mathcal T\),
for every \(n\geq0\).
Its equation set \(E_{\mathcal T}\), in the explicit finite contexts
of Definition~\ref{def:p10-lawvere-presentation}, consists of these
two families:
\[
\begin{aligned}
 \sigma_{p_i^n}(\mathbf x_n)&\approx x_i
       &&\text{in }X_n,\quad 0\leq i<n,\\
 \sigma_f\bigl(\sigma_{g_0}(\mathbf x_n),\ldots,
                    \sigma_{g_{k-1}}(\mathbf x_n)\bigr)
   &\approx \sigma_{\,f\circ\langle g_0,\ldots,g_{k-1}\rangle}
                      (\mathbf x_n)
       &&\text{in }X_n.
\end{aligned}
\]
The second family ranges over \(n,k\geq0\), \(f:k\to1\), and
all \(k\)-tuples \(g_i:n\to1\).
At \(k=0\) it reads
\[
                  \sigma_f()\approx\sigma_{f\circ!_n}(\mathbf x_n)
                        \quad\text{in }X_n,\qquad f:0\to1.
\]
Thus it also specifies how a nullary term is regarded in a
context with unused variables.  There are no projection equations
in \(X_0\).
Smallness of \(\mathcal T\) makes the signature and these indexed
equation families small sets.

For a term \(t\) in \(X_n\), define its interpretation
\(\llbracket t\rrbracket_n:n\to1\) recursively by
\[
 \llbracket x_i\rrbracket_n=p_i^n,\qquad
 \llbracket\sigma_f(t_0,\ldots,t_{k-1})\rrbracket_n
       =f\circ\langle\llbracket t_0\rrbracket_n,\ldots,
                              \llbracket t_{k-1}\rrbracket_n\rangle .
\]
The empty tuple gives
\(\llbracket\sigma_f()\rrbracket_n=f\circ!_n\).
In this section abbreviate the relation
\(\equiv_{E_{\mathcal T}}^n\) of
Definition~\ref{def:p10-lawvere-presentation} by
\(\equiv_{E_{\mathcal T},n}\).
\index{Lawvere theory!named presentation}
\end{hbdefinition}

\begin{hbtheorem}{thm:p10-lawvere-presentation-recovery}{proved=annex:p10.lawvere-present.1}{Recovery of an abstract theory from its named equations}
Every term \(t\) in \(X_n\) satisfies
\[
       t\equiv_{E_{\mathcal T},n}
                         \sigma_{\llbracket t\rrbracket_n}(\mathbf x_n).
\]
Consequently, for any terms \(s,t\) in the same context,
\[
 s\equiv_{E_{\mathcal T},n}t
       \quad\Longleftrightarrow\quad
                  \llbracket s\rrbracket_n=\llbracket t\rrbracket_n.
\]
Interpretation induces an isomorphism of categories
\[
 \varepsilon_{\mathcal T}:
       T_{\Sigma_{\mathcal T},E_{\mathcal T}}
                  \longrightarrow\mathcal T
\]
which is the identity on objects and preserves the selected
finite products and the structural finite-set maps.
For a tuple arrow its value is
\[
       \varepsilon_{\mathcal T}([t_0],\ldots,[t_{m-1}])
          =\langle\llbracket t_0\rrbracket_n,\ldots,
                                     \llbracket t_{m-1}\rrbracket_n\rangle.
\]
The inverse sends \(h:n\to m\) to
\[
       \bigl([\sigma_{p_j^m\circ h}(\mathbf x_n)]\bigr)_{0\leq j<m}.
\]
For \(m=0\) this is the empty tuple.
All assertions include \(n=0\), nullary symbols, and composition
through the terminal object \(0\).  Faithfulness of the structural
finite-set functor is not required.
\end{hbtheorem}

The named signature and quotient of its term theory by evaluation
are given in
\cite[Definition~13.5 and Proposition~13.7, p.~123;
Remark~13.8, p.~124]{AdamekRosickyVitale2010Draft}.
The proof here identifies that evaluation kernel using the two
explicit equation families above.
Lawvere's thesis treats presentations and substitution in
\cite[Chapter~II, \S2, pp.~69--73]{Lawvere2004FunctorialReprint}.
Its theory category uses finite coproducts and operations \(1\to n\);
the product convention used here takes its opposite
\cite[Chapter~II, \S1, pp.~61--63]{Lawvere2004FunctorialReprint}.
The latter locators are pages of the authorized 2004 reprint;
the former are pages of the February 2010 author draft.

\begin{hbdefinition}{def:p10-lawvere-theory-morphism}{Morphisms of one-sorted theories}
A theory morphism \(H:\mathcal T\to\mathcal U\) is a functor
which is the identity on the objects \(n\) and satisfies
\[
                         H(p_i^{n,\mathcal T})=p_i^{n,\mathcal U}.
\]
Equivalently, it commutes with the structural functors from the
opposite of the finite-ordinal category.
Composition and identity functors give the category of these
theories and morphisms.
Write \(\operatorname{Mod}(\mathcal T)\) for the models and all
natural transformations of Definition~\ref{def:p10-lawvere-model}.
This morphism convention is
\cite[Definition~11.3 and Remark~11.4, p.~98]{AdamekRosickyVitale2010Draft}.
\end{hbdefinition}

\begin{hbtheorem}{thm:p10-lawvere-model-restriction}{proved=annex:p10.lawvere-present.2}{Theory morphisms and restriction of models}
Every theory morphism \(H:\mathcal T\to\mathcal U\) preserves
the selected finite products.  It induces a functor
\[
 H^*:\operatorname{Mod}(\mathcal U)\longrightarrow
                         \operatorname{Mod}(\mathcal T),
 \qquad H^*(M)=M\circ H,\qquad
              (H^*\alpha)_n=\alpha_n.
\]
This restricts to the full subcategories with nonempty carrier
\(M(1)\).  For \(K:\mathcal U\to\mathcal V\),
\[
             (K\circ H)^*=H^*\circ K^*,\qquad (1_{\mathcal T})^*=1.
\]
On named presentations, replacing \(\sigma_f\) by \(\sigma_{H(f)}\)
and fixing variables induces a theory morphism \(\widehat H\)
between the corresponding term theories, satisfying
\[
       \varepsilon_{\mathcal U}\circ\widehat H
                        =H\circ\varepsilon_{\mathcal T}.
\]
\end{hbtheorem}

%% Free Lawvere models and the all-carrier/nonempty-carrier boundary.

\subsection{Representable models and free generators}
\label{sec:p10-lawvere-free-algebras}

Generator sets and algebra carriers below follow the fixed-universe
convention of Section~\ref{sec:p10-lawvere-theories}.

Use the product-oriented Lawvere theories and models of
Definitions~\ref{def:p10-lawvere-theory} and~\ref{def:p10-lawvere-model}.
Let \(\mathbb T\) be such a theory.
Write \(\operatorname{Mod}(\mathbb T)\) for their category, with all
natural transformations as morphisms, and \(U(M)=M(1)\).
For \(n\in\mathbb N\), let
\[
 R_n=\mathbb T(n,-):\mathbb T\longrightarrow\mathsf{Set},\qquad
 R_n(t)(a)=t\circ a.
\]
The comparison \(c_n^M:M(n)\to M(1)^n\) has coordinates
\(M(p_i^n)\), \(0\leq i<n\).
Powers here are sets of finite tuples; in particular the canonical
map \(A\to A^1\) sends \(a\) to the singleton tuple \((a)\).

\begin{hbtheorem}{thm:p10-lawvere-finite-free}{proved=annex:p10.lawvere-free.1}{Finite free models are representable}
Each \(R_n\) is a model of \(\mathbb T\).
The map
\[
 \eta^n:n\longrightarrow U(R_n)=\mathbb T(n,1),
 \qquad \eta^n(i)=p_i^n,
\]
makes it free on the finite set \(n\): for every model \(M\),
restriction to these generators is a bijection
\[
 \operatorname{Nat}(R_n,M)\longrightarrow\mathsf{Set}(n,U(M)),
 \qquad \alpha\longmapsto\alpha_1\circ\eta^n.
\]
It is natural in \(M\), including arbitrary model homomorphisms.
For \(n=0\), \(R_0\) is initial and its carrier is \(\mathbb T(0,1)\).

Precomposition gives a fully faithful functor
\[
 \mathbb T^{\mathrm{op}}\longrightarrow\operatorname{Mod}(\mathbb T),
 \qquad n\longmapsto R_n,
\]
with
\[
 \operatorname{Nat}(R_n,R_m)\cong\mathbb T(m,n).
\]
Precisely, \(\mathbb T^{\mathrm{op}}\) is isomorphic to the category
whose objects are the indexed chosen models \((n,R_n)\) and whose
arrows are all natural transformations between them.
Every model free on a finite set is isomorphic to one of these;
the isomorphism is unique once the correspondence of generators is fixed.
The generator map is not required to be injective.
\end{hbtheorem}

The opposite-theory and representable-free-model descriptions are
\cite[1963 thesis, Chapter~III, \S1, pp.~74--75;
1968 article, \S1, p.~110]{Lawvere2004FunctorialReprint} and
\cite[Remark~11.21 and Corollary~11.22, pp.~102--103]{AdamekRosickyVitale2010Draft}.
The latter locators are for the February 2010 author draft.
The context-category formulation also appears in
\cite[Remark~4.2.3, p.~33]{Pitts1995Draft}.

\subsection{Free models on arbitrary sets}
\label{sec:p10-lawvere-free-adjunction}

For this construction use the all-carrier category
\(\operatorname{Alg}_{\rm all}(\Sigma,E)\) of
Theorem~\ref{thm:p10-lawvere-models}.
An empty carrier is allowed when all operations can be interpreted.
This convention is local to the Lawvere-model treatment;
Definition~\ref{def:p10-term-algebra} and
Theorem~\ref{thm:p10-free-algebra} retain their nonempty-algebra hypotheses.

For a set \(X\), use a tagged copy of its elements as formal variables,
disjoint from \(\Sigma\), and let \(T_\Sigma(X)\) be the term set.
For terms \(s,t\), put
\[
 s\mathrel{\theta_X}t
 \quad\Longleftrightarrow\quad
 \forall A\in\operatorname{Alg}_{\rm all}(\Sigma,E)\
 \forall v:X\to A\quad s^A(v)=t^A(v),
 \qquad Q_X=T_\Sigma(X)/\theta_X.
\]
The quotient operations are induced by term formation.
When there are no terms, both the relation and the quotient are empty.
Each equation of \(E\) is interpreted universally in its declared
finite variable context, as in Definition~\ref{def:p10-lawvere-presentation}.

\begin{hbtheorem}{thm:p10-lawvere-free-adjunction}{proved=annex:p10.lawvere-free.2}{Free models on sets and the underlying-set adjunction}
For every finitary equational presentation \((\Sigma,E)\) and every
set \(X\), \(Q_X\) is an all-carrier \((\Sigma,E)\)-algebra.
The map \(x\mapsto[x]\) makes it free on \(X\).
For \(v:X\to A\), its unique extending homomorphism is
\[
 \bar v:Q_X\longrightarrow A,\qquad \bar v([t])=t^A(v).
\]
The free construction extends to a functor on sets.

For any Lawvere theory \(\mathbb T\), take the recovery presentation
and identity-on-objects theory isomorphism
\[
 \rho:T_{\Sigma_{\mathbb T},E_{\mathbb T}}\longrightarrow\mathbb T
\]
of Theorem~\ref{thm:p10-lawvere-presentation-recovery}.
With \(Q_X\) formed for this presentation, define
\[
 \mathcal F_{\mathbb T}(X)=M_{Q_X}\circ\rho^{-1},
 \qquad \mathcal F_{\mathbb T}(X)(n)=Q_X^n.
\]
Here \(M_{Q_X}\) is the evaluation model of
Theorem~\ref{thm:p10-lawvere-models}.
There are natural bijections
\[
 \operatorname{Nat}(\mathcal F_{\mathbb T}(X),M)
       \cong\mathsf{Set}(X,U(M)),
\]
so \(\mathcal F_{\mathbb T}\dashv U\).
Its generator map is
\[
 \eta_X:X\longrightarrow U\mathcal F_{\mathbb T}(X)=Q_X^1,
 \qquad \eta_X(x)=([x]).
\]
On a finite set \(n\), this free model is isomorphic to \(R_n\)
by the unique isomorphism respecting the displayed generators.
All assertions include \(X=\varnothing\).
\end{hbtheorem}

The term and presentation constructions are
\cite[\S13, pp.~121--125]{AdamekRosickyVitale2010Draft}; free models on
all sets are discussed in Remark~11.21, pp.~102--103.
The proof below verifies the empty term case separately and selects
representatives only for the finitely many variables in one equation.

\subsection{Empty carriers and the nonempty restriction}
\label{sec:p10-lawvere-empty-carriers}

\begin{hbtheorem}{thm:p10-lawvere-empty-carriers}{proved=annex:p10.lawvere-free.3}{The exact empty-carrier boundary}
Every model \(M\) has \(M(0)\) a singleton.
A model with \(M(1)=\varnothing\) exists exactly when
\[
                         \mathbb T(0,1)=\varnothing.
\]
When it exists it is unique up to its unique natural isomorphism.
It can be written as \(M(0)=\{*\}\), \(M(n)=\varnothing\) for \(n>0\),
with the uniquely possible functions on arrows.

For a presentation \((\Sigma,E)\), arrows \(0\to1\) are singleton tuples
of ground-term classes, hence are canonically in bijection with those
classes.  Thus a basic nullary operation excludes empty
carriers, while the absence of all nullary symbols permits the empty
algebra and empty-carrier model.  Equations identify existing terms;
they do not create a ground term when the grammar has none.

The full subcategory of models with \(M(1)\ne\varnothing\) corresponds
to \(\operatorname{Alg}_{\rm ne}(\Sigma,E)\).
Free models on nonempty generator sets lie in this subcategory.
If \(\mathbb T(0,1)\ne\varnothing\), all models have nonempty carriers
and \(R_0\) is also initial there.
If \(\mathbb T(0,1)=\varnothing\), the all-carrier initial model is
excluded; the nonempty subcategory may or may not have an initial object.
\end{hbtheorem}

For the empty signature with no equations, nonempty algebras are
nonempty sets and have no initial object.  Nonempty semigroups also
have no initial object: every such semigroup has two distinct constant
homomorphisms to the two-element left-zero semigroup, whose product is
\(a b=a\).
In contrast, for one unary symbol \(d\) and the equation
\(d(x)=d(y)\), the one-element algebra is initial among nonempty models,
although there are no ground terms.  The empty algebra is initial in
the all-carrier category.
These examples distinguish an expressible nullary term from a value
defined by a positive-arity operation independently of its argument,
as in \cite[1963 thesis, Chapter~II, \S1, p.~64]{Lawvere2004FunctorialReprint}.
They also explain why the all-carrier adjunction cannot simply be
restricted at the zero-generator input without checking its carrier.
\index{Lawvere theory!free models}\index{free algebra!zero generators}
\index{empty carrier}\index{definable constant}


%% Explicit first-order predicate doctrines with quantifiers and equality.

\subsection{Predicates, substitution, and quantifier adjunctions}
\label{sec:p10-doctrine-quantifiers}

For a fixed context, predicates form an ordered algebra: entailment
is its order, and conjunction is its meet.  A substitution changes
the context and pulls predicates back to it.  Quantifiers are
characterized by how entailment interacts with that pullback.
The following definition states the precise first-order convention
used here.

Lawvere's original hyperdoctrine is category-valued, with cartesian
closed categories of attributes and adjoints to substitution along
all terms.  His substitution functors are not required to preserve
implication in that generality
\cite[pp.~11--13]{Lawvere2006AdjointnessReprint}.
The posetal first-order convention below uses the indexed orders,
Heyting operations, projection adjoints, and equality axioms
specified separately in
\cite[Definitions~5.2.1, 5.4.1, 5.5.1, and 5.6.1,
pp.~41--56]{Pitts1995Draft}.
Their assumptions are part of this presentation.

\begin{hbdefinition}{def:p10-first-order-hyperdoctrine}{A first-order posetal hyperdoctrine}
Let \(\mathcal C\) be a locally small category with chosen finite
products.  Its objects and indexed families are formed in the
containing ambient universe of
Section~\ref{sec:p10-categorical-toolkit}.
Let \(\mathsf{Heyt}\) have small Heyting algebras as objects and
maps preserving \(0,1,\wedge,\vee,\Rightarrow\) as arrows.
An indexed predicate algebra is a functor
\[
                       P:\mathcal C^{\mathrm{op}}\longrightarrow
                         \mathsf{Heyt}.
\]
The Heyting algebra \(P(I)\) is the fibre of predicates in context
\(I\).  Write its greatest and least elements as
\(\top_I,\bot_I\).
For \(f:J\to I\), write
\[
                    f^*=P(f):P(I)\longrightarrow P(J).
\]
This reindexing preserves all the stated operations and satisfies
\(1_I^*=1_{P(I)}\) and \((f g)^*=g^*f^*\).

For every object \(A\) and product projection
\(\pi_{I,A}:I\times A\to I\), require monotone maps
\[
        \exists_{\pi_{I,A}},\forall_{\pi_{I,A}}:
                    P(I\times A)\longrightarrow P(I)
\]
with the order adjunctions
\[
 \begin{aligned}
 \exists_\pi\alpha\leq\beta
       &\quad\Longleftrightarrow\quad\alpha\leq\pi^*\beta,\\
 \pi^*\beta\leq\alpha
       &\quad\Longleftrightarrow\quad\beta\leq\forall_\pi\alpha.
 \end{aligned}
\]
Here \(\alpha\in P(I\times A)\), \(\beta\in P(I)\).
For \(f:J\to I\), require the two Beck--Chevalley identities
\[
 \begin{aligned}
 f^*\exists_{\pi_{I,A}}
     &=\exists_{\pi_{J,A}}(f\times1_A)^*,\\
 f^*\forall_{\pi_{I,A}}
     &=\forall_{\pi_{J,A}}(f\times1_A)^*.
 \end{aligned}
\]
They express stability under parameter substitution in the product
square with upper arrow \(f\times1_A:J\times A\to I\times A\)
and lower arrow \(f:J\to I\).
These squares are pullbacks using only finite products.
No adjoints along arbitrary base arrows, or other base pullbacks,
are included in this definition.
\index{hyperdoctrine!first-order posetal}
\index{predicate!indexed algebra}\index{reindexing}
\index{Beck--Chevalley condition}
\end{hbdefinition}

\begin{hbtheorem}{thm:p10-doctrine-quantifier-laws}{proved=annex:p10.quant.1}{Order laws of the two quantifiers}
The parameter squares in the preceding definition are pullbacks.
For each projection, both adjoints are uniquely determined by
reindexing.  They satisfy
\[
 \begin{gathered}
 \alpha\leq\pi^*\exists_\pi\alpha,\qquad
 \exists_\pi\pi^*\beta\leq\beta,\\
 \pi^*\forall_\pi\alpha\leq\alpha,\qquad
 \beta\leq\forall_\pi\pi^*\beta .
 \end{gathered}
\]
Existential quantification preserves finite joins and bottom,
while universal quantification preserves finite meets and top:
\[
 \begin{aligned}
 \exists_\pi(\alpha\vee\alpha')
       &=\exists_\pi\alpha\vee\exists_\pi\alpha',
 &\exists_\pi\bot_{I\times A}&=\bot_I,\\
 \forall_\pi(\alpha\wedge\alpha')
       &=\forall_\pi\alpha\wedge\forall_\pi\alpha',
 &\forall_\pi\top_{I\times A}&=\top_I.
 \end{aligned}
\]
If the fibres are Boolean algebras, then
\[
                   \forall_\pi\alpha
                         =\neg\exists_\pi\neg\alpha .
\]
The adjunction hypotheses alone supply these order laws;
the two Beck--Chevalley equations are additional substitution
conditions in the definition.
\end{hbtheorem}

\begin{hbtheorem}{thm:p10-doctrine-frobenius}{proved=annex:p10.quant.2}{Frobenius from implication-preserving reindexing}
Let \(H,K\) be Heyting algebras, and let monotone maps
\(L:H\to K\), \(R:K\to H\) satisfy \(L\dashv R\).
If \(R\) preserves Heyting implication, then
\[
                         L(a\wedge R(b))=L(a)\wedge b
                         \quad(a\in H,\ b\in K).
\]
In particular, the first-order doctrine above satisfies
Frobenius reciprocity:
\[
             \exists_\pi(\alpha\wedge\pi^*\beta)
                    =(\exists_\pi\alpha)\wedge\beta.
\]
This conclusion uses preservation of implication.
In a presentation with only finite-meet fibres, Frobenius
must instead be required or proved from other assumptions.
\index{Frobenius reciprocity}
\end{hbtheorem}

The two projection adjunctions and their parameter naturality are
equations (5.6)--(5.9), pp.~52--53, of
\cite{Pitts1995Draft}.  Frobenius is equation (5.10) and
Remark~5.5.2 on p.~53.  The source treats these assumptions
modularly; the present Heyting-preserving convention implies
Frobenius by the displayed theorem.

\subsection{Equality with parameters}
\label{sec:p10-doctrine-equality}

\begin{hbdefinition}{def:p10-doctrine-equality}{Parameterized equality predicates}
For each object \(A\), specify a predicate
\(\delta_A\in P(A\times A)\).
For any parameter object \(I\), use the left-associated product
\((I\times A)\times A\).
Write \(\pi_I:I\times A\to I\) and \(\pi_A:I\times A\to A\).
Define
\[
 \begin{aligned}
 e_{I,A}
    &=\langle1_{I\times A},\pi_A\rangle:
                        I\times A\longrightarrow(I\times A)\times A,\\
 p_{I,A}
    &=\pi_1:(I\times A)\times A\longrightarrow I\times A,\\
 q_{I,A}
    &=\langle\pi_I\pi_1,\pi_2\rangle:
                        (I\times A)\times A\longrightarrow I\times A,\\
 r_{I,A}
    &=\langle\pi_A\pi_1,\pi_2\rangle:
                        (I\times A)\times A\longrightarrow A\times A.
 \end{aligned}
\]
Thus \(e\) duplicates the last variable, \(p\) retains its first
copy, \(q\) retains its second copy, and \(r\) selects the two
copies while discarding the parameters.
Require that the monotone map
\[
 L_e:P(I\times A)\longrightarrow P((I\times A)\times A),
 \qquad L_e(\alpha)=p^*\alpha\wedge r^*\delta_A
\]
be left adjoint to \(e^*\).  Explicitly,
\[
              p^*\alpha\wedge r^*\delta_A\leq\beta
                  \quad\Longleftrightarrow\quad
              \alpha\leq e^*\beta .
\]
This is required for every parameter object \(I\).
An equality predicate is an element of a fibre; a representation
by actual subobjects is not assumed here.
\index{equality!in an indexed predicate algebra}
\index{diagonal!equality adjunction}
\end{hbdefinition}

The parameterized adjunction is equation (5.12) and
Definition~5.6.1, pp.~55--56, of
\cite{Pitts1995Draft}.
The coordinate maps above state it for the chosen left-associated
products.  Rebracketing uses the canonical product isomorphisms.

\begin{hbtheorem}{thm:p10-doctrine-equality-laws}{proved=annex:p10.quant.3}{Reflexivity and replacement of equals}
Let \(\Delta_A=\langle1_A,1_A\rangle:A\to A\times A\) and
\(\tau_A=\langle\pi_2,\pi_1\rangle:A\times A\to A\times A\).
The parameterized equality axiom implies
\[
                    \Delta_A^*\delta_A=\top_A,\qquad
                    \tau_A^*\delta_A=\delta_A.
\]
For every parameter object \(I\) and \(\gamma\in P(I\times A)\),
\[
                     p_{I,A}^*\gamma\wedge r_{I,A}^*\delta_A
                                  \leq q_{I,A}^*\gamma .
\]
In particular, in the fibre over \((A\times A)\times A\),
\[
                     p_{A,A}^*\delta_A\wedge r_{A,A}^*\delta_A
                                  \leq q_{A,A}^*\delta_A,
\]
which is transitivity.  For every arrow \(f:A\to B\),
\[
                             \delta_A\leq(f\times f)^*\delta_B.
\]
Thus reindexing equality along a pair of terms supports reflexivity,
symmetry, transitivity, and substitution in predicates and terms.
The assertion is internal equality preservation; no reflection
of base-arrow equality is included.
\end{hbtheorem}

\subsection{Inhabited sorts and unused variables}
\label{sec:p10-doctrine-inhabited}

\begin{hbdefinition}{def:p10-inhabited-sort}{An internally inhabited sort}
An object \(A\) is inhabited in this doctrine when
\[
                  \exists_{\pi_{1,A}}\top_{1\times A}=\top_1,
\]
where \(1\) is the chosen terminal object of \(\mathcal C\).
This is an assertion about predicates in the terminal context.
A chosen global element \(a:1\to A\) is a different datum.
\index{inhabited sort}\index{nonempty domain!categorical semantics}
\end{hbdefinition}

\begin{hbtheorem}{thm:p10-doctrine-inhabited}{proved=annex:p10.quant.4}{Inhabitance and reflection along context extension}
An object \(A\) is inhabited exactly when
\(\exists_{\pi_{I,A}}\top_{I\times A}=\top_I\) for every
parameter object \(I\).
A global element \(1\to A\) suffices for inhabitance.

For an inhabited \(A\), every projection \(\pi=\pi_{I,A}\)
satisfies
\[
                   \exists_\pi\pi^*\beta=\beta
                   =\forall_\pi\pi^*\beta
                         \quad(\beta\in P(I)),
\]
and reindexing along \(\pi\) reflects order:
\[
                    \pi^*\alpha\leq\pi^*\beta
                         \quad\Longleftrightarrow\quad
                    \alpha\leq\beta.
\]
Consequently, an unused variable of an inhabited sort may be
removed from the retained variable context of a semantic sequent.
For several such variables the argument applies one extension
at a time.  Without inhabitance this removal is not supplied by
the doctrine axioms.
\end{hbtheorem}

The powerset construction below permits empty sorts and makes
this distinction concrete.  The subsequent soundness theorem
retains variable contexts, then states the additional bridge to
ordinary nonempty-domain first-order semantics.

%% Quantification along arbitrary set maps, including empty fibres.

\subsection{Powersets: substitution and quantifiers along a function}
\label{sec:p10-powerset-quantifiers}

Work in the category \(\mathsf{Set}\) of sets in the fixed universe
of Section~\ref{sec:p10-categorical-toolkit}.
All sets, maps, and indexed families below use that convention.
No nonemptiness hypothesis is imposed on these sets.

\begin{hbdefinition}{def:p10-powerset-quantifiers}{The powerset fibre and quantifiers}
For a set \(X\), let \(P(X)=\mathcal P(X)\), ordered by inclusion.
Its Boolean operations are intersection, union, complement in \(X\),
bottom \(\varnothing\), and top \(X\).
Indexed intersections are taken in the indicated powerset:
the empty intersection in \(P(X)\) is \(X\), and the empty
union is \(\varnothing\).
Its Heyting implication is
\[
                     U\Rightarrow_X V=(X\setminus U)\cup V.
\]
When \(X=\varnothing\), this is the one-element Boolean algebra.

For a function \(f:X\to Y\), define the reindexing map
\[
 f^*:P(Y)\longrightarrow P(X),\qquad
 f^*(V)=f^{-1}(V)=\{x\in X:f(x)\in V\}.
\]
Define two maps in the other direction by
\[
\begin{aligned}
 \exists_f(U)&=\{y\in Y:\exists x\in U\ (f(x)=y)\},\\
 \forall_f(U)&=\{y\in Y:\forall x\in X\
                                  (f(x)=y\Longrightarrow x\in U)\}.
\end{aligned}
\]
Thus \(\exists_f(U)\) is the image \(f[U]\), whereas
\(\forall_f(U)\) consists of points whose entire \(f\)-fibre
lies in \(U\).
\index{powerset!quantifiers}\index{inverse image!as substitution}
\index{quantifier!along a function}
\end{hbdefinition}

\begin{hbtheorem}{thm:p10-powerset-adjunctions}{proved=annex:p10.quant-set.1}{The two powerset adjunctions and their function laws}
Inverse image preserves the Boolean operations and all unions
and intersections.  The maps \(\exists_f,f^*,\forall_f\) are
monotone, and
\[
                      \exists_f\dashv f^*\dashv\forall_f.
\]
Explicitly, for \(U\subseteq X\) and \(V\subseteq Y\),
\[
 \exists_f(U)\subseteq V\ \Longleftrightarrow\ U\subseteq f^*(V),
 \qquad
 f^*(V)\subseteq U\ \Longleftrightarrow\ V\subseteq\forall_f(U).
\]
These properties determine both quantifier maps uniquely.
For \(f:X\to Y\), \(g:Y\to Z\), the function laws are
\[
\begin{gathered}
 (1_X)^*=\exists_{1_X}=\forall_{1_X}=1_{P(X)},\\
 (g\circ f)^*=f^*\circ g^*,\qquad
 \exists_{g\circ f}=\exists_g\circ\exists_f,\qquad
 \forall_{g\circ f}=\forall_g\circ\forall_f .
\end{gathered}
\]
Moreover,
\[
 \forall_f(U)=Y\setminus\exists_f(X\setminus U),\qquad
 \exists_f\Bigl(\bigcup_j U_j\Bigr)=\bigcup_j\exists_f(U_j),\qquad
 \forall_f\Bigl(\bigcap_j U_j\Bigr)=\bigcap_j\forall_f(U_j).
\]
The family may be empty.  In particular,
\[
 \exists_f(\varnothing)=\varnothing,\quad
 \forall_f(X)=Y,\quad
 \exists_f(X)=f[X],\quad
 \forall_f(\varnothing)=Y\setminus f[X].
\]
\end{hbtheorem}

\begin{hbexample}{def:p10-powerset-fibres}{A universal quantifier includes an empty fibre}
Let \(X=\{0,1,2\}\), \(Y=\{a,b,c\}\), and
\(f(0)=f(1)=a\), \(f(2)=b\).
The fibres at \(a,b,c\) are respectively \(\{0,1\},\{2\},\varnothing\).
For \(U=\{0,2\}\),
\[
                    \exists_f(U)=\{a,b\},\qquad
                    \forall_f(U)=\{b,c\}.
\]
The fibre at \(a\) meets \(U\) but is not contained in it;
the empty fibre at \(c\) is contained in \(U\).
For a projection \(\pi:I\times A\to I\), the same definitions give
\[
 \exists_\pi(U)=\{i:\exists a\in A\ ((i,a)\in U)\},\qquad
 \forall_\pi(U)=\{i:\forall a\in A\ ((i,a)\in U)\}.
\]
If \(A=\varnothing\), its only domain predicate is
\(\varnothing\), with
\(\exists_\pi(\varnothing)=\varnothing\) and
\(\forall_\pi(\varnothing)=I\).
\end{hbexample}

The powerset and inverse-image construction is
\cite[Example~5.2.2, p.~42]{Pitts1995Draft};
the projection quantifiers are Example~5.5.3, pp.~53--54.
Lawvere's powerset example uses inverse image for substitution
and direct image for existential quantification
\cite[p.~13]{Lawvere2006AdjointnessReprint}.
These locators refer respectively to the 18 May 1995 draft
and the authorized 2006 reprint.

\subsection{Pullback stability and conjunction}
\label{sec:p10-powerset-base-change}

\begin{hbtheorem}{thm:p10-powerset-base-change}{proved=annex:p10.quant-set.2}{Beck--Chevalley and Frobenius for powersets}
Suppose the following is a pullback square of sets:
\[
\begin{array}{ccc}
 X'&\xrightarrow{\ g'\ }&X\\
 {\scriptstyle f'}\downarrow&&\downarrow{\scriptstyle f}\\
 Y'&\xrightarrow{\ g\ }&Y.
\end{array}
\]
Then both Beck--Chevalley identities hold, as maps
\(P(X)\to P(Y')\):
\[
              g^*\circ\exists_f=\exists_{f'}\circ(g')^*,
              \qquad
              g^*\circ\forall_f=\forall_{f'}\circ(g')^*.
\]
For an arbitrary function \(f:X\to Y\), Frobenius reciprocity is
\[
             \exists_f\bigl(U\cap f^*(V)\bigr)
                         =\exists_f(U)\cap V
                \quad(U\subseteq X,\ V\subseteq Y).
\]
For a square which is only assumed to commute, the following
inclusions always hold for \(U\subseteq X\):
\[
 \exists_{f'}((g')^*U)\subseteq g^*(\exists_fU),
 \qquad
 g^*(\forall_fU)\subseteq\forall_{f'}((g')^*U).
\]
Commutativity alone does not imply the reverse inclusions.
\end{hbtheorem}

\begin{hbexample}{def:p10-powerset-commuting-square}{A commuting square without base change}
In the displayed shape, take
\(X'=\varnothing\), \(X=Y=Y'=\{*\}\), \(f=g=1_{\{*\}}\),
and let \(f',g'\) be the unique maps from the empty set.
The square commutes, but is not a pullback: the compatible
pair of points of \(X,Y'\) has no point over it in \(X'\).
For \(U=X\),
\[
             g^*(\exists_fU)=\{*\},\qquad
             \exists_{f'}((g')^*U)=\varnothing.
\]
For \(U=\varnothing\),
\[
             g^*(\forall_fU)=\varnothing,\qquad
             \forall_{f'}((g')^*U)=\{*\}.
\]
Thus each automatic inclusion can be strict.
\end{hbexample}

The general-square stability equations and the automatic
comparison directions are discussed in
\cite[Remark~5.6.6, pp.~58--59]{Pitts1995Draft};
Frobenius is equation~(5.10) and Remark~5.5.2, p.~53.
The proof below establishes both identities for every pullback
of set maps.

\subsection{Diagonal equality and the Boolean doctrine}
\label{sec:p10-powerset-equality}

\begin{hbtheorem}{thm:p10-powerset-equality}{proved=annex:p10.quant-set.3}{Powerset equality with parameters}
For a set \(A\), let \(\Delta_A:A\to A\times A\) send
\(a\) to \((a,a)\), and put
\[
              \delta_A=\exists_{\Delta_A}(A)
                         =\{(a,a):a\in A\}\in P(A\times A).
\]
For any \(s,t:I\to A\),
\[
             \langle s,t\rangle^*\delta_A
                        =\{i\in I:s(i)=t(i)\}.
\]
In particular \(\Delta_A^*\delta_A=A\).

Use the left-associated product \((I\times A)\times A\)
and the maps
\[
\begin{aligned}
 e:I\times A&\longrightarrow(I\times A)\times A,
                         &e(i,a)&=((i,a),a),\\
 p:(I\times A)\times A&\longrightarrow I\times A,
                         &p((i,a),b)&=(i,a),\\
 r:(I\times A)\times A&\longrightarrow A\times A,
                         &r((i,a),b)&=(a,b).
\end{aligned}
\]
For \(U\subseteq I\times A\),
\[
                    \exists_e(U)=p^*(U)\cap r^*(\delta_A).
\]
Consequently, for \(V\subseteq(I\times A)\times A\),
\[
       p^*(U)\cap r^*(\delta_A)\subseteq V
                   \quad\Longleftrightarrow\quad U\subseteq e^*(V).
\]
The powerset functor with these predicates is therefore a
Boolean first-order hyperdoctrine with equality, in the precise
sense of Definitions~\ref{def:p10-first-order-hyperdoctrine}
and~\ref{def:p10-doctrine-equality}.
It satisfies the parameter-product Beck--Chevalley identities
and the full parameterized diagonal adjunction, including
empty parameter sets and empty sorts.
\end{hbtheorem}

Pitts gives the parameterized equality condition and the Set
diagonal in
\cite[\S5.6, equation~(5.12), Definition~5.6.1, and
Example~5.6.3, pp.~55--56]{Pitts1995Draft}.
Lawvere describes equality through existential quantification
along the diagonal in
\cite[p.~12]{Lawvere2006AdjointnessReprint}.

The terminal set interpreting a context with no variables is
a singleton \(1\).  It is distinct from an empty sort \(A=\varnothing\).
For such a sort, quantification along \(A\to1\) sends its top
predicate \(A=\varnothing\) to \(\varnothing\) existentially
and to \(1\) universally.
The doctrine above permits these empty sorts; interpretations
of ordinary first-order logic with nonempty domains impose
nonemptiness as an additional condition.

%% Contextual predicate logic: finite variables and finite hypotheses differ.

\subsection{Contextual first-order syntax and its interpretation}
\label{sec:p10-first-order-categorical-semantics}

\begin{hbdefinition}{def:p10-contextual-first-order}{The contextual first-order language and calculus}
Fix a many-sorted signature with a set of sorts, function symbols
\(F:\sigma_1,\ldots,\sigma_n\to\tau\), and relation symbols
\(R\subseteq\sigma_1,\ldots,\sigma_n\), each of finite arity,
including zero. Fix a countably infinite supply of variable names;
their sorts are assigned by declarations. A variable context \(\Gamma\) is a finite list of
distinct variables with declared sorts. Write \(M:\sigma\,[\Gamma]\)
for a well-sorted term: a declared variable has its declared sort,
and \(F(M_1,\ldots,M_n)\) has sort \(\tau\) when its arguments have
the specified sorts in that same context.

Following Pitts, use the formula constructors
\[
 R(\mathbf M),\quad (M=_\sigma N),\quad\mathsf{false},\quad
 \varphi\mathbin{\&}\psi,\quad\varphi\vee\psi,\quad
 \varphi\Rightarrow\psi,\quad
 \forall_\sigma(\varphi),\quad\exists_\sigma(\varphi).
\]
Atomic arguments must have their declared sorts, and equality requires
two terms of the same sort. The propositional constructors preserve the
variable context. In the last two constructors, \(\varphi\) is a
unary formula abstraction: if \(e\) is a formula in
\(\Gamma,x:\sigma\), with \(x\notin\operatorname{var}(\Gamma)\),
write \(\varphi=(x)e\) and \(\varphi(M)=e[M/x]\).
Then \(\varphi(x)\) is its body in \(\Gamma,x:\sigma\), while
\(\forall_\sigma(\varphi)\) and \(\exists_\sigma(\varphi)\) are
formulas in \(\Gamma\). This is Pitts's metalevel abstraction
notation \((x)e\), introduced in \cite[\S3, p.~13]{Pitts1995Draft};
the abstraction is notation for binding and substitution, not an
additional object-language sort or connective.
Put \(\neg\varphi=(\varphi\Rightarrow\mathsf{false})\) and
\(\mathsf{true}=(\mathsf{false}\Rightarrow\mathsf{false})\).
Formulas are taken up to sort-preserving bound alpha-renaming.
Every substitution is simultaneous and capture-avoiding; first freshen
bound variables outside the finitely many contexts and terms involved.

The judgments are term equations \(M=N:\sigma\,[\Gamma]\) and sequents
\[
                         \Phi\vdash\psi\,[\Gamma].
\]
Here \(\Phi\) is a finite list of assumptions, and every assumption
and the conclusion is a formula in \(\Gamma\).
The variable context and the hypothesis list have different roles.
In particular, the judgment retains unused variable declarations.
A theory \(\mathcal T\) is any set of well-formed term-equation and
sequent axioms. A derivation is a finite tree using those axioms and
Sheet~\ref{sheet:contextual-first-order}.
The sheet records the rule package of Pitts's Figures~1--3, 5,
and~7 with variable contexts and side conditions explicit.
Its displayed witness typing \(M:\sigma\,[\Gamma]\) in
\(\exists\)-Intr states the source's implicit well-sorted
instantiation condition; Figure~5 omits that premise from the printed
fraction. The equivalent adjoint presentations in Figures~4, 6,
and~8 are displayed separately in
Section~\ref{sec:p10-adjoint-first-order-calculus}, with their
syntactic interderivability proofs; they are not additional
primitive rules of this presentation.
\index{first-order logic!contextual}\index{variable context}\index{hypothesis list}
\end{hbdefinition}

\showcalculus{contextual-first-order}

The term-equation judgment and the formula \((M=_\sigma N)\) remain
distinct. There is no rule reflecting a proved equality formula into
an equation of term arrows. This is the distinction made in
\cite[\S5.6, Remark~5.6.7, p.~59]{Pitts1995Draft}.

\begin{hbdefinition}{def:p10-contextual-doctrine-interpretation}{Interpreting a contextual first-order signature}
Let \(P:\mathcal C^{\mathrm{op}}\to\mathsf{Heyt}\) be a first-order
hyperdoctrine with equality as in
Definitions~\ref{def:p10-first-order-hyperdoctrine} and
\ref{def:p10-doctrine-equality}.
Interpret each sort \(\sigma\) by an object \(A_\sigma\).
For a finite sort list use the chosen left-associated product starting
with \(1\); likewise put
\[
 I_{()}=1,\qquad I_{\Gamma,x:\sigma}=I_\Gamma\times A_\sigma.
\]
A unary variable context therefore denotes \(1\times A_\sigma\).
Interpret an \(n\)-ary function symbol by an arrow from its argument
product to its result object, and an \(n\)-ary relation symbol by a
predicate in the fibre over that argument product.
At arity zero these are respectively an arrow from \(1\) and a
predicate in \(P(1)\).

Terms denote arrows from \(I_\Gamma\): variables are projections and
\[
 \llbracket F(\mathbf M)\rrbracket_\Gamma
       =F^{\mathcal I}\circ\langle\llbracket M_i\rrbracket_\Gamma\rangle_i.
\]
The tuple into a leading-terminal product includes its unique terminal
component. In particular the empty tuple is \(!:I_\Gamma\to1\).
Formulas denote predicates in \(P(I_\Gamma)\):
\[
\begin{aligned}
 \llbracket R(\mathbf M)\rrbracket_\Gamma
   &=\langle\llbracket M_i\rrbracket_\Gamma\rangle_i^*R^{\mathcal I},\\
 \llbracket M=_\sigma N\rrbracket_\Gamma
   &=\langle\llbracket M\rrbracket_\Gamma,
                         \llbracket N\rrbracket_\Gamma\rangle^*\delta_{A_\sigma},\\
 \llbracket\mathsf{false}\rrbracket_\Gamma&=\bot_{I_\Gamma},\\
 \llbracket\varphi\mathbin{\&}\psi\rrbracket_\Gamma
   &=\llbracket\varphi\rrbracket_\Gamma\wedge\llbracket\psi\rrbracket_\Gamma,\\
 \llbracket\varphi\vee\psi\rrbracket_\Gamma
   &=\llbracket\varphi\rrbracket_\Gamma\vee\llbracket\psi\rrbracket_\Gamma,\\
 \llbracket\varphi\Rightarrow\psi\rrbracket_\Gamma
   &=\llbracket\varphi\rrbracket_\Gamma
                        \Rightarrow\llbracket\psi\rrbracket_\Gamma,\\
 \llbracket\forall_\sigma(\varphi)\rrbracket_\Gamma
   &=\forall_{\pi_{I_\Gamma,A_\sigma}}
                      \llbracket\varphi(x)\rrbracket_{\Gamma,x:\sigma},\\
 \llbracket\exists_\sigma(\varphi)\rrbracket_\Gamma
   &=\exists_{\pi_{I_\Gamma,A_\sigma}}
                      \llbracket\varphi(x)\rrbracket_{\Gamma,x:\sigma}.
\end{aligned}
\]
The implication on the right acts in a Heyting fibre; the one on the
left is the formula constructor.
A term equation is satisfied when its arrows are equal.
A sequent is satisfied when
\[
          \bigwedge_{\phi\in\Phi}\llbracket\phi\rrbracket_\Gamma
                                  \leq\llbracket\psi\rrbracket_\Gamma,
\]
with empty meet \(\top_{I_\Gamma}\).
\end{hbdefinition}

\begin{hbtheorem}{thm:p10-doctrine-substitution}{proved=annex:p10.quant-fol.1}{Terms, formulas, and contextual substitution}
Let \(\mathbf s:\Delta\to\Gamma\) be a tuple of well-sorted terms
in \(\Delta\), one for each variable of \(\Gamma\), and let
\(s:I_\Delta\to I_\Gamma\) be its interpreted tuple arrow.
Substitution preserves well-formedness, respects bound alpha-renaming,
and satisfies
\[
 \llbracket M[\mathbf s]\rrbracket_\Delta
       =\llbracket M\rrbracket_\Gamma\circ s,\qquad
 \llbracket\varphi[\mathbf s]\rrbracket_\Delta
       =s^*\llbracket\varphi\rrbracket_\Gamma.
\]
Adding unused variable declarations is reindexing along their product
projection; consistent renaming and permutation use the corresponding
variable-tuple isomorphisms.
\end{hbtheorem}

\begin{hbtheorem}{thm:p10-doctrine-soundness}{proved=annex:p10.quant-fol.2}{Soundness of contextual intuitionistic first-order logic}
Every interpretation in the stated doctrine satisfying every axiom of
\(\mathcal T\) satisfies every term equation and every contextual
sequent derived from \(\mathcal T\) on the displayed sheet.
The set of theory axioms may be arbitrary; each derivation uses only
finitely many. Sort objects need not be inhabited.
\end{hbtheorem}

The interpretations and stability conditions are
\cite[\S\S5.2--5.6, pp.~40--59, especially Lemma~5.2.6,
Theorems~5.3.2 and~5.4.2]{Pitts1995Draft}.
The proof below treats the displayed calculus directly, without
identifying it with a context-erasing presentation of first-order logic.

\begin{hbtheorem}{thm:p10-doctrine-tarski}{proved=annex:p10.quant-fol.3}{The powerset interpretation and nonempty sorts}
In the powerset doctrine over \(\mathsf{Set}\), the interpretation of
each formula in \(\Gamma\) is exactly the set of finite assignments
to \(\Gamma\) satisfying its usual many-sorted Tarskian clauses,
transported through the canonical product-to-tuple bijection.
This assertion allows empty sort carriers.
Term-equation and sequent satisfaction agree with universal validity
over these finite assignments.

When every interpreted sort is nonempty, these are the ordinary
nonempty-sort Tarskian semantics. An unused variable declaration can
then be removed from the semantic context of either kind of judgment.
For predicate sequents this removal is also valid in any doctrine
for inhabited sort objects, by Theorem~\ref{thm:p10-doctrine-inhabited}.
That abstract statement reflects predicate order; it does not assert
reflection of equality between base arrows.
Without this hypothesis, neither context removal nor validity of
the closed formula \(\exists_\sigma((x)(x=_\sigma x))\) is supplied.
\end{hbtheorem}

For an empty sort in Set, the context \((x:\sigma)\) has empty
interpretation, so even \(\vdash\mathsf{false}\,[x:\sigma]\) is satisfied
in that interpretation. The empty context denotes a singleton, where
\(\vdash\mathsf{false}\,[()]\) is not satisfied.
This is why the retained context is part of the semantic judgment.
The powerset comparison is Example~5.5.3, pp.~53--54, together with
Example~5.6.3, p.~56, of \textcite{Pitts1995Draft}; no completeness
claim for the current handbook LJ calculus is made here.

%% Pitts's alternative contextual first-order presentations, with full proofs.

\subsection{Adjoint presentations of the contextual calculus}
\label{sec:p10-adjoint-first-order-calculus}

\begin{hbdefinition}{def:p10-contextual-adjoint-presentations}{The common basis and two presentations}
Use the language, finite variable contexts, finite hypothesis lists,
and arbitrary theory \(\mathcal T\) of
Definition~\ref{def:p10-contextual-first-order}.
For a binary formula abstraction \(\psi=(x)(x')e\), write
\(\psi(M,N)=e[M/x,N/x']\), with simultaneous capture-avoiding
substitution; \(\psi(x,x)\) denotes its diagonal instance.
The \emph{common basis} consists of the two theory-axiom rules,
Eq-Refl, Eq-Sym, Eq-Trans, Eq-Subst, Id, Cut, Contract, Exchange,
Weaken, and Subst of Sheet~\ref{sheet:contextual-first-order}.
In particular, Exchange exchanges hypotheses, Weaken appends variable
declarations, and the two substitution rules retain their displayed
shared prefix. Contexts remain ordered lists.

For this comparison only, write \(\mathrm{ND}\) for the common basis
together with the propositional, quantifier, and equality-formula
rules of Sheet~\ref{sheet:contextual-first-order}.
Write \(\mathrm{ADJ}\) for the same basis together with all fifteen
directions of Sheet~\ref{sheet:contextual-adjoint-first-order}.
These are local abbreviations for two presentations of derivability.
The notation \(\mathcal T\vdash_{\mathrm{ND}}J\), or
\(\mathcal T\vdash_{\mathrm{ADJ}}J\), means that the judgment \(J\)
has a finite derivation from that same set of theory axioms in the
indicated presentation.
\end{hbdefinition}

Pitts displays seven \emph{bi-rules}: all sequents above a double
line are derivable if and only if the sequent below it is derivable.
This is a statement about derivability, distinct from the formula
constructor \(\Rightarrow\) and from object-language bi-implication.
A bi-rule with two upper sequents gives one downward inference and
two separate upward inferences. The empty upper side of
\(\mathsf{false}\)-Adj gives just one zero-premise inference.
The following sheet displays every such direction of
\cite[Figures~4, 6, and~8, pp.~48, 52, and~55]{Pitts1995Draft}.
The words ``down'' and ``up'' in its rule names refer to the source
double line.

\showcalculus{contextual-adjoint-first-order}

The retained list \(\Phi\) in \(\vee\)-Adj and
\(\exists\)-Adj is part of the displayed source presentation.
Its presence expresses the stability needed when additional
hypotheses are retained. In \(=\)-Adj, \(\Phi\) may depend on
the first equality variable \(x\); only the added variable \(x'\)
is fresh for it. These features agree with the respective semantic
conditions of distributivity, Frobenius, and parameterized equality
in \cite[pp.~48--49, 52--53, 55--56]{Pitts1995Draft}.

\begin{hbtheorem}{thm:p10-contextual-predicate-transport}{proved=annex:p10.adjoint.1}{Syntactic transport of predicate sequents}
Consider any presentation containing the common basis and the
natural-deduction introduction and elimination rules for
\(\Rightarrow\) and \(\forall\).
Let \(\Gamma=(x_1:\sigma_1,\ldots,x_n:\sigma_n)\), let
\(\Phi\vdash\psi\,[\Gamma]\) be derivable from \(\mathcal T\),
and let \(t_i:\sigma_i\,[\Delta]\) be a well-sorted term for
each declaration of \(\Gamma\), in an arbitrary finite context
\(\Delta\). Then
\[
       \Phi[\mathbf t/\mathbf x]\vdash
       \psi[\mathbf t/\mathbf x]\,[\Delta]
\]
is derivable from the same \(\mathcal T\).
Substitution is simultaneous and capture-avoiding.
In particular, consistent variable renaming and permutation of
variable contexts are admissible for predicate sequents.
The statement permits empty \(\Phi\), empty \(\Gamma\), and
arbitrary sets of theory axioms. It asserts no corresponding
transport rule for external term-equation judgments.
\end{hbtheorem}

This syntactic theorem uses implication discharge and universal
closure; it does not treat semantic product isomorphisms as
inference rules or assume that theory axioms are closed under
renaming. The distinction matters because the source contexts and
shared-prefix substitution rules are explicit in
\cite[Definition~2.1.1, p.~6; Figure~1, p.~7;
Remark~3.0.2, p.~14; \S5.1, pp.~40--41]{Pitts1995Draft}.

\begin{hbtheorem}{thm:p10-contextual-adjoint-equivalence}{proved=annex:p10.adjoint.2}{Interderivability of the two contextual presentations}
For every theory \(\mathcal T\) and every well-formed contextual
sequent or external term-equation judgment \(J\),
\[
       \mathcal T\vdash_{\mathrm{ND}}J
       \quad\Longleftrightarrow\quad
       \mathcal T\vdash_{\mathrm{ADJ}}J.
\]
Every direction of each displayed adjoint bi-rule is a derived rule
of \(\mathrm{ND}\), and every logical rule of \(\mathrm{ND}\)
is a derived rule of \(\mathrm{ADJ}\).
Both presentations retain Cut. The correspondence concerns
derivable judgments, without identifying their proof trees.
\end{hbtheorem}

The equivalence asserted in
\cite[pp.~49, 51, and~55]{Pitts1995Draft} is proved here by
explicit finite derivations. Consequently the soundness result
of Theorem~\ref{thm:p10-doctrine-soundness} also applies to
\(\mathrm{ADJ}\), including interpretations with empty sorts.

\begin{hbdefinition}{def:p10-contextual-parameter-free-equality}{The parameter-free equality presentation}
The nearby variant in \cite[Remark~5.6.2, p.~56]{Pitts1995Draft}
restricts the equality bi-rule to an empty retained hypothesis list
and no additional variable context. Denote its two directions
locally by \(\mathrm{Eq0}\), as displayed on the next sheet.
Its binary formula abstraction \(\psi\) has body
\(\psi(x,x')\) in exactly \((x:\sigma,x':\sigma)\), so it
has no additional free parameters. Quantified variables inside
that body are allowed.
\end{hbdefinition}

\showcalculus{contextual-parameter-free-equality}

\begin{hbtheorem}{thm:p10-contextual-equality-reduction}{proved=annex:p10.adjoint.3}{Recovering parameterized equality from the reduced rule}
Over the common basis together with the natural-deduction rules
for implication and universal quantification, the two directions
of full \(=\)-Adj are interderivable with the two directions of
\(\mathrm{Eq0}\).
All other included rules and every theory axiom may be held fixed.
Thus either equality presentation gives the same derivable
judgments under these hypotheses, including when the parameter
context or retained hypothesis list is empty.
No conjunction, disjunction, existential-quantifier rule, or
reflection of formula equality to external term equality is
required in this reduction.
\end{hbtheorem}

The extra implication and universal-quantifier rules in this
theorem are the hypotheses stated in Pitts's remark. They permit
the parameters and hypotheses of the full rule to be represented
inside a formula before the parameter-free rule is applied.


\subsection{Topoi and internal logic}

\begin{hbdefinition}{def:p10-topos}{Elementary topos and subobject classifier}
An elementary topos is a locally small category with finite limits,
exponentials, and a subobject classifier.
The classifier is an object \(\Omega\) with an arrow
\(\mathsf{true}:1\to\Omega\) from a terminal object such that
every monomorphism \(m:S\to X\) has a unique characteristic map
\(\chi_m:X\to\Omega\) making
\[
\begin{array}{ccc}
 S&\xrightarrow{\ !_S\ }&1\\
 {\scriptstyle m}\downarrow&&\downarrow{\scriptstyle\mathsf{true}}\\
 X&\xrightarrow{\ \chi_m\ }&\Omega
\end{array}
\]
a pullback.
Thus \(\chi_m m=\mathsf{true}\,!_S\), and for every
\(f:Y\to X\) with \(\chi_m f=\mathsf{true}\,!_Y\) there is
exactly one \(g:Y\to S\) with \(m g=f\).
The condition classifies the subobject, independently of its
isomorphic domain representation.
This is the finite-limit, exponential, and classifier package in
\cite[Definition~1.2, p.~8]{MoerdijkVanOosten2007Topos};
see also \cite{MacLaneMoerdijk1992,Johnstone2002}.
\index{topos}\index{subobject classifier}\index{Kripke--Joyal semantics}
\end{hbdefinition}

Lawvere's \emph{Quantifiers and Sheaves}, presented at the 1970
congress and published in 1971, lists finite colimits as well
\cite[p.~329]{Lawvere1970}.
The presheaf construction below supplies all small colimits
explicitly, so it satisfies that listed condition too.

%% Small-base presheaves, pointwise constructions, and subobjects.

\subsection{Presheaves and their pointwise constructions}
\label{sec:p10-presheaf-categories}

\begin{hbdefinition}{def:p10-presheaf-category}{The category of presheaves on a small category}
Fix a category \(\mathcal C\) whose object and arrow sets belong
to the fixed universe.  A presheaf on \(\mathcal C\) is a functor
\[
                         X:\mathcal C^{\mathrm{op}}\to\mathsf{Set}.
\]
Thus \(u:d\to c\) induces a restriction function
\(X(u):X(c)\to X(d)\).
For \(v:e\to d\), the functor laws are
\[
                    X(1_c)=1_{X(c)},\qquad
                    X(u v)=X(v)X(u).
\]
A map of presheaves \(\alpha:X\to Y\) is a natural
transformation: its functions \(\alpha_c:X(c)\to Y(c)\) satisfy
\[
                         Y(u)\alpha_c=\alpha_d X(u).
\]
Write
\(\widehat{\mathcal C}=[\mathcal C^{\mathrm{op}},\mathsf{Set}]\)
for their category, with componentwise identities and composition.
The category itself may be large in the fixed universe.

For \(c\in\mathcal C\), the representable presheaf is
\(\mathrm y c=\mathcal C(-,c)\).
An arrow \(u:d\to c\) induces
\[
 \mathrm y(u):\mathrm y d\longrightarrow\mathrm y c,\qquad
 (\mathrm y(u))_e(v)=u v\quad(v:e\to d).
\]
This is the covariant Yoneda assignment on representing objects;
each \(\mathrm y c\) is contravariant in its argument.
\index{presheaf}\index{presheaf category}\index{Yoneda embedding}
\end{hbdefinition}

\begin{hbdefinition}{def:p10-subpresheaf}{Subpresheaves and their order}
A subpresheaf \(S\) of \(X\) is a family of subsets
\(S(c)\subseteq X(c)\) such that
\[
                         X(u)[S(c)]\subseteq S(d)
                         \quad(u:d\to c).
\]
Its restriction maps are the restrictions of those of \(X\).
Write \(\operatorname{Sub}(X)\) for the set of such families,
ordered by componentwise inclusion.
This is a small set: it is a subset of
\[
                       \prod_{c\in\operatorname{Ob}\mathcal C}
                                      \mathcal P(X(c)).
\]
The compatibility conditions select a subset of this product.
The theorem below identifies these families with representatives
of categorical subobjects of \(X\).
\index{subpresheaf}\index{subobject!of a presheaf}
\end{hbdefinition}

\begin{hbtheorem}{thm:p10-presheaf-limits}{proved=annex:p10.psh.1}{Small limits and colimits of presheaves}
The category \(\widehat{\mathcal C}\) is locally small.
It has all small limits and colimits, computed pointwise.
More precisely, for a small diagram
\(D:J\to\widehat{\mathcal C}\), the component at \(c\) of
its limiting cone or colimiting cocone is the corresponding
universal cone or cocone for the set diagram \(j\mapsto D(j)(c)\).
Canonical set constructions supply the restriction maps without
choosing an unrelated limit or colimit at every object.

In particular, the terminal presheaf \(1\) has value \(\{*\}\)
at every object, and the initial presheaf \(0\) has value
\(\varnothing\).
Products have
\((X\times Y)(c)=X(c)\times Y(c)\), and the equalizer of
\(\alpha,\beta:X\to Y\) has component
\[
                         \{x\in X(c):\alpha_c(x)=\beta_c(x)\}.
\]
Coproducts and coequalizers are the corresponding pointwise
disjoint unions and quotients.
If \(\mathcal C\) is empty, its presheaf category has one object
and one arrow; the same constructions give its initial and
terminal object.
\end{hbtheorem}

\begin{hbtheorem}{thm:p10-presheaf-mono-epi}{proved=annex:p10.psh.2}{Maps and subobjects of presheaves}
For a natural map \(m:X\to Y\):
\[
 \begin{aligned}
 m\text{ is monic}
      &\quad\Longleftrightarrow\quad
                    m_c\text{ is injective for every }c,\\
 m\text{ is epic}
      &\quad\Longleftrightarrow\quad
                    m_c\text{ is surjective for every }c,\\
 m\text{ is an isomorphism}
      &\quad\Longleftrightarrow\quad
                    m_c\text{ is bijective for every }c.
 \end{aligned}
\]
Its componentwise images form a subpresheaf
\[
                    \operatorname{Im}(m)(c)=m_c[X(c)]
\]
and give a factorization
\[
                     X\xrightarrow{q}\operatorname{Im}(m)
                                      \xrightarrow{i}Y,
\]
where \(q\) is epic and \(i\) is the inclusion and is monic.
If \(m\) is monic, \(q\) is an isomorphism over \(Y\).

Two monomorphisms into \(Y\) are isomorphic over \(Y\) exactly
when their image subpresheaves agree.  One such monomorphism
factors through another exactly when its image is contained
in the other's image.
Thus \(\operatorname{Sub}(Y)\) gives a canonical small set of
representatives for the ordered categorical subobjects.
A map that is both monic and epic is an isomorphism.
No assertion that every epimorphism has a natural section is made.
\end{hbtheorem}

The pointwise constructions and componentwise map criteria are
stated in \cite[\S1, p.~3]{MoerdijkVanOosten2007Topos}.
The annex supplies the complete constructions and cancellation
proofs, including the converse epimorphism argument through a
pointwise pushout.
For the empty base there are no components; the component
criteria are vacuous and hold for the unique arrow.

%% Exponentials, sieve classifier, and power objects for a small presheaf base.

\subsection{Exponentials of presheaves}
\label{sec:p10-presheaf-exponentials}

Fix the small category \(\mathcal C\) and fixed-universe convention of
Definition~\ref{def:p10-presheaf-category}.
Write \(\widehat{\mathcal C}=[\mathcal C^{\mathrm{op}},\mathsf{Set}]\)
and \(\mathrm y c=\mathcal C(-,c)\).
Products are the pointwise products of
Theorem~\ref{thm:p10-presheaf-limits}.
The presheaves \(0\) and \(1\) have respectively the empty set
and a fixed singleton as every component.

\begin{hbtheorem}{thm:p10-presheaf-exponentials}{proved=annex:p10.psh.3}{The exponential presheaf and its evaluation}
For presheaves \(X,Y\), define
\[
                    (Y^X)(c)=
                    \operatorname{Nat}(\mathrm y c\times X,Y).
\]
For \(u:d\to c\), its restriction is precomposition:
\[
 (Y^X)(u)(\alpha)=\alpha\circ(\mathrm y u\times1_X),\qquad
 [(Y^X)(u)(\alpha)]_e(v,x)=\alpha_e(u\circ v,x)
\]
for \(v:e\to d\) and \(x\in X(e)\).
This defines a presheaf with small component sets.
Its evaluation \(\operatorname{ev}:Y^X\times X\to Y\) has components
\[
                         \operatorname{ev}_c(\alpha,x)
                                      =\alpha_c(1_c,x).
\]
For \(h:Z\times X\to Y\), its unique transpose
\(\Lambda h:Z\to Y^X\) is given by
\[
             [(\Lambda h)_c(z)]_d(u,x)=h_d(Z(u)(z),x)
               \quad(u:d\to c,\ z\in Z(c),\ x\in X(d)).
\]
These are natural transformations.  For every \(k:Z\to Y^X\),
the inverse identities are
\[
 \operatorname{ev}\circ(\Lambda h\times1_X)=h,\qquad
 \Lambda\bigl(\operatorname{ev}\circ(k\times1_X)\bigr)=k.
\]
Thus evaluation and transposition give inverse bijections
\[
 \widehat{\mathcal C}(Z\times X,Y)
                   \cong\widehat{\mathcal C}(Z,Y^X),
\]
natural in \(Z\) and \(Y\).  In particular
\(\widehat{\mathcal C}\) is cartesian closed.
The construction includes empty component sets and the empty base.
\end{hbtheorem}

This is the exponential construction of
\cite[I\S6, Proposition~1, pp.~46--47, formulas (5)--(8)]{MacLaneMoerdijk1992}.
The same component and restriction formulas occur in
\cite[\S1, pp.~4--5]{MoerdijkVanOosten2007Topos}.
The element \(\alpha\in(Y^X)(c)\) specifies compatible functions
at every object mapping into \(c\); its data are not in general
just a function \(X(c)\to Y(c)\).

\subsection{The sieve classifier}
\label{sec:p10-presheaf-classifier}

\begin{hbdefinition}{def:p10-presheaf-sieves}{Sieves and the presheaf of truth values}
A sieve \(S\) on \(c\) is a subset of the arrows with codomain \(c\)
closed under precomposition: if \(u:d\to c\) belongs to \(S\)
and \(v:e\to d\), then \(u\circ v\in S\).
Let \(\Omega(c)\) be the set of all sieves on \(c\).
For \(u:d\to c\), put
\[
             \Omega(u)(S)=
                \{v:e\to d:u\circ v\in S\}.
\]
Write \(t(c)\) for the maximal sieve of all arrows into \(c\).
It is the unique sieve containing \(1_c\).
Define the truth components
\[
                    \mathsf{true}_c:1(c)\to\Omega(c),
                    \qquad \mathsf{true}_c(*)=t(c).
\]
The sieve description is equivalently a subpresheaf of
\(\mathrm y c\): its component at \(d\) is
\(S\cap\mathcal C(d,c)\).
\index{sieve}\index{subobject classifier!presheaf}
\end{hbdefinition}

\begin{hbtheorem}{thm:p10-presheaf-topos}{proved=annex:p10.psh.4}{The sieve classifier makes a presheaf category a topos}
The displayed restrictions make \(\Omega\) a presheaf, and
\(\mathsf{true}:1\to\Omega\) a monomorphism.
For a subpresheaf \(A\subseteq X\), define
\[
       (\chi_A)_c(x)=
              \{u:d\to c:X(u)(x)\in A(d)\}.
\]
This is a natural transformation \(X\to\Omega\), and the square
\[
\begin{array}{ccc}
 A&\longrightarrow&1\\
 \downarrow&&\downarrow{\scriptstyle\mathsf{true}}\\
 X&\xrightarrow{\ \chi_A\ }&\Omega
\end{array}
\]
is a pullback.  It is the unique characteristic map for this
subobject.
For an arbitrary monomorphism \(m:B\to X\), use its canonical
image subpresheaf \(A(c)=m_c[B(c)]\); the same characteristic
map gives a pullback square with \(m\) as its left side.

Hence \(\mathsf{true}\) is a subobject classifier.
Together with the finite limits and exponentials already proved,
this makes \(\widehat{\mathcal C}\) an elementary topos in the
sense of Definition~\ref{def:p10-topos}.
For \(\mathcal C=\varnothing\), it is the degenerate category
with one object and one arrow, and still satisfies this definition.
\end{hbtheorem}

Sieves, their restriction, and the characteristic map are
\cite[I\S4, pp.~37--38, formulas (2)--(8)]{MacLaneMoerdijk1992}.
The characteristic-map uniqueness argument is also given in
\cite[\S1, pp.~6--7]{MoerdijkVanOosten2007Topos}.
Smallness of the base is used throughout: each \(\Omega(c)\)
is a subset of the powerset of the set of arrows into \(c\).

\subsection{Power objects and membership}
\label{sec:p10-presheaf-power-objects}

\begin{hbdefinition}{def:p10-presheaf-power-object}{The power object and its membership relation}
For a presheaf \(W\), write \(\operatorname{Sub}(W)\) for its
subobjects, identified with actual subpresheaves using the
canonical images of
Theorem~\ref{thm:p10-presheaf-mono-epi}.
For \(X\), put
\[
                              P(X)=\Omega^X .
\]
Let \(\in_X\hookrightarrow P(X)\times X\) be the pullback
of \(\mathsf{true}:1\to\Omega\) along evaluation.
It is called the membership subobject.
A power object represents subobjects of products with \(X\):
the following theorem gives that universal property explicitly.
\index{power object!presheaf}\index{membership!power object}
\end{hbdefinition}

\begin{hbtheorem}{thm:p10-presheaf-power-objects}{proved=annex:p10.psh.5}{Classification by a power object}
There are bijections, natural in \(Z\),
\[
 \widehat{\mathcal C}(Z,P(X))
       \cong\widehat{\mathcal C}(Z\times X,\Omega)
       \cong\operatorname{Sub}(Z\times X).
\]
The first is uncurrying, and the second pulls back
\(\mathsf{true}\).  Thus \(k:Z\to P(X)\) classifies exactly
\[
                              (k\times1_X)^*(\in_X).
\]
Conversely, the unique classifying map of a subobject
\(S\hookrightarrow Z\times X\) is \(k=\Lambda(\chi_S)\).
Here a superscript \({}^*\) on subobjects denotes pullback.

At each base object \(c\), classification gives a canonical bijection
\[
                       P(X)(c)\cong
                       \operatorname{Sub}(\mathrm y c\times X).
\]
Under this identification, the restriction along \(u:d\to c\)
sends \(S\subseteq\mathrm y c\times X\) to
\[
 S|u=(\mathrm y u\times1_X)^*S,\qquad
 (S|u)(e)=\{(v,x)\in\mathcal C(e,d)\times X(e):
                                      (u\circ v,x)\in S(e)\}.
\]
The membership components are
\[
 (\in_X)(c)=
   \{(S,x):S\in\operatorname{Sub}(\mathrm y c\times X),\
                         x\in X(c),\ (1_c,x)\in S(c)\},
\]
where the first coordinate uses the same canonical identification.
For \(S\subseteq Z\times X\), its classifying map sends
\(z\in Z(c)\) to the subpresheaf
\[
        S_z(d)=\{(u,x)\in\mathcal C(d,c)\times X(d):
                                      (Z(u)(z),x)\in S(d)\}.
\]
These formulas include \(X=0\) and \(\mathcal C=\varnothing\).
All displayed component and subobject collections are small sets.
\end{hbtheorem}

Power objects as \(\Omega^X\), their natural classification, and
the membership formula are described in
\cite[\S1, pp.~7--8, Exercise~4]{MoerdijkVanOosten2007Topos}.
The annex supplies the restriction, membership, and classification
proofs.

\begin{hbexample}{def:p10-presheaf-not-pointwise}{Exponential and power objects depend on restrictions}
Let \(\mathcal C\) have objects \(0,1\) and one nonidentity arrow
\(a:0\to1\).
Let \(X(0)=\{*\}\), \(X(1)=\varnothing\), with its unique restriction.
Then \(\mathrm y1\) is the terminal presheaf.
There is no natural map \(\mathrm y1\times X\to0\), because its
component at \(0\) would map a singleton to the empty set.
Consequently
\[
                     (0^X)(1)=\varnothing,
\]
whereas the external function set
\(\mathsf{Set}(X(1),\varnothing)\) is a singleton.
Also \(\mathrm y1\times X\) has exactly two subpresheaves,
so \(P(X)(1)\) has two elements, whereas
\(\mathcal P(X(1))=\mathcal P(\varnothing)\) has one.
\end{hbexample}


The internal logic of these topoi is generally intuitionistic.
At base-category stages, presheaf disjunction and existence use
respectively a disjunct or a witness at the current stage, whereas
implication and universal quantification test all restriction arrows.
Generalized-element Kripke--Joyal semantics instead permits
epimorphic locality, and site/sheaf semantics uses covering sieves
\cite[\S1.2, p.~13; \S2.2, p.~32; \S3.3, pp.~58--59]
{MoerdijkVanOosten2007Topos}.
Sheaf and Boolean-valued
models can establish relative independence results, including presentations
of the continuum hypothesis, only after the base topos, site, internal
statement, and transfer back to set theory are specified. This pointer does not
replace the forcing development in Part VI.

%% ED26-REQ-N-d-023-01: complete original source account; proof residuals retained.
\subsection{Lawvere's quantifiers and sheaves: Nice 1970, published 1971}
\label{sec:p10-n-d-023-source-01}

This is the presentation of \emph{Quantifiers and Sheaves}, the six-page
Nice congress article published in the 1971 proceedings
\cite[pp.~329--334]{Lawvere1970}.  Its account of joint work with Myles
Tierney puts geometry first, treats a topology as an internal modality,
and transports constructions for abstract sets to a variable geometric
setting.  Its two applications are Boolean-valued models, including a
continuum-hypothesis argument without transfinite induction in the
displayed construction, and the Chevalley--Hakim global spectrum without
choosing an external site.  These are the article's announced program and
source-specific constructions, not a replacement of set-theoretic
independence by the assertion that every topos is a model of ZFC.

The introductory motivation concerns sets varying with time, space, or a
ring of definition, rather than sets characterized solely by their
elements over the terminal object.  Lawvere interprets opposing
mathematical tendencies through adjoint functors and criticizes reliance
on an externally accumulated power-set hierarchy.  He emphasizes that
such geometric sets usually have intuitionistic logic, need not satisfy
choice, and need not be determined by their global points
\cite[p.~329]{Lawvere1970}.  This motivation does not remove any of the
structural hypotheses in the following definition.

\subsubsection{The original topos package and its internal adjoints}

\begin{hbdefinition}{def:p10-lawvere-sheaves-topos}{Lawvere's original topos package}
In this presentation a topos \(\mathcal E\) has finite limits, finite
colimits, exponentials, and a subobject classifier
\(\mathsf{true}:1\to T\).  For each object \(A\), the functor
\((-)^A\) is right adjoint to \((-)\times A\).  For each monomorphism
\(m:X'\hookrightarrow X\), there is exactly one
\(\chi_m:X\to T\) for which
\[
\begin{array}{ccc}
X'&\longrightarrow&1\\
\downarrow{\scriptstyle m}&&\downarrow{\scriptstyle\mathsf{true}}\\
X&\xrightarrow{\ \chi_m\ }&T
\end{array}
\]
is a pullback.  For \(\varphi:X\to T\), Lawvere writes
\(\{X/\varphi\}\) for the classified subobject.  The source's
truth-value object is \(T\), not the \(\Omega\) used in
Definition~\ref{def:p10-topos}; finite colimits are explicitly part of
this original list of assumptions \cite[pp.~329--330]{Lawvere1970}.
\index{Lawvere, F. William!quantifiers and sheaves}
\index{topos!Lawvere's original package}
\end{hbdefinition}

The article states that \(T\) is an internal Heyting algebra.  Its
completeness is indexed by objects and maps of \(\mathcal E\), not an
assertion that \(\mathcal E\) has every externally indexed limit.  For
\(f:X\to Y\), precomposition induces
\[
 T^f:T^Y\longrightarrow T^X,
 \qquad \exists_f\dashv T^f\dashv\forall_f,
 \qquad \exists_f,\forall_f:T^X\longrightarrow T^Y.
\]
These are order adjunctions between the internal predicate objects.
On classified subobjects the corresponding conditions are
\[
 \exists_f A\leq B\ \Longleftrightarrow\ A\leq f^*B,
 \qquad f^*B\leq A\ \Longleftrightarrow\ B\leq\forall_f A.
\]
For a topological space the sheaf \(T\) has, over an open \(U\),
the open subsets of \(U\), with restriction by intersection.  For
set-valued presheaves on a small \(\mathcal C\), its value at \(c\)
is the set of cribles, or sieves, on \(c\).  Neither example is
generally Boolean.  Formula classification gives actual subobjects for
the appropriate higher-order formulas; it is not a claim that their
truth is determined by a Boolean value at every global point
\cite[p.~330]{Lawvere1970}.

\begin{hbdefinition}{def:p10-lawvere-sheaves-families}{Families, support, and internal presheaves}
An object of \(\mathcal E/X\) is a map \(p:E\to X\), regarded as
an \(X\)-indexed family.  Pullback along \(f:X\to Y\) is
\(f^*:\mathcal E/Y\to\mathcal E/X\).  Its right adjoint, when
written in the source's family notation, is
\[
 \Pi_f:\mathcal E/X\longrightarrow\mathcal E/Y.
\]
The support \(\sigma_X(p)\) is the characteristic predicate of the
image of \(p\); it belongs to \(\mathcal E(X,T)\).

An internal category has an object of objects \(C_0=X\), an object of
arrows \(C_1\), domain, codomain, identity, and composition maps
satisfying the category equations.  A contravariant action on a family
\(E\to X\) gives restriction along each internal arrow, with identity
and composition laws.  Lawvere denotes the resulting category of
internal \(\mathcal E\)-valued presheaves by
\(\widehat C=C^{\mathrm{op}}\mathcal E\).  An internal functor
\(f:C\to C'\) induces restriction
\(f^*:\widehat{C'}\to\widehat C\); the source denotes its
left and right adjoints by \(f_0\) and \(f_*\).
\end{hbdefinition}

\begin{hbtheorem}{thm:p10-lawvere-sheaves-internal-adjoints}{proved=annex:p10.lawvere-sheaves.internal}{Lawvere's announced internal completeness}
For a topos with the original package, the slice pullback functors have
the displayed right adjoints \(\Pi_f\).  The internal presheaf
category \(\widehat C\) is again a topos, and restriction along an
internal functor has both adjoints \(f_0\) and \(f_*\).
The source uses these constructions to obtain the usual exactness
properties and to explain how even an externally countable topos can
have this form of internal completeness \cite[p.~330]{Lawvere1970}.
The original article announces these conclusions without giving their
general constructions or complete proofs.  The annex records that
historical boundary while supplying separate editorial deductions
through the general constructions below.  The concrete
external-small-base case also retains its distinct detailed proofs.
\end{hbtheorem}

\begin{hbdefinition}{def:p10-lawvere-sheaves-functional-graphs}{Power objects and fiberwise functional graphs}
For any object \(W\), the power object is \(T^W\), with
membership the pullback of \(\mathsf{true}\) along evaluation.
Classifier and exponential transposition give, naturally in \(Z\),
\[
 \mathcal E(Z,T^W)\cong\mathcal E(Z\times W,T)
                         \cong\operatorname{Sub}(Z\times W).
\]
These bijections use the original package, not external powersets.
Let \(\delta_W:W\times W\to T\) classify equality.  Let
\(\land:T\times T\to T\) classify the intersection of the
two truth subobjects.  Thus \(a\leq b\) for predicates means
\(a\land b=a\); this equation can be imposed by an equalizer.

For \(p:E\to X\) and \(f:X\to Y\), put
\(D_0=Y\times T^{X\times E}\), with universal membership
predicate \(r(y,R,x,e)\).  Select \(D_1\hookrightarrow D_0\)
by the following equations of predicates, imposed after transposing
in the displayed \(x,e,e'\) variables:
\[
\begin{aligned}
 r(y,R,x,e)\land
   \delta_X(x,p(e))\land\delta_Y(f(x),y)&=r(y,R,x,e),\\
 r(y,R,x,e)\land r(y,R,x,e')\land\delta_E(e,e')
                    &=r(y,R,x,e)\land r(y,R,x,e').
\end{aligned}
\]
They say that the relation lies over the fiber of \(f\), lies
over \(p\), and is single-valued.  They use finite equalizers,
not an assumed quantifier or an assumed unique-choice theorem.
Let \(R_1\hookrightarrow D_1\times X\times E\) be the
restricted membership relation.  Its projection
\(d:R_1\to D_1\times X\) is monic.  Write \(\chi_d\)
for its classifier, and write
\(\gamma(y,R,x)=\delta_Y(f(x),y)\) for the classifier of
\(D_1\times_Y X\hookrightarrow D_1\times X\).
Define \(D\hookrightarrow D_1\) as the equalizer of their
transposes into \(T^X\):
\[
                    D=\operatorname{Eq}
                       (\Lambda\chi_d,\Lambda\gamma).
\]
This final condition makes the functional graph's domain exactly
the fiber, not merely a dense or inhabited part of it.
\end{hbdefinition}

\begin{hbtheorem}{thm:p10-lawvere-sheaves-dependent-product}{proved=annex:p10.lawvere-sheaves.dependent-product}{Dependent products from the original power-object package}
The object \(D\to Y\) above, denoted \(\Pi_f(p)\), has
a canonical evaluation or counit
\[
             \varepsilon_p:D\times_Y X\longrightarrow E,
             \qquad p\varepsilon_p=\operatorname{pr}_X.
\]
For every \(a:Z\to Y\), taking the graph of a map and
evaluating a classified graph give inverse bijections
\[
 \mathcal E/Y(Z,\Pi_fp)
             \cong\mathcal E/X(Z\times_Y X,E),
\]
natural in both the family \(p\) and the test family \(Z\).
Thus \(f^*\dashv\Pi_f\), without an image, exactness, or
choice hypothesis in the construction.

For every pullback square
\[
\begin{array}{ccc}
 X'&\xrightarrow{\ a\ }&X\\
 \downarrow{\scriptstyle f'}&&\downarrow{\scriptstyle f}\\
 Y'&\xrightarrow{\ b\ }&Y,
\end{array}
\]
there is a canonical natural isomorphism
\[
                  b^*\Pi_f(p)\cong\Pi_{f'}(a^*p).
\]
It is the mate whose evaluation is the pulled-back counit, and
it is compatible with identity squares and pasting pullback
squares.  The conclusion concerns actual pullbacks, not arbitrary
commuting squares.  For composable maps,
\(\Pi_{gf}\cong\Pi_g\Pi_f\), and \(\Pi_{1_X}\cong1\).
Every slice \(\mathcal E/X\) satisfies the original topos
package; in particular, its exponential of \(F\to X\) by
\(E\xrightarrow{p}X\) is \(\Pi_p(p^*F)\).
This is an editorial complete construction of the slice assertion
announced on p.~330, not a proof printed in the original article
\cite[p.~330]{Lawvere1970}.
\end{hbtheorem}

\begin{hbtheorem}{thm:p10-lawvere-sheaves-predicate-adjoints}{proved=annex:p10.lawvere-sheaves.predicate-adjoints}{Internal Heyting predicates, stable images, and both quantifiers}
From the original package and the constructed dependent products,
the following statements hold without Boolean negation or choice.
Every map has an epimorphism--monomorphism image factorization,
and these images are stable under arbitrary pullback.  Subobjects
of each \(X\) form a Heyting algebra: intersections are meets,
images of \(A+B\to X\) are joins, and implication is classified
by the truth-value predicate of
\[
       \operatorname{Eq}(\land,\operatorname{pr}_1)
                                        \hookrightarrow T\times T.
\]
Thus \(D\leq(A\Rightarrow B)\) exactly when
\(D\cap A\leq B\), and \(T\) itself is an internal
Heyting algebra.  No complement operation is assumed.

For \(f:X\to Y\) and \(A\hookrightarrow X\),
\(\exists_fA\) is the image of \(A\to Y\), and
\(\forall_fA\) is \(\Pi_f(A\to X)\), regarded as a
subobject of \(Y\).  They satisfy
\[
 \exists_fA\leq B\ \Longleftrightarrow\ A\leq f^*B,
 \qquad f^*B\leq A\ \Longleftrightarrow\ B\leq\forall_fA.
\]
For every pullback square of the preceding theorem,
\[
 b^*\exists_fA=\exists_{f'}a^*A,\qquad
 b^*\forall_fA=\forall_{f'}a^*A,
\]
and \(\exists_f(A\cap f^*B)=\exists_fA\cap B\).

These operations are represented internally by actual maps
\(\exists_f,\forall_f:T^X\to T^Y\).  If
\(r:Z\to T^X\) classifies \(R\hookrightarrow Z\times X\),
their composites with \(r\) classify, respectively,
\[
 \operatorname{Im}(R\to Z\times Y),\qquad
 \Pi_{1_Z\times f}(R\hookrightarrow Z\times X).
\]
The internal order adjunctions therefore hold at every parameter
object \(Z\), not just on the global subobject posets.  Empty
fibers are retained: universal quantification of top is top,
and existential image of the initial subobject is initial.
This proves editorially the indexed completeness of \(T\)
announced on p.~330 \cite[p.~330]{Lawvere1970}.
\end{hbtheorem}

\begin{hbtheorem}{thm:p10-lawvere-sheaves-effective-quotients}{proved=annex:p10.lawvere-sheaves.effective-quotients}{Regular epimorphisms and effective internal equivalence relations}
Every epimorphism in a topos with the original package is the
coequalizer of its kernel pair.  The image factorizations are
therefore regular-epi--mono factorizations stable under pullback.
For any internal equivalence relation
\(R\hookrightarrow X\times X\), transpose its classifier
to the row map \(r:X\to T^X\), and factor it by its
constructed image as
\[
                       X\twoheadrightarrow Q\hookrightarrow T^X.
\]
The first map has kernel pair exactly \(R\) and coequalizes
it universally.  Thus the category is regular and its internal
equivalence relations are effective.  These are explicit exactness
deductions from the constructed package, not exactness assumed in
order to obtain \(\Pi_f\) or its images
\cite[p.~330]{Lawvere1970}.
\end{hbtheorem}

\begin{hbdefinition}{def:p10-lawvere-sheaves-internal-actions}{Internal presheaves with their typed action data}
Let \(C\) be an internal category, with objects \(C_0\), arrows
\(C_1\), source and target \(s,t:C_1\to C_0\), identities
\(i:C_0\to C_1\), and composition \(a b=a\circ b\)
defined when \(t(b)=s(a)\).  All these are maps in
\(\mathcal E\) satisfying the category equations.  An internal
presheaf is a family \(E\to C_0\) with a map
\[
  E\times_{C_0,t}C_1\longrightarrow E,
             \qquad(e,a)\longmapsto e\cdot a,
\]
lying over \(s(a)\), with
\(e\cdot i(c)=e\) and
\((e\cdot a)\cdot b=e\cdot(a b)\).
A morphism is a family map preserving this action.  Write
\(U_C:C^{\mathrm{op}}\mathcal E\to\mathcal E/C_0\)
for the forgetful functor.

The generalized-element notation in the following constructions
denotes actual dependent families: a family of values over the
fiber of \(q\) is \(\Pi_q\) of that family, and a universal
compatibility condition is \(\forall_q\) of the corresponding
equality subobject.  The two preceding proofs construct both.
Thus none of these expressions presupposes an external infinite
product or quantification only over global points.
\end{hbdefinition}

\begin{hbtheorem}{thm:p10-lawvere-sheaves-internal-presheaf-topos}{proved=annex:p10.lawvere-sheaves.internal-presheaf-topos}{The topos of internal presheaves}
For an internal category \(C\), the category
\(C^{\mathrm{op}}\mathcal E\) satisfies the original topos
package.  The functor \(U_C\) creates finite limits and finite
colimits and detects monomorphisms.  It has a free left adjoint
\(L_C\) and a cofree right adjoint \(R_C\), with underlying
families
\[
             U_CL_CB=\Sigma_s(t^*B),\qquad
             U_CR_CB=\Pi_t(s^*B),
\]
where \(\Sigma_s\) means composition of a family map with
\(s\).  On generalized fibers, \(L_CB\) consists of
\((a:d\to c,b\in B_c)\) at \(d\), with precomposition
action; \(R_CB\) at \(c\) consists of all assignments
\((b_a\in B_d)_{a:d\to c}\), also restricted by
precomposition.

For presheaves \(G,F\), their exponential at an internal object
\(c\) consists of the compatible families
\[
 \bigl(\phi_a:G_d\to F_d\bigr)_{a:d\to c},
 \qquad
 \phi_{a b}(g\cdot b)=\phi_a(g)\cdot b
       \quad(b:e\to d).
\]
Evaluation is \((\phi,g)\mapsto\phi_{i(c)}(g)\);
restriction along \(u:c'\to c\) replaces \(\phi_a\)
by \(\phi_{u a}\).
The classifier \(T_C\) is the subfamily of
\(R_C(C_0\times T)\) consisting of assignments
\((v_a\in T)_{a:d\to c}\) with
\[
                         v_a\leq v_{a b}
                              \quad(b:e\to d).
\]
Its truth assignment is constantly true.  A subpresheaf
\(A\hookrightarrow E\) has unique characteristic assignment
\[
                    (\chi_A^C(e))_a=\chi_A(e\cdot a),
\]
where the right-hand classifier is that of its underlying family.
These formulas define internal objects, restrictions, evaluation,
and unique pullback classification, not only fiberwise candidates.
Power objects are \(T_C^G\) with evaluation membership.
This is an editorial proof of the internal-presheaf assertion
announced on p.~330 \cite[p.~330]{Lawvere1970}.
\end{hbtheorem}

\begin{hbtheorem}{thm:p10-lawvere-sheaves-internal-kan}{proved=annex:p10.lawvere-sheaves.internal-kan}{Both adjoints of restriction along an internal functor}
Let \(F:C\to D\) be an internal functor.  Restriction
\(F^*:D^{\mathrm{op}}\mathcal E\to C^{\mathrm{op}}\mathcal E\)
pulls underlying families back along \(F_0\) and acts along
\(F_1\).  For \(A\) on \(C\), its right adjoint has
generalized fibers
\[
 (F_*A)_d=
 \left\{(a_{c,b}\in A_c)_{b:F(c)\to d}:
 a_{c,bF(u)}=a_{c',b}\cdot u
            \text{ for }u:c\to c',\ b:F(c')\to d\right\}.
\]
Dependent products construct the assignments, and equality
subobjects and universal quantification impose compatibility.
Restriction along \(v:d'\to d\) sends the coordinate
\((c,b:F(c)\to d')\) to the old \((c,vb)\) coordinate.

For the left adjoint, form internal \(D\)-presheaves \(B_0,B_1\)
whose generalized elements over \(d\) are respectively
\[
\begin{aligned}
 B_0(d)&:\quad(c,x\in A_c,b:d\to F(c)),\\
 B_1(d)&:\quad(u:c\to c',x'\in A_{c'},b:d\to F(c)).
\end{aligned}
\]
Their action is precomposition of \(b\).  The two maps
\(B_1\rightrightarrows B_0\) send the latter tuple to
\[
              (c,x'\cdot u,b),\qquad(c',x',F(u)b).
\]
Define \(F_0A\) as their coequalizer in the internal presheaf
topos.  This uses its proved finite-colimit construction, not an
assumed externally indexed coend or an assumed effective quotient.
The presheaf argument distinguishes this left-extension notation
from the object component \(F_0:C_0\to D_0\) of the internal
functor.
There are natural inverse bijections
\[
 \operatorname{Hom}_D(F_0A,B)\cong\operatorname{Hom}_C(A,F^*B),
 \qquad
 \operatorname{Hom}_C(F^*B,A)\cong\operatorname{Hom}_D(B,F_*A).
\]
Thus \(F_0\dashv F^*\dashv F_*\), with actual unit/counit
maps and both triangle identities furnished by the inverse
constructions in the annex.  Empty internal index families are
included.  This supplies editorially the internal-functor adjoints
announced on p.~330 \cite[p.~330]{Lawvere1970}.
\end{hbtheorem}

\begin{hbtheorem}{thm:p10-lawvere-sheaves-presheaf-bridge}{proved=annex:p10.lawvere-sheaves.presheaf}{The concrete presheaf instance of the original package}
Fix the small-base and universe hypotheses of
Definition~\ref{def:p10-presheaf-category}.  In
\(\mathcal E=[\mathcal C^{\mathrm{op}},\mathsf{Set}]\), finite limits
and finite colimits are pointwise; \(T=\Omega\), with
\[
 T(c)=\{\text{sieves on }c\},\qquad
 T(u)(S)=\{v:e\to d:u v\in S\}\quad(u:d\to c).
\]
For \(A\subseteq X\), the characteristic map is uniquely
\[
 (\chi_A)_c(x)=\{u:d\to c:X(u)x\in A(d)\}.
\]
The exponential and power object are
\[
 (Y^X)(c)=\operatorname{Nat}(\mathrm y c\times X,Y),\qquad
 T^X(c)\cong\operatorname{Sub}(\mathrm y c\times X).
\]
Exponential restriction is precomposition by
\(\mathrm y u\times1_X\); under the power-object identification
the restriction is pullback by that map.  Evaluation is
\(\operatorname{ev}_c(\alpha,x)=\alpha_c(1_c,x)\), and membership
is \((S,x)\in\in_X(c)\) exactly when \((1_c,x)\in S(c)\).
All restriction, naturality, classification, and uniqueness conditions
hold, including for the empty base.  Thus this category satisfies the
finite-colimit condition in Lawvere's package as well as the host
elementary-topos package.  This is a source-to-host identification of
the concrete example, not attribution of the host's later detailed
proofs to the six-page article.
\end{hbtheorem}

\subsubsection{Why support, dependent product, and choice are different}

\begin{hbdefinition}{def:p10-lawvere-sheaves-choice-conditions}{The two source conditions}
Lawvere separates the following assertions.
\[
\begin{aligned}
(\Pi):\quad&
 \sigma_Y(\Pi_f p)=\forall_f\sigma_X(p)
 &&\text{for every }f:X\to Y\text{ and }p:E\to X;\\
(\exists):\quad&
 E\twoheadrightarrow\{1/\sigma_1(E\to1)\}
 &&\text{has a section for every }E.
\end{aligned}
\]
The first says that taking support commutes with dependent product.
The second splits the epimorphism onto an object's support, not every
epimorphism onto every base.  In the category of formulas of a higher
theory the source calls the latter a Skolem condition.  In geometry an
existential predicate may be true only locally; it need not furnish a
global section \cite[pp.~330--331]{Lawvere1970}.
\index{axiom of choice!support and dependent product}
\index{Skolem condition!Lawvere}
\end{hbdefinition}

\begin{hbtheorem}{thm:p10-lawvere-sheaves-choice}{proved=annex:p10.lawvere-sheaves.choice}{The conjunction, not either condition alone, expresses choice}
In a topos with the slice and predicate adjoints above, the conjunction
of \((\Pi)\) and \((\exists)\) is equivalent to every epimorphism
having a section.  Lawvere calls this geometrically
zero-dimensionality and logically the axiom of choice
\cite[p.~330]{Lawvere1970}.  The annex gives an editorial complete
proof of this equivalence, rather than asserting that the article
prints that proof.
\end{hbtheorem}

\begin{hbexample}{def:p10-lawvere-sheaves-choice-examples}{The source's group and ordered-diagram comparisons}
For permutation representations of a group \(G\), condition
\((\Pi)\) holds: at a base point the dependent product consists of
sections over its fiber, and its support says all the fiberwise
supports hold.  This comparison uses choice in the external category
of sets.  If \(G\) is nontrivial, its regular action has support
\(1\) but no fixed point; its map to \(1\) has no equivariant
section.  Hence \((\exists)\) fails.

For covariant set-valued diagrams on a well-ordered poset \(P\),
\((\exists)\) holds: an inhabited diagram first becomes inhabited
at a least stage, and one element there propagates to all later
stages in its support.  These are the source's \(P\mathsf{Set}\),
equivalently presheaves on \(P^{\mathrm{op}}\), not presheaves
on \(P\) with the opposite restrictions.

The source specifically says \((\Pi)\) fails for the arrow category
\(2\mathsf{Set}\).  Here is an explicit reconstruction.  In the
covariant diagram category \(0\to1\), put
\[
 X=(\{a,b\}\to\{c\}),\qquad
 E=(\{e_a,e_b\}\to\{u,v\}),
\]
where the first map sends both elements to \(c\), the second sends
\(e_a\) to \(u\) and \(e_b\) to \(v\), and the family
\(p:E\to X\) has \(p_0(e_a)=a\), \(p_0(e_b)=b\),
\(p_1(u)=p_1(v)=c\).  It is pointwise surjective.  A section
would have to choose both \(u\) and \(v\) as the image of \(c\),
so none exists.  For \(f:X\to1\), the dependent product has no
element at stage \(0\) but has two at stage \(1\), whereas
\(\forall_f\sigma_X(p)=\mathsf{true}\) at both stages.
The annex verifies these comparisons and their variance
\cite[p.~330]{Lawvere1970}.
\end{hbexample}

\subsubsection{Recursion, geometric morphisms, and associated sheaves}

Lawvere's additional condition \((\omega)\) concerns the forgetful
functor from objects with an endomap to objects:
\[
 \mathcal E^{\Omega}\longrightarrow\mathcal E.
\]
Here the source's \(\Omega\) encodes an endomap; it is not its
truth-value object \(T\).  The functor is required not to be an
equivalence and to have the free-recursion left adjoint
\((-)\times\omega\).  The endomap on the free object is successor
on the \(\omega\) factor.  Truth-value-valued sequences give the
induction case.  The article invokes this condition for coequalizer
analysis and free group or free ring constructions, but intentionally
does not put it in the topos definition.  It announces that it lifts
to a topos defined geometrically over a base satisfying it
\cite[p.~331]{Lawvere1970}.  It supplies no general free-object
construction proof here.

\begin{hbdefinition}{def:p10-lawvere-sheaves-morphisms}{The source's geometric and logical morphisms}
A geometric morphism \(u:\mathcal E'\to\mathcal E''\) is given
by a functor \(u_*:\mathcal E'\to\mathcal E''\) with an exact
left adjoint \(u^*:\mathcal E''\to\mathcal E'\); the left
adjoint preserves finite limits.  A logical morphism preserves, up
to isomorphism, the structures in the original topos package.  The
article uses \emph{local homeomorphism} for a geometric morphism
whose inverse-image part is logical.  This records its stated
usage, without importing a different modern characterization under
the same name \cite[p.~331]{Lawvere1970}.
\index{geometric morphism}\index{logical morphism}
\end{hbdefinition}

\begin{hbtheorem}{thm:p10-lawvere-sheaves-factorization}{proved=annex:p10.lawvere-sheaves.sheaves}{The announced geometric-morphism factorization}
Every geometric morphism
\(u:\mathcal E'\to\mathcal E''\) factors through a topos
\(\mathcal E\) as \(u''u'\), where
\((u'')_*:\mathcal E\to\mathcal E''\) is full and faithful
and \((u')^*:\mathcal E\to\mathcal E'\) reflects
isomorphisms.  The full-and-faithful geometric part is determined
by a topology \(j_u:T''\to T''\), and its domain is the
category of \(j_u\)-sheaves \cite[p.~331]{Lawvere1970}.
The source gives the associated-sheaf route below, not a complete
factorization proof.  The annex supplies a separate editorial
construction using the proved reflector and subobject adjunction
of the given morphism; that does not alter the historical source
standing.  No uniqueness statement stronger than the source's
topology determination is added.
\end{hbtheorem}

\begin{hbdefinition}{def:p10-lawvere-sheaves-topology}{An internal topology and its closure}
The source's topology is a map \(j:T\to T\) satisfying
\[
 j j=j,\qquad j\mathsf{true}=\mathsf{true},\qquad
 j\circ\land=\land\circ(j\times j).
\]
For \(A\hookrightarrow X\) classified by \(\chi_A\), its
\(j\)-closure is classified by \(j\chi_A\).  It is
\(j\)-closed when \(j\chi_A=\chi_A\), and \(j\)-dense
when its closure is \(X\).  An object is separated when its
diagonal is \(j\)-closed.  A \(j\)-sheaf has a unique extension
of each map along a \(j\)-dense monomorphism.  This unpacks the
usual sheaf condition used by the article.

The closure is not Kuratowski closure of subsets of a space.  In
presheaves on the open-set poset, the relevant operator sends an
order ideal of opens to the principal ideal generated by their
union \cite[p.~331]{Lawvere1970}.
\index{Lawvere--Tierney topology}\index{associated sheaf}
\end{hbdefinition}

\begin{hbtheorem}{thm:p10-lawvere-sheaves-closure}{proved=annex:p10.lawvere-sheaves.closure}{The closure properties of the source topology}
The closure defined above is extensive, monotone, idempotent,
preserves finite intersections, and commutes with pullback.
In particular the source's three displayed axioms yield a genuine
universal closure operator on subobjects; an additional arbitrary
external powerset operation is not being assumed.
\end{hbtheorem}

The article's associated-sheaf argument has four explicit observations:
the image \(T_j\) of \(j\) is a sheaf; \(Y^X\) is a sheaf
when \(Y\) is; the closure of \(\Delta_X\) in \(X\times X\)
is an equivalence relation; and, for a monomorphism of \(X\) into
a sheaf \(Y\), the closure of \(X\) in \(Y\) is its associated
sheaf.  Before applying them it uses the singleton map
\(\{\}:X\to T^X\).  It also characterizes sheaves as the
objects that are separated and closed in every separated object
into which they embed.  The reflector is thus obtained without
infinite direct limits, and its exactness is studied through the
maps it inverts.  The annex retains this sequence and explains
the complete editorial deductions below supply those steps; the
source's printed list is not itself a full proof
\cite[p.~331]{Lawvere1970}.

\begin{hbdefinition}{def:p10-lawvere-sheaves-associated-construction}{The closed-singleton construction}
Write \(T_j\hookrightarrow T\) for the equalizer of \(j\)
and \(1_T\); the next theorem identifies it with the source's
image of \(j\).  Put
\[
 R_X=\overline{\Delta_X}\hookrightarrow X\times X,
 \qquad \kappa_X:X\longrightarrow T_j^X,
\]
where \(\kappa_X\) transposes the closed equality predicate
\(j\delta_X\).  It is the \(j\)-closed companion of the
source's singleton map.  Factor it by its already constructed image:
\[
 X\xrightarrow{q_X}S_X\hookrightarrow T_j^X,
 \qquad a_jX=\overline{S_X}^{T_j^X}.
\]
Write \(\eta_X:X\to a_jX\) for the image map followed by
the closure inclusion.  These are finite/internal constructions
using the proved power objects, images and closure, not an
assumed sheaf reflector or external transfinite iteration.
\end{hbdefinition}

\begin{hbtheorem}{thm:p10-lawvere-sheaves-associated-sheaf}{proved=annex:p10.lawvere-sheaves.associated-sheaf}{The four observations and the actual associated sheaf}
The object \(T_j\) is the image of \(j\) and is a
\(j\)-sheaf.  If \(Y\) is a sheaf, \(Y^X\) is a
sheaf for every \(X\).  The relation \(R_X\) is an
equivalence relation and is the kernel pair of \(\kappa_X\)
and \(q_X\); \(S_X\) is separated and is the universal
separated image of \(X\).

The object \(a_jX\) is a sheaf, and precomposition with
\(\eta_X\) gives natural bijections
\[
       \mathcal E(a_jX,Y)\cong\mathcal E(X,Y)
                           \quad(Y\text{ a }j\text{-sheaf}).
\]
Thus the construction defines a functor left adjoint to the
full inclusion of sheaves.  In particular, if
\(X\hookrightarrow Y\) is a mono into any sheaf, its
closure in \(Y\) is an associated sheaf of \(X\).
An object is a sheaf exactly when it is separated and closed
in every separated object into which it embeds.  The assertions
and all their extension/uniqueness arguments hold at arbitrary
generalized parameter objects; no global-point sufficiency or
choice principle is used.  This is a complete editorial proof
of the source's four observations and characterization, whose
printed treatment is a sketch \cite[p.~331]{Lawvere1970}.
\end{hbtheorem}

\begin{hbtheorem}{thm:p10-lawvere-sheaves-sheaf-topos}{proved=annex:p10.lawvere-sheaves.sheaf-topos}{Left exactness, inverted maps, and the sheaf topos}
For \(f:X\to Y\), call the following two conditions
\(j\)-bijectivity: its image is \(j\)-dense in \(Y\),
and \(\Delta_X\hookrightarrow X\times_YX\) is
\(j\)-dense.  Then
\[
 a_j(f)\text{ is an isomorphism}
 \quad\Longleftrightarrow\quad f\text{ is }j\text{-bijective}
 \quad\Longleftrightarrow\quad
 \mathcal E(Y,Z)\to\mathcal E(X,Z)\text{ is bijective
 for every sheaf }Z.
\]
Both conditions are stable under every pullback.  Every unit
\(\eta_X\) is \(j\)-bijective, and the reflector preserves
terminal objects and all pullbacks, hence all finite limits.

The full category \(\mathcal E_j\) of sheaves has finite
limits computed in \(\mathcal E\), finite colimits obtained
by reflecting the base colimits, exponentials computed in
\(\mathcal E\), and subobject classifier
\(\mathsf{true}_j:1\to T_j\).  A mono between sheaves
is an underlying closed mono, with the unique base classifier
factoring through \(T_j\).  Consequently \(\mathcal E_j\)
satisfies the complete original topos package.  The adjunction
\(a_j\dashv i_j\), with full-and-faithful inclusion \(i_j\),
is a geometric morphism whose inverse image is the constructed
left-exact reflector.  No sheafification or sheaf-topos theorem
is assumed as a premise \cite[p.~331]{Lawvere1970}.
For every parameter object \(B\), the same construction in
\(\mathcal E/B\) with topology \(1_B\times j\) has
relative associated sheaf
\[
                    B\times_{a_jB}a_jX\longrightarrow B
                         \quad(X\to B),
\]
using \(\eta_B:B\to a_jB\).  These relative reflectors
commute canonically with every pullback \(B'\to B\),
including noninhabited parameter objects.
\end{hbtheorem}

The geometric example is Godement's construction lifted to an arbitrary
base.  A topology basis has an object of points \(X\), an object of
indices \(A\), and a pairing \(X\times A\to T\).  The
comparison relation on indices is obtained by pulling back inclusion
of subobjects of \(X\) along \(A\to T^X\).  A directness
condition makes the realization/section adjunction
\[
 \mathsf{lim}:A^{\mathrm{op}}\mathcal E
       \rightleftarrows\mathcal E/X:\Gamma
 \qquad (\mathsf{lim}\dashv\Gamma)
\]
a geometric morphism.  Its image is the sheaf category, described
either by a topology on the presheaf side or by a left-exact
cotriple on the family side.  The source does not spell out the
directness predicate or give the complete adjunction construction.
The following reconstruction derives a necessary-and-sufficient
intrinsic criterion for the displayed geometric adjunction, with
separately identified primary comparisons; it is not a claim that
the original article prints the formula
\cite[pp.~331--332]{Lawvere1970}.

\begin{hbdefinition}{def:p10-lawvere-sheaves-godement-basis}{The intrinsic local directness criterion}
Let \(P\hookrightarrow X\times A\) be the pairing's membership
relation, and write \(U_a\hookrightarrow X\) for its indicated
subobject.  The relation \(a\leq b\) means
\(U_a\subseteq U_b\), defined by the internal predicate
\(\forall x\,(P(x,a)\Rightarrow P(x,b))\).
It gives a thin internal category on \(A\).  The source calls
it a poset; if two distinct indices name the same subobject,
the literal pulled-back order is a preorder.  The construction
keeps those indices and their thin category instead of silently
discarding them.

The local directness conditions are the following two epimorphisms:
\[
 P\longrightarrow X,
 \qquad
 \{(x,a,b,c):P(x,c),\ c\leq a,\ c\leq b\}
      \longrightarrow\{(x,a,b):P(x,a)\land P(x,b)\}.
\]
Equivalently, internally every \(x\) lies in some basis member,
and membership in two members has a local common lower member
containing \(x\).  These are existential image/epi conditions,
not choices of globally defined indices.  They remain valid under
every parameter pullback.  For a thin category the parallel-arrow
part of filtering is automatic.  This is the basis special case
of the explicitly epimorphic-local filtering conditions of
\cite[VII\S8, Definition~1 and Lemma~3, pp.~394--397]
{MacLaneMoerdijk1992}; the theorem below proves necessity as well
as sufficiency here, without assuming that book's externally
cocomplete setting.
\end{hbdefinition}

\begin{hbtheorem}{thm:p10-lawvere-sheaves-godement}{proved=annex:p10.lawvere-sheaves.godement}{The intrinsic Godement realization, image, and cotriple}
For the pairing and thin category above, put
\[
       (\Gamma M)_a=\Pi_{U_a\to1}(M|_{U_a})
                  \quad(M\to X),
\]
interpreted internally as \(\Pi_{P\to A}(P\times_XM)\)
over \(A\), with restriction by precomposition on the basis
inclusions.  For a presheaf \(F\) on \(A\), form in
\(\mathcal E/X\) the coequalizer
\[
 H_1(F)\rightrightarrows H_0(F)\longrightarrow L(F),
\]
where generalized elements are
\[
 H_0(F): (x,a,v\in F_a),\ P(x,a),
 \qquad H_1(F):(x,a\leq b,v\in F_b),\ P(x,a),
\]
and the parallel maps give \((x,b,v)\) and
\((x,a,v|a)\).  Then \(L\dashv\Gamma\), with actual
section/germ unit and evaluation counit.  It is the source's
\(\mathsf{lim}\dashv\Gamma\) adjunction.  The two local
directness conditions are equivalent to \(L\) being left exact.
Under them, two germs agree exactly when their representatives
agree on a local common lower basis member.

Its geometric image is the sheaf topos for the internal coverage
in which a sieve \(S\) on \(a\) covers precisely when
\[
                           U_a=\bigvee_{b\in S}U_b.
\]
Here the join is an internally constructed image, not an assumed
external infinite union.  Equivalently the image is the coalgebra
category of the actual left-exact cotriple \(L\Gamma\) on
\(\mathcal E/X\), with counit evaluation and comultiplication
\(L\) of the adjunction unit.  The finite equalizer proof of
this comparison is supplied in the annex, not assumed as a
comonadicity theorem.
Here a coalgebra is a family with \(\theta:M\to L\Gamma M\)
satisfying \(\varepsilon_M\theta=1_M\) and
\((L\Gamma\theta)\theta=(L\zeta_{\Gamma M})\theta\),
where \(\zeta\) is the adjunction unit; its morphisms commute
with these structure maps.

For ordinary sets and a genuine topological basis, these local
conditions are the usual cover/common-refinement conditions and
the coequalizer is the germ construction.  The primary comparison
is Godement's \emph{espace \'etal\'e} construction and generated
sheaf \cite[II\S1.2, pp.~110--112]{Godement1998Printing}.
The held book is a 1998 printing with 1958 copyright; no first-edition
byte identity or printed 1970 directness formula is inferred.
This reconstruction lifts the source's named construction to the
arbitrary base and retains its original announcement standing
\cite[pp.~331--332]{Lawvere1970}.
\end{hbtheorem}

\begin{hbexample}{def:p10-lawvere-sheaves-godement-failures}{Both directness clauses matter}
For \(X\) a singleton and \(A\) empty, realization sends the
terminal presheaf to the empty family, not to \(X\); it does
not preserve the terminal object.  For
\(X=\{x,y,z\}\), \(U_a=\{x,y\}\),
\(U_b=\{x,z\}\), with only these two incomparable indices,
the basis covers \(X\) but has no member refining their
intersection at \(x\).  Realization of the two representables
is \(U_a,U_b\), whose product over \(X\) is \(\{x\}\),
whereas their presheaf product and its realization are empty.
Thus cover alone does not give left exactness.  In a non-set base,
the repaired condition requires epimorphic local witnesses, not
a choice of one index over each global point.
\end{hbexample}

\subsubsection{Double negation and generalized ultraproducts}

Double negation \(j=\neg\neg\) is a topology in any topos.  The
source interprets it as cofinal locality.  Its sheaves have classical
logic: their classifier satisfies \(1+1\cong T\).  This is
equivalent to \(T\) being an internal Boolean algebra and to
each mono having a unique complementary summand
\(X'+(\neg X')\cong X\).  These are source assertions about
the double-negation sheaf topos, not assertions that the original
\(\mathcal E\) was already classical
\cite[p.~332]{Lawvere1970}.

\begin{hbtheorem}{thm:p10-lawvere-sheaves-double-negation}{proved=annex:p10.lawvere-sheaves.double-negation}{Double-negation sheaves have classical internal logic}
The operation \(j(p)=\neg\neg p\) on the internal Heyting
classifier is idempotent, preserves truth and binary meets, and
is therefore a topology.  Its fixed-predicate object \(T_j\)
is the classifier of the constructed sheaf topos.  For closed
subobjects \(A,B\) of a sheaf \(X\), their meet is ambient
intersection, their join is \(\overline{A\lor B}\), and
their complement is the ambient Heyting negation \(\neg A\),
which is itself closed.  Thus every internal predicate in
\(\mathcal E_{\neg\neg}\) satisfies excluded middle,
at every generalized parameter base.

In a topos with the original package, the following conditions
are equivalent: the classifier is an internal Boolean algebra;
every mono has a unique complementary summand; and the canonical
truth/false map \(1+1\to T\) is an isomorphism.  Applying
this to \(\mathcal E_{\neg\neg}\) gives the exact classicality
condition \(1+1\cong T_j\) announced by Lawvere
\cite[p.~332]{Lawvere1970}.  Its coproduct is the one in the
sheaf topos, not an assumed unchanged ambient coproduct.
Classicality alone does not supply the source's additional
Skolem condition in its stronger definition of a Boolean topos,
nor an axiom of choice.
\end{hbtheorem}

\begin{hbdefinition}{def:p10-lawvere-sheaves-ultraproduct}{The source's generalized ultraproduct data}
Take a finite-limit-preserving functor
\(u_*:\mathcal E\to\mathcal E_0\), not necessarily a
geometric morphism, and a homomorphism of internal Heyting algebras
\(h:u_*(T)\to T_0\).  Form \(\mathcal E_h\) by formally
inverting the monomorphisms \(m:X'\hookrightarrow X\) for which
\[
 h\bigl(u_*(\sigma_1(\Pi_{X\to1}(m)))\bigr)
                  =\mathsf{true}_0.
\]
The expression inside \(u_*\) is a predicate \(1\to T\):
it is support of the family of sections of \(m\), or equivalently
universal quantification of its membership predicate.  Finite-limit
preservation identifies \(u_*(1)\) with \(1_0\), making the
display typed \cite[p.~332]{Lawvere1970}.
\index{ultraproduct!Lawvere's generalized construction}
\end{hbdefinition}

\begin{hbtheorem}{thm:p10-lawvere-sheaves-ultraproduct}{sketched=annex:p10.lawvere-sheaves.ultraproduct}{The generalized-ultraproduct announcement}
The localized category \(\mathcal E_h\) is a topos, and
\(\mathcal E\to\mathcal E_h\) is a logical morphism.  Its
relation to \(\mathcal E_0\) is one of closed categories, not
in general a geometric morphism of topoi.  Lawvere invokes this
construction to collapse a Boolean-valued model of abstract sets
to a two-valued model, while doing the principal independence work
before choosing \(h\) \cite[p.~332]{Lawvere1970}.
This is an announced theorem without a localization proof in the
article; no geometric direct-image structure is inferred merely
from preservation of finite limits.
The full ordinary localization and logical-preservation proofs are
\cref{thm:p10-lawvere-localization-denominators,
thm:p10-lawvere-localization-fractions,thm:p10-lawvere-localization-logical}.
The actual full-datum closed functor is constructed in
\cref{thm:p10-lawvere-closed-base-partials,thm:p10-lawvere-closed-base-closed};
it is not automatically strong or normal, and its enriched global
underlying category is not silently identified with this ordinary
localization. The nondegenerate Boolean-valued collapse, including
existence of the required \(h\) and preservation of choice, is
\cref{thm:p10-lawvere-closed-base-collapse}. A degenerate Boolean
model cannot collapse to nondegenerate \(2\). The original combined
announcement retains its historical sketch boundary and unprinted
closed-base convention, not a claim that these separate editorial
constructions remain missing.
\end{hbtheorem}

For comparison, if \(\mathcal E=\mathsf{Set}^I\), take
\(u_*\) to be product and let a given ultrafilter \(U\) define
\(h:2^I\to2\).  For \(A_i\subseteq X_i\), the displayed
condition becomes
\[
                     \{i:A_i=X_i\}\in U.
\]
This calculation explains the ultrafilter terminology.  It does not
identify a localized category of families with the single ordinary
model-theoretic structure \(\prod_i X_i/U\), nor supply an
ultrafilter without the relevant external existence assumption.

%% New owned fragment; released source01 manuscript files remain frozen.
\subsubsection{The exact selected-mono localization and its logical functor}
\label{sec:p10-lawvere-localization}

The original article's datum is the finite-limit-preserving
\(u_*:\mathcal E\to\mathcal E_0\) and the internal Heyting
homomorphism \(h:u_*(T)\to T_0\), not an assumption of a
geometric inverse image or an ordinary model-theoretic ultrafilter
\cite[p.~332]{Lawvere1970}.  The following construction supplies
the ordinary localization and its logical preservation at exactly
that datum.  Its reduction to a truth filter is proved by the
universal property; no finite, principal, Boolean or prime-filter
restriction is imposed.  The base-relative closed-category and
collapse assertions remain distinct below.

\begin{hbdefinition}{def:p10-lawvere-localization-filter}{The full truth filter induced by the original datum}
Identify \(\operatorname{Sub}_{\mathcal E}(1)\) with the
maps \(1\to T\) by classification.  Define
\[
 \theta_h(U)=h\circ u_*(\chi_U):1_0\longrightarrow T_0,
 \qquad
 \mathcal F_h=\{U\hookrightarrow1:\theta_h(U)=\mathsf{true}_0\},
\]
using the finite-limit identification \(u_*(1)\cong1_0\).
The value \(\theta_h(U)\) is not asserted to be the
characteristic map of \(u_*(U)\).  For a mono
\(m:A\hookrightarrow X\), put
\[
         \gamma_m=\forall_{X\to1}(m)\in\operatorname{Sub}(1),
 \qquad \mathcal W_h=\{m:\gamma_m\in\mathcal F_h\}.
\]
The predicate of \(\gamma_m\) is exactly
\(\sigma_1(\Pi_{X\to1}(m))\), by the proved restriction
of the family adjunction to monos.  Thus \(\mathcal W_h\)
is precisely the monomorphism class in the source display,
without interchanging \(u_*\) with universal quantification.
\end{hbdefinition}

\begin{hblemma}{thm:p10-lawvere-localization-denominators}{proved=annex:p10.lawvere-localization.denominators}{The source denominators have uniform cofinal restrictions}
The map \(\theta_h\) is a Heyting homomorphism on global
truth algebras.  The set \(\mathcal F_h\) contains top,
is upward closed, and is closed under binary intersection.
It is proper when \(\mathcal E_0\) is nondegenerate;
the degenerate-base case permits the improper filter.
For every \(U\in\mathcal F_h\), the mono
\(i_{U,X}:U\times X\hookrightarrow X\) lies in
\(\mathcal W_h\).  Conversely every \(m\in\mathcal W_h\)
contains \(\gamma_m\times X\), and its pullback to
that admitted subterminal is an isomorphism.
Hence localization at the source's full mono class and
localization at all \(i_{U,X}\), \(U\in\mathcal F_h\),
have the same universal property.  This is a mathematical
equivalence of the exact classes, not replacement of the
general datum by a special finite filter.
\end{hblemma}

\begin{hbdefinition}{def:p10-lawvere-localization-fractions}{Fractions and their actual equality relation}
Fix the locally small/universe convention, so \(\operatorname{Sub}(1)\)
and each indicated hom collection are sets.  The objects of
\(\mathcal E/\mathcal F_h\) are those of \(\mathcal E\).
A morphism \(X\to Y\) is represented by
\((U,f)\), where \(U\in\mathcal F_h\) and
\(f:U\times X\to Y\).  Set
\[
 (U,f)\sim(V,g)\quad\Longleftrightarrow\quad
 \text{some }W\in\mathcal F_h, W\leq U\cap V,
       \text{ has }f|_{W\times X}=g|_{W\times X}.
\]
Write \([U,f]\) for its class.  Identity has denominator
\(1\).  Compose \([U,f]:X\to Y\) and
\([V,g]:Y\to Z\) on \(W=U\cap V\): restrict
\(f\) to \(W\times X\), pair it with the unique
parameter map to \(V\), then apply \(g\).
No representatives outside the admitted subterminal are
chosen or required.  The canonical functor is
\(Q(X)=X\), \(Q(f)=[1,f]\).

In the 1992 primary comparison, a representative is a
map \(X\times U\to Y\times U\) over \(U\).
Its first coordinate is the present representative;
the other coordinate is forced because \(U\) is
subterminal.  Product symmetry gives the exact typed
bridge.  The ordinary set-sized quotient of hom sets
does not assume an infinite colimit inside \(\mathcal E\)
\cite[V\S9, pp.~258--260]{MacLaneMoerdijk1992}.
That source's filter convention excludes zero; the
improper-filter boundary here is explicitly separate.
\end{hbdefinition}

\begin{hbtheorem}{thm:p10-lawvere-localization-fractions}{proved=annex:p10.lawvere-localization.fractions}{The complete fraction calculus and localization property}
The displayed relation is an equivalence relation and
the composition is independent of representatives and
refinements, associative and unital.  The functor \(Q\)
inverts every source mono in \(\mathcal W_h\).
For any category \(\mathcal D\) and functor
\(H:\mathcal E\to\mathcal D\) inverting that class,
there is a unique extension on the same objects,
\[
 \bar H[U,f]=H(f)\,H(i_{U,X})^{-1}.
\]
Natural transformations extend uniquely by their original
components.  Consequently the constructed category is
the ordinary \(\mathcal E_h\) formally localized at
the exact class stated on p.~332.  Every source fraction
has a uniform representative of this form, and equality
is exactly agreement on a further admitted common
restriction; there is no additional unproved fraction
existence or coherence requirement.
\end{hbtheorem}

\begin{hbtheorem}{thm:p10-lawvere-localization-logical}{proved=annex:p10.lawvere-localization.logical}{The ordinary localized topos and full logical preservation}
The constructed \(\mathcal E_h\) has finite limits,
finite colimits, exponentials and subobject classifier
\(Q(\mathsf{true}):Q1\to QT\), and \(Q\) preserves
all these structures up to their canonical isomorphisms.
In particular \(Q(Y^X)\) is the exponential of
\(QY\) by \(QX\), with its actual evaluation, and
every localized mono has a unique characteristic map,
not merely those arising from initially chosen strict
monos.  Every finite diagram of germs is handled on
a common admitted restriction, including equalizers
and coequalizers of genuinely localized arrows.

Power objects, membership, the graph constructions of
dependent products, and both internal quantifier maps
are preserved as well.  Thus \(Q\) is a logical
morphism in the original source's package.  The
improper filter gives the degenerate topos and is
included explicitly.  The later primary comparison is
\cite[V\S9, Theorem~2 and Lemma~3, pp.~261--263]
{MacLaneMoerdijk1992}; its proof boundaries and proper
filter convention are not silently promoted to a
different 1970 source statement.
\end{hbtheorem}

\begin{hbexample}{def:p10-lawvere-localization-base-functor}{The original base functor need not descend unchanged}
Let \(\mathcal E=\mathsf{Set}\times\mathsf{Set}\),
\(\mathcal E_0=\mathsf{Set}\), \(u_*\) be binary
product, and \(h:2\times2\to2\) project to the first
coordinate.  These satisfy the original hypotheses.
The subterminal \(U=(1,0)\) is admitted, so its mono
into \(1=(1,1)\) is inverted.  But its image under
\(u_*\) is \(0\to1\), not an isomorphism.
Therefore \(u_*\) itself cannot factor through this
localization.  A new base-relative construction is
needed; finite-limit preservation alone does not
produce a geometric direct image or a strong
closed-functor descent.
\end{hbexample}

\begin{proofboundary}{def:p10-lawvere-localization-boundary}{The exact completed and still-owed source assertions}
The ordinary localization and its logical topos preservation are
proved at the full original \(u_* / h\) datum, through the exact
cofinal reduction to its arbitrary induced filter. The article
announces these conclusions. The additional full-datum closed
functor is constructed in
\cref{thm:p10-lawvere-closed-base-partials,thm:p10-lawvere-closed-base-closed},
without asserting unchanged descent of \(u_*\), an invertible
closed comparison, or arbitrary normality. BVM classification and
choice at the declared cell regimes are
\cref{thm:p10-lawvere-bvm-presentation,thm:p10-lawvere-bvm-realization,thm:p10-lawvere-bvm-correspondence};
the nondegenerate collapse, with its actual \(h\)-existence and
choice hypotheses, is \cref{thm:p10-lawvere-closed-base-collapse}.
The degenerate/improper boundary is not a nondegenerate two-valued
collapse. General stronger normal-relative identification is not
entailed by the original datum; the source's unprinted convention
is kept distinct. CH/ccc/no-epi/transfer, spectrum and independent
recursion/free-algebra proofs remain owed. No finite, principal
or improper special case discharges those separate obligations.
\end{proofboundary}


\subsubsection{Abstract sets, Boolean-valued models, and the CH sketch}

\begin{hbdefinition}{def:p10-lawvere-sheaves-bvm}{The source's abstract-set and Boolean packages}
Lawvere takes the category \(\mathcal S\) of abstract sets to
satisfy \((\Pi)\), \((\exists)\), \((\omega)\), classical
logic \((\neg):1+1\cong T\), and irreducibility
\[
 (\lor):\quad\varphi_1,\varphi_2:1\to T,
 \quad\varphi_1\lor\varphi_2=\mathsf{true}
 \quad\Longrightarrow\quad
 \varphi_1=\mathsf{true}\text{ or }\varphi_2=\mathsf{true}.
\]
Classicality plus irreducibility implies that \(1\) has at most the
two truth/false subobjects; they are distinct when the topos is
nondegenerate.  Exact two-valuedness includes that nondegeneracy
condition.  The converse is not asserted: the source uses
\(M^{\mathrm{op}}\mathsf{Set}\) for a monoid that is not a
group to distinguish it.

In this article a \emph{Boolean topos} means a topos satisfying
both \((\exists)\) and \((\neg)\).  Its classifier is usually
written \(2=T\).  A Boolean-valued model
\(\mathcal E\in\mathsf{BVM}/\mathcal S\) is such a Boolean
topos with a specified geometric morphism to \(\mathcal S\).
This is stronger terminology than merely saying that a topos's
internal logic is Boolean \cite[pp.~332--333]{Lawvere1970}.
\index{Boolean-valued model!Lawvere's definition}
\end{hbdefinition}

\begin{hbtheorem}{thm:p10-lawvere-sheaves-generators}{proved=annex:p10.lawvere-sheaves.generators}{Subterminal generators under the source's Boolean package}
If every subobject has a complementary summand and
\((\exists)\) holds, the subobjects of \(1\) form a generating
family: unequal maps \(f,g:E\to F\) are distinguished by
some map \(U\to E\) with \(U\hookrightarrow1\).
This is the generation conclusion stated on p.~333, with an
editorial full proof in the annex.
\end{hbtheorem}

The article further announces that every Boolean-valued model over
\(\mathcal S\) satisfies \((\Pi)\), hence choice, and that
the bicategory \(\mathsf{BVM}/\mathcal S\) is equivalent to
\(\mathsf{CBA}(\mathcal S)\), the category of
\(\mathcal S\)-complete Boolean algebra objects
\cite[p.~333]{Lawvere1970}. These conclusions are announced, not
proved in the short article. The following editorial construction
supplies the actual functors, comparisons and choice proof while
declaring its morphism/2-cell directions explicitly; their names
alone would not supply those data.

% Adopted Lawvere continuation; working spec fragment preserved separately.
\subsubsection{The abstract-set base and the Boolean-valued correspondence}

Lawvere's Nice 1970 article, published in 1971, announces on p.~333
that its Boolean-valued models satisfy choice and correspond to complete
Boolean algebras in the abstract-set base \cite[p.~333]{Lawvere1970}.
Here ``Boolean'' retains the source's conjunction of classical logic and
support splitting, not just the modern classical-logic condition.
The following constructions are independent editorial proofs of that
announcement at an explicitly specified morphism scope. They do not turn
the six-page source into a source-native detailed proof.

\begin{hblemma}{thm:p10-lawvere-bvm-base}{proved=annex:p10.lawvere-bvm.base}{Bounded well-ordering in the source's base}
Let \(\mathcal S\) have the abstract-set package of
Definition~\ref{def:p10-lawvere-sheaves-bvm}. Then \(\mathcal S\)
is nondegenerate, \(1\) is a generator, every epimorphism splits,
and every object admits an internal well-order. The latter construction
uses power objects and a choice function on nonempty subobjects of that
object; it does not require a cumulative universe, Replacement, an
external ordinal, or a comparison with all models of ZFC.
\end{hblemma}

\begin{hbdefinition}{def:p10-lawvere-bvm-matrices}{Complete Boolean algebras and their matrices}
An \(\mathcal S\)-complete Boolean algebra \(B\) is a Boolean
algebra object with \(\bigvee:\mathcal P_{\mathcal S}B\to B\)
satisfying the least-upper-bound property internally. Meets are obtained
by complementation. A morphism preserves Boolean operations and all
\(\mathcal S\)-indexed joins.

The category \(\operatorname{Mat}(B)\) has objects \((I,b,e)\),
where \(I\in\mathcal S\), \(b:I\to B\), and
\(e:I\times I\to B\) satisfy
\[
 e_{ii}=b_i,\quad e_{ij}\le b_i\wedge b_j,\quad
 e_{ij}=e_{ji},\quad e_{ij}\wedge e_{jk}\le e_{ik}.
\]
A map \((I,b,e)\to(J,c,f)\) is a matrix \(r:I\times J\to B\)
with
\[
\begin{gathered}
 r_{ij}\le b_i\wedge c_j,\qquad \bigvee_j r_{ij}=b_i,\qquad
 r_{ij}\wedge r_{ik}\le f_{jk},\\
 e_{ii'}\wedge r_{ij}\le r_{i'j},\qquad
 r_{ij}\wedge f_{jj'}\le r_{ij'}.
\end{gathered}
\]
These are saturated, total, single-valued relations, not functions between
chosen representatives. Identity is \(e\); composition is
\((sr)_{ik}=\bigvee_j r_{ij}\wedge s_{jk}\).
All displayed indexing and inequalities are internal to \(\mathcal S\).
\end{hbdefinition}

\begin{hbtheorem}{thm:p10-lawvere-bvm-presentation}{proved=annex:p10.lawvere-bvm.presentation}{The actual Boolean-algebra presentation of a model}
For a source Boolean-valued model \(p:\mathcal E\to\mathcal S\),
put \(B_{\mathcal E}=p_*T_{\mathcal E}\). This is an
\(\mathcal S\)-complete Boolean algebra, and there is an equivalence
over \(\mathcal S\)
\[
             \mathcal E\simeq\operatorname{Mat}(B_{\mathcal E}).
\]
Each object is represented by an actual effective quotient of a subobject
of \(p^*I\), for an object \(I\) of the base. The functors on maps,
inverse comparison maps, and their triangle identities are supplied in
the annex. No global points of \(\mathcal E\) are assumed sufficient.
\end{hbtheorem}

\begin{hbtheorem}{thm:p10-lawvere-bvm-realization}{proved=annex:p10.lawvere-bvm.realization}{Realization and choice}
For every \(\mathcal S\)-complete Boolean algebra \(B\), let
\(\operatorname{Sh}_{\mathcal S}(B)\) be internal sheaves on the
poset \(B\), where a sieve on \(b\) covers precisely when its join
is \(b\). Its geometric base map \(q\) has
\(q^*=a\Delta\), and
\[
 q_*T\cong B,\qquad
 \operatorname{Sh}_{\mathcal S}(B)\simeq\operatorname{Mat}(B).
\]
It is classical and every epimorphism splits. Therefore it is a Boolean
topos in the source's stronger sense and satisfies \((\Pi)\) and choice.
The same splitting construction applies to every \(\mathcal E\) in
the preceding theorem. It uses a well-order of a base indexing object
and Boolean disjointification, not a choice of a sheaf-global witness.
The one-element Boolean algebra is allowed and realizes the degenerate
topos; excluding it on one side requires excluding it on the other.
\end{hbtheorem}

\begin{hbdefinition}{def:p10-lawvere-bvm-cell-regime}{The specified morphism and 2-cell regimes}
Write \(\mathsf{BVM}^{\leftarrow}(\mathcal S)\) for the
inverse-image-directed bicategory: a 1-cell from \(\mathcal E\) to
\(\mathcal F\) is a finite-limit-preserving left adjoint
\(L:\mathcal E\to\mathcal F\), with its right adjoint and a specified
base isomorphism \(Lp^*\cong q^*\). A 2-cell is a natural
transformation of these left adjoints compatible with that isomorphism.
Composition has the usual composite base isomorphism.
In the conventional geometric direction a 1-cell
\(\mathcal E\to\mathcal F\) instead has inverse image
\(\mathcal F\to\mathcal E\). We keep this opposite direction
separate. The article names a bicategory on p.~333 without printing its
1-cell or 2-cell definitions; the notation here specifies a reconstruction,
not a quotation of an unprinted convention.
\end{hbdefinition}

\begin{hbtheorem}{thm:p10-lawvere-bvm-correspondence}{proved=annex:p10.lawvere-bvm.correspondence}{Maps, adjunctions, and coherence of the correspondence}
There are biequivalences, with the algebra categories regarded as locally
discrete bicategories,
\[
 \mathsf{BVM}^{\leftarrow}(\mathcal S)
       \simeq\mathsf{CBA}(\mathcal S),\qquad
 \mathsf{BVM}_{\mathrm{geom}}(\mathcal S)
       \simeq\mathsf{CBA}(\mathcal S)^{\mathrm{op}}.
\]
A complete homomorphism \(\alpha:B\to C\) acts on Boolean matrices
entry by entry. Its right adjoint on sheaves is the actual functor
\(R_\alpha(Y)(b)=Y(\alpha b)\). Conversely a 1-cell induces its
homomorphism by transport of subobjects of base-indexed objects.
The adjunction bijections, units and counits, reconstruction isomorphisms,
uniqueness of compatible 2-cells, and composition coherence are established
in the annex. Such inverse images are geometric; they need not preserve
power objects or be logical morphisms.
\end{hbtheorem}

For comparison, Mac Lane--Moerdijk explicitly separate the internal
Boolean condition, two-valuedness, and the splitting of all epimorphisms
\cite[pp.~270--276]{MacLaneMoerdijk1992}. Their later choice-independence
construction is not a substitute for the source's stronger Boolean
package \cite[pp.~291--295]{MacLaneMoerdijk1992}; the proof of that later
construction continues beyond this selected reading. The correspondence
above is relative to the given abstract-set base, not a classification
of all elementary topoi or of ordinary set-theoretic models.
The full-datum closed functor and nondegenerate two-valued collapse
follow in Theorems~\ref{thm:p10-lawvere-closed-base-partials}--
\ref{thm:p10-lawvere-closed-base-collapse}, after this BVM construction.
The countable-chain-condition/CH transfer, global-spectrum universal
property, and independently owed recursion/free-algebra dependencies
remain separate; neither choice nor the correspondence proves them.

% Adopted Lawvere continuation; working spec fragment preserved separately.
\subsubsection{The full-datum closed-base construction and collapse}

The original p.~332 theorem distinguishes a closed-category base from a
geometric base \cite[p.~332]{Lawvere1970}; the event is Nice 1970 and the
proceedings publication is 1971. The following construction retains the
given finite-limit functor \(u_*:\mathcal E\to\mathcal E_0\) and
Heyting homomorphism \(h:u_*T\to T_0\). It does not assume that
\(u_*\) itself descends through the selected-mono localization.

For the historical terminology, a closed functor has a functor, a unit
map, and a natural internal-hom comparison satisfying identity, unit and
composition compatibilities. The hom comparison need not be invertible.
Eilenberg--Kelly distinguish this from \emph{normality}, which identifies
the underlying-element functors \cite[pp.~433--434,439]{EilenbergKelly1966}.
Their relative-base paragraph imposes a stronger normal convention
\cite[p.~549]{EilenbergKelly1966}. Lawvere's six-page article does not
print such an additional element-functor requirement. Here we construct
the actual closed functor, its base-valued hom objects and comparison;
we do not silently identify the resulting enriched global underlying
category with the ordinary localization. The countercase below explains
why that stronger identification cannot follow from the full printed datum.

\begin{hbdefinition}{def:p10-lawvere-closed-base-partials}{The h-total partial-element quotient}
Let \(\partial X\subseteq\mathcal P_{\mathcal E}X\) be the object
of subsingleton predicates, with universal partial member
\(D_X\hookrightarrow\partial X\) and evaluation \(D_X\to X\).
Write \(d_X:\partial X\to T\) for its domain classifier. On
\(\partial X\times\partial X\), let
\(e_X(p,q)\) classify the domain on which the two partial members
have a common value. Thus
\[
 e_X(p,p)=d_X(p),\quad e_X(p,q)=e_X(q,p),\quad
 e_X(p,q)\le d_X(p)\wedge d_X(q),\quad
 e_X(p,q)\wedge e_X(q,r)\le e_X(p,r).
\]
Define in \(\mathcal E_0\)
\[
 \widehat d_X=h\,u_*d_X,\qquad \widehat e_X=h\,u_*e_X,
 \qquad A_X=\{p\in u_*\partial X:\widehat d_X(p)=\mathsf{true}_0\}.
\]
The finite-product identifications of \(u_*\) are understood.
On \(A_X\) let \(R_X\) be classified by \(\widehat e_X\),
and put \(H_hX=A_X/R_X\), using the established effective quotient
construction in \(\mathcal E_0\). The hats classify newly defined
predicates, not the images under \(u_*\) of the original subobjects.
Every formula is internal in the base and applies at all generalized
parameter objects.
\end{hbdefinition}

\begin{hbtheorem}{thm:p10-lawvere-closed-base-partials}{proved=annex:p10.lawvere-closed-base.partials}{The quotient functor and its exact finite-limit behavior}
The relation \(R_X\) is an equivalence relation. Forward image of
partial members defines an actual functor
\(H_h:\mathcal E\to\mathcal E_0\), with a natural singleton map
\(\eta_X:u_*X\to H_hX\). The quotient has the universal property
that maps \(H_hX\to Z\) correspond uniquely to maps
\(A_X\to Z\) equal on \(R_X\).
There is also an actual map \(\tau:H_hT\to T_0\) with
\(\tau\eta_T=h\); it is not asserted a classifier isomorphism.

The functor preserves terminal objects, products and equalizers, with
explicit inverse product maps and equalizer lifts obtained by restricting
partial members. It inverts every mono in the original class
\(\mathcal W_h\), not merely principal or finite-filter examples.
Consequently the exact ordinary localization property yields
\[
 \overline H_h:\mathcal E_h\longrightarrow\mathcal E_0,
 \qquad
 \overline H_h[U,f]=H_h(f)H_h(i_{U,X})^{-1}.
\]
This extension preserves finite limits of arbitrary localized arrows.
It is a new base functor, not a claimed factorization of \(u_*\).
\end{hbtheorem}

\begin{hbtheorem}{thm:p10-lawvere-closed-base-closed}{proved=annex:p10.lawvere-closed-base.closed}{Actual hom objects, composition and the closed comparison}
The cartesian closed structures of \(\mathcal E_h\) and
\(\mathcal E_0\) give the unit
\(1_0\cong\overline H_h1\) and natural comparison
\[
 \kappa_{X,Y}:\overline H_h(Y^X)
                  \longrightarrow(\overline H_hY)^{\overline H_hX},
\]
obtained by transposing transported evaluation. They satisfy all three
closed-functor compatibilities. In particular
\(\mathcal C_h(X,Y)=\overline H_h(Y^X)\) are actual base hom
objects, with identities transported from names of identity maps and
composition transported from internal composition. Associativity, unit
laws and the compatibility of \(\kappa\) with composition are proved
as equations of maps in \(\mathcal E_0\).
The closed comparison is not asserted an isomorphism or normal. The
ordinary maps of \(\mathcal E_h\) have natural names in these hom
objects; the naming maps need not be bijections onto their global elements.
The quotient and localization universal properties, and the exponential
transpose property fixing this product-induced comparison, are explicit
in the annex. No adjoint or geometric direct image is assumed for
\(\overline H_h\).
\end{hbtheorem}

\begin{hblemma}{thm:p10-lawvere-closed-base-countercases}{proved=annex:p10.lawvere-closed-base.countercases}{Two different base requirements that cannot be conflated}
For product \(u_*:\mathsf{Set}\times\mathsf{Set}\to\mathsf{Set}\)
and first-coordinate truth homomorphism, the new \(H_h\) is first
projection, whereas \(u_*\) cannot descend unchanged.
For diagonal \(u_*:\mathsf{Set}\to\mathsf{Set}\times\mathsf{Set}\)
and the identity truth homomorphism, \(\mathcal E_h=\mathsf{Set}\)
and \(H_h\) is diagonal. Its base hom object is
\((Y^X,Y^X)\); its global elements are pairs of functions, while
ordinary localized maps give only equal pairs. Thus even an invertible
closed comparison does not make the enriched global underlying category
the specified ordinary localization. Neither stronger requirement is
inserted into the source theorem.
\end{hblemma}

\begin{hbtheorem}{thm:p10-lawvere-closed-base-collapse}{proved=annex:p10.lawvere-closed-base.collapse}{The nondegenerate Boolean-valued two-valued collapse}
Let \(p:\mathcal E\to\mathcal S\) be a \emph{nondegenerate}
source Boolean-valued model over the abstract-set base. Then its internally
complete Boolean algebra \(B=p_*T\) admits a Boolean homomorphism
\(h:B\to2=T_{\mathcal S}\). This need not preserve arbitrary joins.
For this original datum \((p_*,h)\), the exact localization
\(\mathcal E\to\mathcal E_h\) is logical and its target is
nondegenerate, classical, two-valued, and still satisfies choice.
It has the constructed closed functor to \(\mathcal S\); it is
not claimed geometrically defined over \(\mathcal S\). In this
particular collapse case the naming maps are bijective, giving the
normal element-functor identification up to canonical isomorphism;
this uses the geometric original \(p_*\) and the proved choice
in \(\mathcal S\), unlike the unrestricted full-datum case.

The annex constructs \(h\) by bounded well-ordered proper-filter
histories, rather than assuming an ultrafilter lemma or a ZFC universe.
Nondegeneracy is necessary: the one-element Boolean algebra admits no
Boolean homomorphism to the nondegenerate algebra \(2\). The general
full-datum construction above includes the improper-filter/degenerate-base
boundary, but that boundary is not a nondegenerate two-valued model.
\end{hbtheorem}

These are independent editorial proofs. The original closed-base and
collapse statements remain historically announcements on p.~332. They
do not settle the unprinted normal-relative convention by inventing an
element-functor identity. They also do not supply the outstanding suitable
ccc/CH comparison and transfer, spectrum/locality/universality, or the
independent recursion-lifting/free-algebra arguments.


\subsubsection{Double-negation Boolean-valued models and the CH sketch}

The article also gives the construction
\[
           \widetilde P=(P\mathcal S)_{\neg\neg}
\]
for a poset \(P\) internal to \(\mathcal S\)
\cite[p.~333]{Lawvere1970}. The BVM correspondence and collapse
just constructed do not prove the subsequent cardinal-comparison
or chain-condition claims.

The forcing notation in this construction is stagewise and explicitly
relative to \(\widetilde P\).  If \(q\geq p\), \(x\) is
an element of \(X\) defined at \(p\), and
\(\varphi:X\to2\), then \(q\) forces \(\varphi(x)\)
when \(\varphi(x/q)=\mathsf{true}\).  In double-negation
sheaves this holds at \(q\) exactly when it holds on a set of
stages cofinal beyond \(q\).  This is the source's specialization
of sheaf locality; it is not a general current-stage witness clause
for all presheaves.

\begin{hbtheorem}{thm:p10-lawvere-sheaves-ch}{sketched=annex:p10.lawvere-sheaves.ch}{Lawvere's Boolean-valued CH-refutation construction}
In \(\mathcal S\), choose \(I\) with \(2^\omega<I\), where
the inequality means a monomorphism \(2^\omega\to I\) but no
epimorphism in that same direction.  Let \(P\)
be finite partial maps \(\omega\times I\to2\), ordered by
extension, and use its double-negation sheaf topos
\(\widetilde P\).  The source calls this finite-partial-map
object intrinsically definable; its application here is over the
abstract-set base with the stated recursion object.  For the
constant-sheaf inverse image
\(u^*:\mathcal S\to\widetilde P\), the source asserts
\[
                    \omega<u^*(2^\omega)<2^\omega
                      \quad\text{in }\widetilde P.
\]
The rightmost power object is formed in \(\widetilde P\);
the middle one is the transported ground object.  The source
produces a sheaf-local injection \(u^*(I)\hookrightarrow2^\omega\)
and depends on its following lemma: for a poset satisfying a suitable
countable-chain condition and ground objects \(X,Y\) with
\(Y\times\omega\cong Y\),
\[
 \operatorname{Epi}(X,Y)=0\text{ in }\mathcal S
 \quad\Longrightarrow\quad
 \operatorname{Epi}(u^*X,u^*Y)=0\text{ in }\widetilde P.
\]
Here \(\operatorname{Epi}(X,Y)\hookrightarrow Y^X\) is the
object of maps whose image is all of \(Y\), obtained by pulling
back \(\mathsf{true}\) along the predicate of having full image.
The article leaves the word ``suitable'' in the chain-condition
hypothesis and supplies no proof of this lemma
\cite[p.~333]{Lawvere1970}. The following editorial continuation
gives a precise sufficient internal ccc and proves it for this actual
finite-partial-map object. Its no-epi, local-injection and strict
internal comparison proofs are
\cref{thm:p10-lawvere-ch-chain-ccc,thm:p10-lawvere-ch-chain-no-epi,
thm:p10-lawvere-ch-chain-generic,thm:p10-lawvere-ch-chain-comparisons},
with their cardinal and internal-site dependencies supplied separately.
This is not a claim that the article prints that chosen ccc formula
or those complete proofs. The original statement retains its historical
sketch standing. Specified set-theoretic transfer/model interpretation
and CH's positive consistency direction still require separate proofs.
\end{hbtheorem}

% Adopted CH-chain continuation; working spec fragment preserved separately.
\subsubsection{The finite-partial-map chain condition and the source's no-epi lemma}

The original Nice 1970 article, published in 1971, leaves the word
\emph{suitable} in its countable-chain-condition lemma on p.~333
\cite[p.~333]{Lawvere1970}. Here is an explicit sufficient internal
condition and its proof for the very finite-partial-map object used there.
These are independent editorial proofs, not unprinted formulas attributed
to the original source. The base is the declared abstract-set topos
\(\mathcal S\), with classical logic, choice, power objects and the source's
recursion object \(\omega\), not an external universe of ZFC sets.

\begin{hbdefinition}{def:p10-lawvere-ch-chain-poset}{Finite conditions and internally countable antichains}
Let \(D=\omega\times I\). Write \(\operatorname{Fin}(D)\) for its
finite-subset object, the least
subobject of \(\mathcal P D\) containing empty and singleton predicates
and closed under binary union, obtained by intersection of all such
subobjects. In the present classical recursion base these are exactly
subsets bijective with some initial segment \([n]\subseteq\omega\).
Let \(P\) be the object of functional relations \(D\rightharpoonup2\)
whose domain lies in this finite-subset object.

Write \(q\succeq p\) when \(q\) extends \(p\), the source's \(q\ge p\).
The site category has arrows \(q\to p\) in that case, so its presheaf
restriction runs toward stronger conditions. Conditions are compatible
iff they agree on their common domain; their union is then a common
extension. An antichain is a subobject of pairwise incompatible distinct
conditions. Countable means admitting a monomorphism into \(\omega\),
including finite and empty objects. The internal ccc says that every
antichain subobject is countable. This is an internal quantified statement
on the relevant power objects, valid with parameters; it is not merely an
external assertion about global conditions.

On any such poset use the double-negation coverage: a sieve at \(p\)
covers iff for every \(q\succeq p\) there is \(r\succeq q\) in the sieve.
Let \(\widetilde P\) be its internal sheaf topos and
\(u^*=a\Delta:\mathcal S\to\widetilde P\) its constant-sheaf inverse image.
The earlier internal-presheaf and reflector constructions supply this
topos and geometric adjunction. Constant sheaves are not assumed to have
only ground sections at every stage.
\end{hbdefinition}

\begin{hblemma}{thm:p10-lawvere-ch-chain-cardinal-tools}{proved=annex:p10.lawvere-ch-chain.cardinal-tools}{The bounded cardinal tools actually used}
In \(\mathcal S\), countable unions of internally countable families
are countable, images of countable objects are countable, and
\(\omega\times\omega\cong\omega\). Every infinite object \(X\) satisfies
\(X\times\omega\cong X\).
Pairwise-disjoint subfamilies and antichains have maximal subfamilies.
The annex constructs these by bounded well-ordered histories in power
objects and ordinary natural-number recursion, not external ordinal
recursion, an unproved Zorn lemma, or Replacement.
\end{hblemma}

\begin{hbtheorem}{thm:p10-lawvere-ch-chain-ccc}{proved=annex:p10.lawvere-ch-chain.ccc}{The actual finite-partial-map poset is internally ccc}
For every \(I\in\mathcal S\), the object \(P\) above satisfies the
internal ccc. The proof constructs an uncountable delta-system from an
uncountable family of distinct finite supports, then fixes the finite
root assignment. Equal root assignments make two conditions compatible.
Finite support fibers and countable unions are proved, not assumed to
preserve cardinal size.
\end{hbtheorem}

\begin{hblemma}{thm:p10-lawvere-ch-chain-constant-sections}{proved=annex:p10.lawvere-ch-chain.constant-sections}{Constant sections, deciding covers and number objects}
For a base object \(A\), the constant presheaf is separated and its
unit into \(u^*A\) is a dense mono. Thus every section of \(u^*A\)
is locally a ground section, and unequal ground sections never become
equal. These assertions hold in every generalized parameter context.
The image \(u^*\omega\), with transported zero and successor, is the
natural-number object of \(\widetilde P\).
This classical sheaf topos is a Boolean-valued model in the source's
stronger sense: its canonical free-action covers are subobjects of
base-indexed constant sheaves, and the proved matrix disjointification
therefore supplies choice and support splitting.
No enough-points assumption or logical preservation of arbitrary
function objects is made.
\end{hblemma}

\begin{hbtheorem}{thm:p10-lawvere-ch-chain-no-epi}{proved=annex:p10.lawvere-ch-chain.no-epi}{A precise sufficient no-epi preservation lemma}
Let an internal poset \(P\) satisfy the internal ccc just defined, and
let \(X,Y\in\mathcal S\), retaining the source's hypothesis
\(Y\times\omega\cong Y\). Then
\[
 \operatorname{Epi}_{\mathcal S}(X,Y)=0
 \quad\Longrightarrow\quad
 \operatorname{Epi}_{\widetilde P}(u^*X,u^*Y)=0.
\]
Here the epi objects are the full internally parametrized predicates of
surjectivity, not the absence of one chosen global map. For each ground
input, an actual maximal deciding antichain gives a countable family of
possible output values. A purported local epi yields a ground epi from
\(X\times\omega\); infinite \(X\) absorbs that factor. Finite \(X\)
is handled by deciding all its values on a common extension and choosing
a missing ground output. The proof does not make the invalid inference
that a surjection from \(X\times\omega\) always gives one from \(X\).
\end{hbtheorem}

\begin{hbexample}{def:p10-lawvere-ch-chain-naive-countercase}{Why the finite-input branch is necessary}
For \(X=1,Y=\omega\), the projection \(X\times\omega\to Y\) is
an isomorphism, while no map \(X\to Y\) is epi: a singleton image
cannot contain both zero and successor of zero. Also
\(Y\times\omega\cong Y\). Thus even the source's absorption
hypothesis does not justify the naive intermediate cardinal inference.
The proved finite-input deciding argument is what repairs that step.
\end{hbexample}

\begin{hbtheorem}{thm:p10-lawvere-ch-chain-generic}{proved=annex:p10.lawvere-ch-chain.generic}{Actual generic evaluation and local injection descent}
For the finite-partial-map \(P\), the domain where a condition assigns a
value at \((n,i)\) is a dense subpresheaf of \(\Delta(\omega\times I)\).
Its evaluation into \(\Delta2\), after sheafification, gives an actual map
\[
 u^*I\times u^*\omega\longrightarrow u^*2.
\]
Its transpose \(\alpha:u^*I\to2^\omega\), using the new sheaf topos's
power object and number object, is monic. The proof uses a fresh finite
coordinate to distinguish two unequal indices and proves the required
density and sheaf descent at arbitrary parameter bases.
\end{hbtheorem}

\begin{hbtheorem}{thm:p10-lawvere-ch-chain-comparisons}{proved=annex:p10.lawvere-ch-chain.comparisons}{The displayed strict comparisons at their internal scope}
Put \(C=2^\omega\) in \(\mathcal S\), and take \(I\) with a mono
\(C\to I\) and \(\operatorname{Epi}_{\mathcal S}(C,I)=0\), as in the
source. Such an \(I\) is available, for example \(I=2^C\), by the
singleton injection and Cantor's diagonal argument.
In the finite-partial-map sheaf topos the actual comparisons are
\[
             \omega<u^*C<2^\omega.
\]
Each strict inequality means an exhibited mono in that direction and
zero internally parametrized epi object in the same direction.
The middle-to-right mono is the transported-evaluation comparison, not
an assumed isomorphism \(u^*(2^\omega)\cong2^\omega\).
The zero epi object on the right follows from the generic injection and a complement
retraction, together with the proved no-epi lemma; no extra global sheaf
choice is assumed for that retraction.

This supplies the source's internal cardinal comparison and its omitted
ccc/no-epi/local-injection dependencies. It does not specify or prove
a translation to a separately chosen set-theoretic axiom system or model,
nor CH's positive consistency direction. Spectrum and the independently
owed general recursion/free-algebra constructions remain separate.
\end{hbtheorem}


\subsubsection{The ringed-topos global spectrum}

The third source section first takes a commutative multiplicative
monoid object \(A\) and a homomorphism
\(u:A\to T^X\), where multiplication in \(T^X\) is
intersection.  Its pairing defines a membership subobject
\(P\hookrightarrow X\times A\), and the inclusion relation
between basis elements defines \(F\rightrightarrows A\).
The latter gives a submonoid of the constant monoid presheaf
\(\overline A\) on the index category; the former gives a
submonoid of the constant family \(\dot A\) over \(X\).
Forming fractions gives \((\overline A)_F\) on the presheaf
side and \((\dot A)_P\) on the family side.  The intermediate
sheaf category has the reflected monoid \(\widetilde A\)
between them \cite[p.~333]{Lawvere1970}.

\begin{hbdefinition}{def:p10-lawvere-sheaves-prime}{The source's prime subobject and local ring}
For a commutative ring object \(A\), a prime \(x\) in the
source's terminology is a subobject of \(A\), not an ideal,
whose membership predicate satisfies, internally,
\[
\begin{aligned}
 [1\in x]&=\mathsf{true},\\
 [fg\in x]&=[f\in x]\land[g\in x],\\
 [0\in x]&=\mathsf{false},\\
 [f+g\in x]&\leq[f\in x]\lor[g\in x].
\end{aligned}
\]
The multiplication clause is an equivalence, not one implication.
The disjunction is the join of subobjects, and may hold only
locally.  The source notes that this intuitionistic distinction
already matters in \(2\mathsf{Set}\).  A ring is local when
its subobject of units is prime in
this sense.  The object \(X\hookrightarrow T^A\) of such
subobjects is obtained by imposing these conditions with finite
limits on predicate objects and their quantified maps
\cite[pp.~333--334]{Lawvere1970}.
\index{prime subobject!Lawvere's spectrum}
\index{local ring!internal definition}
\end{hbdefinition}

\begin{hbtheorem}{thm:p10-lawvere-sheaves-spectrum}{sketched=annex:p10.lawvere-sheaves.spectrum}{The announced global-spectrum universal construction}
The resulting pairing gives a topology basis whose sheaves form
\(\operatorname{Spec}(A)\), the global spectrum of the ringed
topos \((\mathcal E,A)\).  In it \(\widetilde A\) is a local
ring object and is the universal local \(A\)-algebra among
topoi defined over \(\mathcal E\).  The article describes the
membership relation as being transformed into its opposite in this
process; the fractions invert the selected multiplicative subobjects.
This is the source's intrinsic
Chevalley--Hakim construction \cite[p.~334]{Lawvere1970}.
The article outlines the monoid, fraction, and sheaf steps; it does
not supply their universal-property proofs.  The annex keeps those
dependencies explicit.
\end{hbtheorem}

\begin{hbexample}{def:p10-lawvere-sheaves-classical-prime}{Comparison with the complement of a prime ideal}
In ordinary sets with classical logic, the four conditions say that
\(x\) is the complement of a proper prime ideal.  For
\(A=\mathbb Z\) and a prime number \(\ell\), take
\(x=\{n:\ell\nmid n\}\).  Its localization inverts just
these elements and gives \(\mathbb Z_{(\ell)}\); the units
are its localized image of \(x\).  This worked comparison is
editorial.  In intuitionistic internal logic, replacing the source's
prime subobject by its putative ideal complement would erase the
local-disjunction qualification.
\end{hbexample}

The closing section describes the algebraic-geometric application as
only beginning.  It attributes an interpretation of intuitionistic
semantics in a suitable varying-set topos to unpublished work of
George Rousseau, and the analogous interpretation of L\"auchli's
proof-theoretic semantics to Anders Kock.  It proposes parameters for
actual time rather than only imagined construction stages.  It says
that each definition of the reals determines an object in a topos
satisfying \((\omega)\), but records the analytic theorems then
available as unresolved.  These are historical attributions and the
article's 1970 questions, not assertions of a current unresolved
state \cite[p.~334]{Lawvere1970}.

The source bibliography identifies Mao's 1966 philosophical work;
Grothendieck--Giraud--Verdier's \emph{Topologie et faisceaux}; Hakim's
1967 Orsay thesis \emph{Sch\'emas relatifs}; Solovay--Scott's
Boolean-valued-model work as forthcoming; and Lawvere's 1969
\emph{Adjointness in Foundations} and 1970 work on equality and
comprehension in hyperdoctrines.  Those references locate the
article's dependencies; they are not substituted here for acquired
full proofs of the announced claims.

\begin{proofboundary}{def:p10-lawvere-sheaves-proof-boundary}{Source standing and the independent proof obligation}
The original article's complete account is distinct from a complete
modern development of its announced results.  The concrete and
internal presheaf packages, dependent products, internal
Heyting/predicate adjoints, stable images and effective quotients,
both internal-functor adjoints, choice equivalence, closure laws,
subterminal-generator argument, all four associated-sheaf
observations, sheaf characterization, left-exact reflector,
sheaf-topos/classifier package, geometric factorization,
double-negation classicality, and the intrinsic Godement
realization/image/cotriple reconstruction have
full, separately attributed editorial annex proofs. The exact ordinary
localization and logical preservation, Boolean-matrix classification
and choice at the declared cell regimes, and full-datum closed functor
with nondegenerate two-valued collapse likewise have their actual
editorial constructions above. General strong or normal closed-base
identification is not inferred from the arbitrary original datum;
the diagonal countercase and the separate normality proof for the
collapse case retain that distinction.

The explicit sufficient internal ccc, its finite-partial-map verification,
cardinal tools, internal-site bridge, complete no-epi predicate argument,
local injection descent and both strict internal CH comparisons now have
the six reciprocal editorial proofs in the continuation above. The
unprinted word suitable is not retroactively assigned a unique original
formal definition, and the proofs retain all declared base parameters
and cardinal/choice/NNO hypotheses.

Specified set-theoretic transfer/model interpretation, CH's positive
consistency direction, spectrum locality/universality and the independent
general recursion-lifting/free-algebra dependencies remain owed. Arbitrary
strong or normal-base obligations are not discharged by this internal
forcing argument. Historical source announcement or sketch does not
waive a separately promised full proof. The reciprocal source sketch
targets retain these actual remaining and historical boundaries.
\end{proofboundary}

% ED26-REQ-N-d-023-03: Mac Lane--Moerdijk 1992, source-specific semantics.
\subsection{Mac Lane--Moerdijk: internal predicates and three forcing scopes}
\label{sec:p10-n-d-023-source-03}

Mac Lane and Moerdijk construct presheaf examples before interpreting typed
language and its generalized-element and site forcing semantics
\cite[I\S4, pp.~35--38; I\S6, pp.~44--47; VI\S5--7, pp.~296--318]
{MacLaneMoerdijk1992}. Here \(\mathcal C\) is small in the fixed universe,
\(\widehat{\mathcal C}=\mathsf{Set}^{\mathcal C^{\mathrm{op}}}\), and
\(x\cdot u=X(u)(x)\) for \(x\in X(c)\), \(u:d\to c\).
The source's \(\Rightarrow\) denotes Heyting implication, not material
implication; \(f^{-1}\) on subobjects is the host's \(f^*\).

The elementary-topos package is Definition~\ref{def:p10-topos}.
Its presheaf realization is supplied by Theorems~\ref{thm:p10-presheaf-limits},
\ref{thm:p10-presheaf-mono-epi}, \ref{thm:p10-presheaf-exponentials},
\ref{thm:p10-presheaf-topos} and \ref{thm:p10-presheaf-power-objects}, with
reciprocal proofs in Annex items~\ref{annex:p10.psh.1}--\ref{annex:p10.psh.5}.
These are the foundation interfaces below, not proofs of the following
semantic claims. The source writes \(Q^P(c)=\operatorname{Nat}(\mathrm yc\times
P,Q)\), with evaluation \((\theta,x)\mapsto\theta_c(1_c,x)\), and sends \(x\)
to the characteristic sieve \(\{u:x\cdot u\in A\}\). These are the same
host constructions with \(P,Q\) renamed \(X,Y\), not pointwise function sets.

\begin{hbtheorem}{thm:p10-mlm-presheaf-predicates}{proved=annex:p10.mlm.1}{Heyting predicates and all-arrow quantifiers}
For subpresheaves \(A,B\subseteq X\), top and bottom are \(X\) and \(0\),
and meets and joins are componentwise intersections and unions. Put
\[
 (A\Rightarrow B)(c)=\{x\in X(c):\forall u:d\to c,\
          x\cdot u\in A(d)\ \Longrightarrow\ x\cdot u\in B(d)\}.
\]
Then \(R\subseteq(A\Rightarrow B)\) iff \(R\cap A\subseteq B\).
Thus \(\operatorname{Sub}(X)\) is a Heyting algebra, with
\((\neg A)(c)=\{x:\forall u:d\to c,\ x\cdot u\notin A(d)\}\).
For any natural map \(f:X\to Y\), inverse image preserves these Heyting
operations and has both order adjoints \(\exists_f\dashv f^*\dashv\forall_f\):
\[
\begin{aligned}
 (\exists_f A)(c)&=\{y\in Y(c):\exists x\in A(c),\ f_c(x)=y\},\\
 (\forall_f A)(c)&=\{y\in Y(c):\forall u:d\to c\ \forall x\in X(d),\
                  f_d(x)=y\cdot u\ \Longrightarrow\ x\in A(d)\}.
\end{aligned}
\]
Frobenius says \(\exists_f(A\cap f^*B)=(\exists_f A)\cap B\).
For a pullback \(P=X\times_Y Z\) with projections \(p:P\to X\), \(q:P\to Z\),
and \(g:Z\to Y\), both Beck--Chevalley identities hold:
\[
             g^*\exists_f=\exists_qp^*,\qquad
             g^*\forall_f=\forall_qp^*.
\]
All formulas include empty components and the empty base.
\end{hbtheorem}

Implication and negation are I\S8 Proposition~5 (18)--(19), p.~56, where
further verification is left to the reader. The correctly typed universal
image is III\S8 (15), p.~148. Frobenius and Beck--Chevalley are IV\S9 (2),
Proposition~1 and (13), pp.~204--209 \cite{MacLaneMoerdijk1992}.
The annex supplies direct componentwise proofs; the source's general-topos
route instead uses cartesian-closed slices and stable epi--mono images.

\begin{hbdefinition}{def:p10-mlm-internal-operations}{Internal operations on power objects}
Give \(P(X)(c)=\operatorname{Sub}(\mathrm yc\times X)\) its Heyting structure.
Restrictions are pullback along \(\mathrm yu\times1_X\), so preserve that
structure. Thus component operations define natural internal maps
\(\wedge,\vee,\Rightarrow:P(X)\times P(X)\to P(X)\), \(\top,\bot:1\to P(X)\),
and \(\neg S=(S\Rightarrow\bot)\). Internal order is the equalizer of
\(\wedge\) and first projection: \(S\leq T\) means \(S\wedge T=S\).
For \(f:X\to Y\), define \(P(f):P(Y)\to P(X)\) and
\(\exists_f,\forall_f:P(X)\to P(Y)\), at component \(c\), using the external
inverse image and its adjoints along \(1_{\mathrm yc}\times f\).
On \(\operatorname{Nat}(Z,P(X))\cong\operatorname{Sub}(Z\times X)\), they act
by the corresponding external operations along \(1_Z\times f\).
This separates a map of power objects from an operation on an external
subobject poset. In particular \(P(1)\cong\Omega\) is an internal Heyting algebra.
\end{hbdefinition}

\begin{hbtheorem}{thm:p10-mlm-internal-representation}{proved=annex:p10.mlm.2}{Representation of the predicate operations}
These maps are natural, satisfy the Heyting identities and both internal
order adjunctions, and act on every parameter object exactly as stated.
Frobenius and both Beck--Chevalley identities are identities of natural maps,
not merely of global elements. This is the representability route of IV\S8
Theorem~1 (Internal), IV\S9 Theorem~2 and Proposition~4, pp.~201--209
\cite{MacLaneMoerdijk1992}.
\end{hbtheorem}

\begin{hbdefinition}{def:p10-mlm-language}{The Mitchell--B\'enabou language and interpretation}
Fix an elementary topos \(\mathcal E\). Its objects are types, including
\(1\), products, function types \(Y^X\), \(\Omega\), and \(P(X)=\Omega^X\).
A finite declared context \(\Gamma=(x_1:X_1,\ldots,x_n:X_n)\) denotes
\(G=\prod_iX_i\), with \(G=1\) if \(n=0\). A term of type \(Y\) denotes an
arrow \(G\to Y\); a variable denotes its projection. Keep contexts explicit,
including unused variables. In a common context the inductive clauses are:
\begin{itemize}
\item Pairing denotes \(\langle s,t\rangle\); projections and the unit term
      use their categorical arrows.
\item Each arrow \(h:X\to Y\) supplies \(h(s)\), denoting \(h\circ s\).
      In particular \(h:1\to Y\) supplies a closed term.
\item For \(r:G\to Y^X\), \(s:G\to X\), application \(r(s)\) denotes
      \(\operatorname{ev}\circ\langle r,s\rangle\).
\item For \(t:G\times X\to Y\), \(\lambda x.t\) denotes
      \(\Lambda t:G\to Y^X\); binding is capture-avoiding.
\item Equality of \(s,t:G\to X\) denotes
      \(\delta_X\circ\langle s,t\rangle\), where \(\delta_X\) classifies
      the diagonal \(X\to X\times X\).
\item For \(a:G\to X\), \(r:G\to P(X)\), membership \(a\in r\) denotes
      evaluation on \(\langle r,a\rangle\).
\item Formulas are \(\Omega\)-valued terms. Constants, connectives and
      negation use the internal Heyting maps on \(\Omega\).
\item For \(\phi\) in \(\Gamma,x:X\), its extension \([\phi]\hookrightarrow
      G\times X\) is the pullback of truth. Along \(\pi:G\times X\to G\),
      the extensions of \(\exists x\,\phi\), \(\forall x\,\phi\) are
      \(\exists_\pi[\phi]\), \(\forall_\pi[\phi]\). Equivalently transpose
      \(\phi:G\times X\to\Omega\) to \(G\to P(X)\), then apply internal
      quantifiers along \(X\to1\).
\item Comprehension \(\{x:X\mid\phi\}\) as a predicate-valued term denotes
      \(\Lambda\phi:G\to P(X)\). The same braces as an object denote
      \([\phi]\hookrightarrow G\times X\), not that arrow.
\end{itemize}
These are VI\S5 pp.~298--301 \cite{MacLaneMoerdijk1992}, expressed in one
declared context. The source merges ordered lists of free variables using
projections; repeated variables occur once, different variables of the same
type remain different factors. This context convention makes those
projections explicit without deleting unused possibly empty types.
Natural-number, recursion and arithmetic constructions require an NNO when
invoked; they are not consequences of the elementary-topos definition alone.
\end{hbdefinition}

\begin{hbtheorem}{thm:p10-mlm-contextual-semantics}{proved=annex:p10.mlm.3}{Substitution, comprehension and contextual validity in presheaves}
All language clauses are defined in \(\widehat{\mathcal C}\).
For capture-avoiding simultaneous typed substitution denoting \(a:H\to G\),
term interpretation is composition with \(a\), and \([\phi[a]]=a^*[\phi]\).
In \(\Gamma,x:X\), membership in the comprehension
\(\{x:X\mid\phi\}\) has exactly extension \([\phi]\).
Evaluation and transposition validate typed \(\beta\)- and \(\eta\)-identities.
Equality denotes the equalizer of the compared terms and supports replacing
equal terms throughout expressions. Predicate equality is extensional:
\(r=s\) has the same extension as
\(\forall x:X\,(x\in r\Leftrightarrow x\in s)\), with
\(\Leftrightarrow\) abbreviating the two Heyting implications.
Function equality is likewise equivalent to quantified pointwise equality.
For finite hypotheses \(\Delta\), semantic consequence in context \(G\)
means \(\bigcap_{\theta\in\Delta}[\theta]\subseteq[\psi]\); the empty
intersection is \(G\). Validity means \([\phi]=G\), or characteristic arrow
\(\mathsf{true}\,!_G\), and is equivalent to validity of universal closure
in context \(1\). No type inhabitation or choice principle is presumed.
\end{hbtheorem}

This theorem supplies semantic construction, not soundness of an unspecified
formal higher-order calculus. VI\S5 p.~302 sends specification of the proof
rules to Boileau--Joyal or Lambek--Scott \cite{MacLaneMoerdijk1992}.
The separately attributed Lambek--Scott1980 specification below supplies an
explicit dependent rule surface and soundness proof; it does not stand in
for the distinct 1986 book or Boileau--Joyal1981 presentation.

\begin{hbtheorem}{thm:p10-mlm-topos-structural-package}{cited=MacLaneMoerdijk1992}{Dependencies of generalized forcing}
In an elementary topos, subobject posets are Heyting algebras; pullback
preserves their Heyting structure and has both adjoints, with Beck--Chevalley.
Joins and existential images are images of coproduct and projection maps.
Epi--mono images and finite coproducts are stable under pullback; a monic epi
is an isomorphism; the initial object is strict. These are the earlier
structural results used by VI\S6: IV\S6--7, IV\S8 Theorem~1, and IV\S9.
They are not assumptions asserting the forcing clauses. For presheaves their
relevant constructions are provided by the five foundation proofs and
Theorem~\ref{thm:p10-mlm-presheaf-predicates}.
\end{hbtheorem}

\begin{hbdefinition}{def:p10-mlm-generalized-forcing}{Forcing at a generalized element}
For \(\alpha:U\to G\), \(U\forces\phi(\alpha)\) means that \(\alpha\) factors
through \([\phi]\hookrightarrow G\), equivalently
\(\operatorname{Im}(\alpha)\subseteq[\phi]\).
For a closed formula use \(U\to1\); for multiple variables use their tuple.
This is VI\S6 (1)--(2), p.~303 \cite{MacLaneMoerdijk1992}.
Here \(\forces\) is semantic, not derivability or set-theoretic forcing.
\end{hbdefinition}

\begin{hbtheorem}{thm:p10-mlm-generalized-clauses}{proved=annex:p10.mlm.4}{Generalized-element Kripke--Joyal clauses}
For \(\alpha:U\to G\), forcing persists under every \(p:V\to U\) and
descends along an epimorphism. Its clauses are:
\begin{itemize}
\item Top always holds; bottom holds iff \(U\cong0\).
\item Conjunction holds iff both conjuncts hold at \(\alpha\).
\item Disjunction holds iff there are \(p:V\to U\), \(q:W\to U\) with
      \([p,q]:V+W\to U\) epic, \(V\forces\phi(\alpha p)\) and
      \(W\forces\psi(\alpha q)\).
\item Implication holds iff every \(p:V\to U\) with
      \(V\forces\phi(\alpha p)\) also has \(V\forces\psi(\alpha p)\).
\item Negation holds iff every \(p:V\to U\) with
      \(V\forces\phi(\alpha p)\) has \(V\cong0\).
\item Existence \(U\forces\exists x:X\,\phi\) holds iff there are an epic
      \(p:V\to U\) and \(b:V\to X\) with \(V\forces\phi(\alpha p,b)\).
\item Universality \(U\forces\forall x:X\,\phi\) holds iff every \(p:V\to U\)
      and \(b:V\to X\) has \(V\forces\phi(\alpha p,b)\), equivalently
      \(U\times X\forces\phi(\alpha\pi_1,\pi_2)\).
\end{itemize}
Equality of terms \(s,t:G\to Y\) is forced iff \(s\alpha=t\alpha\).
Membership is forced iff the interpreted element--predicate tuple factors
through membership. These statements use the preceding structural package,
not enough points. They are VI\S6 Theorem~1 and Propositions~2--3,
pp.~304--310 \cite{MacLaneMoerdijk1992}; the annex expands the source's reader
verification for negation and membership too.
\end{hbtheorem}

\subsubsection{An explicit dependent rule specification: Lambek--Scott (1980)}
\label{sec:p10-ls1980-dependent-calculus}

Lambek and Scott's \(\mathcal L_1\) is a precise product/power type theory,
not literally the preceding exponential-term syntax. We use the actual
1980 article's \S1 pp.~216--218 and \S5 pp.~240--242
\cite{LambekScott1980ITT} as a separately qualified dependent specification.
The 1980 logical connectives are primitive and equality is Leibniz-defined;
neither the primitive-equality1983 version nor the 1986 book is merged into it.

\begin{hbdefinition}{def:p10-ls1980-language}{The source's product/power grammar and its declared fragments}
The full pure source language has basic types \(1,N,\Omega\), products
\(A\times B\), and power types \(PA\), with countably many variables of each
type. Its ten term clauses are: (1) a typed variable; (2) \(*:1\);
(3) \(\langle a,b\rangle:A\times B\) from \(a:A,b:B\);
(4) \(a\in u:\Omega\) from \(a:A,u:PA\);
(5) \(\{x:A\mid\phi(x)\}:PA\) from an \(\Omega\)-term, with binding of \(x\);
(6) \(0:N\); (7) \(Sn:N\) from \(n:N\);
(8) \(\top,\bot:\Omega\); (9) \(p\wedge q,p\vee q,p\Rightarrow q:\Omega\);
(10) \(\forall x:A\,\phi\), \(\exists x:A\,\phi:\Omega\).
Formulas are the terms of type \(\Omega\), and
\[
 \neg p=(p\Rightarrow\bot),\qquad
 p\Leftrightarrow q=(p\Rightarrow q)\wedge(q\Rightarrow p),\qquad
 a=_Ab=\forall u:PA\,(a\in u\Leftrightarrow b\in u),
\]
where the last \(u\) is fresh. There is no primitive equality term,
projection term, function type, lambda binder or application term in this
pure grammar. All binding and substitution use the source's free-for
condition, here made explicitly capture-avoiding.

The nonarithmetic fragment \(\mathcal L_1^-\) omits \(N,0,S\) and Peano
axioms; it is an explicitly identified restriction, not the full source
language renamed. The arithmetic extension retains exactly those source
constructors and axioms and requires an interpreted NNO. The optional
Boolean axiom is independent of that arithmetic extension.
The source's applied internal-language extension (\S5 p.~242) permits named
objects as types and named arrows as term operations: for an interpreted
\(h:A\to B\), \(h(a):B\) denotes composition. This extension allows arbitrary
object types, including uninhabited ones; it does not silently add a
primitive function/lambda grammar.
\end{hbdefinition}

\begin{hbdefinition}{def:p10-ls1980-contextual-validity}{Source-indexed entailment and interpretation}
The source relation \(p\vdash_Xq\) is indexed by a declared set \(X\) of
typed variables, containing both formulas' free variables; it is the least
relation satisfying the displayed structural, logical and nonlogical rows
of \ref{sheet:p10-internal-hol-ls1980}. The sheet splits each source logical
``if and only if'' into its rule directions and makes quantifier freshness
explicit. \(\vdash_Xp\) abbreviates \(\top\vdash_Xp\), not truth in a smaller
undeclared context.

Interpret types as the corresponding objects, products and power objects;
interpret the ten constructors by projections, pairing, membership,
comprehension, internal Heyting operations and quantification as above.
For possibly infinite \(X\), a generalized assignment at \(U\) is a family
\((a_x:U\to A_x)_{x\in X}\). Each term has finite support, so is evaluated
using the finite tuple for its free variables; no infinite context product
is presumed. Define \(p\models_Xq\) to mean: at every object \(U\) and every
such assignment, forcing of \(p\) implies forcing of \(q\).
For a finite \(X\), this is exactly
\([p]_G\subseteq[q]_G\) in its full declared context \(G=\prod_{x\in X}A_x\).
Generalized assignments retain unused declarations; if no assignment exists
at a noninitial \(U\), the corresponding test is vacuous, not an instruction
to delete its empty types.
\end{hbdefinition}

\begin{hbtheorem}{thm:p10-ls1980-leibniz-equality}{proved=annex:p10.mlm.7}{The Leibniz--diagonal interpretation bridge}
In an elementary topos, using the explicit structural package of
Theorem~\ref{thm:p10-mlm-topos-structural-package}, the source's Leibniz
formula \(a=_Ab\) has exactly the diagonal-equality extension for every pair
of contextual terms. The substitution and comprehension interpretations
therefore agree with the source's \S5 diagonal interpretation, pp.~241--242
\cite{LambekScott1980ITT}, without changing the source's defined syntax.
The same statement holds concretely in the already constructed presheaf
models, where the structural ingredients have the linked component proofs.
\end{hbtheorem}

% Lambek--Scott 1980 L_1, section 1 pp.216--218; biconditionals split into rules.
\calculussheet{sheet:p10-internal-hol-ls1980}{Lambek--Scott 1980 product/power type theory}{Source: \cite[\S1, Definitions~1.1--1.2 and axioms, pp.~216--218]{LambekScott1980ITT}. Types are \(1,N,\Omega\), closed under \(A\times B\) and \(PA\); terms are exactly the ten clauses in Definition~\ref{def:p10-ls1980-language}. The explicitly nonarithmetic fragment omits \(N,0,S\) and Peano rows; the applied extension names additional interpreted object types and arrows as in source \S5 p.~242. Judgments \(p\vdash_Xq\) have formula free variables contained in the declared set \(X\); \(\vdash_Xp\) abbreviates \(\top\vdash_Xp\). All substitutions are typed and free-for/capture-avoiding. Write \(\neg p=(p\Rightarrow\bot)\), \(p\Leftrightarrow q=(p\Rightarrow q)\wedge(q\Rightarrow p)\), and \(a=_Ab\) for \(\forall u:PA\,(a\in u\Leftrightarrow b\in u)\), with \(u\) fresh. Equality is defined, not primitive; neither function/lambda types nor projections are primitive terms of pure \(\mathcal L_1\). Each source logical iff is displayed as its forward and reverse rules. Quantifier freshness is explicit; optional Peano rows require the NNO interpretation and the Boolean row requires Booleanness.}
\hbrule{rule:p10-internal-hol-ls1980-refl}{Reflexivity}{}{p\vdash_Xp}{\none}
\hbrule{rule:p10-internal-hol-ls1980-trans}{Transitivity}{p\vdash_Xq\qquad q\vdash_Xr}{p\vdash_Xr}{\none}
\hbrule{rule:p10-internal-hol-ls1980-weaken}{Variable weakening}{p\vdash_Xq}{p\vdash_Yq}{\(X\subseteq Y\); declared types are unchanged. No unused-variable deletion rule.}
\hbrule{rule:p10-internal-hol-ls1980-subst}{Typed substitution}{\phi(y)\vdash_{X\cup\{y:B\}}\psi(y)}{\phi(b)\vdash_X\psi(b)}{\(b:B\), \(\mathrm{FV}(b)\subseteq X\); \(b\) is free for \(y\) in both formulas. If \(y\in X\), its original assigned value is used when evaluating \(b\), and the premise receives the substituted value at \(y\).}
\hbrule{rule:p10-internal-hol-ls1980-top}{Logical (1)}{}{p\vdash_X\top}{\none}
\hbrule{rule:p10-internal-hol-ls1980-bottom}{Logical (2)}{}{\bot\vdash_Xp}{\none}
\hbrule{rule:p10-internal-hol-ls1980-meet-intro}{Logical (3), reverse}{p\vdash_Xq\qquad p\vdash_Xr}{p\vdash_Xq\wedge r}{\none}
\hbrule{rule:p10-internal-hol-ls1980-meet-left}{Logical (3), forward 1}{p\vdash_Xq\wedge r}{p\vdash_Xq}{\none}
\hbrule{rule:p10-internal-hol-ls1980-meet-right}{Logical (3), forward 2}{p\vdash_Xq\wedge r}{p\vdash_Xr}{\none}
\hbrule{rule:p10-internal-hol-ls1980-join-elim}{Logical (4), reverse}{p\vdash_Xr\qquad q\vdash_Xr}{p\vee q\vdash_Xr}{\none}
\hbrule{rule:p10-internal-hol-ls1980-join-left}{Logical (4), forward 1}{p\vee q\vdash_Xr}{p\vdash_Xr}{\none}
\hbrule{rule:p10-internal-hol-ls1980-join-right}{Logical (4), forward 2}{p\vee q\vdash_Xr}{q\vdash_Xr}{\none}
\hbrule{rule:p10-internal-hol-ls1980-residual-intro}{Logical (5), forward}{p\wedge q\vdash_Xr}{p\vdash_Xq\Rightarrow r}{\none}
\hbrule{rule:p10-internal-hol-ls1980-residual-elim}{Logical (5), reverse}{p\vdash_Xq\Rightarrow r}{p\wedge q\vdash_Xr}{\none}
\hbrule{rule:p10-internal-hol-ls1980-forall-intro}{Logical (6), reverse}{p\vdash_{X\cup\{y:B\}}\psi(y)}{p\vdash_X\forall y:B\,\psi(y)}{\(y\notin\mathrm{FV}(p)\); other free variables of \(\psi\) lie in \(X\). If \(y\in X\), its declared value is overridden when testing the bound variable.}
\hbrule{rule:p10-internal-hol-ls1980-forall-elim}{Logical (6), forward}{p\vdash_X\forall y:B\,\psi(y)}{p\vdash_{X\cup\{y:B\}}\psi(y)}{\(y\notin\mathrm{FV}(p)\); other free variables of \(\psi\) lie in \(X\).}
\hbrule{rule:p10-internal-hol-ls1980-exists-elim}{Logical (7), reverse}{\psi(y)\vdash_{X\cup\{y:B\}}p}{\exists y:B\,\psi(y)\vdash_Xp}{\(y\notin\mathrm{FV}(p)\); other free variables of \(\psi\) lie in \(X\).}
\hbrule{rule:p10-internal-hol-ls1980-exists-intro}{Logical (7), forward}{\exists y:B\,\psi(y)\vdash_Xp}{\psi(y)\vdash_{X\cup\{y:B\}}p}{\(y\notin\mathrm{FV}(p)\); other free variables of \(\psi\) lie in \(X\).}
\hbrule{rule:p10-internal-hol-ls1980-comprehension}{Nonlogical (a)}{}{\vdash_X\forall x:A\,(x\in\{x:A\mid\phi(x)\}\Leftrightarrow\phi(x))}{Parameters of \(\phi\) other than its displayed \(x\) lie in \(X\); binders are renamed to avoid capture, not to delete declared variables.}
\hbrule{rule:p10-internal-hol-ls1980-power-ext}{Nonlogical (b), power}{}{\vdash_{\varnothing}\forall u:PA\,\forall v:PA\,((\forall x:A\,(x\in u\Leftrightarrow x\in v))\Rightarrow u=_{PA}v)}{\(u,v,x\) are distinct bound variables.}
\hbrule{rule:p10-internal-hol-ls1980-proposition-ext}{Nonlogical (b), propositions}{}{\vdash_{\varnothing}\forall s:\Omega\,\forall t:\Omega\,((s\Leftrightarrow t)\Rightarrow s=_{\Omega}t)}{\(s,t\) are distinct bound variables.}
\hbrule{rule:p10-internal-hol-ls1980-unit}{Nonlogical (c), unit}{}{\vdash_{\varnothing}\forall x:1\,x=_1*}{\none}
\hbrule{rule:p10-internal-hol-ls1980-product-surj}{Nonlogical (c), decomposition}{}{\vdash_{\varnothing}\forall z:A\times B\,\exists x:A\,\exists y:B\,z=_{A\times B}\langle x,y\rangle}{\(z,x,y\) are distinct bound variables.}
\hbrule{rule:p10-internal-hol-ls1980-product-inj}{Nonlogical (c), injectivity}{}{\vdash_{\varnothing}\forall x:A\,\forall x':A\,\forall y:B\,\forall y':B\,(\langle x,y\rangle=_{A\times B}\langle x',y'\rangle\Rightarrow(x=_Ax'\wedge y=_By'))}{All four bound variables are distinct.}
\hbrule{rule:p10-internal-hol-ls1980-peano-zero}{Optional nonlogical (d), zero}{}{\vdash_{\varnothing}\forall x:N\,(Sx=_N0\Rightarrow\bot)}{Arithmetic extension only; \(N\) is interpreted as an NNO with its zero and successor.}
\hbrule{rule:p10-internal-hol-ls1980-peano-inj}{Optional nonlogical (d), successor}{}{\vdash_{\varnothing}\forall x:N\,\forall y:N\,(Sx=_NSy\Rightarrow x=_Ny)}{Arithmetic extension only; same NNO interpretation.}
\hbrule{rule:p10-internal-hol-ls1980-peano-induction}{Optional nonlogical (d), induction}{}{\vdash_{\varnothing}\forall u:PN\,((0\in u\wedge\forall x:N\,(x\in u\Rightarrow Sx\in u))\Rightarrow\forall y:N\,y\in u)}{Arithmetic extension only; same NNO interpretation. Bound \(u,x,y\) are distinct.}
\hbrule{rule:p10-internal-hol-ls1980-boolean}{Optional Boolean axiom}{}{\vdash_{\varnothing}\forall t:\Omega\,(t\vee\neg t)}{Classical extension only; the interpreting topos is Boolean. Independent of whether Peano rows are included.}


\begin{hbtheorem}{thm:p10-ls1980-hol-soundness}{proved=annex:p10.mlm.8}{Rule-by-rule intuitionistic higher-order soundness}
Every derivation using the nonarithmetic intuitionistic rows (excluding both
optional extensions) of
\ref{sheet:p10-internal-hol-ls1980}, in either the pure or the declared
object/arrow-named applied language, satisfies
\[
                 p\vdash_Xq\quad\Longrightarrow\quad p\models_Xq.
\]
In finite contexts this is the subobject inclusion of
Definition~\ref{def:p10-ls1980-contextual-validity}; in sheaf models it uses
the site forcing clauses, not pointwise Boolean satisfaction. The proof
validates every rule direction and nonlogical axiom before induction on the
derivation. An arbitrary elementary-topos interpretation uses only the
identified earlier structural package, not soundness as an assumed premise.

The full arithmetic rows are also sound when \(N,0,S\) are interpreted as an
NNO and its zero and successor; the annex derives the three Peano properties
from that universal property. Presheaves have the explicitly constructed
constant natural-number NNO. The Boolean row is sound precisely under
Booleanness, by Theorem~\ref{thm:p10-mlm-internal-epi-boolean-choice}.
Neither NNO nor Booleanness is required for the nonarithmetic intuitionistic
fragment. No completeness, normalization, realizability or free-topos
metatheorem is claimed here.
\end{hbtheorem}

\begin{hbtheorem}{thm:p10-ls1980-exponential-graph-bridge}{proved=annex:p10.mlm.9}{The graph bridge to exponential terms}
Inside \(P(A\times B)\), let \(\operatorname{Fun}(A,B)\) be the subobject
defined, in the source's product/power language, by
\[
 \forall x:A\,\exists y:B\left(\langle x,y\rangle\in R\wedge
       \forall y':B\,(\langle x,y'\rangle\in R\Rightarrow y'=_{B}y)\right).
\]
There is a canonical isomorphism \(\operatorname{Fun}(A,B)\cong B^A\),
natural on all parameter objects. A function is sent to its graph, and a
total single-valued relation is sent to its unique arrow. Under it membership
in the graph is equality with evaluation. A lambda term's graph is the
source comprehension
\[
 \{z:A\times B\mid\exists x:A\,\exists y:B\,
             (z=_{A\times B}\langle x,y\rangle\wedge y=_Bt(x))\}.
\]
It represents exactly the exponential transpose of the interpreted body.
Graph extensionality and the inverse representation give the typed beta and
eta conversions, including empty types. Thus the larger exponential-term
interpretation has an actual object/arrow bridge to the applied source
language. This is not a claim that function types or lambdas were primitive
in \(\mathcal L_1\), nor a proof of syntactic conservativity of a pure-language
extension. If exponential-term macros or their beta/eta conversions are used
with the source logical rules, their denotation-preserving expansion and the
same rule validation preserve soundness.
\end{hbtheorem}

Freshness is consequential. Allowing \(y\) to remain a parameter of \(p\)
in logical (6) would infer \(p\vdash_X\forall y\,p\) from reflexivity when
\(p\) depends on \(y\); in a two-element set with \(p(y)\) selecting one
element that conclusion is false. Conversely deleting an unused variable
of empty type changes the contextual meaning. Neither maneuver is a rule
of the declared interpretation.

\begin{hbdefinition}{def:p10-mlm-site-setup}{Site setup and its separate construction boundary}
A Grothendieck topology \(J\) assigns covering sieves \(J(c)\), contains the
maximal sieve, is stable under inverse image, and is transitive: if \(S\in J(c)\)
and \(u^*R\in J(d)\) for every \(u:d\to c\) in \(S\), then \(R\in J(c)\).
A sheaf uniquely amalgamates matching families on covers. A subsheaf has
local membership: a section belongs iff its restrictions on a cover belong.
Take the already established \(\mathcal E=\operatorname{Sh}(\mathcal C,J)\)
and associated-sheaf adjunction \(a\dashv i\); their construction is outside
this overview. Yoneda and the adjunction give
\[
 X(c)\cong\operatorname{Nat}(\mathrm yc,iX)\cong\mathcal E(a(\mathrm yc),X).
\]
For \(\gamma\in G(c)\), define \(c\forces\phi(\gamma)\) by
\(\gamma\in[\phi](c)\). It is generalized forcing from \(a(\mathrm yc)\),
not from a possibly nonsheaf \(\mathrm yc\). This identification is VI\S7
(1)--(3), pp.~315--316 \cite{MacLaneMoerdijk1992}; no sheafification or
geometric-morphism construction is asserted here.
\end{hbdefinition}

\begin{hbtheorem}{thm:p10-mlm-site-clauses}{proved=annex:p10.mlm.5}{Site and presheaf stage clauses}
At a site stage, equality and membership use their interpreted components,
conjunction means both at \(c\), top always holds, and bottom holds iff
\(\varnothing\in J(c)\). For \(\gamma\in G(c)\):
\begin{itemize}
\item Disjunction holds iff some covering sieve \(S\) has, for every
      \(u:d\to c\) in \(S\), either \(d\forces\phi(\gamma\cdot u)\) or
      \(d\forces\psi(\gamma\cdot u)\).
\item Implication holds iff every \(u:d\to c\) forcing the antecedent at
      \(\gamma\cdot u\) also forces the consequent there.
\item Negation holds iff every positive-forcing \(u:d\to c\) has
      \(\varnothing\in J(d)\).
\item Existence holds iff some covering sieve \(S\) has for each \(u:d\to c\)
      in \(S\) a \(b\in X(d)\) with \(d\forces\phi(\gamma\cdot u,b)\).
\item Universality holds iff for every \(u:d\to c\) and \(b\in X(d)\),
      \(d\forces\phi(\gamma\cdot u,b)\).
\end{itemize}
Forcing persists under restriction and is local on covers. At the trivial
topology, only the maximal sieve covers: disjunction then has a disjunct at
\(c\), existence a witness in \(X(c)\), and negation excludes every
positive-forcing restriction. Bottom is never forced at a base object.
Implication and universality still test all arrows, including identity.
No products in \(\mathcal C\) are required.
\end{hbtheorem}

These are VI\S7 Theorem~1 and (ii\(^{\prime}\)), (iv\(^{\prime}\)),
(v\(^{\prime}\)), pp.~316--318 \cite{MacLaneMoerdijk1992}.
The source's first proof route uses explicit subsheaf operations from III\S8;
the annex expands it. The other uses generalized forcing and the generation
and cover-detection properties of sheafified representables (III\S6,
III\S7 Corollary~7). Those are separate dependencies, not an enough-points
assumption. The generic-element test \(U\times X\) does not become one base
stage without further hypotheses; VI\S7 Remark~2 assumes products in
\(\mathcal C\) and a quantified type \(a(\mathrm yB)\).

\begin{hbexample}{def:p10-mlm-arrow-failures}{Two stages and the false pointwise substitute}
Let \(a:0\to1\) be the only nonidentity arrow of \(\mathcal C\).
In the terminal presheaf take \(A(0)=\{*\}\), \(A(1)=\varnothing\).
Then \(\neg A=0\), since restriction along \(a\) reaches \(A(0)\), and
\(\neg\neg A=1\neq A\). At \(1\), neither \(A\) nor \(\neg A\) holds.
Pointwise complement would put \(*\) at \(1\) but not \(0\), so is not a
subpresheaf. For \(f:A\to1\), \(\exists_f A=A\) as a subobject of \(1\),
whereas \(\forall_f0=0\). A component-only universal formula would falsely
put \(*\) at \(1\), because \(A(1)\) is empty; \(a\) and \(*\in A(0)\)
refute it. The classifier has two sieves at \(0\) and three at \(1\):
\(\varnothing\), \(\{a\}\), and \(\{a,1_1\}\). The middle sieve becomes
true after restriction but is not true at \(1\).
This matches the source's arrow-category example after reversing its arrow
for the host's contravariant convention \cite[I\S4, pp.~35--36; I\S8, p.~56]
{MacLaneMoerdijk1992}.
\end{hbexample}

\begin{hbexample}{def:p10-mlm-locality-failures}{Locality, choice, and unused types}
In \(\mathsf{Set}\), let \(U=\{0,1\}\), with predicates supported on the
respective singletons. Their disjunction holds at generalized element
\(1_U\), although neither disjunct does: the two inclusions have epic
copairing. The one presheaf base-stage representable for \(\mathsf{Set}\)
is a singleton, not \(U\), so its disjunction clause is different.
For a two-element group \(H\), the map from its regular right action to
the trivial singleton action is epic but has no equivariant section:
such a section would select a point fixed by the nonidentity element.
Internal existence is therefore not natural choice. Invariant subsets
nevertheless have invariant complements, so the subobject Heyting algebras
of \(H\)-sets are Boolean, as in I\S8 p.~56 \cite{MacLaneMoerdijk1992}.
Finally \(\bot\) is valid in declared context \(x:0\), whose context is
initial, but not in the empty context of a nondegenerate topos. Deleting
unused variables is not validity-preserving without inhabitation.
If the empty sieve covers a site object \(c\), every sheaf has one section
there and bottom is forced; bottom is not universally absent at site stages.
\end{hbexample}

\begin{hbtheorem}{thm:p10-mlm-internal-epi-boolean-choice}{proved=annex:p10.mlm.6}{Three internal-to-external translations}
For \(r:U\to Y^X\) with transpose \(\widetilde r:U\times X\to Y\),
\[
 U\forces\forall y:Y\,\exists x:X\,r(x)=y
 \ \Longleftrightarrow\
 \langle\pi_1,\widetilde r\rangle:U\times X\to U\times Y\text{ is epic}.
\]
A topos is Boolean iff \(\forall p:\Omega\,(p\vee\neg p)\) is valid.
For all types \(X,Y\), validity of
\[
 \forall f:Y^X\left((\forall y:Y\,\exists x:X\,f(x)=y)
       \Rightarrow\exists s:X^Y\,\forall y:Y\,f(s(y))=y\right)
\]
is equivalent to every \((-)^E\) preserving epimorphisms, for every object
\(E\). This is the source's internal axiom of choice (IAC), not the external
assertion that every epi has a global section, and is not assumed here.
The translations are VI\S5 (7)--(10), VI\S6 (10)--(26), pp.~301--315
\cite{MacLaneMoerdijk1992}.
\end{hbtheorem}

For the dense topology on a poset, a sieve covers \(p\) iff every \(q\leq p\)
has an \(r\leq q\) in it. Thus disjunction and existence require, after
each refinement, a further refinement with a disjunct or witness.
Implication and universality test all \(q\leq p\); negation allows no
positive-forcing refinement. This is VI\S7's dense/double-negation example
on p.~318, not a proof of continuum-hypothesis independence or a transfer
back to set theory.

% Source-bound reconstruction of the corrected/expanded Abramsky--Jung author PDF.
\section{Domain theory and recursive semantic models}
\label{sec:p10-domain-theory-recursive-models}

Recursive values and recursive types pose different fixed-point problems.
The first asks for an element \(d=f(d)\) in a given ordered structure;
the second asks for a structure and an actual isomorphism
\(D\cong F(D)\), possibly with contravariant occurrences of \(D\).
This treatment follows Abramsky and Jung's \emph{Corrected and expanded
version} of \emph{Domain Theory}, with its own page and result numbering
\cite[pp.~1,5--7]{AbramskyJungCorrectedExpandedDomainTheory}.
That author presentation retains numbering while changing and explaining
the published text; it is not identified with the printed 1994 chapter,
the separate PostScript manifestation, or other semantic-domain sources.

\subsection{Directed convergence and approximation}

\begin{hbconvention}{def:p10-domain-order-convention}{Information order and nonempty directed sets}
Write \(x\sqsubseteq y\) for the information order of a poset, not for
an independently given numerical order. A subset \(A\) is
\emph{directed} when it is nonempty and every two members have an upper
bound belonging to \(A\). Consequently every finite subset of \(A\)
has an upper bound in \(A\). Write \(\bigsqcup^{\uparrow}A\) for its
supremum when this exists, and \(\downarrow x=\{y:y\sqsubseteq x\}\).
A \emph{dcpo} is a poset with suprema of all directed subsets. A
\emph{pointed} dcpo has a least element \(\bot\). Pointedness is an
additional condition: a nonempty discrete poset is a dcpo but generally
has no least element. The empty poset is allowed as a dcpo, but cannot
be pointed. All constructions here use set-sized objects and maps;
claims explicitly involving countable unions or enumerations use ordinary
set theory with choice. No reduction of arbitrary directed sets to
countable chains is assumed
\cite[Definitions~2.1.1,2.1.3,2.1.8,2.1.13, pp.~10--15]
{AbramskyJungCorrectedExpandedDomainTheory}.
\end{hbconvention}

\begin{hbdefinition}{def:p10-domain-continuous-map}{Scott-continuous and strict maps}
For dcpo's \(D,E\), a map \(f:D\to E\) is Scott-continuous when it is
monotone and
\[
 f(\bigsqcup^{\uparrow} A)=\bigsqcup^{\uparrow}_{a\in A}f(a)
                    \qquad(A\subseteq D\text{ directed}).
\]
A map between pointed dcpo's is \emph{strict} when
\(f(\bot_D)=\bot_E\). Let \([D\longrightarrow E]\) be all
Scott-continuous maps, ordered pointwise, and
\([D\mathrel{\mathop{\longrightarrow}^{\bot!}}E]\) the strict ones.
The category \(\mathsf{DCPO}\) has all dcpo's and Scott-continuous maps;
\(\mathsf{DCPO}_{\bot}\) restricts objects to pointed dcpo's;
\(\mathsf{DCPO}_{\bot!}\) also restricts maps to strict ones.
These are different categories, not different spellings of one category
\cite[Definition~2.1.17; pp.~15,41]{AbramskyJungCorrectedExpandedDomainTheory}.
\index{dcpo}\index{Scott continuity}\index{strict map!domain theory}
\end{hbdefinition}

\begin{hbdefinition}{def:p10-domain-way-below}{Approximation, bases and domains}
For \(x,y\in D\), put \(x\ll y\) when, for every directed
\(A\subseteq D\),
\[
 y\sqsubseteq\bigsqcup^{\uparrow}A
       \quad\Longrightarrow\quad
       (\exists a\in A)\ x\sqsubseteq a.
\]
Write \(\mathord{\downarrow\!\downarrow}y=\{x:x\ll y\}\),
\(\mathord{\uparrow\!\uparrow}x=\{y:x\ll y\}\), and
\(K(D)=\{x:x\ll x\}\). Members of \(K(D)\) are compact elements.
A subset \(B\subseteq D\) is a basis when, for every \(x\),
\(B_x=B\cap\mathord{\downarrow\!\downarrow}x\) contains a directed
subset with supremum \(x\). This is the source's basis definition;
directedness of the entire \(B_x\) is a consequence, not a second
unmentioned premise. A continuous domain is a dcpo with a basis; an
algebraic domain has a basis of compact elements. The prefixes
\(\omega\)-continuous and \(\omega\)-algebraic mean that a countable
basis exists, respectively that \(K(D)\) is a countable basis.
Here ``domain'' includes approximation as well as convergence
\cite[Definitions~2.2.1,2.2.3,2.2.6, pp.~18--19]
{AbramskyJungCorrectedExpandedDomainTheory}.
\index{way-below relation}\index{continuous domain}\index{algebraic domain}
\end{hbdefinition}

\begin{hbtheorem}{thm:p10-domain-basis}{proved=annex:p10.domain.basis}{Basis properties and interpolation}
In any dcpo, \(x\ll y\) implies \(x\sqsubseteq y\), and
\(x'\sqsubseteq x\ll y\sqsubseteq y'\) implies \(x'\ll y'\).
For a basis \(B\), every \(B_x\) is directed with supremum \(x\),
and \(K(D)\subseteq B\). Hence continuity is equivalent to
\(x=\bigsqcup^{\uparrow}\mathord{\downarrow\!\downarrow}x\) for
every \(x\); algebraicity is equivalent to the same assertion with
\(K(D)\cap\downarrow x\).
If \(D\) is continuous, \(M\subseteq D\) is finite, and every
\(m\in M\) satisfies \(m\ll y\), there is \(b\in B\) with
\(m\ll b\ll y\) for all \(m\in M\), including the empty \(M\).
In particular,
\[
 \mathord{\downarrow\!\downarrow}\bigsqcup^{\uparrow}A
       =\bigcup_{a\in A}\mathord{\downarrow\!\downarrow}a
\]
for directed \(A\) in a continuous domain
\cite[Propositions~2.2.2,2.2.4,2.2.7; Lemma~2.2.15;
Corollary~2.2.16, pp.~18--25]{AbramskyJungCorrectedExpandedDomainTheory}.
\end{hbtheorem}

\begin{hbtheorem}{thm:p10-domain-basic-examples}{proved=annex:p10.domain.examples}{Finite information and genuinely directed examples}
The powerset \(\mathcal P(S)\), ordered by inclusion, is algebraic;
directed suprema are unions, compact elements are exactly finite subsets,
and \(A\ll B\) exactly when \(A\) is finite and \(A\subseteq B\).
For uncountable \(S\), its directed family of finite subsets is not a
chain and has no cofinal chain.
The partial functions \([X\rightharpoonup Y]\), ordered by graph
inclusion, form a pointed algebraic dcpo. Their directed suprema are
graph unions and their compact elements are the finite partial functions.
The interval \([0,1]\), with its numerical order, is continuous but
not algebraic: \(x\ll y\) exactly when \(x=0\) or \(x<y\), and
\(K([0,1])=\{0\}\). A countable basis is the rational points in
\([0,1]\). A dcpo need not be continuous: two incomparable increasing
\(\omega\)-chains with a common supremum, even after adding bottom,
have only bottom as an approximant to their nonbottom elements.
These reconstruct the source examples and Exercises~2.3.9(6),(19),
with the additional pointed two-chain comparison made explicit
\cite[pp.~13,18--21,31--32]{AbramskyJungCorrectedExpandedDomainTheory}.
\end{hbtheorem}

\begin{hbdefinition}{def:p10-domain-abstract-basis}{Abstract bases and round ideals}
An abstract basis \((B,\prec)\) has a transitive relation such that
for every finite \(M\subseteq B\) and \(x\in B\),
\[
  (\forall m\in M)\ m\prec x
       \quad\Longrightarrow\quad
  (\exists y\in B)\ (\forall m\in M)\ m\prec y\prec x.
\]
An ideal is a nonempty lower subset directed with respect to \(\prec\):
for \(a,b\) in it some \(c\) in it satisfies \(a\prec c\) and
\(b\prec c\). Taking \(a=b\) enforces roundness. Let
\(\operatorname{Idl}(B,\prec)\) be its ideals ordered by inclusion,
and \(i(b)=\{a:a\prec b\}\). Neither reflexivity nor
\(b\in i(b)\) is assumed
\cite[Definitions~2.2.20--21, pp.~26--27]
{AbramskyJungCorrectedExpandedDomainTheory}.
\end{hbdefinition}

\begin{hbtheorem}{thm:p10-domain-ideal-completion}{proved=annex:p10.domain.ideals}{Actual ideal completion and recovery from a basis}
The ideal completion is a continuous dcpo with directed unions as
suprema and basis \(i(B)\). Its approximation relation is
\[
 A\ll A'\quad\Longleftrightarrow\quad
       (\exists b\in A')\ A\subseteq i(b).
\]
For reflexive \(\prec\) it is algebraic. In particular the ideal
completion of a poset \(B\) has compact elements exactly its
principal ideals. If \(B\) is a basis of a continuous domain \(D\),
with \(\prec\) the restricted \(\ll\), the inverse isomorphisms are
\[
 A\longmapsto\bigsqcup^{\uparrow}A,
       \qquad x\longmapsto B_x.
\]
If \(f:D\to E\) is continuous, its values on that basis determine it:
\(f(x)=\bigsqcup^{\uparrow}_{b\in B_x}f(b)\)
\cite[Propositions~2.2.22,2.2.25; Corollary~2.2.26, pp.~27--28]
{AbramskyJungCorrectedExpandedDomainTheory}.
\end{hbtheorem}

\begin{hbdefinition}{def:p10-domain-scott-topology}{Scott-open observations}
A Scott-open subset \(O\subseteq D\) is upper and inaccessible by
directed suprema: \(\bigsqcup^{\uparrow}A\in O\) implies
\(A\cap O\ne\varnothing\). Equivalently its complement is lower
and closed under directed suprema. Write \(\sigma D\) for these
opens. This topology need not be Hausdorff
\cite[Definition~2.3.1, p.~29]{AbramskyJungCorrectedExpandedDomainTheory}.
\end{hbdefinition}

\begin{hbtheorem}{thm:p10-domain-scott-topology}{proved=annex:p10.domain.topology}{Topology, order and continuity agree at the stated scope}
The Scott opens form a topology and recover the order:
\(x\sqsubseteq y\) iff every Scott-open neighborhood of \(x\)
contains \(y\). Thus this topology is \(T_0\).
A map of dcpo's is Scott-continuous exactly when it is continuous
between their Scott topologies. In a continuous domain with basis
\(B\), the sets \(\mathord{\uparrow\!\uparrow}b\), \(b\in B\),
form a basis for \(\sigma D\).
\cite[Propositions~2.3.2,2.3.4--6, pp.~29--31]
{AbramskyJungCorrectedExpandedDomainTheory}.
\end{hbtheorem}

\subsection{Least fixed points as meanings of recursion}

\begin{hbtheorem}{thm:p10-domain-function-space}{proved=annex:p10.domain.functions}{Function dcpo's, evaluation and currying}
For dcpo's \(D,E\), \([D\longrightarrow E]\) is a dcpo whose
directed suprema are pointwise. If \(E\) is pointed it has bottom
the constant \(\bot_E\) map. For pointed \(D,E\), the strict function space is also a
dcpo with pointwise directed suprema.
Finite products are coordinatewise dcpo's, and their approximation
relation is coordinatewise; products of bases are bases.
Evaluation, currying and composition are Scott-continuous, with
\[
 [D\times E\longrightarrow F]\cong[D\longrightarrow[E\longrightarrow F]].
\]
The functorial action on function spaces is
\([f,g](h)=g\circ h\circ f\), contravariant in \(f\)'s object.
This proves cartesian closure of \(\mathsf{DCPO}\), not of all
continuous or algebraic domains
\cite[Proposition~2.1.18; Section~3.2.1--2, pp.~16,42--44]
{AbramskyJungCorrectedExpandedDomainTheory}.
\end{hbtheorem}

\begin{hbtheorem}{thm:p10-domain-least-fixpoint}{proved=annex:p10.domain.fixpoint}{Least fixed point and continuity of its operator}
Let \(D\) be a pointed dcpo and \(f:D\to D\) Scott-continuous.
Then
\[
 \operatorname{fix}(f)=\bigsqcup^{\uparrow}_{n\in\mathbb N}f^n(\bot)
\]
is a fixed point and is below every pre-fixed point
\(x\) with \(f(x)\sqsubseteq x\). The actual assignment
\(\operatorname{fix}:[D\longrightarrow D]\to D\) is Scott-continuous.
For continuous \(F:P\times D\to D\), the parametrized solution
\(p\mapsto\operatorname{fix}(F(p,-))\) is therefore continuous and
satisfies its recursive equation
\cite[Theorem~2.1.19; Exercise~2.3.9(17), pp.~16,32]
{AbramskyJungCorrectedExpandedDomainTheory}.
\end{hbtheorem}

\begin{hbtheorem}{thm:p10-domain-fixpoint-principles}{proved=annex:p10.domain.induction}{Induction and strict change of representation}
If \(P\subseteq D\) contains \(\bot\), is closed under suprema
of increasing \(\omega\)-chains, and \(f(P)\subseteq P\), then
\(\operatorname{fix}(f)\in P\). If \(h:D\to E\) is strict
continuous, \(f,g\) are continuous endomaps of pointed dcpo's, and
\(hf=gh\), then \(h(\operatorname{fix}(f))=\operatorname{fix}(g)\).
The assumptions do not reduce to monotonicity, or to an arbitrary
not-necessarily-strict representation map
\cite[Lemmas~2.1.20--21, p.~17]{AbramskyJungCorrectedExpandedDomainTheory}.
\end{hbtheorem}

\begin{hbexample}{def:p10-domain-recursive-factorial}{A partial-function recursive meaning}
On \([\mathbb N\rightharpoonup\mathbb N]\), define \(\Phi(p)\)
by \(\Phi(p)(0)=1\), and, for \(n>0\), define
\(\Phi(p)(n)=n\,p(n-1)\) exactly when \(p(n-1)\) is defined.
This is continuous: every produced graph pair uses at most one graph
pair from its input. Its \(k\)-th approximation defines exactly
\(0,\ldots,k-1\) with their factorial values, so its least fixed
point is total factorial. In contrast, the recursive equation
\(p=p\) has every partial function as a solution and the least
meaning is the empty function. These are editorial worked instances
of the source's partial-function domain and least-fixed-point theorem;
their construction and approximation calculation are supplied in
Annex~\ref{annex:p10.domain.examples} and
Annex~\ref{annex:p10.domain.fixpoint}.
\end{hbexample}

\begin{hbtheorem}{thm:p10-domain-failed-premises}{proved=annex:p10.domain.failures}{Failures when the fixed-point premises are removed}
On the discrete dcpo \(\mathbb N\), successor is continuous and
has no fixed point: bottom is missing. On the pointed chain
\(0<1<\cdots<\omega<\omega+1\), the monotone map sending
\(n\) to \(n+1\), \(\omega\) to \(\omega+1\), and
\(\omega+1\) to itself has an \(\omega\)-iteration supremum
\(\omega\) which is not fixed: continuity is missing.
On \(\omega+1\), successor extended by fixing \(\omega\) leaves
the finite elements invariant, but its least fixed point is not finite:
that predicate is not admissible. Finally, a non-strict constant map
from the one-point dcpo to the top of \(\{0<1\}\) intertwines
identity maps but does not preserve their least fixed points.
The monotone discontinuous example answers the actual
Exercise~2.3.9(13), rather than denying the separate existence theorem
for monotone endomaps on pointed dcpo's
\cite[pp.~15--17,32]{AbramskyJungCorrectedExpandedDomainTheory}.
\end{hbtheorem}

\subsection{Domain constructors and expanding sequences}

\begin{hbdefinition}{def:p10-domain-type-constructors}{Lifting, coalesced sum and smash product}
Let \(I\) denote the one-point dcpo. The lifting \(D_{\bot}\)
adds a \emph{new} bottom to \(D\);
write \(\operatorname{up}(d)\) for its old elements. If \(D\)
already has bottom, \(\operatorname{up}(\bot_D)\ne\bot_{D_{\bot}}\).
For pointed \(D,E\), the coalesced sum \(D\oplus E\) identifies
their bottoms and otherwise retains tagged elements \((d:1)\),
\((e:2)\), with no cross-branch order. The smash product
\(D\otimes E\) identifies all product pairs having a bottom
coordinate; other pairs keep their coordinatewise order.
\cite[Sections~3.2.3--5, pp.~44--45]
{AbramskyJungCorrectedExpandedDomainTheory}.
\end{hbdefinition}

\begin{hbtheorem}{thm:p10-domain-constructors}{proved=annex:p10.domain.constructors}{Exact constructor objects and universal properties}
These constructions are pointed dcpo's. Lifting preserves continuity
and algebraicity; its basis is new bottom together with the lifted old
basis. Coalesced sum and smash product preserve those properties for
pointed factors, with their induced tagged and paired bases.
Coalesced sum is a coproduct for \emph{strict} continuous maps, not
for all continuous maps between pointed objects. Lifting has the
natural bijection
\[
 [D\longrightarrow E]\cong
 [D_{\bot}\mathrel{\mathop{\longrightarrow}^{\bot!}}E]
                    \qquad(E\text{ pointed}),
\]
and, for pointed \(F\), smash product has the strict currying bijection
\[
 [D\otimes E\mathrel{\mathop{\longrightarrow}^{\bot!}}F]
 \cong
 [D\mathrel{\mathop{\longrightarrow}^{\bot!}}
      [E\mathrel{\mathop{\longrightarrow}^{\bot!}}F]].
\]
Thus the strict category is monoidal closed with smash product,
unit \(I_{\bot}\), and strict function space; these are not its
cartesian exponentials.
\cite[Propositions~3.2.8--9; Sections~3.2.4--5, pp.~44--45]
{AbramskyJungCorrectedExpandedDomainTheory}.
\end{hbtheorem}

\begin{hbtheorem}{thm:p10-domain-algebraic-functions}{proved=annex:p10.domain.steps}{A genuine algebraic function domain}
Let \(D\) be algebraic and \(E\) pointed, algebraic and
bounded-complete: every bounded pair in \(E\) has a join.
For compact \(a\in D,b\in E\), define
\[
 (a\mathbin{\searrow}b)(x)=
 \begin{cases}b,&a\sqsubseteq x,\\\bot_E,&a\not\sqsubseteq x.\end{cases}
\]
The compact elements of \([D\longrightarrow E]\) are exactly
finite pointwise joins of these step functions for which the join
exists; the empty join is the bottom function. This is a pointed
bounded-complete algebraic domain. Countable compact bases in \(D,E\)
give a countable compact basis of the function domain.
This is the algebraic bounded-complete case of the source's
lattice-like function-domain construction, not an assertion that
every algebraic-domain function space is algebraic
\cite[pp.~54--57, Proposition~4.1.5]
{AbramskyJungCorrectedExpandedDomainTheory}.
\end{hbtheorem}

\begin{hbdefinition}{def:p10-domain-expanding-sequence}{Embedding--projection pairs and expanding sequences}
An embedding--projection pair \(e:D\to E,p:E\to D\) consists
of continuous maps with \(pe=1_D\) and \(ep\sqsubseteq1_E\).
These equations imply that both maps are strict when their objects
are pointed. An expanding sequence has such pairs
\((e_{mn}:D_n\to D_m,p_{nm}:D_m\to D_n)\), \(n\le m\),
with identity maps at equal indices and
\(e_{kn}=e_{km}e_{mn}\) for \(n\le m\le k\).
Consequently \(p_{nk}=p_{nm}p_{mk}\).
\cite[Definition~3.1.15; Definition~3.3.6, pp.~39,48]
{AbramskyJungCorrectedExpandedDomainTheory}.
\end{hbdefinition}

\begin{hbtheorem}{thm:p10-domain-bilimit}{proved=annex:p10.domain.bilimit}{The actual chain bilimit and its two universal properties}
For an expanding sequence of pointed dcpo's set
\[
 D=\{(x_n)_n\in\prod_nD_n:p_{nm}(x_m)=x_n\ (n\le m)\}.
\]
Its order and directed suprema are coordinatewise. Let \(p_n\)
be coordinate projection and define \(e_n:D_n\to D\) by
\[
 (e_nx)_i=\begin{cases}p_{in}(x),&i\le n,\\
                         e_{in}(x),&n\le i.\end{cases}
\]
These definitions agree at \(i=n\). They form embedding--projection
pairs and satisfy \(\bigsqcup_n^{\uparrow}e_np_n=1_D\).
The \(p_n\)'s give the limit of the projection diagram; the
\(e_n\)'s give the colimit of the embedding diagram, even for
cocones consisting of arbitrary continuous maps. Its mediators are
\[
 z\longmapsto(g_n(z))_n,
       \qquad (x_n)_n\longmapsto\bigsqcup_n^{\uparrow}h_n(x_n),
\]
for a projection cone \((g_n)\) and embedding cocone \((h_n)\).
A cocone of embedding--projection pairs \((u_n,v_n)\) is
colimiting exactly when \(\bigsqcup_n^{\uparrow}u_nv_n=1\).
If each \(D_n\) is continuous with basis \(B_n\), then
\(\bigcup_ne_n(B_n)\) is a basis of \(D\); if each is algebraic,
\(K(D)=\bigcup_ne_n(K(D_n))\). Countability is preserved for a
countable sequence of countable bases.
\cite[Theorem~3.3.7; Lemma~3.3.8; Theorem~3.3.11;
Exercise~3.3.12(15), pp.~48--53]
{AbramskyJungCorrectedExpandedDomainTheory}.
\end{hbtheorem}

\subsection{Recursive type equations with actual fold and unfold}

\begin{hbdefinition}{def:p10-domain-continuous-functor}{Local continuity and continuity on embeddings}
A functor on categories of dcpo's is locally continuous when its
action on each function dcpo is Scott-continuous, hence monotone
and preserving directed suprema. For a
mixed-variance functor \(F:\mathcal D^{\mathrm{op}}\times
\mathcal D'\to\mathcal E\), this means
\[
 F(\bigsqcup^{\uparrow}A,\bigsqcup^{\uparrow}A')
   =\bigsqcup^{\uparrow}_{f\in A,\,g\in A'}F(f,g)
\]
for correctly typed directed families of maps. A functor on
embedding categories is continuous when it preserves the actual
colimit cocones of expanding sequences. This is a preservation
condition to prove, not a synonym for the desired type solution
\cite[Definitions~5.2.1,5.2.3,5.2.5, pp.~70--72]
{AbramskyJungCorrectedExpandedDomainTheory}.
\end{hbdefinition}

\begin{hbtheorem}{thm:p10-domain-mixed-embedding}{proved=annex:p10.domain.mixed}{The typed mixed-variance passage to embeddings}
Let \(F:\mathsf{DCPO}_{\bot}^{\mathrm{op}}\times
\mathsf{DCPO}_{\bot}\to\mathsf{DCPO}_{\bot}\) be locally
continuous. On embedding pairs define
\[
 \widehat F(e,e')=F(p,e'),\qquad
                   \widehat F(e,e')^*=F(e,p').
\]
These are an embedding--projection pair, define a covariant functor
on embedding categories, and preserve synchronized chain bilimits.
The diagonal functor \(G(D)=F(D,D)\),
\(G(e)=F(p,e)\), therefore preserves expanding-sequence colimits.
The function-space functor \(F(D,E)=[D\longrightarrow E]\)
satisfies this local-continuity hypothesis by the actual composition
and pointwise-supremum laws. This passage does not assert an
ordinary covariant function-space functor on all maps
\cite[Proposition~5.2.6, pp.~72--73]
{AbramskyJungCorrectedExpandedDomainTheory}.
\end{hbtheorem}

\begin{hbtheorem}{thm:p10-domain-recursive-solution}{proved=annex:p10.domain.fold}{An equation solution constructed from an expanding sequence}
Let \(G\) be a continuous functor on pointed-dcpo embedding
categories, and choose an actual embedding \(a:A\to G(A)\).
Form \(D_n=G^n(A)\), with adjacent embeddings \(G^n(a)\),
and their bilimit \(D\). With bilimit embeddings \(e_n\) and
projections \(p_n\), the maps
\[
 \operatorname{fold}=\bigsqcup_n^{\uparrow}e_{n+1}G(e_n)^*:G(D)\to D,
 \qquad
 \operatorname{unfold}=\bigsqcup_n^{\uparrow}G(e_n)p_{n+1}:D\to G(D)
\]
are actual inverse continuous maps. In particular
\(\operatorname{unfold}e_{n+1}=G(e_n)\) and
\(G(e_n)^*=p_{n+1}\operatorname{fold}\).
The one-point pointed dcpo \(I\) has an embedding into every
pointed dcpo by selecting bottom. Starting at \(A=I\) gives
the source's \(\operatorname{FIX}(G)\). A different seed is
specified rather than silently called that canonical solution.
An isomorphism alone does not prove an operational adequacy or
initial/final universal property for an unspecified language.
\cite[Lemma~5.2.2, pp.~70--71]{AbramskyJungCorrectedExpandedDomainTheory}.
\end{hbtheorem}

\begin{hbtheorem}{thm:p10-domain-recursive-examples}{proved=annex:p10.domain.recursive-examples}{Two explicit recursive semantic structures}
For a nonempty discrete dcpo \(A\), finite and infinite words
\(A^{\le\omega}\), ordered by prefix with empty word as bottom,
form an algebraic dcpo with the finite words as its compact basis.
There is an actual isomorphism
\[
 A^{\le\omega}\cong(A\times A^{\le\omega})_{\bot}:
 \quad\epsilon\mapsto\bot,\qquad aw\mapsto\operatorname{up}(a,w).
\]
Truncation at depth \(n\) and finite-prefix inclusion are
embedding--projection pairs whose supremum is identity; they
identify this object with the actual unfolding bilimit. Finite
words here are partial prefixes, not a separately tagged terminated
list value.

For the flat natural-number domain \(A=\mathbb N_{\bot}\), the
maps \(a(x)=\text{constant }x\) and \(p(f)=f(\bot_A)\)
form an embedding--projection pair
\(A\rightleftarrows[A\longrightarrow A]\).
Iterating the mixed function-space construction yields an
\(\omega\)-algebraic bilimit \(D\) with explicit inverse maps
\(D\rightleftarrows[D\longrightarrow D]\).
This solution is nontrivial because it contains the seed \(A\)
as a subdomain. Starting with the one-point seed instead gives the
one-point solution of that same equation. Application in the
nontrivial reflexive structure is
\(\operatorname{unfold}(d)(x)\), and abstraction is
\(\operatorname{fold}\) of the continuous curried denotation;
their beta/eta equations follow from the stated inverse maps.
These are editorial worked constructions using the source's
constructors, bilimit theorem and genuine-equation example,
not a claim that a solution of an equation is unique without
its specified seed or universal property
\cite[pp.~20,43,48--53,68--72]{AbramskyJungCorrectedExpandedDomainTheory}.
\end{hbtheorem}

% Continuation of the frozen domain-theory-and-recursive-models packet.
\section{Canonical recursion, invariance and relational reasoning}
\label{sec:p10-domain-recursion-invariance}

An inverse pair solving a domain equation does not by itself establish
minimality, initiality or finality. These conclusions require the
canonical construction and its least-endomorphism property. This
continuation develops that dependency from the corrected and expanded
Abramsky--Jung author presentation, including its separate strictness
regimes \cite[Sections~5.2.3--5.4, pp.~73--84]
{AbramskyJungCorrectedExpandedDomainTheory}.
The underlying directed, function-space and bilimit constructions are
Theorems~\ref{thm:p10-domain-function-space},
\ref{thm:p10-domain-least-fixpoint},\ref{thm:p10-domain-bilimit}
and~\ref{thm:p10-domain-recursive-solution}.

\begin{hbconvention}{def:p10-domain-invariance-category}{Objects, hom dcpo's and the canonical seed}
The default category \(\mathcal D\) is
\(\mathsf{DCPO}_{\bot}\), with set-sized pointed dcpo's and all
Scott-continuous maps under the previously declared universe convention.
One may instead use a full subcategory containing the one-point
\(I\) and closed under the actual expanding-sequence bilimits;
its hom objects still contain all continuous maps between its objects.
Write \(\mathcal D_{\bot!}\) for the same objects with strict maps.
Its hom dcpo's are the strict function spaces, not an assumption that
an unrestricted functor preserves strict maps.
All functors below are locally continuous on the displayed hom dcpo's.
For a mixed functor this includes correctly typed joint continuity in
the contravariant and covariant hom variables. Products of finitely many
such categories use coordinatewise hom orders and bilimits.
Every \(\operatorname{FIX}\) in this section starts at \(I\),
not at an arbitrary seed from Theorem~\ref{thm:p10-domain-recursive-solution}.
All colimits used are countable expanding-sequence bilimits; no
unasserted arbitrary colimit, universe enlargement or effective
presentation is required \cite[pp.~70--74]
{AbramskyJungCorrectedExpandedDomainTheory}.
\end{hbconvention}

\subsection{Least comparisons and covariant universal properties}

\begin{hbdefinition}{def:p10-domain-invariance-algebra}{Algebras, coalgebras and solution comparisons}
For a covariant \(F\), an \(F\)-algebra has structure
\(a:F(A)\to A\); an algebra homomorphism \(h:A\to B\)
satisfies \(ha=bF(h)\). A coalgebra has structure
\(c:A\to F(A)\); a coalgebra homomorphism into
\(d:B\to F(B)\) satisfies \(dh=F(h)c\).
A solution \((A,i)\) has an actual isomorphism \(i:F(A)\to A\).
An isomorphism of solutions includes preservation of these structure
maps; it is not merely an untyped bijection of carriers.
For the canonical solution \(D\), write
\(i_D=\operatorname{fold}\), \(i_D^{-1}=\operatorname{unfold}\),
and \(q_n=e_np_n\) for its bilimit approximating idempotents.
\cite[Definition~5.1.1; pp.~69--70,74--77]
{AbramskyJungCorrectedExpandedDomainTheory}.
\end{hbdefinition}

\begin{hbtheorem}{thm:p10-domain-invariance-least-hom}{proved=annex:p10.domain-invariance.least-hom}{Least homomorphisms and their composition}
Let \(F:\mathcal D\to\mathcal D\) be locally continuous,
\(i:F(A)\to A\) an isomorphism, and \(a:F(C)\to C\)
any continuous algebra structure. Define
\[
 \phi_{C,A}(u)=aF(u)i^{-1},\qquad
 h_{C,A}=\operatorname{fix}(\phi_{C,A})
                       \quad\text{in }[A\longrightarrow C].
\]
This is the least algebra homomorphism. For another solution \((B,j)\),
\(h_{C,A}=h_{C,B}h_{B,A}\). Consequently
\(\mu_A=h_{A,A}\) is an idempotent below identity, and is strict.
\cite[Lemma~5.3.1, pp.~74--75]
{AbramskyJungCorrectedExpandedDomainTheory}.
\end{hbtheorem}

\begin{hbtheorem}{thm:p10-domain-invariance-covariant}{proved=annex:p10.domain-invariance.covariant}{Canonical identity, invariance and minimality}
For the canonical \(D=\operatorname{FIX}(F)\),
\[
 \theta_D(u)=i_DF(u)i_D^{-1},\qquad
 \theta_D^n(c_{\bot})=q_n,\qquad
 \operatorname{fix}(\theta_D)=1_D.
\]
For any solution \((A,i)\), its least comparisons
\(e=h_{A,D}:D\to A\), \(p=h_{D,A}:A\to D\)
satisfy \(pe=1_D\) and \(ep=\mu_A\sqsubseteq1_A\).
Thus the canonical solution is a subdomain of every solution.
The following are equivalent: \((A,i)\) is isomorphic as a solution
to the canonical solution; \(\mu_A=1_A\); and identity is the
only strict algebra endomorphism of \((A,i)\).
\cite[Theorems~5.3.2--3, pp.~75--76]
{AbramskyJungCorrectedExpandedDomainTheory}.
\end{hbtheorem}

\begin{hbtheorem}{thm:p10-domain-invariance-initial-final}{proved=annex:p10.domain-invariance.initial-final}{Initiality and finality in their different regimes}
If \(F:\mathcal D_{\bot!}\to\mathcal D_{\bot!}\) is
locally continuous, canonical fold is an initial \(F\)-algebra
in the strict-map category. Its unique strict homomorphism to a
strict algebra \(a:F(A)\to A\) is the least fixed point of
\(u\mapsto aF(u)\operatorname{unfold}\) on the strict hom dcpo.
If instead \(F:\mathcal D\to\mathcal D\) is locally continuous
on unrestricted continuous maps, canonical unfold is a final
coalgebra for \emph{all} continuous coalgebra structures
\(c:A\to F(A)\), including non-strict ones. Its unique
homomorphism is
\[
 \operatorname{fix}\bigl(u\mapsto\operatorname{fold}F(u)c\bigr)
                              \quad\text{in }[A\longrightarrow D].
\]
The second assertion does not extend a functor known only on strict
maps to non-strict maps. In the strict-only regime it gives finality
for strict coalgebras and strict homomorphisms.
\cite[Theorems~5.3.4--5; Exercise~5.4.11(6), pp.~77,83]
{AbramskyJungCorrectedExpandedDomainTheory}.
\end{hbtheorem}

\subsection{Mixed variance and the free comparison property}

\begin{hbdefinition}{def:p10-domain-invariance-mixed-comparison}{Paired comparisons for mixed solutions}
For \(F:\mathcal D^{\mathrm{op}}\times\mathcal D\to\mathcal D\),
solutions have \(i:F(A,A)\to A\), \(j:F(B,B)\to B\).
A comparison pair \(h:A\to B,k:B\to A\) satisfies
\[
 hi=jF(k,h),\qquad kj=iF(h,k).
\]
These are two typed equations, not an ordinary covariant
\(F\)-algebra homomorphism for the diagonal object assignment.
A mixed endomorphism of \((A,i)\) satisfies \(hi=iF(h,h)\).
\cite[Lemma~5.3.6, pp.~77--78]
{AbramskyJungCorrectedExpandedDomainTheory}.
\end{hbdefinition}

\begin{hbtheorem}{thm:p10-domain-invariance-least-pairs}{proved=annex:p10.domain-invariance.least-pairs}{Least pairs, exact composition and the mixed kernel}
The least comparison pair is the least fixed point of
\[
 \Phi_{B,A}(u,v)
       =(jF(v,u)i^{-1},\ iF(u,v)j^{-1})
\]
on \([A\longrightarrow B]\times[B\longrightarrow A]\).
For a third solution, composing forward maps forward and backward
maps backward gives the least comparison pair between the outside
solutions. For \(A=B\) the two least maps agree and equal
\[
 \mu_A=\operatorname{fix}\bigl(u\mapsto iF(u,u)i^{-1}\bigr).
\]
This mixed kernel is idempotent, strict and below identity.
\cite[Lemma~5.3.6 and its corrected-author comment, p.~78]
{AbramskyJungCorrectedExpandedDomainTheory}.
\end{hbtheorem}

\begin{hbtheorem}{thm:p10-domain-invariance-mixed}{proved=annex:p10.domain-invariance.mixed}{Mixed canonical identity and minimality}
Let \(D=\operatorname{FIX}(\widehat F\circ\Delta)\), using the
actual embedding functor of Theorem~\ref{thm:p10-domain-mixed-embedding}.
Then for \(\theta_D(u)=\operatorname{fold}F(u,u)\operatorname{unfold}\),
\(\theta_D^n(c_{\bot})=q_n\) and \(\mu_D=1_D\).
Its least comparison pair with any mixed solution \((A,i)\)
is an embedding--projection pair with composite on \(A\) equal
to \(\mu_A\). Moreover canonical-solution isomorphism,
\(\mu_A=1_A\), and uniqueness of the strict mixed endomorphism
are equivalent. In the uniqueness proof, \(u\mapsto huh\)
first gives \(h^2=1\); the least-identity order comparison is
also used to conclude \(h=1\).
\cite[Theorems~5.3.7--8, pp.~78--79]
{AbramskyJungCorrectedExpandedDomainTheory}.
\end{hbtheorem}

\begin{hbtheorem}{thm:p10-domain-invariance-free-mixed}{proved=annex:p10.domain-invariance.free-mixed}{The actual free mixed-variant universal property}
Suppose \(F:\mathcal D_{\bot!}^{\mathrm{op}}\times
\mathcal D_{\bot!}\to\mathcal D_{\bot!}\) is locally
continuous and \(D\) its canonical mixed solution. For strict
\(f:A\to F(B,A)\), \(g:F(A,B)\to B\), there is a
\emph{unique strict} pair \(h:A\to D,k:D\to B\) with
\[
 \operatorname{unfold}h=F(k,h)f,\qquad
 k\operatorname{fold}=gF(h,k).
\]
It is the least fixed point on the strict product hom dcpo of
\[
 (u,v)\longmapsto
 (\operatorname{fold}F(v,u)f,\ gF(u,v)\operatorname{unfold}).
\]
Uniqueness is proved by the canonical approximants \(q_n\),
not by the mere invertibility of fold and unfold.
\cite[Theorem~5.3.9, p.~79]{AbramskyJungCorrectedExpandedDomainTheory}.
\end{hbtheorem}

\subsection{Parameters, simultaneous equations and resumptions}

\begin{hbtheorem}{thm:p10-domain-invariance-parameters}{proved=annex:p10.domain-invariance.parameters}{The parameter solution is an actual locally continuous functor}
Let \(F:\mathcal C_{\bot!}\times\mathcal D_{\bot!}\to
\mathcal D_{\bot!}\) be locally continuous. Put
\(F_A(X)=F(A,X)\), \(M(A)=\operatorname{FIX}(F_A)\), and
\(X_n(A)=F_A^n(I)\). For a strict \(u:A\to A'\), define
\[
 u_0=1_I,\quad u_{n+1}=F(u,u_n),\qquad
 M(u)=\bigsqcup_n^{\uparrow}e_n^{A'}u_np_n^A.
\]
These maps are strict, the supremum is directed, and this gives
a locally continuous functor. Its actual structure isomorphisms
\(\alpha_A:F(A,M(A))\to M(A)\) satisfy
\[
 M(u)\alpha_A=\alpha_{A'}F(u,M(u)),\qquad
 \alpha_{A'}^{-1}M(u)=F(u,M(u))\alpha_A^{-1}.
\]
\cite[Proposition~5.2.7, pp.~73--74]
{AbramskyJungCorrectedExpandedDomainTheory}.
\end{hbtheorem}

\begin{hbtheorem}{thm:p10-domain-invariance-mixed-parameters}{proved=annex:p10.domain-invariance.mixed-parameters}{Explicitly paired mixed parameters}
Let \(F:\mathcal C_{\bot!}^{\mathrm{op}}\times
\mathcal C_{\bot!}\times\mathcal D_{\bot!}^{\mathrm{op}}
\times\mathcal D_{\bot!}\to\mathcal D_{\bot!}\) be locally
continuous. For parameter objects \(P,Q\), start \(X_0=Y_0=I\)
and iterate
\[
 X_{n+1}=F(P,Q,Y_n,X_n),\quad
 Y_{n+1}=F(Q,P,X_n,Y_n).
\]
The mixed embedding action gives actual expanding sequences.
Their bilimits \(M(P,Q)\), \(M(Q,P)\) satisfy
\[
 M(P,Q)\cong F(P,Q,M(Q,P),M(P,Q)).
\]
For \(u:P'\to P,v:Q\to Q'\), define
\[
 r_0=s_0=1_I,\quad
 r_{n+1}=F(u,v,s_n,r_n),\quad
 s_{n+1}=F(v,u,r_n,s_n).
\]
The first sequence induces \(M(u,v):M(P,Q)\to M(P',Q')\);
the second induces the reversed map \(M(v,u)\).
The induced maps use the same bilimit supremum formula as the
preceding theorem and define a locally continuous mixed functor.
The structure maps are natural with the typed paired action
\(F(u,v,M(v,u),M(u,v))\). This makes the negative recursive
slot explicit instead of substituting a covariant action there.
This is the four-slot extension of the source's parameter
construction and its prescribed variance handling
\cite[pp.~72--74; Exercise~5.4.11(9), p.~84]
{AbramskyJungCorrectedExpandedDomainTheory}.
\end{hbtheorem}

\begin{hbtheorem}{thm:p10-domain-invariance-bekic}{proved=annex:p10.domain-invariance.bekic}{Beki\v c elimination, for values and for canonical domains}
For pointed dcpo's \(D,E\) and continuous
\(f:D\times E\to D,g:D\times E\to E\), set
\[
 h(y)=\operatorname{fix}(x\mapsto f(x,y)),\quad
 k(y)=g(h(y),y),\quad d=\operatorname{fix}(k),\quad c=h(d).
\]
Then \((c,d)=\operatorname{fix}\langle f,g\rangle\).
For locally continuous strict functors
\(P:\mathcal C_{\bot!}\times\mathcal E_{\bot!}\to\mathcal C_{\bot!}\)
and \(Q:\mathcal C_{\bot!}\times\mathcal E_{\bot!}\to\mathcal E_{\bot!}\),
put \(H(B)=\operatorname{FIX}(P(-,B))\),
\(K(B)=Q(H(B),B)\), \(B_* =\operatorname{FIX}(K)\),
\(A_*=H(B_*)\). Their two structure maps make \((A_*,B_*)\)
an initial algebra for \((A,B)\mapsto(P(A,B),Q(A,B))\).
Hence the nested and simultaneous canonical solutions are
isomorphic by actual structure-preserving comparison maps.
The value assertion needs no strictness; the domain initiality
assertion uses the displayed strict regime.
\cite[pp.~73--74; Exercise~5.4.11(4), p.~83]
{AbramskyJungCorrectedExpandedDomainTheory}.
\end{hbtheorem}

\begin{hbtheorem}{thm:p10-domain-invariance-resumptions}{proved=annex:p10.domain-invariance.resumptions}{The source resumption functor and observable recursive behavior}
For a pointed state dcpo \(S\), define
\[
 T(R)=[S\longrightarrow(S\oplus(S\times R))],\qquad
 T(u)(r)=(1_S\oplus(1_S\times u))\circ r
\]
for \emph{strict} continuous \(u\). This is an actual locally
continuous functor on \(\mathsf{DCPO}_{\bot!}\). Its canonical
solution has inverse fold/unfold, initial strict-algebra and
final strict-coalgebra properties. No unrestricted-map extension
is inferred from the coalesced-sum notation.
For a nonbottom state \(s\), define
\[
 \operatorname{return}_s=\operatorname{fold}(x\mapsto\operatorname{inl}(s)),
 \quad
 \operatorname{step}(s,r)=\operatorname{fold}
              (x\mapsto\operatorname{inr}(s,r)).
\]
The continuous functional \(r\mapsto\operatorname{step}(s,r)\)
has a least fixed point whose \(n\)-th approximation exposes
\(n\) successive intermediary states and then an unknown remainder.
Its limit never has a final-return observation; return and step
remain distinct coproduct-branch observations. This example
retains the source's final-state versus state-plus-resumption
interpretation without asserting an unproved programming-language
adequacy theorem \cite[Section~5.1.2, pp.~68--69]
{AbramskyJungCorrectedExpandedDomainTheory}.
\end{hbtheorem}

\subsection{Admissible relations, induction and coinduction}

\begin{hbdefinition}{def:p10-domain-invariance-relations}{Admissibility and typed relation actions}
For finite arity \(n\), \(R\subseteq D^n\) is admissible when
it contains the constantly-bottom tuple and is closed under
coordinatewise suprema of increasing \(\omega\)-chains. Let
\(\mathcal R_n(D)\) be these relations ordered by inclusion.
For strict \(u:D\to E\), write
\(u^{\mathrm{rel}}R=\{(u(x_1),\ldots,u(x_n)):(x_i)_i\in R\}\).
An admissible action for mixed \(F\) assigns
\(F^{\mathrm{rel}}(R,S)\in\mathcal R_n(F(D,E))\).
It must satisfy the typed condition: if
\(u:D_2\to D_1,v:E_1\to E_2\) are strict and
\(u^{\mathrm{rel}}R_2\subseteq R_1\),
\(v^{\mathrm{rel}}S_1\subseteq S_2\), then
\[
 F(u,v)^{\mathrm{rel}}F^{\mathrm{rel}}(R_1,S_1)
                       \subseteq F^{\mathrm{rel}}(R_2,S_2).
\]
It is therefore antitone in the negative relation and monotone
in the positive relation. These actions are not obtained merely
by treating the relations as domain objects and applying \(F\).
\cite[Definitions~5.4.1,5.4.3; Proposition~5.4.4, pp.~80--81]
{AbramskyJungCorrectedExpandedDomainTheory}.
\end{hbdefinition}

\begin{hbtheorem}{thm:p10-domain-invariance-relational}{proved=annex:p10.domain-invariance.relational}{The general admissible-relation comparison principle}
The predicate \(\{u:u^{\mathrm{rel}}R\subseteq S\}\) is
admissible on the strict function dcpo whenever \(R,S\) are
admissible. If \(D\) is the canonical mixed solution and
\[
 \operatorname{unfold}^{\mathrm{rel}}R\subseteq F^{\mathrm{rel}}(S,R),
 \qquad
 \operatorname{fold}^{\mathrm{rel}}F^{\mathrm{rel}}(R,S)\subseteq S,
\]
then \(R\subseteq S\). This follows from the least-identity
endomorphism and an actual invariant admissible predicate.
\cite[Proposition~5.4.2; Theorem~5.4.5, pp.~80--81]
{AbramskyJungCorrectedExpandedDomainTheory}.
\end{hbtheorem}

\begin{hbtheorem}{thm:p10-domain-invariance-actions}{proved=annex:p10.domain-invariance.actions}{Actual constructor actions and variance-compatible composition}
For unary predicates the logical actions on lifting, product,
sum and function space are
\[
 R_{\bot}=\{\bot\}\cup\operatorname{up}(R),\quad
 R\times S,\quad\operatorname{inl}(R)\cup\operatorname{inr}(S),
 \quad [R\longrightarrow S]=\{f:f(R)\subseteq S\}.
\]
Smash and strict function space use the corresponding bottom
quotient and strict-map restriction. For binary relations, lifting
uses pairs whose first coordinate is new bottom or whose two
coordinates are old lifts related by \(R\); products act
coordinatewise; sums use a bottom first coordinate or a pair
in the same branch related by its component relation. Function
space uses the pointwise action
\[
 [R\longrightarrow S]=\{(f,g):(\forall x)\ (f(x),g(x))\in S\},
\]
independent of \(R\). These are admissible actions with their
declared constructor variances. The unary mixed actions are
logical: \(F^{\mathrm{rel}}(R,E)=F(D,E)\); the binary ones
are extensional: \(F^{\mathrm{rel}}(R,\sqsubseteq_E)
=\sqsubseteq_{F(D,E)}\). For
\(H(X,Y)=F(G_1(Y,X),G_2(X,Y))\), their composition is
\[
 H^{\mathrm{rel}}(R,S)=
 F^{\mathrm{rel}}(G_1^{\mathrm{rel}}(S,R),G_2^{\mathrm{rel}}(R,S)).
\]
The binary lifting relation lives on \((D_{\bot})^2\), with
explicit old-element injections and a separate new-bottom clause.
\cite[pp.~81--83]{AbramskyJungCorrectedExpandedDomainTheory}.
\end{hbtheorem}

\begin{hbtheorem}{thm:p10-domain-invariance-induction-coinduction}{proved=annex:p10.domain-invariance.induction-coinduction}{Structural logical induction and extensional simulation}
For a logical unary action, an admissible \(S\subseteq D\)
with \(\operatorname{fold}(F^{\mathrm{rel}}(D,S))\subseteq S\)
equals \(D\). For an extensional binary action, an admissible
\(R\subseteq D^2\) with
\[
 \operatorname{unfold}^{\mathrm{rel}}R
                       \subseteq F^{\mathrm{rel}}(\sqsubseteq_D,R)
\]
is contained in \(\sqsubseteq_D\). Such an \(R\) is a
simulation; \(x\sqsubseteq y\) iff some simulation relates
them. These conclusions use the declared action, not a slogan
that every arbitrary relation is inductive or coinductive.
If a canonical solution is \(\omega\)-algebraic and its
compact basis is given by finite constructor expressions, induction
on those expressions plus admissibility
also establishes a predicate on the whole domain. Without
countable approximation, containing all compact points and
being only omega-chain closed is insufficient.
\cite[Sections~5.4.1--4; Theorems~5.4.7,5.4.9;
Corollary~5.4.10; Proposition~2.2.13, pp.~22,79--83]
{AbramskyJungCorrectedExpandedDomainTheory}.
\end{hbtheorem}

\begin{hbexample}{def:p10-domain-invariance-stream-simulation}{Two recursive descriptions of the same infinite stream}
In the canonical prefix-word domain of
Theorem~\ref{thm:p10-domain-recursive-examples}, fix a letter \(a\).
If \(r,q\) both satisfy
\(\operatorname{unfold}(r)=\operatorname{up}(a,r)\) and
\(\operatorname{unfold}(q)=\operatorname{up}(a,q)\), then
\(\{(\bot,\bot),(r,q)\}\) is an admissible simulation for
the lifted product constructor. Coinduction gives \(r\sqsubseteq q\);
the reversed relation gives \(q\sqsubseteq r\), hence equality.
Annex~\ref{annex:p10.domain-invariance.induction-coinduction}
checks its actual action and admissibility.
\end{hbexample}

\begin{hbtheorem}{thm:p10-domain-invariance-failures}{proved=annex:p10.domain-invariance.failures}{Universal-property and induction boundaries that have counterexamples}
For the identity functor, the two-point pointed chain with
identity structure solves the domain equation but is not initial;
identity and the constant-bottom strict map are distinct
endomorphisms. Its canonical solution is the one-point domain.
Allowing non-strict algebras/maps destroys strict initiality:
every constant map from the one-point domain to a two-point
identity algebra is a homomorphism. The same example with
\(F(X,Y)=Y\) defeats uniqueness of the free mixed comparison
when its backward map is not required strict.
For uncountable \(U\), the predicate of countable subsets of
\(\mathcal P(U)\) contains bottom and every compact element
and is omega-chain closed, but is not the whole domain.
Finally initial algebra and final coalgebra for \(X\mapsto1+X\)
in sets have carriers \(\mathbb N\) and
\(\mathbb N\cup\{\infty\}\); their algebra structures are
not isomorphic. Domain algebraic compactness is not inferred
for an arbitrary category from a recursive equation.
\cite[Exercises~5.4.11(5)--(7), p.~83; pp.~13,18--23]
{AbramskyJungCorrectedExpandedDomainTheory}.
\end{hbtheorem}

\section{Retracts and a constructed nondeterministic domain}
\label{sec:p10-domain-constructions-retractions}

The chain bilimit and canonical recursion arguments of
Theorems~\ref{thm:p10-domain-bilimit} and
\ref{thm:p10-domain-invariance-mixed} already construct recursive objects.
Here the additional tasks are to split continuous idempotents, to distinguish
retracts from subdomains, and to construct a free domain for finite
nondeterministic choice. The source is the corrected and expanded
Abramsky--Jung author presentation, Sections~3.1 and~6.1.2--6.2.1
\cite[pp.~34--41,87--95]{AbramskyJungCorrectedExpandedDomainTheory}.
Its printed chapter and the Gunter--Scott, Gunter and Mosses presentations
use their own conventions; equivalence with them requires a matched comparison.

\subsection{Retractions are mathematical representation data}

\begin{hbdefinition}{def:p10-dc-retraction}{Sections, retractions and idempotent images}
A continuous section--retraction pair consists of Scott-continuous
\(s:D\to E\), \(r:E\to D\) with \(rs=1_D\).
It need not satisfy \(sr\sqsubseteq1_E\); that additional inequality
makes it the embedding--projection pair of
Definition~\ref{def:p10-domain-expanding-sequence}.
A continuous idempotent \(f:E\to E\) satisfies \(f^2=f\).
Give \(\operatorname{im}(f)=\{x:f(x)=x\}\) its inherited order.
All objects and bases here are set-sized under the existing universe
convention; directed sets are nonempty. A section--retraction pair is
not by itself an effective encoding of a programming language
\cite[Definitions~3.1.1,3.1.15, pp.~34,39]
{AbramskyJungCorrectedExpandedDomainTheory}.
\end{hbdefinition}

\begin{hbtheorem}{thm:p10-dc-splitting}{proved=annex:p10.dc.splitting}{Actual splitting, topology and directed idempotents}
The image of a continuous idempotent on a dcpo is a dcpo with
directed suprema computed in its parent. Its inclusion \(s\) and
corestriction \(r(x)=f(x)\) are continuous and satisfy \(rs=1\),
\(sr=f\). Its Scott topology is the subspace topology inherited from
the parent. A directed family of continuous idempotents, ordered
pointwise, has a continuous idempotent supremum.
If \(E\) is continuous with basis \(B\), every continuous retract
\(D\) is continuous with basis \(r(B)\); countability of a basis is
preserved. This does not say that \(r\) preserves every way-below pair
or that algebraicity is preserved
\cite[Lemma~3.1.3; Theorem~3.1.4; Propositions~3.1.7--8, pp.~34--36]
{AbramskyJungCorrectedExpandedDomainTheory}.
\end{hbtheorem}

\begin{hbtheorem}{thm:p10-dc-algebraic-cover}{proved=annex:p10.dc.algebraic-cover}{Every continuous domain has an algebraic retract presentation}
Let \(B\) be a basis of a continuous domain \(D\). In this theorem
\(\operatorname{Idl}(B,\sqsubseteq)\) uses the restricted information
order, not the usually nonreflexive way-below relation. Put
\[
 s(x)=B_x=\{b\in B:b\ll x\},\qquad
 r(J)=\bigsqcup^{\uparrow}J.
\]
These are continuous maps and \(rs=1_D\).
The ideal domain is algebraic, with compact principal ideals
\(\downarrow_B b\). A countable \(B\) gives an omega-algebraic
cover of an omega-continuous \(D\).
For \(D=[0,1]\) and its rational basis, this is an algebraic domain
retracting onto a nonalgebraic continuous domain. Thus ordinary continuous
retractions genuinely fail to preserve algebraicity
\cite[Theorem~3.1.5; Proposition~3.1.6, p.~35]
{AbramskyJungCorrectedExpandedDomainTheory}.
\end{hbtheorem}

\begin{hbtheorem}{thm:p10-dc-function-retract}{proved=annex:p10.dc.function-retract}{Retraction transport for continuous functions}
For section--retraction pairs \((s_D,r_D)\), \((s_E,r_E)\),
from \(D,E\) into \(D',E'\), respectively, the maps
\[
 S(h)=s_Ehr_D,
 \quad R(k)=r_Eks_D,
 \quad [D\longrightarrow E]\mathrel{\mathop{\rightleftarrows}^{S}_{R}}
       [D'\longrightarrow E']
\]
are continuous and satisfy \(RS=1\).
If the original pairs are embedding--projection pairs, so is \((S,R)\).
Products transport retractions coordinatewise as well.
In particular this gives an actual retract representation of a function
domain, not merely a bijection of function names. Computability of those
names and maps still requires effective basis and representation data.
This is an explicit construction using the source's retraction and
function-space machinery \cite[pp.~34--36,43]
{AbramskyJungCorrectedExpandedDomainTheory}.
\end{hbtheorem}

\subsection{The Plotkin powerdomain from its approximation basis}

\begin{hbdefinition}{def:p10-dc-plotkin-basis}{Two-sided approximation and formal union}
Let \(D\) be a nonempty continuous domain with basis \(B\), and let
\(\mathcal P_f(B)\) mean its \emph{nonempty finite} subsets.
The source's Egli--Milner approximation relation is
\[
 M\ll_{\mathrm{EM}}N\quad\Longleftrightarrow\quad
 (\forall a\in M)(\exists b\in N)\ a\ll b
 \ \text{and}\
 (\forall b\in N)(\exists a\in M)\ a\ll b.
\]
Define \(P^{\mathrm P}(D)=
\operatorname{Idl}(\mathcal P_f(B),\ll_{\mathrm{EM}})\), using round
ideals as in Definition~\ref{def:p10-domain-abstract-basis}.
Write \(i(M)=\{L:L\ll_{\mathrm{EM}}M\}\) and
\[
 \eta_D(x)=\{M\in\mathcal P_f(B):(\forall b\in M)\ b\ll x\}.
\]
The symbol \(\cup\) has two expressly typed readings below.
For finite basis sets \(M,N\), it is ordinary set union.
For powerdomain ideals or elements of a dcpo-semilattice, it is
\emph{formal union}: a continuous commutative, associative, idempotent
binary operation. It is not assumed to be a join in the information order.
On powerdomain ideals define the latter operation by
\[
 I\cup J=\{L:(\exists M\in I)(\exists N\in J)\
                         L\ll_{\mathrm{EM}}(M\cup N)\}.
\]
It generally differs from their ordinary union as sets of finite sets
\cite[Definitions~6.2.1--2; Theorem~6.2.3, pp.~93--94]
{AbramskyJungCorrectedExpandedDomainTheory}.
\end{hbdefinition}

\begin{hbtheorem}{thm:p10-dc-plotkin-object}{proved=annex:p10.dc.plotkin-object}{Existence, approximation, and the generator map}
\((\mathcal P_f(B),\ll_{\mathrm{EM}})\) is an abstract basis.
Consequently \(P^{\mathrm P}(D)\) is a continuous dcpo with directed
unions as suprema and basis \(i(\mathcal P_f(B))\).
The map \(\eta_D:D\to P^{\mathrm P}(D)\) is Scott-continuous.
If \(D\) is algebraic, choosing \(B=K(D)\) makes the powerdomain
algebraic. Its compact elements are the principal ideals of the
Egli--Milner preorder, with equal principal ideals identified.
Countable bases yield countable powerdomain bases. If \(D\) is pointed
and \(\bot_D\in B\), then \(i(\{\bot_D\})\) is its bottom and
equals \(\eta_D(\bot_D)\)
\cite[Theorem~6.2.3; pp.~94--95]
{AbramskyJungCorrectedExpandedDomainTheory}.
\end{hbtheorem}

\begin{hbtheorem}{thm:p10-dc-formal-union}{proved=annex:p10.dc.formal-union}{A genuine continuous semilattice operation}
The displayed formal union produces round ideals and is jointly
Scott-continuous. It obeys commutativity, associativity and idempotence.
For every \(M\in\mathcal P_f(B)\),
\[
 i(M)=\mathop{\bigcup}_{b\in M}^{\mathrm{formal}}\eta_D(b).
\]
Here the finite formal union is unambiguous by the stated equations;
no nullary union or empty nondeterministic choice has been introduced.
Thus \(P^{\mathrm P}(D)\) is a dcpo-semilattice in the source's
sense \cite[pp.~93--94]{AbramskyJungCorrectedExpandedDomainTheory}.
\end{hbtheorem}

\begin{hbtheorem}{thm:p10-dc-free-powerdomain}{proved=annex:p10.dc.free-powerdomain}{The complete free universal property and its continuity}
Let \(A\) be any dcpo with a jointly Scott-continuous commutative,
associative, idempotent operation \(\cup_A\), and let \(g:D\to A\)
be Scott-continuous. There is exactly one continuous semilattice
homomorphism \(\operatorname{ext}(g):P^{\mathrm P}(D)\to A\)
with \(\operatorname{ext}(g)\eta_D=g\). It is explicitly
\[
 \operatorname{ext}(g)(I)=
 \bigsqcup^{\uparrow}_{M\in I}
                 \mathop{\bigcup}_{b\in M}^{A}g(b).
\]
The family inside the supremum is directed in \(A\)'s information
order. The assignment \(g\mapsto\operatorname{ext}(g)\) preserves
directed suprema in the pointwise function orders.
The target \(A\) need not itself be a continuous domain.
This is the Plotkin specialization of the source's free
continuous-domain algebra theorem, not a claim that every semantic
choice structure is a set of possible terminating outcomes
\cite[Theorems~6.1.6--7; Proposition~6.1.8; Theorem~6.2.3,
pp.~89--91,94]{AbramskyJungCorrectedExpandedDomainTheory}.
\end{hbtheorem}

\begin{hbtheorem}{thm:p10-dc-powerdomain-functor}{proved=annex:p10.dc.powerdomain-functor}{Functorial action, retract preservation and flattening}
For continuous \(u:D\to E\) between nonempty continuous domains put
\(P^{\mathrm P}(u)=\operatorname{ext}(\eta_Eu)\).
This is a locally continuous functor. It transports continuous
section--retraction pairs and embedding--projection pairs without
changing their inequalities. The maps
\[
 \mu_D=\operatorname{ext}(1_{P^{\mathrm P}(D)}):
       P^{\mathrm P}(P^{\mathrm P}(D))\to P^{\mathrm P}(D)
\]
and \(\eta_D\) are natural and satisfy
\[
 \mu_D\eta_{P^{\mathrm P}(D)}=1,
 \quad\mu_DP^{\mathrm P}(\eta_D)=1,
 \quad\mu_DP^{\mathrm P}(\mu_D)=\mu_D\mu_{P^{\mathrm P}(D)}.
\]
Thus the generator and flattening form a monad with an actual free
universal property. This conclusion does not assert closure of every
cartesian closed subclass of domains under the powerdomain constructor
\cite[Theorem~6.1.9; pp.~91,93--95]
{AbramskyJungCorrectedExpandedDomainTheory}.
\end{hbtheorem}

\begin{hbtheorem}{thm:p10-dc-choice-example}{proved=annex:p10.dc.choice-example}{A finite nondeterministic example and failed substitutes}
Take the flat finite domain \(D=\{\bot,a,b\}\), with incomparable
\(a,b\). The seven nonempty subsets represent seven distinct compact
powerdomain elements, ordered by \(\sqsubseteq_{\mathrm{EM}}\).
In particular
\[
 \{\bot\}\sqsubseteq_{\mathrm{EM}}\{\bot,a\}
 \sqsubseteq_{\mathrm{EM}}\{a\},\qquad
 \{a\}\not\sqsubseteq_{\mathrm{EM}}\{a,b\},
 \quad\{a,b\}\not\sqsubseteq_{\mathrm{EM}}\{a\}.
\]
Formal union of the represented outcomes \(\{a\}\), \(\{b\}\)
is \(\{a,b\}\), but that element is not an information-order
upper bound of either input. Ordinary union of their principal ideals
is not directed and is not a powerdomain ideal.
If empty finite subsets are admitted without changing the theory,
the empty set has no Egli--Milner relation to a nonempty set and supplies
an extra isolated alternative, not the existing bottom \(\{\bot\}\).
Finally an ordinary retract can create compactness: the constant map
to \(1/2\) on \([0,1]\) has a compact one-point image although
\(1/2\not\ll1/2\) in the parent.
These are explicit construction/failure cases, not operational
adequacy claims for a programming language
\cite[pp.~35--36,93--96]{AbramskyJungCorrectedExpandedDomainTheory}.
\end{hbtheorem}

\subsection{An explicit effective algebraic regime}

\begin{hbdefinition}{def:p10-dc-effective}{Numbered compact bases and effective maps}
In this subsection an object is a nonempty algebraic domain with a
fixed numbering \(K(D)=\{b_n:n\in\mathbb N\}\), allowing repeated
names, whose compact order \(b_i\sqsubseteq b_j\) is recursive in
the indices. The numbering and decision procedure are part of the object.
For a point \(x\), put \(I_x=\{n:b_n\sqsubseteq x\}\); call
\(x\) computable when \(I_x\) is recursively enumerable.
For a continuous \(f:D\to E\), with target compact names \(c_m\),
put
\[
 R_f=\{(n,m):c_m\sqsubseteq f(b_n)\};
\]
call \(f\) effective when this relation is recursively enumerable.
This is an explicit algebraic specialization of the source's numbered
abstract bases and recursively enumerable approximable mappings
\cite[Section~8.1.1, p.~145]{AbramskyJungCorrectedExpandedDomainTheory}.
It does not identify all Scott-continuous maps with computable maps.
All finite choice codes below are nonempty finite lists of compact
indices; repetition merely gives alternative names for the same finite
basis set. No algorithm deciding arbitrary program termination is assumed.
\end{hbdefinition}

\begin{hbtheorem}{thm:p10-dc-effective-choice}{proved=annex:p10.dc.effective-choice}{An actual effective category and effective choice construction}
These objects and effective continuous maps form a category. Effective
maps send computable points to computable points. The Plotkin powerdomain
has an effective compact-basis numbering by finite choice codes:
its compact order is the recursive finite Egli--Milner comparison.
The functor \(P^{\mathrm P}\) preserves effective maps, and its unit
and flattening maps are effective in these numberings. In particular the
monad of Theorem~\ref{thm:p10-dc-powerdomain-functor} restricts to this
explicit effective category.
Continuity alone is insufficient: if \(U\subseteq\mathbb N\) is
not recursively enumerable, the constant map from the one-point domain
to \(\mathcal P(\mathbb N)\) with value \(U\) is continuous but not
effective in the standard finite-subset compact numbering.
These are proved finite-basis lifts of the constructed choice model,
not a claim to have supplied every effective function-space, fixed-point
or universal-domain result mentioned by the source's overview
\cite[pp.~94,145]{AbramskyJungCorrectedExpandedDomainTheory}.
\end{hbtheorem}

\begin{hbmethod}{def:p10-dc-using-the-model}{Using retractions and finite-choice semantics}
Specify the source and target dcpo's, their bases, and the actual maps
\(s,r\); prove \(rs=1\) and, when a subdomain is claimed,
\(sr\sqsubseteq1\). For finite nondeterministic observations, specify
the two-sided approximation relation before forming round ideals.
Interpret a generator by a continuous \(g\), evaluate finite formal
choices with \(\cup_A\), and use the proved directed extension for
infinite information. Keep information refinement distinct from
semilattice accumulation. An effective model additionally needs
enumerated basis codes and computable approximation/interpretation maps;
an operational adequacy theorem additionally needs a specified language,
reduction, observation and logical relation. None of these data follows
from the free universal property alone.
\end{hbmethod}

\begin{proofboundary}{def:p10-dc-source-boundary}{What this construction does and does not establish}
The preceding proofs supply retraction representation and a free
continuous nondeterministic-domain construction. They do not establish
an operational/denotational adequacy or full-abstraction theorem,
the general effective continuous-domain/function-space/recursive-equation
theory beyond the declared algebraic choice regime, the inverse-limit theorem for directed-index
adjunction diagrams, or the full bifinite/FS-domain hierarchy.
The further convex-hull and topological representations of the Plotkin
construction, one-sided powerdomains and probabilistic powerdomains
require their own hypotheses and proofs.
The distinct Gunter--Scott, Mosses and type-semantics presentations are
not identified with this corrected author presentation. The results here
apply to the specified continuous-domain and effective algebraic regimes,
not to every domain model or programming-language interpretation.
\end{proofboundary}

\section{Operational adequacy for call-by-name PCF}
\label{sec:p10-operational-adequacy-domain-models}

A continuous interpretation is not automatically an adequate programming
model. We now specify the language, evaluation and observations before
proving their agreement with the domain interpretation. The language is
Gunter's author-presentation PCF, including its predecessor-at-zero
convention \cite[Sections~4.1--4.4, pp.~29--35]{Gunter1994TypesAuthorOA}.
The adequacy argument supplies the proof omitted at that source's p.~35;
it uses a typed logical relation and the least-fixed-point construction,
not the free-powerdomain property. Plotkin's original computability
argument is the primary proof-method reference
\cite[Theorem~3.1; Lemmas~3.2--3.3, pp.~230--232]{Plotkin1977LCFAuthorOA}.
His original constants and predecessor convention are not silently
identified with this Gunter language. Abramsky--Jung expressly separate
the mathematical domain theory from detailed applications
\cite[pp.~7,148--149]{AbramskyJungCorrectedExpandedDomainTheory}.

\subsection{The language, strategy and observations}

\begin{hbdefinition}{def:p10-oa-language}{Typed PCF with the source's recursion binder}
Types and terms are
\[
 t::=\mathsf{num}\mid\mathsf{bool}\mid t\to t,
\]
\[
 \begin{split}
 M::={}&0\mid\mathsf{true}\mid\mathsf{false}
 \mid\operatorname{succ}(M)\mid\operatorname{pred}(M)
 \mid\operatorname{zero?}(M)\\
 &\mid\mathsf{if}\ M\ \mathsf{then}\ M\ \mathsf{else}\ M
 \mid x\mid\lambda x:t.M\mid M\,M\mid\mu x:t.M.
 \end{split}
\]
A finite context \(H\) assigns types to distinct variables.
The source judgment \(H\vdash M:t\) has variable, application and
abstraction rules, the displayed ground constructors, conditional with
same-typed branches, and recursion: the body and bound variable of
\(\mu x:t.M\) both have type \(t\). This \(\mu\) is a term binder,
not a probability mean or a semantic monad multiplication.
Both binders use capture-avoiding substitution; write \([N/x]M\).
Let \(\bar n=\operatorname{succ}^n(0)\).
The closed well-typed values are numerals, the two booleans and typed
abstractions. Function values need not have terminating bodies.
\cite[pp.~17--18,29--31]{Gunter1994TypesAuthorOA}.
\end{hbdefinition}

\begin{hbdefinition}{def:p10-oa-evaluation}{Weak call-by-name evaluation and ground results}
Write \(M\to N\) for the small-step relation registered below.
The argument of a function is substituted without being evaluated.
No reduction is performed under \(\lambda\), in an argument position,
or in an unselected conditional branch. Evaluation contexts are
\[
 E::=[]\mid\operatorname{succ}(E)\mid\operatorname{pred}(E)
 \mid\operatorname{zero?}(E)\mid E\,L
 \mid\mathsf{if}\ E\ \mathsf{then}\ M\ \mathsf{else}\ N.
\]
They propagate the registered primitive contractions. In particular
\(\operatorname{pred}(0)\) contracts to \(0\), and the recursion
binder unfolds once by substituting itself into its body.
Write \(M\Downarrow V\) for the least relation generated by the
registered natural-evaluation rules, and \(\to^*\) for the finite
reflexive transitive closure. A program is a closed term of ground type.
Its observation is its resulting numeral or boolean, when one exists.
Write \(M\Uparrow\) when no such finite evaluation exists.
The symbols \(\to\) and \(\Downarrow\) here describe operational
steps and evaluation, not denotational order or logical implication
\cite[Tables~7--8; Theorem~4.2.2, pp.~31--33]{Gunter1994TypesAuthorOA}.
\end{hbdefinition}

\calculussheet{sheet:p10-operational-model}{Gunter PCF: typing and call-by-name evaluation}{Simple num/bool/arrow types; source mu binder; weak call-by-name steps and natural evaluation; pred(0)=0.}

\hbrule{rule:p10-oa-ty-var}{Var}{}{H\vdash x:t}{\(x:t\in H\).}
\hbrule{rule:p10-oa-ty-abs}{Abs}{H,x:s\vdash M:t}{H\vdash\lambda x:s.M:s\to t}{\(x\text{ fresh for }H\).}
\hbrule{rule:p10-oa-ty-app}{App}{H\vdash M:s\to t\quad H\vdash N:s}{H\vdash M\,N:t}{\none}
\hbrule{rule:p10-oa-ty-zero}{Zero}{}{H\vdash0:\mathsf{num}}{\none}
\hbrule{rule:p10-oa-ty-true}{True}{}{H\vdash\mathsf{true}:\mathsf{bool}}{\none}
\hbrule{rule:p10-oa-ty-false}{False}{}{H\vdash\mathsf{false}:\mathsf{bool}}{\none}
\hbrule{rule:p10-oa-ty-succ}{Succ}{H\vdash M:\mathsf{num}}{H\vdash\operatorname{succ}(M):\mathsf{num}}{\none}
\hbrule{rule:p10-oa-ty-pred}{Pred}{H\vdash M:\mathsf{num}}{H\vdash\operatorname{pred}(M):\mathsf{num}}{\none}
\hbrule{rule:p10-oa-ty-test}{IsZero}{H\vdash M:\mathsf{num}}{H\vdash\operatorname{zero?}(M):\mathsf{bool}}{\none}
\hbrule{rule:p10-oa-ty-if}{Cond}{H\vdash L:\mathsf{bool}\quad H\vdash M:t\quad H\vdash N:t}{H\vdash\mathsf{if}\ L\ \mathsf{then}\ M\ \mathsf{else}\ N:t}{\none}
\hbrule{rule:p10-oa-ty-rec}{Rec}{H,x:t\vdash M:t}{H\vdash\mu x:t.M:t}{\(x\text{ fresh for }H\).}

\hbrule{rule:p10-oa-step-beta}{Beta}{}{(\lambda x:s.M)N\to[N/x]M}{Capture-avoiding substitution; the argument is not evaluated.}
\hbrule{rule:p10-oa-step-pred-zero}{PredZero}{}{\operatorname{pred}(0)\to0}{\none}
\hbrule{rule:p10-oa-step-pred-succ}{PredSucc}{}{\operatorname{pred}(\operatorname{succ}(\bar n))\to\bar n}{\(n\geq0\).}
\hbrule{rule:p10-oa-step-test-zero}{TestZero}{}{\operatorname{zero?}(0)\to\mathsf{true}}{\none}
\hbrule{rule:p10-oa-step-test-succ}{TestSucc}{}{\operatorname{zero?}(\operatorname{succ}(\bar n))\to\mathsf{false}}{\(n\geq0\).}
\hbrule{rule:p10-oa-step-if-true}{IfTrue}{}{\mathsf{if}\ \mathsf{true}\ \mathsf{then}\ M\ \mathsf{else}\ N\to M}{\none}
\hbrule{rule:p10-oa-step-if-false}{IfFalse}{}{\mathsf{if}\ \mathsf{false}\ \mathsf{then}\ M\ \mathsf{else}\ N\to N}{\none}
\hbrule{rule:p10-oa-step-rec}{Unfold}{}{\mu x:t.M\to[\mu x:t.M/x]M}{Capture-avoiding substitution.}
\hbrule{rule:p10-oa-step-context}{EvalContext}{M\to N}{E[M]\to E[N]}{\(E\) is a typed evaluation context of Definition~\ref{def:p10-oa-evaluation}.}

\hbrule{rule:p10-oa-eval-zero}{ZeroValue}{}{0\Downarrow0}{\none}
\hbrule{rule:p10-oa-eval-true}{TrueValue}{}{\mathsf{true}\Downarrow\mathsf{true}}{\none}
\hbrule{rule:p10-oa-eval-false}{FalseValue}{}{\mathsf{false}\Downarrow\mathsf{false}}{\none}
\hbrule{rule:p10-oa-eval-abs}{AbsValue}{}{\lambda x:s.M\Downarrow\lambda x:s.M}{The abstraction is closed and well typed.}
\hbrule{rule:p10-oa-eval-succ}{SuccEval}{M\Downarrow\bar n}{\operatorname{succ}(M)\Downarrow\operatorname{succ}(\bar n)}{\(n\geq0\).}
\hbrule{rule:p10-oa-eval-pred-zero}{PredZeroEval}{M\Downarrow0}{\operatorname{pred}(M)\Downarrow0}{\none}
\hbrule{rule:p10-oa-eval-pred-succ}{PredSuccEval}{M\Downarrow\operatorname{succ}(\bar n)}{\operatorname{pred}(M)\Downarrow\bar n}{\(n\geq0\).}
\hbrule{rule:p10-oa-eval-test-zero}{TestZeroEval}{M\Downarrow0}{\operatorname{zero?}(M)\Downarrow\mathsf{true}}{\none}
\hbrule{rule:p10-oa-eval-test-succ}{TestSuccEval}{M\Downarrow\operatorname{succ}(\bar n)}{\operatorname{zero?}(M)\Downarrow\mathsf{false}}{\(n\geq0\).}
\hbrule{rule:p10-oa-eval-app}{NameApp}{M\Downarrow\lambda x:s.L\quad[N/x]L\Downarrow V}{M\,N\Downarrow V}{No evaluation premise for \(N\).}
\hbrule{rule:p10-oa-eval-if-true}{TrueBranch}{L\Downarrow\mathsf{true}\quad M\Downarrow V}{\mathsf{if}\ L\ \mathsf{then}\ M\ \mathsf{else}\ N\Downarrow V}{No premise for the unselected branch.}
\hbrule{rule:p10-oa-eval-if-false}{FalseBranch}{L\Downarrow\mathsf{false}\quad N\Downarrow V}{\mathsf{if}\ L\ \mathsf{then}\ M\ \mathsf{else}\ N\Downarrow V}{No premise for the unselected branch.}
\hbrule{rule:p10-oa-eval-rec}{RecEval}{[\mu x:t.M/x]M\Downarrow V}{\mu x:t.M\Downarrow V}{The recursive term is substituted, not evaluated as a premise.}


\begin{hbtheorem}{thm:p10-oa-operational}{proved=annex:p10.oa.operational}{Substitution, preservation, progress and the two evaluations}
Typed capture-avoiding substitution preserves typing. Small steps preserve
typing; a closed well-typed nonvalue has exactly one next step, and a
value has none. For closed well-typed terms,
\[
 M\Downarrow V\quad\Longleftrightarrow\quad M\to^*V.
\]
The result is unique up to consistent bound-variable renaming.
At ground type, absence of a finite result is equivalent to an infinite
sequence of the deterministic evaluation steps; it is not a stuck
type error. Moreover \(M\to N\) implies
\(M\Downarrow V\) iff \(N\Downarrow V\) for every value \(V\).
\end{hbtheorem}

\subsection{The actual domain interpretation}

\begin{hbdefinition}{def:p10-oa-interpretation}{Flat observations and full continuous function spaces}
Put \(D_t=C\llbracket t\rrbracket\), using Gunter's interpretation
\(C\). Specifically
\[
 D_{\mathsf{num}}=\mathbb N_{\bot},\quad
 D_{\mathsf{bool}}=\{\bot,\mathsf{true},\mathsf{false}\},\quad
 D_{s\to t}=[D_s\longrightarrow D_t].
\]
Both ground orders are flat, and function spaces have pointwise order.
An \(H\)-environment \(\rho\) supplies \(\rho(x)\in D_{H(x)}\).
Abbreviate \(C\llbracket H\triangleright M:t\rrbracket\rho\) by
\(\llbracket M\rrbracket_\rho\); this drops redundant typing data,
not a source version. Variables select their environment components,
application is evaluation, and abstraction is
\(d\mapsto\llbracket M\rrbracket_{\rho[x\mapsto d]}\).
Ground constants have their corresponding nonbottom values.
Successor, predecessor and zero-test are strict on bottom and have
their ordinary displayed numeral/boolean results, with
\(\operatorname{pred}(0)=0\).
The conditional is strict in its guard, selects its true/false branch,
and returns \(\bot_t\) for bottom guard; it is not strict in both
branches. Finally
\[
 \llbracket\mu x:t.M\rrbracket_\rho
 =\operatorname{fix}\bigl(d\mapsto
                    \llbracket M\rrbracket_{\rho[x\mapsto d]}\bigr).
\]
The function space is not lifted again: this is the call-by-name
interpretation, not a value/computation model for call-by-value
\cite[pp.~34--35]{Gunter1994TypesAuthorOA}.
\end{hbdefinition}

\begin{hbtheorem}{thm:p10-oa-semantic-soundness}{proved=annex:p10.oa.semantic-soundness}{Well-defined continuous meanings and evaluation soundness}
Every \(D_t\) is a pointed bounded-complete algebraic domain.
For each typed term its interpretation is jointly Scott-continuous in
its finite environment. Interpretation respects bound renaming and
typed simultaneous substitution:
\[
 \llbracket M[\sigma]\rrbracket_\rho
   =\llbracket M\rrbracket_{x\mapsto\llbracket\sigma(x)\rrbracket_\rho}.
\]
Every small step preserves denotation. Consequently
\(M\Downarrow V\) implies \(\llbracket M\rrbracket=\llbracket V\rrbracket\).
This proves one direction only; a separate logical-relation argument
is needed for the converse.
\end{hbtheorem}

\subsection{The logical relation and fixed-point induction}

\begin{hbdefinition}{def:p10-oa-logical-relation}{Semantic information related to an operational term}
Relate \(d\in D_t\) to a closed term \(M:t\) by \(d\mathrel{\mathcal R_t}M\).
At ground type, bottom is related to every term; a nonbottom ground
value is related to \(M\) precisely when \(M\) evaluates to its
corresponding numeral or boolean. At arrow type put
\[
 f\mathrel{\mathcal R_{s\to t}}M\quad\Longleftrightarrow\quad
 (\forall N:s\text{ closed})(\forall a\in D_s)\
 a\mathrel{\mathcal R_s}N\ \Longrightarrow\
 f(a)\mathrel{\mathcal R_t}(M\,N).
\]
This relation is between a semantic element and syntax, not equality
of two denotations. Its arrow clause intentionally does not require
\(M\) itself to evaluate to an abstraction.
\end{hbdefinition}

\begin{hbtheorem}{thm:p10-oa-relation-closure}{proved=annex:p10.oa.relation-closure}{Bottom, directed admissibility and operational conversion}
For every closed \(M:t\), the set
\(\{d:d\mathrel{\mathcal R_t}M\}\) contains \(\bot_t\) and is
closed under directed suprema. If \(M\to^*N\), then for every
\(d\in D_t\),
\(d\mathrel{\mathcal R_t}M\) iff \(d\mathrel{\mathcal R_t}N\).
These statements hold at all simple types, including higher arrows.
\end{hbtheorem}

\begin{hbtheorem}{thm:p10-oa-fundamental}{proved=annex:p10.oa.fundamental}{Fundamental lemma, including arbitrary typed recursion}
Let \(H\vdash M:t\). Suppose \(\sigma\) replaces each free
\(x:s\) in \(H\) by a closed term \(\sigma(x):s\), and
\(\rho(x)\mathrel{\mathcal R_s}\sigma(x)\) for every such variable.
Then
\[
                 \llbracket M\rrbracket_\rho
                       \mathrel{\mathcal R_t}M[\sigma].
\]
For a recursion binder, every semantic iterate from bottom is related
to the same full recursive term, by induction on the iterate; directed
admissibility then relates their supremum. Neither that induction nor
its substitution into the body is replaced by a mere fixed-point name.
\end{hbtheorem}

\begin{hbtheorem}{thm:p10-oa-adequacy}{proved=annex:p10.oa.adequacy}{Both directions of ground-observation adequacy}
For every closed program \(M:g\), with
\(g\in\{\mathsf{num},\mathsf{bool}\}\), and every value \(V:g\),
\[
 M\Downarrow V\quad\Longleftrightarrow\quad
                  \llbracket M\rrbracket=\llbracket V\rrbracket.
\]
Consequently \(M\Uparrow\) iff \(\llbracket M\rrbracket=\bot_g\).
This conclusion uses the flat ground carriers and their distinct
nonbottom observations. It is not asserted as termination iff nonbottom
at every higher type.
\end{hbtheorem}

\begin{hbdefinition}{def:p10-oa-contextual}{Contextual approximation and the separate full-abstraction question}
For \(H\vdash M,N:t\), write \(M\sqsubseteq_o N\) when every
well-typed program context \(C[-]\) closing both terms at ground type
satisfies: \(C[M]\Downarrow V\) implies \(C[N]\Downarrow V\).
As in the source, context binders may close free variables of their
inserted terms; that capture is specified by the context, not confused
with capture-avoiding substitution. Contexts retain their binder names.
Operational equivalence \(M\simeq_o N\) is approximation both ways.
Denotational order means pointwise order in every \(H\)-environment.
The source calls equality-implies-operational-equivalence adequacy;
we derive that conclusion from the ground theorem. Full abstraction
would additionally require operational equivalence to imply denotational
equality \cite[Section~4.3, p.~34]{Gunter1994TypesAuthorOA}.
\end{hbdefinition}

\begin{hbtheorem}{thm:p10-oa-contextual-soundness}{proved=annex:p10.oa.contextual-soundness}{The contextual consequence actually supported by adequacy}
Pointwise denotational order implies contextual approximation.
In particular denotational equality implies operational equivalence
for open as well as closed same-typed terms. Thus the constructed
interpretation is adequate in the source's contextual sense.
No converse is inferred from compositionality or from this implication.
\end{hbtheorem}

\begin{hbtheorem}{thm:p10-oa-strategy-boundaries}{proved=annex:p10.oa.strategy-boundaries}{Call-by-value and higher-type termination are different claims}
Let \(\Omega_t=\mu x:t.x\).
In this call-by-name language \((\lambda x:\mathsf{num}.0)\Omega_{\mathsf{num}}
\Downarrow0\), although \(\Omega_{\mathsf{num}}\Uparrow\).
Changing only application to call-by-value, by requiring its argument
to evaluate to a value before substitution, makes that application
diverge; the unchanged denotation remains zero and is therefore not
adequate for that changed strategy.
Also \(\lambda x:\mathsf{num}.\Omega_{\mathsf{num}}\) is a value
whose denotation is the bottom function, whereas
\(\Omega_{\mathsf{num}\to\mathsf{num}}\) diverges with that same
denotation. Higher-type termination is not a ground observation and
does not satisfy the displayed ground bottom criterion.
\end{hbtheorem}

\begin{hbtheorem}{thm:p10-oa-countdown}{proved=annex:p10.oa.countdown}{A recursive PCF computation, with syntax and approximants}
As a worked abbreviation, not a new primitive, put
\[
 K=\mu f:\mathsf{num}\to\mathsf{num}.
    \lambda x:\mathsf{num}.
    \mathsf{if}\ \operatorname{zero?}(x)\ \mathsf{then}\ 0\
    \mathsf{else}\ \operatorname{succ}(f(\operatorname{pred}(x))).
\]
For every closed \(P:\mathsf{num}\) with \(P\Downarrow\bar n\),
\(K P\Downarrow\bar n\); in particular \(K\bar3\Downarrow\bar3\).
Its semantic functional has iterates from bottom whose \(i\)-th
approximation gives \(n\) on precisely \(n<i\), and bottom on the
remaining inputs. The least fixed point is the strict identity on
\(\mathbb N_{\bot}\). Therefore \(K\Omega_{\mathsf{num}}\Uparrow\),
unlike the constant-zero application of the preceding theorem.
The operational proof permits unevaluated argument expressions \(P\),
not just already evaluated numeral arguments.
\end{hbtheorem}

\subsection{An adequate model that is not fully abstract}

\begin{hbdefinition}{def:p10-oa-sequential-relations}{Semantic relations for the sequentiality dependency}
At ground type let
\(Q_g=\{(a,b,c):a=\bot\text{ or }b=\bot\text{ or }a=b=c\}\).
At arrows a function triple is in \(Q_{s\to t}\) when it sends every
\(Q_s\)-related argument triple to a \(Q_t\)-related result triple.
Put \(A_t=\{d:(d,d,d)\in Q_t\}\). These selected elements are not
claimed to form a new full domain category. On \(A_g\), let \(B_g\)
be equality; put \((f,h)\in B_{s\to t}\) when both endpoints are
in \(A_{s\to t}\) and \((f(a),h(b))\in B_t\) whenever
\((a,b)\in B_s\). These editorial semantic logical relations prove
sequentiality; neither replaces the adequacy relation \(\mathcal R_t\).
\end{hbdefinition}

\begin{hbtheorem}{thm:p10-oa-sequential-relations}{proved=annex:p10.oa.sequential-relations}{Two complete semantic fundamental lemmas}
\(Q_t\) and \(B_t\) contain their constantly-bottom tuples and
are closed under directed suprema in the product carriers. Any triple
whose first or second coordinate is bottom belongs to \(Q_t\).
PCF interpretations preserve both relations on related environments.
Every closed denotation therefore belongs to \(A_t\), and inserting
\(B_t\)-related closed denotations into any typed closing ground context
produces equal ground denotations. This is a theorem about the source
syntax, not a premise that every continuous function is sequential.
\end{hbtheorem}

\begin{hbtheorem}{thm:p10-oa-not-fully-abstract}{proved=annex:p10.oa.not-fully-abstract}{The source's T,F counterexample with its sequentiality proved}
Put \(\Omega=\Omega_{\mathsf{bool}}\). Define
\[
 \begin{split}
 T={}&\lambda f:\mathsf{bool}\to\mathsf{bool}\to\mathsf{bool}.
 \mathsf{if}\ f\,\mathsf{true}\,\Omega\ \mathsf{then}\\
 &\quad\mathsf{if}\ f\,\Omega\,\mathsf{true}\ \mathsf{then}\\
 &\qquad\mathsf{if}\ f\,\mathsf{false}\,\mathsf{false}
    \ \mathsf{then}\ \Omega\ \mathsf{else}\ \mathsf{true}\\
 &\quad\mathsf{else}\ \Omega\ \mathsf{else}\ \Omega.
 \end{split}
\]
Let \(F\) be the identical term with the innermost else-result false
instead of true. Both have type
\((\mathsf{bool}\to\mathsf{bool}\to\mathsf{bool})\to\mathsf{bool}\).
Then \(T\simeq_o F\), although their denotations are different.
The separating argument is the curried continuous parallel disjunction
with the source truth table
\[
 \begin{array}{c|ccc}
 \operatorname{por}&\mathsf{true}&\mathsf{false}&\bot\\\hline
 \mathsf{true}&\mathsf{true}&\mathsf{true}&\mathsf{true}\\
 \mathsf{false}&\mathsf{true}&\mathsf{false}&\bot\\
 \bot&\mathsf{true}&\bot&\bot
 \end{array}
\]
It gives \(\llbracket T\rrbracket(\operatorname{por})=\mathsf{true}\)
and \(\llbracket F\rrbracket(\operatorname{por})=\mathsf{false}\).
Thus the continuous model is not fully abstract for sequential PCF.
The contextual equivalence is proved directly here without importing
stable-model adequacy \cite[Tables~9--10, pp.~36--38]{Gunter1994TypesAuthorOA}
\cite[pp.~234--236]{Plotkin1977LCFAuthorOA}.
\end{hbtheorem}

\begin{hbmethod}{def:p10-oa-use}{Applying the adequacy theorem}
Fix the source language and evaluation strategy; check typing and identify
the ground numeral/boolean observation. Interpret recursion by the least
continuous fixed point, prove substitution and evaluation soundness,
and verify bottom/closure of the typed semantic--syntactic relation.
Use its fundamental lemma for the converse observation direction.
For contextual claims quantify over the actual typed closing contexts.
Changing evaluation to call-by-value, adding effects or parallel operations,
or treating function termination as an observation changes the proof target.
\end{hbmethod}

\begin{proofboundary}{def:p10-oa-source-boundary}{Source-specific scope and stronger language variants}
This supplies the omitted proof for Gunter's continuous bc-domain
call-by-name PCF interpretation, not its stable dI-domain interpretation
or all later record/subtyping languages. The predecessor-zero behavior
differs from Plotkin's original arithmetic constants. The continuous
model's source full-abstraction counterexample now has its sequentiality
dependency proved above. The stronger parallel-extension full-abstraction
and definability theorems do not follow from that failure proof.
Abramsky--Jung's predomain/partial-map and monadic
call-by-value proposals need their own value/computation semantics and
adequacy argument. Printed chapters, the held author presentations and
the whole-source allocations remain distinct.
\end{proofboundary}

\section{Constructing stable function domains and their PCF model}
\label{sec:p10-stable-domains-sequential-models}

The continuous PCF model in Section~\ref{sec:p10-operational-adequacy-domain-models}
uses full function spaces with pointwise order. Gunter's other interpretation
uses stable maps with a different order; its existence is not established by
changing the word ``continuous'' to ``stable''. We construct the stable
function domain, its finite approximants and its evaluation/recursion maps
\cite[Definitions~4.4.3--4; Lemma~4.4.5, pp.~35--36]{Gunter1994TypesAuthorOA}.
The modulus interpretation is from Abramsky--Jung's corrected/expanded
author chapter \cite[Section~8.2, p.~147]{AbramskyJungCorrectedExpandedDomainTheory}.
The construction here concerns pointed dI-domains and unrestricted stable
maps, not a claim that every stable domain theory or bidomain presentation
has these same hypotheses.

\begin{hbconvention}{def:p10-sd-domains}{Distributivity, finite compact ideals and compatibility}
A dI-domain here is a set-sized pointed algebraic dcpo, bounded-complete,
distributive, with property I: \(\downarrow a\) is finite for each
compact \(a\in K(D)\). Bounded-completeness supplies joins of all
bounded subsets; the empty join is bottom. Write \(x\mathrel{\uparrow}y\)
when \(x,y\) have a common upper bound. Nonempty infima exist, and
distributivity means
\[
 x\sqcap(y\sqcup z)=(x\sqcap y)\sqcup(x\sqcap z)
                       \qquad(y\mathrel{\uparrow}z).
\]
The compatibility condition is on \(y,z\), whose join is being used;
it is not replaced by compatibility of \(x,y\).
All maps called stable below are Scott-continuous and preserve meets of
compatible pairs. They need not be strict: constant maps are included.
\cite[pp.~35--36]{Gunter1994TypesAuthorOA}
\cite[Proposition~4.1.2, p.~56]{AbramskyJungCorrectedExpandedDomainTheory}.
\end{hbconvention}

\begin{hbtheorem}{thm:p10-sd-foundations}{proved=annex:p10.sd.foundations}{The finite-ideal facts needed by the construction}
Every element below a compact element of a dI-domain is compact.
Finite bounded joins of compact elements are compact. Binary meet is
jointly Scott-continuous. Every principal ideal is a distributive complete
lattice. Finite products of dI-domains are dI-domains, with coordinatewise
order and operations. These statements include the one-point product.
\end{hbtheorem}

\begin{hbdefinition}{def:p10-sd-stable-order}{Stable functions, moduli and their stronger order}
For dI-domains \(D,E\), let \([D\longrightarrow_s E]\) consist of
Scott-continuous functions satisfying
\(f(x\sqcap y)=f(x)\sqcap f(y)\) whenever \(x\mathrel{\uparrow}y\).
Its source-native stable order is
\[
 f\sqsubseteq_s g\quad\Longleftrightarrow\quad
 (\forall x\sqsubseteq y)\ f(x)=f(y)\sqcap g(x).
\]
Ordinary pointwise order is denoted \(f\sqsubseteq g\).
For compact output information \(b\sqsubseteq f(x)\), the modulus
\(M(f,x,b)\) is the least \(d\sqsubseteq x\) with
\(b\sqsubseteq f(d)\), when it exists. Thus the output must contain
the required information, not be below it. This is the compact algebraic
case of the source's minimum-input notion
\cite[p.~36]{Gunter1994TypesAuthorOA}
\cite[p.~147]{AbramskyJungCorrectedExpandedDomainTheory}.
\end{hbdefinition}

\begin{hbtheorem}{thm:p10-sd-moduli}{proved=annex:p10.sd.moduli}{Actual minimum witnesses and the stable-order equivalence}
A continuous map between dI-domains is stable iff these moduli exist for
every \(x\) and compact \(b\sqsubseteq f(x)\). Each modulus is compact.
Such a map preserves every bounded nonempty infimum.
For stable maps,
\[
 f\sqsubseteq_s g\quad\Longleftrightarrow\quad
 f\sqsubseteq g\ \text{and}\
 M(f,x,b)=M(g,x,b)\quad(b\in K(E),\ b\sqsubseteq f(x)).
\]
The stable order therefore records required input information, not merely
pointwise output refinement.
\end{hbtheorem}

\begin{hbtheorem}{thm:p10-sd-ordered-functions}{proved=annex:p10.sd.ordered-functions}{The actual stable function dcpo and its bounded joins}
The stable order is a partial order with the constant-bottom map least.
Directed suprema are pointwise suprema and remain stable maps.
Every stable-order bounded family has its stable supremum pointwise.
If \(f,g\sqsubseteq_s h\), then their stable meet is pointwise meet.
Within \(\downarrow_s h\), stable order agrees with pointwise order;
this need not hold for an arbitrary pair of stable functions.
\end{hbtheorem}

\begin{hbdefinition}{def:p10-sd-finite-cuts}{Finite cut functions, not an assumed function-space basis}
For stable \(f:D\to E\), compact \(a\in K(D)\), \(b\in K(E)\), put
\[
                     f_{a,b}(x)=f(x\sqcap a)\sqcap b.
\]
Its input cut lies in the finite ideal \(\downarrow a\), and its
range in \(\downarrow b\). A finite join of such cuts is formed only
when they have a common stable upper bound, for example when all are
cuts of the same \(f\). No unbounded pointwise join is stipulated.
\end{hbdefinition}

\begin{hbtheorem}{thm:p10-sd-function-domain}{proved=annex:p10.sd.function-domain}{Finite approximants prove full dI closure}
Each cut \(f_{a,b}\) is stable, compact in the stable function dcpo,
and \(f_{a,b}\sqsubseteq_s f\). The directed family of finite joins
of cuts of \(f\), including the empty bottom join, has supremum \(f\).
Every compact stable function is such a finite join of its own cuts.
Its stable principal ideal is finite. The function dcpo is distributive
and bounded-complete, hence a dI-domain.
Countable compact bases in \(D,E\) give a countable compact basis in
\([D\longrightarrow_s E]\). This proves the source's asserted
closure rather than using it as a premise
\cite[p.~36]{Gunter1994TypesAuthorOA}.
\end{hbtheorem}

\begin{hbtheorem}{thm:p10-sd-cartesian-closed}{proved=annex:p10.sd.cartesian-closed}{Stable evaluation, currying and composition}
Pointed dI-domains and all stable maps form a cartesian closed category.
Products are the actual coordinate products and the exponential is the
stable function domain. Evaluation \((f,x)\mapsto f(x)\) is jointly
continuous and stable. Currying gives inverse bijections of stable maps
\[
 [C\times D\longrightarrow_s E]
       \cong[C\longrightarrow_s[D\longrightarrow_s E]].
\]
Composition is monotone and preserves directed suprema in the displayed
stable hom orders. None of these morphisms is assumed strict.
\end{hbtheorem}

\begin{hbtheorem}{thm:p10-sd-stable-fix}{proved=annex:p10.sd.stable-fix}{The stable fixed-point operator is proved, not postulated}
For a dI-domain \(D\), the ordinary least fixed point of a stable
endomap gives a Scott-continuous stable map
\[
 \operatorname{fix}:[D\longrightarrow_s D]\longrightarrow_s D.
\]
For compatible endomaps with common stable upper bound,
\[
 \operatorname{fix}(f\sqcap_s g)
       =\operatorname{fix}(f)\sqcap\operatorname{fix}(g).
\]
This assertion is about the stable hom order; its proof tracks both
iteration chains under their common upper-bound chain.
\end{hbtheorem}

\begin{hbdefinition}{def:p10-sd-pcf-model}{Gunter's stable interpretation of the already specified PCF}
Keep exactly the language/evaluation of Definitions~\ref{def:p10-oa-language}
and~\ref{def:p10-oa-evaluation}, including \(\operatorname{pred}(0)=0\).
Write \(S_t=D\llbracket t\rrbracket\) for the source's stable
interpretation, distinct from its continuous interpretation \(C\):
\[
 S_{\mathsf{num}}=\mathbb N_{\bot},\quad
 S_{\mathsf{bool}}=\{\bot,\mathsf{true},\mathsf{false}\},\quad
 S_{s\to t}=[S_s\longrightarrow_s S_t].
\]
Use the stable evaluation/currying just constructed; ground constructors
and conditional retain their specified values; interpret \(\mu\) by
the proved stable fixed-point operator. No PCF parallel primitive is added.
The source calls the boolean carrier \(T=S_{\mathsf{bool}}\); here
\(T\) is a domain, not the earlier closed PCF counterexample term.
\cite[Lemma~4.4.5, p.~36]{Gunter1994TypesAuthorOA}.
\end{hbdefinition}

\begin{hbtheorem}{thm:p10-sd-pcf-interpretation}{proved=annex:p10.sd.pcf-interpretation}{Stable environment dependence and sound substitution}
Every stable PCF interpretation is a stable map from its finite environment
product to its result domain. It respects bound renaming and typed
substitution, and each registered call-by-name contraction preserves it.
In particular the body of a typed recursion binder induces an actual
stable endomap and the binder's meaning is well defined.
\end{hbtheorem}

\begin{hbdefinition}{def:p10-sd-relation}{The adequacy relation on stable type carriers}
For \(a\in S_t\) and a closed \(M:t\), write
\(a\mathrel{\mathcal R^s_t}M\). At ground type bottom is related
to every term, and a nonbottom point is related exactly when the term
evaluates to its corresponding numeral or boolean. At an arrow type,
\[
 f\mathrel{\mathcal R^s_{u\to t}}M\quad\Longleftrightarrow\quad
 (\forall N:u\text{ closed})(\forall a\in S_u)\
 a\mathrel{\mathcal R^s_u}N\ \Longrightarrow\
 f(a)\mathrel{\mathcal R^s_t}(M N).
\]
The arrow element is in the stable function carrier. This definition
retains the old operational language but changes the semantic carriers;
its closure and fundamental lemma are proved below.
\end{hbdefinition}

\begin{hbtheorem}{thm:p10-sd-pcf-adequacy}{proved=annex:p10.sd.pcf-adequacy}{Both stable-model observation directions and contextual adequacy}
For closed ground programs and ground values,
\[
 M\Downarrow V\quad\Longleftrightarrow\quad
             D\llbracket M\rrbracket=D\llbracket V\rrbracket,
 \qquad
 M\Uparrow\quad\Longleftrightarrow\quad D\llbracket M\rrbracket=\bot.
\]
Stable-model denotational equality implies contextual operational
equivalence. The typed semantic-to-syntax logical relation is instantiated
on \(S_t\), not assumed equal to the earlier relation on the continuous
type carriers. Its recursion and directed-closure dependencies are checked.
\end{hbtheorem}

\begin{hbtheorem}{thm:p10-sd-counterexample}{proved=annex:p10.sd.counterexample}{Pointwise/stable incomparability and the source's second failure of full abstraction}
On the flat boolean domain let
\(p(x)=\mathsf{true}\) and
\(q(\bot)=\bot\), \(q(\mathsf{true})=q(\mathsf{false})=\mathsf{true}\).
Both are stable; \(q\sqsubseteq p\) pointwise, but they are incomparable
and incompatible in the stable order. Their stable meet is the bottom
function although their pointwise meet is \(q\). On the stable boolean function
domain define
\[
 r(h)=\begin{cases}\mathsf{true},&h=q,\\
                    \mathsf{false},&h=p,\\\bot,&\text{otherwise}.
       \end{cases}
\]
This is continuous and stable. The functions \(p,q\) are represented
by the closed PCF abbreviations
\(p=\lambda x:\mathsf{bool}.\mathsf{true}\) and
\(q=\lambda x:\mathsf{bool}.\mathsf{if}\ x\ \mathsf{then}\ \mathsf{true}\ \mathsf{else}\ \mathsf{true}\).
With \(\Omega=\Omega_{\mathsf{bool}}\), set
\[
 \begin{split}
 M={}&\lambda f:(\mathsf{bool}\to\mathsf{bool})\to\mathsf{bool}.
    \mathsf{if}\ f q\ \mathsf{then}\
    (\mathsf{if}\ f p\ \mathsf{then}\ \Omega\ \mathsf{else}\ \mathsf{true})
       \ \mathsf{else}\ \Omega,\\
 N={}&\lambda f:(\mathsf{bool}\to\mathsf{bool})\to\mathsf{bool}.
    \mathsf{if}\ f q\ \mathsf{then}\
    (\mathsf{if}\ f p\ \mathsf{then}\ \Omega\ \mathsf{else}\ \mathsf{false})
       \ \mathsf{else}\ \Omega.
 \end{split}
\]
The source-native names \(p,q,r,M,N\) denote these typed functions
and abbreviations, not ambient predicate or set variables.
Then \(M\simeq_o N\), while
\(D\llbracket M\rrbracket(r)=\mathsf{true}\ne\mathsf{false}=D\llbracket N\rrbracket(r)\).
Thus the stable model is adequate but not fully abstract. Parallel-or from
the previous counterexample is itself not stable. These are different
failures: \(r\) is not the denotation of any closed term of this PCF.
Stable maps therefore do not coincide with its sequentially definable programs
\cite[p.~38]{Gunter1994TypesAuthorOA}.
\end{hbtheorem}

\begin{proofboundary}{def:p10-sd-scope}{A stable construction is not a bidomain classification}
The proved construction is the pointed dI stable-function model, with
finite-ideal hypotheses used explicitly. General continuous stable
L-domains, concrete-domain sequential algorithms, two-order bidomain
axioms/classifications, and stable recursive type equations have separate
construction obligations. Merely displaying both pointwise and stable
orders does not establish a bidomain theorem. The distinct Berry/Girard/
coherence-space presentations referred to by the sources, their original
papers, and complete printed/author manifestations remain separate duties.
\end{proofboundary}


\section{Theorem and definition register}
\label{sec:p10-theorem-register}
% Generated by tools/audit_tables.py -- do not hand-edit.
\begingroup\scriptsize\setlength{\tabcolsep}{3pt}
\begin{longtable}{@{}p{0.26\textwidth}p{0.40\textwidth}p{0.11\textwidth}p{0.17\textwidth}@{}}
\caption{Theorem and definition register for Part 10.}\label{tab:p10-theorem-definition-register}\\
\toprule
\textbf{ID} & \textbf{Title} & \textbf{Status} & \textbf{Target} \\
\midrule
\endfirsthead
\toprule
\textbf{ID} & \textbf{Title} & \textbf{Status} & \textbf{Target} \\
\midrule
\endhead
\midrule\multicolumn{4}{r}{\emph{continued}}\\\endfoot
\bottomrule\endlastfoot
\hyperref[def:p10-algebra]{\texttt{\detokenize{def:}\allowbreak\detokenize{p10-}\allowbreak\detokenize{algebra}}} & Algebras, homomorphisms, and congruences & -- & -- \\
\hyperref[thm:p10-hsp]{\texttt{\detokenize{thm:}\allowbreak\detokenize{p10-}\allowbreak\detokenize{hsp}}} & Birkhoff HSP theorem & cited & \texttt{\detokenize{Birkhoff1935}} \\
\hyperref[def:p10-boolean-algebra]{\texttt{\detokenize{def:}\allowbreak\detokenize{p10-}\allowbreak\detokenize{boolean-}\allowbreak\detokenize{algebra}}} & Boolean algebra, filter, ideal, and ultrafilter & -- & -- \\
\hyperref[def:p10-lindenbaum-algebra]{\texttt{\detokenize{def:}\allowbreak\detokenize{p10-}\allowbreak\detokenize{lindenbaum-}\allowbreak\detokenize{algebra}}} & Lindenbaum--Tarski algebra & -- & -- \\
\hyperref[def:p10-mv-algebra]{\texttt{\detokenize{def:}\allowbreak\detokenize{p10-}\allowbreak\detokenize{mv-}\allowbreak\detokenize{algebra}}} & MV-algebra & -- & -- \\
\hyperref[def:p10-residuated-lattice]{\texttt{\detokenize{def:}\allowbreak\detokenize{p10-}\allowbreak\detokenize{residuated-}\allowbreak\detokenize{lattice}}} & Residuated lattice & -- & -- \\
\hyperref[def:p10-topos]{\texttt{\detokenize{def:}\allowbreak\detokenize{p10-}\allowbreak\detokenize{topos}}} & Elementary topos and subobject classifier & -- & -- \\
\hyperref[def:p10-term-algebra]{\texttt{\detokenize{def:}\allowbreak\detokenize{p10-}\allowbreak\detokenize{term-}\allowbreak\detokenize{algebra}}} & Term algebra and evaluation & -- & -- \\
\hyperref[thm:p10-term-universal]{\texttt{\detokenize{thm:}\allowbreak\detokenize{p10-}\allowbreak\detokenize{term-}\allowbreak\detokenize{universal}}} & Universal property and substitution law for terms & proved & \texttt{\detokenize{annex:}\allowbreak\detokenize{p10.eq.1}} \\
\hyperref[def:p10-equational-logic]{\texttt{\detokenize{def:}\allowbreak\detokenize{p10-}\allowbreak\detokenize{equational-}\allowbreak\detokenize{logic}}} & Equational consequence & -- & -- \\
\hyperref[def:p10-identities-valuations]{\texttt{\detokenize{def:}\allowbreak\detokenize{p10-}\allowbreak\detokenize{identities-}\allowbreak\detokenize{valuations}}} & An identity and a local equation & -- & -- \\
\hyperref[def:p10-deductive-congruence]{\texttt{\detokenize{def:}\allowbreak\detokenize{p10-}\allowbreak\detokenize{deductive-}\allowbreak\detokenize{congruence}}} & The congruence of provable identities & -- & -- \\
\hyperref[thm:p10-equational-completeness]{\texttt{\detokenize{thm:}\allowbreak\detokenize{p10-}\allowbreak\detokenize{equational-}\allowbreak\detokenize{completeness}}} & Completeness of equational logic & proved & \texttt{\detokenize{annex:}\allowbreak\detokenize{p10.eq.2}} \\
\hyperref[def:p10-free-algebra]{\texttt{\detokenize{def:}\allowbreak\detokenize{p10-}\allowbreak\detokenize{free-}\allowbreak\detokenize{algebra}}} & Free algebra on a set of generators & -- & -- \\
\hyperref[thm:p10-free-algebra]{\texttt{\detokenize{thm:}\allowbreak\detokenize{p10-}\allowbreak\detokenize{free-}\allowbreak\detokenize{algebra}}} & Construction and uniqueness of free algebras & proved & \texttt{\detokenize{annex:}\allowbreak\detokenize{p10.eq.3}} \\
\hyperref[def:p10-subdirect-product]{\texttt{\detokenize{def:}\allowbreak\detokenize{p10-}\allowbreak\detokenize{subdirect-}\allowbreak\detokenize{product}}} & Subdirect products and embeddings & -- & -- \\
\hyperref[thm:p10-subdirect-kernels]{\texttt{\detokenize{thm:}\allowbreak\detokenize{p10-}\allowbreak\detokenize{subdirect-}\allowbreak\detokenize{kernels}}} & The congruence-intersection criterion & proved & \texttt{\detokenize{annex:}\allowbreak\detokenize{p10.subdirect.1}} \\
\hyperref[thm:p10-congruence-correspondence]{\texttt{\detokenize{thm:}\allowbreak\detokenize{p10-}\allowbreak\detokenize{congruence-}\allowbreak\detokenize{correspondence}}} & Congruences above a quotient kernel & proved & \texttt{\detokenize{annex:}\allowbreak\detokenize{p10.subdirect.2}} \\
\hyperref[def:p10-subdirect-irreducible]{\texttt{\detokenize{def:}\allowbreak\detokenize{p10-}\allowbreak\detokenize{subdirect-}\allowbreak\detokenize{irreducible}}} & Subdirect irreducibility and the monolith & -- & -- \\
\hyperref[thm:p10-subdirect-monolith]{\texttt{\detokenize{thm:}\allowbreak\detokenize{p10-}\allowbreak\detokenize{subdirect-}\allowbreak\detokenize{monolith}}} & Characterizing subdirect irreducibility & proved & \texttt{\detokenize{annex:}\allowbreak\detokenize{p10.subdirect.3}} \\
\hyperref[thm:p10-subdirect-representation]{\texttt{\detokenize{thm:}\allowbreak\detokenize{p10-}\allowbreak\detokenize{subdirect-}\allowbreak\detokenize{representation}}} & Birkhoff's subdirect representation theorem & proved & \texttt{\detokenize{annex:}\allowbreak\detokenize{p10.subdirect.4}} \\
\hyperref[thm:p10-ultrafilter-separation]{\texttt{\detokenize{thm:}\allowbreak\detokenize{p10-}\allowbreak\detokenize{ultrafilter-}\allowbreak\detokenize{separation}}} & Ultrafilter dichotomy and separation & proved & \texttt{\detokenize{annex:}\allowbreak\detokenize{p10.stone.1}} \\
\hyperref[def:p10-stone-space]{\texttt{\detokenize{def:}\allowbreak\detokenize{p10-}\allowbreak\detokenize{stone-}\allowbreak\detokenize{space}}} & Stone space of a Boolean algebra & -- & -- \\
\hyperref[thm:p10-stone-representation]{\texttt{\detokenize{thm:}\allowbreak\detokenize{p10-}\allowbreak\detokenize{stone-}\allowbreak\detokenize{representation}}} & Stone representation theorem & proved & \texttt{\detokenize{annex:}\allowbreak\detokenize{p10.stone.2}} \\
\hyperref[thm:p10-stone-space-reconstruction]{\texttt{\detokenize{thm:}\allowbreak\detokenize{p10-}\allowbreak\detokenize{stone-}\allowbreak\detokenize{space-}\allowbreak\detokenize{reconstruction}}} & A Stone space is recovered from its clopen algebra & proved & \texttt{\detokenize{annex:}\allowbreak\detokenize{p10.stone.3}} \\
\hyperref[def:p10-stone-dual-maps]{\texttt{\detokenize{def:}\allowbreak\detokenize{p10-}\allowbreak\detokenize{stone-}\allowbreak\detokenize{dual-}\allowbreak\detokenize{maps}}} & The inverse-image maps of Stone duality & -- & -- \\
\hyperref[thm:p10-stone-duality]{\texttt{\detokenize{thm:}\allowbreak\detokenize{p10-}\allowbreak\detokenize{stone-}\allowbreak\detokenize{duality}}} & Functorial Stone duality & proved & \texttt{\detokenize{annex:}\allowbreak\detokenize{p10.stone.4}} \\
\hyperref[def:p10-atomless-countability]{\texttt{\detokenize{def:}\allowbreak\detokenize{p10-}\allowbreak\detokenize{atomless-}\allowbreak\detokenize{countability}}} & Atomlessness and countability & -- & -- \\
\hyperref[thm:p10-atomless-extension]{\texttt{\detokenize{thm:}\allowbreak\detokenize{p10-}\allowbreak\detokenize{atomless-}\allowbreak\detokenize{extension}}} & Extending an isomorphism of finite Boolean subalgebras & proved & \texttt{\detokenize{annex:}\allowbreak\detokenize{p10.atomless.1}} \\
\hyperref[thm:p10-countable-atomless-uniqueness]{\texttt{\detokenize{thm:}\allowbreak\detokenize{p10-}\allowbreak\detokenize{countable-}\allowbreak\detokenize{atomless-}\allowbreak\detokenize{uniqueness}}} & Uniqueness of the countable atomless Boolean algebra & proved & \texttt{\detokenize{annex:}\allowbreak\detokenize{p10.atomless.2}} \\
\hyperref[def:p10-cantor-cylinder-algebra]{\texttt{\detokenize{def:}\allowbreak\detokenize{p10-}\allowbreak\detokenize{cantor-}\allowbreak\detokenize{cylinder-}\allowbreak\detokenize{algebra}}} & Cantor space and its cylinder algebra & -- & -- \\
\hyperref[thm:p10-cantor-cylinder-algebra]{\texttt{\detokenize{thm:}\allowbreak\detokenize{p10-}\allowbreak\detokenize{cantor-}\allowbreak\detokenize{cylinder-}\allowbreak\detokenize{algebra}}} & Existence through clopen cylinders & proved & \texttt{\detokenize{annex:}\allowbreak\detokenize{p10.atomless.3}} \\
\hyperref[thm:p10-countable-atomless-stone]{\texttt{\detokenize{thm:}\allowbreak\detokenize{p10-}\allowbreak\detokenize{countable-}\allowbreak\detokenize{atomless-}\allowbreak\detokenize{stone}}} & The Stone space is Cantor space & proved & \texttt{\detokenize{annex:}\allowbreak\detokenize{p10.atomless.4}} \\
\hyperref[thm:p10-atoms-isolated]{\texttt{\detokenize{thm:}\allowbreak\detokenize{p10-}\allowbreak\detokenize{atoms-}\allowbreak\detokenize{isolated}}} & Atoms and isolated Stone points & proved & \texttt{\detokenize{annex:}\allowbreak\detokenize{p10.atomless.5}} \\
\hyperref[def:p10-distributive-prime-filter]{\texttt{\detokenize{def:}\allowbreak\detokenize{p10-}\allowbreak\detokenize{distributive-}\allowbreak\detokenize{prime-}\allowbreak\detokenize{filter}}} & Distributive-lattice filters and prime filters & -- & -- \\
\hyperref[thm:p10-prime-filter-separation]{\texttt{\detokenize{thm:}\allowbreak\detokenize{p10-}\allowbreak\detokenize{prime-}\allowbreak\detokenize{filter-}\allowbreak\detokenize{separation}}} & Prime-filter separation & proved & \texttt{\detokenize{annex:}\allowbreak\detokenize{p10.priestley.1}} \\
\hyperref[def:p10-priestley-space]{\texttt{\detokenize{def:}\allowbreak\detokenize{p10-}\allowbreak\detokenize{priestley-}\allowbreak\detokenize{space}}} & Priestley space and clopen upsets & -- & -- \\
\hyperref[thm:p10-priestley-representation]{\texttt{\detokenize{thm:}\allowbreak\detokenize{p10-}\allowbreak\detokenize{priestley-}\allowbreak\detokenize{representation}}} & Priestley representation & proved & \texttt{\detokenize{annex:}\allowbreak\detokenize{p10.priestley.2}} \\
\hyperref[thm:p10-priestley-reconstruction]{\texttt{\detokenize{thm:}\allowbreak\detokenize{p10-}\allowbreak\detokenize{priestley-}\allowbreak\detokenize{reconstruction}}} & Recovering a Priestley space & proved & \texttt{\detokenize{annex:}\allowbreak\detokenize{p10.priestley.3}} \\
\hyperref[thm:p10-priestley-functors]{\texttt{\detokenize{thm:}\allowbreak\detokenize{p10-}\allowbreak\detokenize{priestley-}\allowbreak\detokenize{functors}}} & Functorial Priestley duality & proved & \texttt{\detokenize{annex:}\allowbreak\detokenize{p10.priestley.4}} \\
\hyperref[def:p10-dense-generic-filter]{\texttt{\detokenize{def:}\allowbreak\detokenize{p10-}\allowbreak\detokenize{dense-}\allowbreak\detokenize{generic-}\allowbreak\detokenize{filter}}} & Dense sets and generic filters in a poset & -- & -- \\
\hyperref[def:p10-dependent-choice]{\texttt{\detokenize{def:}\allowbreak\detokenize{p10-}\allowbreak\detokenize{dependent-}\allowbreak\detokenize{choice}}} & Dependent choice for a prescribed start & -- & -- \\
\hyperref[thm:p10-rasiowa-sikorski]{\texttt{\detokenize{thm:}\allowbreak\detokenize{p10-}\allowbreak\detokenize{rasiowa-}\allowbreak\detokenize{sikorski}}} & Rasiowa--Sikorski lemma & proved & \texttt{\detokenize{annex:}\allowbreak\detokenize{p10.rs.1}} \\
\hyperref[def:p10-heyting-algebra]{\texttt{\detokenize{def:}\allowbreak\detokenize{p10-}\allowbreak\detokenize{heyting-}\allowbreak\detokenize{algebra}}} & Heyting algebra & -- & -- \\
\hyperref[thm:p10-heyting-examples]{\texttt{\detokenize{thm:}\allowbreak\detokenize{p10-}\allowbreak\detokenize{heyting-}\allowbreak\detokenize{examples}}} & Open-set and upset Heyting algebras & proved & \texttt{\detokenize{annex:}\allowbreak\detokenize{p10.heyting.1}} \\
\hyperref[def:p10-ipc-lindenbaum]{\texttt{\detokenize{def:}\allowbreak\detokenize{p10-}\allowbreak\detokenize{ipc-}\allowbreak\detokenize{lindenbaum}}} & IPC derivability and its Lindenbaum algebra & -- & -- \\
\hyperref[thm:p10-heyting-completeness]{\texttt{\detokenize{thm:}\allowbreak\detokenize{p10-}\allowbreak\detokenize{heyting-}\allowbreak\detokenize{completeness}}} & Algebraic completeness of IPC & proved & \texttt{\detokenize{annex:}\allowbreak\detokenize{p10.heyting.2}} \\
\hyperref[def:p10-heyting-prime-frame]{\texttt{\detokenize{def:}\allowbreak\detokenize{p10-}\allowbreak\detokenize{heyting-}\allowbreak\detokenize{prime-}\allowbreak\detokenize{frame}}} & The prime-filter order of a Heyting algebra & -- & -- \\
\hyperref[thm:p10-heyting-prime-representation]{\texttt{\detokenize{thm:}\allowbreak\detokenize{p10-}\allowbreak\detokenize{heyting-}\allowbreak\detokenize{prime-}\allowbreak\detokenize{representation}}} & Prime-filter representation and the implication witness & proved & \texttt{\detokenize{annex:}\allowbreak\detokenize{p10.heyting.3}} \\
\hyperref[thm:p10-ipc-prime-kripke-completeness]{\texttt{\detokenize{thm:}\allowbreak\detokenize{p10-}\allowbreak\detokenize{ipc-}\allowbreak\detokenize{prime-}\allowbreak\detokenize{kripke-}\allowbreak\detokenize{completeness}}} & Strong local Kripke completeness of IPC & proved & \texttt{\detokenize{annex:}\allowbreak\detokenize{p10.heyting.4}} \\
\hyperref[def:p10-esakia-space]{\texttt{\detokenize{def:}\allowbreak\detokenize{p10-}\allowbreak\detokenize{esakia-}\allowbreak\detokenize{space}}} & Esakia spaces & -- & -- \\
\hyperref[thm:p10-esakia-objects]{\texttt{\detokenize{thm:}\allowbreak\detokenize{p10-}\allowbreak\detokenize{esakia-}\allowbreak\detokenize{objects}}} & Heyting algebras and Esakia spaces & proved & \texttt{\detokenize{annex:}\allowbreak\detokenize{p10.esakia.1}} \\
\hyperref[def:p10-esakia-morphism]{\texttt{\detokenize{def:}\allowbreak\detokenize{p10-}\allowbreak\detokenize{esakia-}\allowbreak\detokenize{morphism}}} & Esakia morphisms & -- & -- \\
\hyperref[thm:p10-heyting-prime-lift]{\texttt{\detokenize{thm:}\allowbreak\detokenize{p10-}\allowbreak\detokenize{heyting-}\allowbreak\detokenize{prime-}\allowbreak\detokenize{lift}}} & Lifting prime filters along a Heyting homomorphism & proved & \texttt{\detokenize{annex:}\allowbreak\detokenize{p10.esakia.2}} \\
\hyperref[thm:p10-esakia-functors]{\texttt{\detokenize{thm:}\allowbreak\detokenize{p10-}\allowbreak\detokenize{esakia-}\allowbreak\detokenize{functors}}} & Functorial Esakia duality & proved & \texttt{\detokenize{annex:}\allowbreak\detokenize{p10.esakia.3}} \\
\hyperref[thm:p10-heyting-equational]{\texttt{\detokenize{thm:}\allowbreak\detokenize{p10-}\allowbreak\detokenize{heyting-}\allowbreak\detokenize{equational}}} & The equational class of Heyting algebras & proved & \texttt{\detokenize{annex:}\allowbreak\detokenize{p10.intermediate.1}} \\
\hyperref[def:p10-intermediate-logic]{\texttt{\detokenize{def:}\allowbreak\detokenize{p10-}\allowbreak\detokenize{intermediate-}\allowbreak\detokenize{logic}}} & Intermediate theorem logics and their quotients & -- & -- \\
\hyperref[def:p10-intermediate-variety-maps]{\texttt{\detokenize{def:}\allowbreak\detokenize{p10-}\allowbreak\detokenize{intermediate-}\allowbreak\detokenize{variety-}\allowbreak\detokenize{maps}}} & The logic--variety assignments & -- & -- \\
\hyperref[thm:p10-intermediate-algebra-completeness]{\texttt{\detokenize{thm:}\allowbreak\detokenize{p10-}\allowbreak\detokenize{intermediate-}\allowbreak\detokenize{algebra-}\allowbreak\detokenize{completeness}}} & Algebraic completeness of an intermediate logic & proved & \texttt{\detokenize{annex:}\allowbreak\detokenize{p10.intermediate.2}} \\
\hyperref[thm:p10-intermediate-variety-correspondence]{\texttt{\detokenize{thm:}\allowbreak\detokenize{p10-}\allowbreak\detokenize{intermediate-}\allowbreak\detokenize{variety-}\allowbreak\detokenize{correspondence}}} & The lattice correspondence for intermediate logics & proved & \texttt{\detokenize{annex:}\allowbreak\detokenize{p10.intermediate.3}} \\
\hyperref[def:p10-bao]{\texttt{\detokenize{def:}\allowbreak\detokenize{p10-}\allowbreak\detokenize{bao}}} & Boolean algebra with operators & -- & -- \\
\hyperref[def:p10-bao-complex-algebra]{\texttt{\detokenize{def:}\allowbreak\detokenize{p10-}\allowbreak\detokenize{bao-}\allowbreak\detokenize{complex-}\allowbreak\detokenize{algebra}}} & Relational frames and complex algebras & -- & -- \\
\hyperref[def:p10-bao-ultrafilter-frame]{\texttt{\detokenize{def:}\allowbreak\detokenize{p10-}\allowbreak\detokenize{bao-}\allowbreak\detokenize{ultrafilter-}\allowbreak\detokenize{frame}}} & The ultrafilter frame of a Boolean algebra with operators & -- & -- \\
\hyperref[def:p10-bao-binary-nullary]{\texttt{\detokenize{def:}\allowbreak\detokenize{p10-}\allowbreak\detokenize{bao-}\allowbreak\detokenize{binary-}\allowbreak\detokenize{nullary}}} & Binary meet and a nullary operator & -- & -- \\
\hyperref[thm:p10-bao-ultrafilter-witness]{\texttt{\detokenize{thm:}\allowbreak\detokenize{p10-}\allowbreak\detokenize{bao-}\allowbreak\detokenize{ultrafilter-}\allowbreak\detokenize{witness}}} & Simultaneous ultrafilter witnesses & proved & \texttt{\detokenize{annex:}\allowbreak\detokenize{p10.bao.1}} \\
\hyperref[thm:p10-jonsson-tarski]{\texttt{\detokenize{thm:}\allowbreak\detokenize{p10-}\allowbreak\detokenize{jonsson-}\allowbreak\detokenize{tarski}}} & J\'onsson--Tarski representation & proved & \texttt{\detokenize{annex:}\allowbreak\detokenize{p10.bao.2}} \\
\hyperref[def:p10-bao-canonical-extension]{\texttt{\detokenize{def:}\allowbreak\detokenize{p10-}\allowbreak\detokenize{bao-}\allowbreak\detokenize{canonical-}\allowbreak\detokenize{extension}}} & The canonical extension and its Boolean completion & -- & -- \\
\hyperref[thm:p10-bao-canonical-completion]{\texttt{\detokenize{thm:}\allowbreak\detokenize{p10-}\allowbreak\detokenize{bao-}\allowbreak\detokenize{canonical-}\allowbreak\detokenize{completion}}} & The dense and compact canonical completion & proved & \texttt{\detokenize{annex:}\allowbreak\detokenize{p10.bao.3}} \\
\hyperref[def:p10-bao-general-frame]{\texttt{\detokenize{def:}\allowbreak\detokenize{p10-}\allowbreak\detokenize{bao-}\allowbreak\detokenize{general-}\allowbreak\detokenize{frame}}} & A general frame for a finitary operator signature & -- & -- \\
\hyperref[def:p10-bao-descriptive]{\texttt{\detokenize{def:}\allowbreak\detokenize{p10-}\allowbreak\detokenize{bao-}\allowbreak\detokenize{descriptive}}} & Descriptive general frames & -- & -- \\
\hyperref[thm:p10-bao-descriptive-representation]{\texttt{\detokenize{thm:}\allowbreak\detokenize{p10-}\allowbreak\detokenize{bao-}\allowbreak\detokenize{descriptive-}\allowbreak\detokenize{representation}}} & Descriptive representation and reconstruction & proved & \texttt{\detokenize{annex:}\allowbreak\detokenize{p10.bao.4}} \\
\hyperref[def:p10-bao-general-morphism]{\texttt{\detokenize{def:}\allowbreak\detokenize{p10-}\allowbreak\detokenize{bao-}\allowbreak\detokenize{general-}\allowbreak\detokenize{morphism}}} & Bounded morphisms respecting admissible sets & -- & -- \\
\hyperref[thm:p10-bao-descriptive-duality]{\texttt{\detokenize{thm:}\allowbreak\detokenize{p10-}\allowbreak\detokenize{bao-}\allowbreak\detokenize{descriptive-}\allowbreak\detokenize{duality}}} & Duality with descriptive general frames & proved & \texttt{\detokenize{annex:}\allowbreak\detokenize{p10.bao.5}} \\
\hyperref[def:p10-normal-modal-algebra-logic]{\texttt{\detokenize{def:}\allowbreak\detokenize{p10-}\allowbreak\detokenize{normal-}\allowbreak\detokenize{modal-}\allowbreak\detokenize{algebra-}\allowbreak\detokenize{logic}}} & The fixed modal language and its algebras & -- & -- \\
\hyperref[def:p10-modal-lindenbaum]{\texttt{\detokenize{def:}\allowbreak\detokenize{p10-}\allowbreak\detokenize{modal-}\allowbreak\detokenize{lindenbaum}}} & The normal modal Lindenbaum algebra & -- & -- \\
\hyperref[thm:p10-modal-algebra-completeness]{\texttt{\detokenize{thm:}\allowbreak\detokenize{p10-}\allowbreak\detokenize{modal-}\allowbreak\detokenize{algebra-}\allowbreak\detokenize{completeness}}} & Algebraic completeness for every normal modal logic & proved & \texttt{\detokenize{annex:}\allowbreak\detokenize{p10.modal-}\allowbreak\detokenize{algebra.1}} \\
\hyperref[def:p10-modal-canonical-general-frame]{\texttt{\detokenize{def:}\allowbreak\detokenize{p10-}\allowbreak\detokenize{modal-}\allowbreak\detokenize{canonical-}\allowbreak\detokenize{general-}\allowbreak\detokenize{frame}}} & The canonical general frame and its underlying frame & -- & -- \\
\hyperref[thm:p10-modal-general-frame-completeness]{\texttt{\detokenize{thm:}\allowbreak\detokenize{p10-}\allowbreak\detokenize{modal-}\allowbreak\detokenize{general-}\allowbreak\detokenize{frame-}\allowbreak\detokenize{completeness}}} & General-frame completeness and the canonical truth lemma & proved & \texttt{\detokenize{annex:}\allowbreak\detokenize{p10.modal-}\allowbreak\detokenize{algebra.2}} \\
\hyperref[thm:p10-modal-general-frame-strong]{\texttt{\detokenize{thm:}\allowbreak\detokenize{p10-}\allowbreak\detokenize{modal-}\allowbreak\detokenize{general-}\allowbreak\detokenize{frame-}\allowbreak\detokenize{strong}}} & Strong completeness for local general-frame consequence & proved & \texttt{\detokenize{annex:}\allowbreak\detokenize{p10.modal-}\allowbreak\detokenize{algebra.3}} \\
\hyperref[def:p10-modal-canonicity]{\texttt{\detokenize{def:}\allowbreak\detokenize{p10-}\allowbreak\detokenize{modal-}\allowbreak\detokenize{canonicity}}} & Canonicity and unrestricted valuations & -- & -- \\
\hyperref[thm:p10-modal-canonicity-completeness]{\texttt{\detokenize{thm:}\allowbreak\detokenize{p10-}\allowbreak\detokenize{modal-}\allowbreak\detokenize{canonicity-}\allowbreak\detokenize{completeness}}} & Completeness from canonical-frame validity & proved & \texttt{\detokenize{annex:}\allowbreak\detokenize{p10.modal-}\allowbreak\detokenize{algebra.4}} \\
\hyperref[def:p10-polyadic-normal-logic]{\texttt{\detokenize{def:}\allowbreak\detokenize{p10-}\allowbreak\detokenize{polyadic-}\allowbreak\detokenize{normal-}\allowbreak\detokenize{logic}}} & Finitary normal modal language and logic & -- & -- \\
\hyperref[def:p10-polyadic-modal-algebra]{\texttt{\detokenize{def:}\allowbreak\detokenize{p10-}\allowbreak\detokenize{polyadic-}\allowbreak\detokenize{modal-}\allowbreak\detokenize{algebra}}} & Algebraic interpretation and the finitary Lindenbaum algebra & -- & -- \\
\hyperref[thm:p10-polyadic-algebra-completeness]{\texttt{\detokenize{thm:}\allowbreak\detokenize{p10-}\allowbreak\detokenize{polyadic-}\allowbreak\detokenize{algebra-}\allowbreak\detokenize{completeness}}} & Algebraic completeness for a finitary normal logic & proved & \texttt{\detokenize{annex:}\allowbreak\detokenize{p10.polyadic.1}} \\
\hyperref[def:p10-polyadic-canonical-frame]{\texttt{\detokenize{def:}\allowbreak\detokenize{p10-}\allowbreak\detokenize{polyadic-}\allowbreak\detokenize{canonical-}\allowbreak\detokenize{frame}}} & The descriptive canonical frame & -- & -- \\
\hyperref[thm:p10-polyadic-canonical-truth]{\texttt{\detokenize{thm:}\allowbreak\detokenize{p10-}\allowbreak\detokenize{polyadic-}\allowbreak\detokenize{canonical-}\allowbreak\detokenize{truth}}} & Descriptive representation and canonical truth & proved & \texttt{\detokenize{annex:}\allowbreak\detokenize{p10.polyadic.2}} \\
\hyperref[thm:p10-polyadic-local-completeness]{\texttt{\detokenize{thm:}\allowbreak\detokenize{p10-}\allowbreak\detokenize{polyadic-}\allowbreak\detokenize{local-}\allowbreak\detokenize{completeness}}} & Strong local completeness over descriptive general frames & proved & \texttt{\detokenize{annex:}\allowbreak\detokenize{p10.polyadic.3}} \\
\hyperref[def:p10-categorical-size]{\texttt{\detokenize{def:}\allowbreak\detokenize{p10-}\allowbreak\detokenize{categorical-}\allowbreak\detokenize{size}}} & Size conventions for categorical constructions & -- & -- \\
\hyperref[def:p10-category]{\texttt{\detokenize{def:}\allowbreak\detokenize{p10-}\allowbreak\detokenize{category}}} & Category, functor, and natural transformation & -- & -- \\
\hyperref[def:p10-category-equivalence]{\texttt{\detokenize{def:}\allowbreak\detokenize{p10-}\allowbreak\detokenize{category-}\allowbreak\detokenize{equivalence}}} & Fullness, faithfulness, and equivalence & -- & -- \\
\hyperref[def:p10-limit]{\texttt{\detokenize{def:}\allowbreak\detokenize{p10-}\allowbreak\detokenize{limit}}} & Diagrams, cones, limits, and colimits & -- & -- \\
\hyperref[def:p10-products-equalizers]{\texttt{\detokenize{def:}\allowbreak\detokenize{p10-}\allowbreak\detokenize{products-}\allowbreak\detokenize{equalizers}}} & Products, equalizers, and their duals & -- & -- \\
\hyperref[thm:p10-limits-unique]{\texttt{\detokenize{thm:}\allowbreak\detokenize{p10-}\allowbreak\detokenize{limits-}\allowbreak\detokenize{unique}}} & Uniqueness of universal constructions & proved & \texttt{\detokenize{annex:}\allowbreak\detokenize{p10.cat.1}} \\
\hyperref[thm:p10-limits-products-equalizers]{\texttt{\detokenize{thm:}\allowbreak\detokenize{p10-}\allowbreak\detokenize{limits-}\allowbreak\detokenize{products-}\allowbreak\detokenize{equalizers}}} & Constructing limits and colimits & proved & \texttt{\detokenize{annex:}\allowbreak\detokenize{p10.cat.2}} \\
\hyperref[thm:p10-yoneda]{\texttt{\detokenize{thm:}\allowbreak\detokenize{p10-}\allowbreak\detokenize{yoneda}}} & Yoneda lemma & proved & \texttt{\detokenize{annex:}\allowbreak\detokenize{p10.cat.3}} \\
\hyperref[def:p10-adjunction]{\texttt{\detokenize{def:}\allowbreak\detokenize{p10-}\allowbreak\detokenize{adjunction}}} & Adjunction by natural hom-set bijections & -- & -- \\
\hyperref[thm:p10-adjunction-data]{\texttt{\detokenize{thm:}\allowbreak\detokenize{p10-}\allowbreak\detokenize{adjunction-}\allowbreak\detokenize{data}}} & Units, counits, and triangle identities & proved & \texttt{\detokenize{annex:}\allowbreak\detokenize{p10.cat.4}} \\
\hyperref[thm:p10-adjoint-preservation]{\texttt{\detokenize{thm:}\allowbreak\detokenize{p10-}\allowbreak\detokenize{adjoint-}\allowbreak\detokenize{preservation}}} & Adjoints preserve universal constructions & proved & \texttt{\detokenize{annex:}\allowbreak\detokenize{p10.cat.5}} \\
\hyperref[def:p10-ccc]{\texttt{\detokenize{def:}\allowbreak\detokenize{p10-}\allowbreak\detokenize{ccc}}} & Cartesian closed category & -- & -- \\
\hyperref[thm:p10-exponential-adjunction]{\texttt{\detokenize{thm:}\allowbreak\detokenize{p10-}\allowbreak\detokenize{exponential-}\allowbreak\detokenize{adjunction}}} & Exponential adjunction and examples & proved & \texttt{\detokenize{annex:}\allowbreak\detokenize{p10.cat.6}} \\
\hyperref[def:p10-stlc-beta-eta]{\texttt{\detokenize{def:}\allowbreak\detokenize{p10-}\allowbreak\detokenize{stlc-}\allowbreak\detokenize{beta-}\allowbreak\detokenize{eta}}} & Typed \(\beta\eta\) equality and simultaneous substitution & -- & -- \\
\hyperref[thm:p10-stlc-equational-substitution]{\texttt{\detokenize{thm:}\allowbreak\detokenize{p10-}\allowbreak\detokenize{stlc-}\allowbreak\detokenize{equational-}\allowbreak\detokenize{substitution}}} & Typing and equality under substitution & proved & \texttt{\detokenize{annex:}\allowbreak\detokenize{p10.lambda.1}} \\
\hyperref[def:p10-stlc-context-category]{\texttt{\detokenize{def:}\allowbreak\detokenize{p10-}\allowbreak\detokenize{stlc-}\allowbreak\detokenize{context-}\allowbreak\detokenize{category}}} & The category of typed contexts & -- & -- \\
\hyperref[thm:p10-stlc-context-ccc]{\texttt{\detokenize{thm:}\allowbreak\detokenize{p10-}\allowbreak\detokenize{stlc-}\allowbreak\detokenize{context-}\allowbreak\detokenize{ccc}}} & Cartesian closed structure on contexts & proved & \texttt{\detokenize{annex:}\allowbreak\detokenize{p10.lambda.2}} \\
\hyperref[def:p10-stlc-ccc-interpretation]{\texttt{\detokenize{def:}\allowbreak\detokenize{p10-}\allowbreak\detokenize{stlc-}\allowbreak\detokenize{ccc-}\allowbreak\detokenize{interpretation}}} & Interpretation of arrow-only STLC in a CCC & -- & -- \\
\hyperref[thm:p10-stlc-ccc-completeness]{\texttt{\detokenize{thm:}\allowbreak\detokenize{p10-}\allowbreak\detokenize{stlc-}\allowbreak\detokenize{ccc-}\allowbreak\detokenize{completeness}}} & Substitution, soundness, and completeness for CCC equality & proved & \texttt{\detokenize{annex:}\allowbreak\detokenize{p10.lambda.3}} \\
\hyperref[def:p10-internal-ccc-language]{\texttt{\detokenize{def:}\allowbreak\detokenize{p10-}\allowbreak\detokenize{internal-}\allowbreak\detokenize{ccc-}\allowbreak\detokenize{language}}} & The named internal language of a small CCC & -- & -- \\
\hyperref[def:p10-internal-ccc-diagram]{\texttt{\detokenize{def:}\allowbreak\detokenize{p10-}\allowbreak\detokenize{internal-}\allowbreak\detokenize{ccc-}\allowbreak\detokenize{diagram}}} & The full true equational diagram & -- & -- \\
\hyperref[thm:p10-internal-ccc-comparisons]{\texttt{\detokenize{thm:}\allowbreak\detokenize{p10-}\allowbreak\detokenize{internal-}\allowbreak\detokenize{ccc-}\allowbreak\detokenize{comparisons}}} & Named type and context comparisons & proved & \texttt{\detokenize{annex:}\allowbreak\detokenize{p10.internal.1}} \\
\hyperref[thm:p10-internal-ccc-equivalence]{\texttt{\detokenize{thm:}\allowbreak\detokenize{p10-}\allowbreak\detokenize{internal-}\allowbreak\detokenize{ccc-}\allowbreak\detokenize{equivalence}}} & Evaluation recovers the given category & proved & \texttt{\detokenize{annex:}\allowbreak\detokenize{p10.internal.2}} \\
\hyperref[thm:p10-internal-ccc-structure]{\texttt{\detokenize{thm:}\allowbreak\detokenize{p10-}\allowbreak\detokenize{internal-}\allowbreak\detokenize{ccc-}\allowbreak\detokenize{structure}}} & Preservation of cartesian closed structure & proved & \texttt{\detokenize{annex:}\allowbreak\detokenize{p10.internal.3}} \\
\hyperref[def:p10-lawvere-theory]{\texttt{\detokenize{def:}\allowbreak\detokenize{p10-}\allowbreak\detokenize{lawvere-}\allowbreak\detokenize{theory}}} & A one-sorted Lawvere theory & -- & -- \\
\hyperref[def:p10-lawvere-presentation]{\texttt{\detokenize{def:}\allowbreak\detokenize{p10-}\allowbreak\detokenize{lawvere-}\allowbreak\detokenize{presentation}}} & Equational presentations with specified finite contexts & -- & -- \\
\hyperref[thm:p10-lawvere-term-theory]{\texttt{\detokenize{thm:}\allowbreak\detokenize{p10-}\allowbreak\detokenize{lawvere-}\allowbreak\detokenize{term-}\allowbreak\detokenize{theory}}} & The Lawvere theory of an equational presentation & proved & \texttt{\detokenize{annex:}\allowbreak\detokenize{p10.lawvere.1}} \\
\hyperref[def:p10-lawvere-model]{\texttt{\detokenize{def:}\allowbreak\detokenize{p10-}\allowbreak\detokenize{lawvere-}\allowbreak\detokenize{model}}} & Models and their underlying sets & -- & -- \\
\hyperref[thm:p10-lawvere-models]{\texttt{\detokenize{thm:}\allowbreak\detokenize{p10-}\allowbreak\detokenize{lawvere-}\allowbreak\detokenize{models}}} & Equational algebras and functorial models & proved & \texttt{\detokenize{annex:}\allowbreak\detokenize{p10.lawvere.2}} \\
\hyperref[def:p10-lawvere-structural-collapse]{\texttt{\detokenize{def:}\allowbreak\detokenize{p10-}\allowbreak\detokenize{lawvere-}\allowbreak\detokenize{structural-}\allowbreak\detokenize{collapse}}} & A structural functor can identify projections & -- & -- \\
\hyperref[def:p10-lawvere-named-presentation]{\texttt{\detokenize{def:}\allowbreak\detokenize{p10-}\allowbreak\detokenize{lawvere-}\allowbreak\detokenize{named-}\allowbreak\detokenize{presentation}}} & The named operation presentation & -- & -- \\
\hyperref[thm:p10-lawvere-presentation-recovery]{\texttt{\detokenize{thm:}\allowbreak\detokenize{p10-}\allowbreak\detokenize{lawvere-}\allowbreak\detokenize{presentation-}\allowbreak\detokenize{recovery}}} & Recovery of an abstract theory from its named equations & proved & \texttt{\detokenize{annex:}\allowbreak\detokenize{p10.lawvere-}\allowbreak\detokenize{present.1}} \\
\hyperref[def:p10-lawvere-theory-morphism]{\texttt{\detokenize{def:}\allowbreak\detokenize{p10-}\allowbreak\detokenize{lawvere-}\allowbreak\detokenize{theory-}\allowbreak\detokenize{morphism}}} & Morphisms of one-sorted theories & -- & -- \\
\hyperref[thm:p10-lawvere-model-restriction]{\texttt{\detokenize{thm:}\allowbreak\detokenize{p10-}\allowbreak\detokenize{lawvere-}\allowbreak\detokenize{model-}\allowbreak\detokenize{restriction}}} & Theory morphisms and restriction of models & proved & \texttt{\detokenize{annex:}\allowbreak\detokenize{p10.lawvere-}\allowbreak\detokenize{present.2}} \\
\hyperref[thm:p10-lawvere-finite-free]{\texttt{\detokenize{thm:}\allowbreak\detokenize{p10-}\allowbreak\detokenize{lawvere-}\allowbreak\detokenize{finite-}\allowbreak\detokenize{free}}} & Finite free models are representable & proved & \texttt{\detokenize{annex:}\allowbreak\detokenize{p10.lawvere-}\allowbreak\detokenize{free.1}} \\
\hyperref[thm:p10-lawvere-free-adjunction]{\texttt{\detokenize{thm:}\allowbreak\detokenize{p10-}\allowbreak\detokenize{lawvere-}\allowbreak\detokenize{free-}\allowbreak\detokenize{adjunction}}} & Free models on sets and the underlying-set adjunction & proved & \texttt{\detokenize{annex:}\allowbreak\detokenize{p10.lawvere-}\allowbreak\detokenize{free.2}} \\
\hyperref[thm:p10-lawvere-empty-carriers]{\texttt{\detokenize{thm:}\allowbreak\detokenize{p10-}\allowbreak\detokenize{lawvere-}\allowbreak\detokenize{empty-}\allowbreak\detokenize{carriers}}} & The exact empty-carrier boundary & proved & \texttt{\detokenize{annex:}\allowbreak\detokenize{p10.lawvere-}\allowbreak\detokenize{free.3}} \\
\hyperref[def:p10-first-order-hyperdoctrine]{\texttt{\detokenize{def:}\allowbreak\detokenize{p10-}\allowbreak\detokenize{first-}\allowbreak\detokenize{order-}\allowbreak\detokenize{hyperdoctrine}}} & A first-order posetal hyperdoctrine & -- & -- \\
\hyperref[thm:p10-doctrine-quantifier-laws]{\texttt{\detokenize{thm:}\allowbreak\detokenize{p10-}\allowbreak\detokenize{doctrine-}\allowbreak\detokenize{quantifier-}\allowbreak\detokenize{laws}}} & Order laws of the two quantifiers & proved & \texttt{\detokenize{annex:}\allowbreak\detokenize{p10.quant.1}} \\
\hyperref[thm:p10-doctrine-frobenius]{\texttt{\detokenize{thm:}\allowbreak\detokenize{p10-}\allowbreak\detokenize{doctrine-}\allowbreak\detokenize{frobenius}}} & Frobenius from implication-preserving reindexing & proved & \texttt{\detokenize{annex:}\allowbreak\detokenize{p10.quant.2}} \\
\hyperref[def:p10-doctrine-equality]{\texttt{\detokenize{def:}\allowbreak\detokenize{p10-}\allowbreak\detokenize{doctrine-}\allowbreak\detokenize{equality}}} & Parameterized equality predicates & -- & -- \\
\hyperref[thm:p10-doctrine-equality-laws]{\texttt{\detokenize{thm:}\allowbreak\detokenize{p10-}\allowbreak\detokenize{doctrine-}\allowbreak\detokenize{equality-}\allowbreak\detokenize{laws}}} & Reflexivity and replacement of equals & proved & \texttt{\detokenize{annex:}\allowbreak\detokenize{p10.quant.3}} \\
\hyperref[def:p10-inhabited-sort]{\texttt{\detokenize{def:}\allowbreak\detokenize{p10-}\allowbreak\detokenize{inhabited-}\allowbreak\detokenize{sort}}} & An internally inhabited sort & -- & -- \\
\hyperref[thm:p10-doctrine-inhabited]{\texttt{\detokenize{thm:}\allowbreak\detokenize{p10-}\allowbreak\detokenize{doctrine-}\allowbreak\detokenize{inhabited}}} & Inhabitance and reflection along context extension & proved & \texttt{\detokenize{annex:}\allowbreak\detokenize{p10.quant.4}} \\
\hyperref[def:p10-powerset-quantifiers]{\texttt{\detokenize{def:}\allowbreak\detokenize{p10-}\allowbreak\detokenize{powerset-}\allowbreak\detokenize{quantifiers}}} & The powerset fibre and quantifiers & -- & -- \\
\hyperref[thm:p10-powerset-adjunctions]{\texttt{\detokenize{thm:}\allowbreak\detokenize{p10-}\allowbreak\detokenize{powerset-}\allowbreak\detokenize{adjunctions}}} & The two powerset adjunctions and their function laws & proved & \texttt{\detokenize{annex:}\allowbreak\detokenize{p10.quant-}\allowbreak\detokenize{set.1}} \\
\hyperref[def:p10-powerset-fibres]{\texttt{\detokenize{def:}\allowbreak\detokenize{p10-}\allowbreak\detokenize{powerset-}\allowbreak\detokenize{fibres}}} & A universal quantifier includes an empty fibre & -- & -- \\
\hyperref[thm:p10-powerset-base-change]{\texttt{\detokenize{thm:}\allowbreak\detokenize{p10-}\allowbreak\detokenize{powerset-}\allowbreak\detokenize{base-}\allowbreak\detokenize{change}}} & Beck--Chevalley and Frobenius for powersets & proved & \texttt{\detokenize{annex:}\allowbreak\detokenize{p10.quant-}\allowbreak\detokenize{set.2}} \\
\hyperref[def:p10-powerset-commuting-square]{\texttt{\detokenize{def:}\allowbreak\detokenize{p10-}\allowbreak\detokenize{powerset-}\allowbreak\detokenize{commuting-}\allowbreak\detokenize{square}}} & A commuting square without base change & -- & -- \\
\hyperref[thm:p10-powerset-equality]{\texttt{\detokenize{thm:}\allowbreak\detokenize{p10-}\allowbreak\detokenize{powerset-}\allowbreak\detokenize{equality}}} & Powerset equality with parameters & proved & \texttt{\detokenize{annex:}\allowbreak\detokenize{p10.quant-}\allowbreak\detokenize{set.3}} \\
\hyperref[def:p10-contextual-first-order]{\texttt{\detokenize{def:}\allowbreak\detokenize{p10-}\allowbreak\detokenize{contextual-}\allowbreak\detokenize{first-}\allowbreak\detokenize{order}}} & The contextual first-order language and calculus & -- & -- \\
\hyperref[def:p10-contextual-doctrine-interpretation]{\texttt{\detokenize{def:}\allowbreak\detokenize{p10-}\allowbreak\detokenize{contextual-}\allowbreak\detokenize{doctrine-}\allowbreak\detokenize{interpretation}}} & Interpreting a contextual first-order signature & -- & -- \\
\hyperref[thm:p10-doctrine-substitution]{\texttt{\detokenize{thm:}\allowbreak\detokenize{p10-}\allowbreak\detokenize{doctrine-}\allowbreak\detokenize{substitution}}} & Terms, formulas, and contextual substitution & proved & \texttt{\detokenize{annex:}\allowbreak\detokenize{p10.quant-}\allowbreak\detokenize{fol.1}} \\
\hyperref[thm:p10-doctrine-soundness]{\texttt{\detokenize{thm:}\allowbreak\detokenize{p10-}\allowbreak\detokenize{doctrine-}\allowbreak\detokenize{soundness}}} & Soundness of contextual intuitionistic first-order logic & proved & \texttt{\detokenize{annex:}\allowbreak\detokenize{p10.quant-}\allowbreak\detokenize{fol.2}} \\
\hyperref[thm:p10-doctrine-tarski]{\texttt{\detokenize{thm:}\allowbreak\detokenize{p10-}\allowbreak\detokenize{doctrine-}\allowbreak\detokenize{tarski}}} & The powerset interpretation and nonempty sorts & proved & \texttt{\detokenize{annex:}\allowbreak\detokenize{p10.quant-}\allowbreak\detokenize{fol.3}} \\
\hyperref[def:p10-contextual-adjoint-presentations]{\texttt{\detokenize{def:}\allowbreak\detokenize{p10-}\allowbreak\detokenize{contextual-}\allowbreak\detokenize{adjoint-}\allowbreak\detokenize{presentations}}} & The common basis and two presentations & -- & -- \\
\hyperref[thm:p10-contextual-predicate-transport]{\texttt{\detokenize{thm:}\allowbreak\detokenize{p10-}\allowbreak\detokenize{contextual-}\allowbreak\detokenize{predicate-}\allowbreak\detokenize{transport}}} & Syntactic transport of predicate sequents & proved & \texttt{\detokenize{annex:}\allowbreak\detokenize{p10.adjoint.1}} \\
\hyperref[thm:p10-contextual-adjoint-equivalence]{\texttt{\detokenize{thm:}\allowbreak\detokenize{p10-}\allowbreak\detokenize{contextual-}\allowbreak\detokenize{adjoint-}\allowbreak\detokenize{equivalence}}} & Interderivability of the two contextual presentations & proved & \texttt{\detokenize{annex:}\allowbreak\detokenize{p10.adjoint.2}} \\
\hyperref[def:p10-contextual-parameter-free-equality]{\texttt{\detokenize{def:}\allowbreak\detokenize{p10-}\allowbreak\detokenize{contextual-}\allowbreak\detokenize{parameter-}\allowbreak\detokenize{free-}\allowbreak\detokenize{equality}}} & The parameter-free equality presentation & -- & -- \\
\hyperref[thm:p10-contextual-equality-reduction]{\texttt{\detokenize{thm:}\allowbreak\detokenize{p10-}\allowbreak\detokenize{contextual-}\allowbreak\detokenize{equality-}\allowbreak\detokenize{reduction}}} & Recovering parameterized equality from the reduced rule & proved & \texttt{\detokenize{annex:}\allowbreak\detokenize{p10.adjoint.3}} \\
\hyperref[def:p10-presheaf-category]{\texttt{\detokenize{def:}\allowbreak\detokenize{p10-}\allowbreak\detokenize{presheaf-}\allowbreak\detokenize{category}}} & The category of presheaves on a small category & -- & -- \\
\hyperref[def:p10-subpresheaf]{\texttt{\detokenize{def:}\allowbreak\detokenize{p10-}\allowbreak\detokenize{subpresheaf}}} & Subpresheaves and their order & -- & -- \\
\hyperref[thm:p10-presheaf-limits]{\texttt{\detokenize{thm:}\allowbreak\detokenize{p10-}\allowbreak\detokenize{presheaf-}\allowbreak\detokenize{limits}}} & Small limits and colimits of presheaves & proved & \texttt{\detokenize{annex:}\allowbreak\detokenize{p10.psh.1}} \\
\hyperref[thm:p10-presheaf-mono-epi]{\texttt{\detokenize{thm:}\allowbreak\detokenize{p10-}\allowbreak\detokenize{presheaf-}\allowbreak\detokenize{mono-}\allowbreak\detokenize{epi}}} & Maps and subobjects of presheaves & proved & \texttt{\detokenize{annex:}\allowbreak\detokenize{p10.psh.2}} \\
\hyperref[thm:p10-presheaf-exponentials]{\texttt{\detokenize{thm:}\allowbreak\detokenize{p10-}\allowbreak\detokenize{presheaf-}\allowbreak\detokenize{exponentials}}} & The exponential presheaf and its evaluation & proved & \texttt{\detokenize{annex:}\allowbreak\detokenize{p10.psh.3}} \\
\hyperref[def:p10-presheaf-sieves]{\texttt{\detokenize{def:}\allowbreak\detokenize{p10-}\allowbreak\detokenize{presheaf-}\allowbreak\detokenize{sieves}}} & Sieves and the presheaf of truth values & -- & -- \\
\hyperref[thm:p10-presheaf-topos]{\texttt{\detokenize{thm:}\allowbreak\detokenize{p10-}\allowbreak\detokenize{presheaf-}\allowbreak\detokenize{topos}}} & The sieve classifier makes a presheaf category a topos & proved & \texttt{\detokenize{annex:}\allowbreak\detokenize{p10.psh.4}} \\
\hyperref[def:p10-presheaf-power-object]{\texttt{\detokenize{def:}\allowbreak\detokenize{p10-}\allowbreak\detokenize{presheaf-}\allowbreak\detokenize{power-}\allowbreak\detokenize{object}}} & The power object and its membership relation & -- & -- \\
\hyperref[thm:p10-presheaf-power-objects]{\texttt{\detokenize{thm:}\allowbreak\detokenize{p10-}\allowbreak\detokenize{presheaf-}\allowbreak\detokenize{power-}\allowbreak\detokenize{objects}}} & Classification by a power object & proved & \texttt{\detokenize{annex:}\allowbreak\detokenize{p10.psh.5}} \\
\hyperref[def:p10-presheaf-not-pointwise]{\texttt{\detokenize{def:}\allowbreak\detokenize{p10-}\allowbreak\detokenize{presheaf-}\allowbreak\detokenize{not-}\allowbreak\detokenize{pointwise}}} & Exponential and power objects depend on restrictions & -- & -- \\
\end{longtable}
\endgroup

